\documentclass[11pt,reqno]{amsart}
\usepackage{amssymb,mathtools}
\usepackage{enumitem}
\usepackage{microtype}
\emergencystretch=1em
\usepackage{geometry}
\geometry{}
\usepackage[hidelinks]{hyperref}
\let\affil\address

\newtheorem{theorem}{Theorem}[section]
\newtheorem{lemma}[theorem]{Lemma}
\newtheorem{proposition}[theorem]{Proposition}
\newtheorem{corollary}[theorem]{Corollary}
\newtheorem{conjecture}[theorem]{Conjecture}
\theoremstyle{definition}
\newtheorem{definition}[theorem]{Definition}
\theoremstyle{remark}
\newtheorem{remark}[theorem]{Remark}
\numberwithin{equation}{section}

\newcommand{\R}{\mathbb R}
\newcommand{\E}{\mathbb E}
\newcommand{\Pp}{\mathbb P}
\newcommand{\cF}{\mathfrak F}
\newcommand{\Lip}{\operatorname{Lip}}
\newcommand{\sgn}{\operatorname{sgn}}

\title[Empirical Selection on Finsler Manifolds]
{Deterministic Empirical Selection on Finsler Manifolds}
\author{Wentao Lyu}
\affil{Shanghai Center for Mathematical Sciences,
Fudan University}
\subjclass[2020]{Primary 49Q22; Secondary 60D05, 60G55}
\keywords{Optimal transport, empirical matching, Palm measures}
\thanks{The accompanying Lean development is available at
\url{https://github.com/LvDAO/OT-Selector-Finsler}. The general empirical selector theorem is formally verified in Lean 4
under the original smooth Finsler assumptions; see the repository for the exact
statements and verification commands.}
\date{}

\begin{document}

\begin{abstract}
We prove deterministic empirical selection for genuine smooth strongly
convex length distances in dimension \(d\ge6\). On a connected complete
smooth reversible strongly convex Finsler manifold, let the two marginal
probabilities have nonnegative Lipschitz densities and disjoint compact
supports. Independent iid sampling makes the entire empirical optimal
set converge in probability to one intrinsic deterministic population
optimum, including fractional empirical plans. No prescribed projection
balance, density lower bound, support-boundary regularity, curvature sign
or topology restriction is required. The theorem includes Euclidean
space, curved Riemannian manifolds and Finsler tori.

The continuous optimal face decomposes into forward conditional
couplings on intrinsic active geodesic arcs. A genuine Jacobi expansion
and original-input Poisson estimates control complete path currents even
at focal endpoints. Their centered prefixes and nonnegative cut
occupations permit stagewise legal repair with independent reserves,
without requiring a displacement second moment. This identifies the
weighted and finite-second-moment stationary pressures for the full
interleaved conditional populations.

An exact finite cavity identity relates iid assignment to semidiscrete
transport with Dirichlet target weights. Comparing opposite perturbations
in the same diffuse-source environment gives differentiability of the
local paired pressure. Rough-chart pressure localization and countably
many test functions then identify the common limit of all empirical
optima. The general-population problem in dimensions four and five
remains open.
\end{abstract}

\maketitle

\section{Introduction}
\label{sec:introduction}

Optimal transport with a length cost can have many minimizers, even for
compactly supported measures with densities. We ask whether independent
empirical sampling selects an intrinsic deterministic optimal plan on a
complete smooth reversible strongly convex Finsler manifold, with
disjoint compact source and target regions. Theorem~\ref{thm:general-finsler-selector}
answers this question in dimension \(d\ge6\).
The geometric and fixed-root sampling inputs hold for
every \(d\ge2\). The cost is the genuine length distance, and the question
concerns the entire empirical optimal set.

\begin{theorem}[General Finsler empirical distance selection]
\label{thm:general-finsler-selector}
Let \((M,F)\) be a connected complete smooth reversible strongly
convex Finsler manifold of dimension \(d\ge6\), with smooth positive
volume \(m\). Let \(\mu=\rho_A m\), \(\nu=\rho_B m\) be probability
measures with nonnegative Lipschitz densities and compact disjoint
supports \(A,B\). For independent iid empirical probabilities
\(\mu_n,\nu_n\), let \(\operatorname{Opt}_n\) be the entire optimal
coupling set for the original length cost \(d_F\).

There is a deterministic population distance-optimal plan \(\pi_*\)
such that, for every metric \(d_{\rm weak}\) inducing weak convergence
on probabilities on \(A\times B\),
\begin{equation}
 \sup_{\gamma\in\operatorname{Opt}_n}
       d_{\rm weak}(\gamma,\pi_*)\longrightarrow0
                         \qquad\text{in probability}.
 \label{eq:general-finsler-entire-selector}
\end{equation}
Fractional optima are included. The physical limit depends only on
\((F,\mu,\nu)\), independently of optimal potentials, coordinates,
labels, and finite-sample tie breaking.
No density lower bound, support-boundary regularity, predetermined
mass balance, geometric separator, endpoint nonfocusing, or
raywise interior-support condition is assumed.
\end{theorem}

The limit is a transport plan; no Monge representation or identification
with an entropy selector is asserted. The theorem includes every allowed
Euclidean population, complete Riemannian manifolds of arbitrary curvature,
and reversible strongly convex Finsler tori.

\subsection{Intrinsic quantities and proof organization}
\label{subsec:general-proof-organization}

Fix a population optimal \(1\)-Lipschitz potential \(\phi\), oriented
so that transport increases \(\phi\). Its nonnegative calibrated slack is
\begin{equation}
 e(x,y)=d_F(x,y)-\phi(y)+\phi(x)\ge0.
 \label{eq:general-finsler-slack}
\end{equation}
The continuous optimal face consists exactly of couplings with zero
slack integral. For \(q=d-1\), a test \(f\in C(A\times B)\), and
\(t\in\R\), define the original empirical pressure by
\begin{equation}
 V_n(t)=\min_{\gamma\in\Pi(\mu_n,\nu_n)}
       \left\{n^{2/q}\int e\,d\gamma+t\int f\,d\gamma\right\}.
 \label{eq:general-finsler-pressure}
\end{equation}
At \(t=0\), its minimizers are exactly the original distance-optimal
set: the potential integral depends only on the marginals. The quadratic
costs below describe the local expansion of \(e\); the transport problem
always uses \(d_F\).

The proof has four stages.
\begin{enumerate}[label=\textup{(\roman*)},leftmargin=*]
\item \emph{Continuous structure and local geometry.}
Propositions~\ref{prop:general-regular-pairs}
and~\ref{prop:general-optimal-face} disintegrate the entire optimal
face along calibrated nonbranching geodesics. Jacobi second variation
identifies the quadratic slack, and a second Jacobi solution controls
focal endpoints (Lemma~\ref{lem:curved-jacobi-second-solution}).
Finite measurable charts and compatible compact cores provide the
localization; no global smooth foliation is assumed.

\item \emph{Recovery of the local stationary value.}
Complete path currents, centered prefixes and nonnegative cut
occupations give legal repair with independent reserves
(Theorem~\ref{thm:interleaved-focal-traffic} and
Proposition~\ref{prop:interleaved-path-screening}). This equates the
weighted and finite-second-moment pressures on the full conditional
populations. Count defects in a box of volume \(L^q\) have size
\(O(L^{q/2})\); quadratic repair has normalized cost
\(O(L^{2-q/2})\), which vanishes for \(q>4\).

\item \emph{Convergence of the original pressure.}
Section~\ref{sec:global-weighted-pressure} gives a lower bound from
actual optimizer scenes and an expectation upper bound from admissible
recovery. The local value equality in
Theorem~\ref{thm:general-focal-value-equality} closes this bracket,
including interleaved populations and focal endpoints.

\item \emph{Differentiation and the entire empirical face.}
The exact iid cavity identity and mean-plan formula
(Theorem~\ref{thm:iid-cavity-dirichlet} and
Proposition~\ref{prop:iid-cavity-mean}) use semidiscrete transport
with Dirichlet target weights. Opposite perturbations in the same
diffuse-source environment give pressure differentiability. Fixed
finite-difference bounds control every empirical optimizer at once;
countably many continuous tests reconstruct the intrinsic limit in
Section~\ref{sec:finsler-extension}.
\end{enumerate}

Regular separated tubes and homogeneous matching costs below are local
proof ingredients. Their temporary regularity and ordering conditions
are obtained by localization, clipping and finite active-block
decomposition; they add no assumptions to
Theorem~\ref{thm:general-finsler-selector}.
The separated-square question in \cite[Section~7, Problem~4]{delBarrioEtAl2025}
was treated by Lyu~\cite{Lyu2026EmpiricalW1}: for independent uniform
samples from \([0,1]^2\) and \([2,3]\times[0,1]\), the empirical
optimal plans converge weakly in probability to a deterministic plan,
identified through a stationary local matching and its source-Palm law.


\section{Continuous optimal transport and intrinsic active components}\label{sec:continuous-geometry}

\subsection{The Finsler distance and the empirical problem}

Let \(M\) be a connected smooth manifold of dimension \(d\ge6\).
A smooth reversible strongly convex Finsler structure is a function
\(F:TM\to[0,\infty)\), smooth away from the zero section, whose
restriction to each tangent space is a norm and for which
\(D_v^2(F_x^2/2)\) is positive definite at nonzero vectors. Assume
the induced length distance
\[
d_F(x,y)=\inf_{\gamma(0)=x,\,\gamma(1)=y}
\int_0^1F_{\gamma(r)}(\dot\gamma(r))\,dr
\]
is complete. Fix a smooth positive volume measure \(m\). Let
\(\mu=\rho_A m\) and \(\nu=\rho_Bm\) be probability measures
with nonnegative Lipschitz densities and compact disjoint supports
\(A=\operatorname{supp}\mu\), \(B=\operatorname{supp}\nu\).
No boundary regularity or finite-component condition is imposed on
these supports. Their metric separation satisfies
\[
\delta_0:=\min_{x\in A,\,y\in B}d_F(x,y)>0.
\]
In particular every calibrated source--target segment has length
at least \(\delta_0\). On a common unit-speed calibrated ray,
source and target mark supports are separated by at least
\(\delta_0\) in the longitudinal coordinate. This removes
arbitrarily short cross-marginal transports. It does not require
that all source marks precede all target marks: the two regions
can have several interleaving components along the same ray.

Write \(\mathcal O_n^F=\operatorname{Opt}_n\) for the set of exact empirical
distance-optimal couplings, and let \(d_{\rm BL}\) be the bounded Lipschitz
distance induced by the product length metric on \(A\times B\).

\begin{proposition}[Consistency with the continuum optimal set]
For any continuous length distance and compactly supported marginals,
independent iid samples satisfy, almost surely,
\[
\sup_{\pi\in\mathcal O_n^F}\inf_{\gamma\in\mathcal O^F}
d_{\rm BL}(\pi,\gamma)\longrightarrow0,
\]
where \(\mathcal O^F\) is the continuum optimal set.
\end{proposition}

\begin{proof}
For each bounded continuous test, Hoeffding's inequality gives
summable probabilities for every fixed positive deviation of its
empirical average. Borel--Cantelli, on a countable uniformly dense
family of tests and rational positive deviations, gives almost-sure
weak convergence of each empirical marginal. This argument needs
no independence between different sample sizes.
On compact supports this also gives convergence in \(W_1\). The
triangle inequality yields
\[
|W_1(\mu_n,\nu_n)-W_1(\mu,\nu)|
\le W_1(\mu_n,\mu)+W_1(\nu_n,\nu)\longrightarrow0.
\]
If the assertion failed, compactness would supply a subsequence of
optimal empirical plans bounded away from \(\mathcal O^F\) and
converging weakly. Its limit has the prescribed marginals, and
continuity of the cost makes its cost \(W_1(\mu,\nu)\), a
contradiction.
\end{proof}

\subsection{General geometric reductions}
\label{subsec:general-geometric-reductions}

The following reductions hold under the geometric assumptions above
already in dimension \(d\ge2\). They require only absolute continuity
of the two marginals and disjoint compact supports.

\begin{proposition}[Finite uniqueness for the genuine distance]
\label{prop:general-finite-uniqueness}
Almost surely, for every \(n\), independent iid sampling from
\(\mu\) and \(\nu\) gives a unique minimum-cost permutation for
\(d_F\), and \(\mathcal O_n^F\) consists of its empirical coupling.
\end{proposition}

\begin{proof}
For fixed \(y\), the function \(x\mapsto d_F(x,y)\) is locally
Lipschitz. Rademacher's theorem and Fubini imply that for
\(\mu\)-almost every \(x\), the covector
\(P_x(y):=D_xd_F(x,y)\) exists for \(\nu\)-almost every \(y\).
At such a pair, every initial unit velocity \(v\) of a minimizing
geodesic from \(x\) to \(y\) satisfies
\[
F_x(v)=1,\qquad P_x(y)(v)=-1.
\]
Strong convexity makes this velocity unique. Thus, for fixed \(x\)
and fixed covector \(p\), the fiber \(\{y:P_x(y)=p\}\) lies on
one geodesic starting from \(x\). Its image is a countable union
of smooth finite arcs and has \(m\)-measure zero since \(d\ge2\).
Consequently the conditional law of \(P_x(Y)\) is non-atomic.
For independent \(X\sim\mu\) and \(Y,Z\sim\nu\), it follows that
\(P_X(Y)\ne P_X(Z)\) almost surely.

For distinct permutations \(\sigma,\tau\), consider the locally
Lipschitz function of all sample points
\[
H(X,Y)=\sum_{i=1}^n
\bigl[d_F(X_i,Y_{\sigma(i)})-d_F(X_i,Y_{\tau(i)})\bigr].
\]
Choose \(i\) with \(\sigma(i)\ne\tau(i)\). The preceding argument
gives \(D_{X_i}H\ne0\) almost surely. A locally Lipschitz function
has derivative zero at almost every point of its zero level set:
at a differentiability and density point of that set, a nonzero
derivative would exclude a fixed cone of nearby points from it.
Absolute continuity of the joint input law therefore gives
\(\mathbb P(H=0)=0\). A finite union over permutation pairs,
the Birkhoff representation of empirical couplings, and a countable
intersection over \(n\) prove the assertion.
\end{proof}

\begin{proposition}[Common regular rays and empirical cut exclusion]
\label{prop:general-regular-pairs}
Fix one optimal global \(1\)-Lipschitz potential \(\phi\), with
\[
W_1(\mu,\nu)=\int\phi\,d(\nu-\mu),\qquad
s(x,y):=d_F(x,y)-\phi(y)+\phi(x)\ge0.
\]
Every continuum optimal plan is concentrated on the common calibrated
relation \(\Gamma_\phi:=\{s=0\}\), and almost every such pair joins
two interior points of the same nonbranching calibrated ray.
Let \(\mathcal R\) be the open set of distinct pairs joined by a
unique nonconjugate minimizing geodesic, and put
\(C=(A\times B)\setminus\mathcal R\). The true distance is smooth
on \(\mathcal R\), and every continuum optimal plan gives \(C\)
zero mass. For the closed product-metric neighborhoods \(C_r\)
of \(C\), with \(C_r=\varnothing\) if \(C=\varnothing\),
\[
\lim_{r\downarrow0}\limsup_{n\to\infty}
\sup_{\pi\in\mathcal O_n^F}\pi(C_r)=0
\qquad\text{almost surely}.
\]
More precisely, for each \(\epsilon>0\) there are compact
\(K_A\subset A\), \(K_B\subset B\), each losing at most
\(\epsilon\) of its marginal mass, and constants \(r,\eta>0\)
such that, for every \(n\) and \(\pi\in\mathcal O_n^F\),
\begin{equation}
\label{eq:general-cut-slack}
\begin{aligned}
\pi(C_r)&\le
\mu_n(A\setminus K_A)+\nu_n(B\setminus K_B)
+\frac{a_n}{\eta},\\
a_n&:=W_1(\mu_n,\nu_n)-\int\phi\,d(\nu_n-\mu_n).
\end{aligned}
\end{equation}
\end{proposition}

\begin{proof}
Compact-support metric duality supplies \(\phi\). For every
population optimum \(\pi\), \(\int s\,d\pi=0\), so the same
\(\Gamma_\phi\) supports all of them. Along each calibrated
unit-speed minimizing segment, \(\phi\) increases at unit speed.
At a differentiability source point its initial velocity is the
unique unit vector \(v\) with \(d\phi(v)=1\). Moreover, calibrated
segments cannot meet in their interiors and then branch: calibration
makes their splice minimizing, whereas a corner cannot minimize a
smooth strongly convex length. Geodesic uniqueness then identifies
the oriented rays.

By \cite[Theorem~1.1(i),(ii)]{OhtaNeedles}, the branch set is
\(m\)-null and the volume disintegration on maximal calibrated
rays has conditionals absolutely continuous in arc length.
These parts of Ohta's theorem impose no global Ricci lower bound;
the curvature hypothesis belongs to part~(iii).
One may use an equivalent smooth positive finite reference volume
without changing the relevant null sets. Since each nonbranching
ray has at most two endpoints, their union is also null.
Measurability follows, for example, from the forward and backward
calibrated reach functions \(L_\pm\): for \(t>0\),
\(\{L_\pm\ge t\}\) is closed by properness and continuity, and
the one-sided endpoint sets are Borel. Let \(T_\phi^\circ\) be
the Borel set with nontrivial calibrated predecessor and successor.
Disjointness of the supports and absolute continuity give
\(\mu(T_\phi^\circ)=\nu(T_\phi^\circ)=1\).

If \((x,y)\in\Gamma_\phi\) and \(x,y\in T_\phi^\circ\),
extend the ray past \(y\) to a calibrated point \(z\).
The segment from \(x\) through \(y\) to \(z\) is minimizing,
so \(y\) lies strictly before its cut time from \(x\).
An alternative route to \(y\) would give a minimizing corner
after splicing with the tail to \(z\); a conjugate point at \(y\)
would prevent the longer segment from minimizing by second variation.
Thus \((x,y)\in\mathcal R\). Invertibility of the exponential map
gives the local smooth branch, and compactness of minimizing
geodesics excludes competing branches at converging nearby pairs.
This proves openness of \(\mathcal R\) and smoothness of the
global distance there, as well as \(\pi(C)=0\) for every optimum.

Choose compact \(K_A\subset A\cap T_\phi^\circ\) and
\(K_B\subset B\cap T_\phi^\circ\) losing at most
\(\epsilon\) of their marginal masses. The compact set
\(C\cap(K_A\times K_B)\) is disjoint from \(\Gamma_\phi\).
Continuity and compactness give \(r,\eta>0\) with
\(s\ge\eta\) on \(C_r\cap(K_A\times K_B)\).
If this set is empty for all sufficiently small \(r\), any
positive \(\eta\) works. Since every empirical optimum satisfies
\(\int s\,d\pi=a_n\), splitting according to whether both
endpoints lie in \(K_A\times K_B\) proves
\eqref{eq:general-cut-slack}.
Compact-support empirical convergence gives convergence of
\(W_1(\mu_n,\nu_n)\) and of the two \(\phi\)-averages, hence
\(a_n\to0\) almost surely. The strong law for the fixed complements gives
a limsup of at most \(2\epsilon\); apply this for a countable
sequence \(\epsilon\downarrow0\) to finish.
\end{proof}

By monotonicity of \(C_r\), the conclusion also gives
\(\sup_{\pi\in\mathcal O_n^F}\pi(C_{r_n})\to0\) almost surely
for every sequence \(r_n\to0\): eventually \(C_{r_n}\subset C_r\)
for each fixed \(r>0\). No rate is obtained. These reductions
leave the macroscopic selection of pairings along the common rays
unresolved.

\subsection{The continuous optimal face on the actual ray quotient}
\label{subsec:general-optimal-face}

Assume here only that \(\mu=\rho_A m\) and \(\nu=\rho_B m\)
have nonnegative Lipschitz densities and disjoint compact supports
\(A,B\). No regularity of their support boundaries is required.
Retain the complete smooth reversible strongly convex Finsler geometry,
and put \(\delta_0=\operatorname{dist}_F(A,B)>0\) and
\(D=\operatorname{diam}_F(A\cup B)\). Fix one optimal
1-Lipschitz potential with orientation
\(\phi(y)-\phi(x)=d_F(x,y)\) on calibrated pairs.

\begin{proposition}[The full optimal face on the ray quotient]
\label{prop:general-optimal-face}
There are a common full-marginal-measure Borel interior-ray set \(N\),
a standard Borel ray-label space \(Z\), and a Borel quotient map
\(\ell:N\to Z\), with
\[
 \Lambda=\ell_\#\mu=\ell_\#\nu.
\]
There are absolutely continuous probability kernels
\(\alpha_z,\beta_z\) in the potential coordinate \(r\), and
jointly Borel ray coordinates \(\gamma_z(r)\), such that
\(\phi(\gamma_z(r))=r\) and
\[
 \mu=\int(\gamma_z)_\#\alpha_z\,\Lambda(dz),\qquad
 \nu=\int(\gamma_z)_\#\beta_z\,\Lambda(dz).
\]
The population optimal plans are exactly
\begin{equation}
\label{eq:general-optimal-face}
 \pi=\int(\gamma_z\times\gamma_z)_\#\kappa_z\,\Lambda(dz),
 \qquad \kappa_z\in\Pi(\alpha_z,\beta_z),\quad
 \kappa_z\{a<b\}=1,
\end{equation}
with measurable coupling kernels. Their conditional support spans are
at most \(D\), opposite colors are at least \(\delta_0\) apart,
and every forward pair satisfies \(\delta_0\le b-a\le D\).
The common quotient marginal is derived from optimality.
\end{proposition}
\begin{proof}
Proposition~\ref{prop:general-regular-pairs} gives \(N\); every
calibrated pair in \(N\times N\) belongs to one common oriented ray.
The Borel selector of \cite[Lemmas~4.10--4.11]{OhtaNeedles} is
constant on rays and idempotent. Its image is its Borel fixed-point
set, giving a standard Borel realization of \(Z\).
The injective Borel map \(x\mapsto(\ell(x),\phi(x))\) has a
Borel inverse on its image, which supplies \(\gamma_z(r)\).
Duality says that a coupling is optimal exactly when supported on
\(\Gamma_\phi\). Hence every optimum satisfies
\(\ell(x)=\ell(y)\) almost surely, proving the equality of quotient
marginals independently of which optimum is chosen.

For the conditional absolute continuity, if necessary replace \(m\)
by an equivalent smooth positive probability measure \(m_0=w m\).
Disintegrate it along rays as \(m_{0,z}\,v(dz)\).
By \cite[Theorem~1.1(i),(ii)]{OhtaNeedles},
\(m_{0,z}=h_z(r)\,dr\). Writing \(f_A=d\mu/dm_0\), the
source quotient density is \(a(z)=\int f_A\,dm_{0,z}\), and its
conditional measure is \(f_A m_{0,z}/a(z)\) on \(a(z)>0\).
The target quotient density equals \(a(z)\) by the equality just
proved. This gives the asserted absolutely continuous kernels.

Disintegrating an optimum over its common ray label gives precisely
the forward kernels in \eqref{eq:general-optimal-face}. Conversely
every displayed reconstruction is calibrated and hence optimal.
On a calibrated ray \(d_F(\gamma_z(r),\gamma_z(s))=|r-s|\).
Apply this to marginal-support points and then their limits, using
completeness and compactness, to obtain the span and separation bounds.
Conditional endpoints have zero mass, so extending the coordinate to
such limits creates no ambiguity. All assertions about a given
disintegration hold for \(\Lambda\)-almost every label.
\end{proof}

\subsection{General conditional order bands}

\begin{proposition}[Uniformly finite monochromatic order bands]
\label{prop:ray-finite-bands}
On almost every conditional ray, the combined support has at most
\(1+\lfloor D/\delta_0\rfloor\) maximal monochromatic order-convex
classes. Their compact interval hulls alternate in color and consecutive
hulls are separated by at least \(\delta_0\). A class may have
infinitely many connected components and empty gaps.
Moreover, the quotient admits a finite Borel partition into types on
each of which fixed compact interval containers cover the conditional
supports, alternate in color, and have successive gaps at least
\(\delta_0/2\), with the same bound on their number.
\end{proposition}
\begin{proof}
Starting at the first occupied coordinate, collect points of its color
until the first point of the other color, and repeat. Compactness and
cross-color separation make each resulting class compact. First points
\(t_1<\cdots<t_k\) of successive classes have alternating colors,
so \((k-1)\delta_0\le t_k-t_1\le D\). The process therefore
terminates and gives exactly the maximal order-convex classes. Applying
separation to adjacent last and first points proves the hull-gap claim.

For measurability put \(h=\delta_0/4\) and cover the fixed interval
\([\min_{A\cup B}\phi,\max_{A\cup B}\phi]\) by finitely many
half-open grid cells of width \(h\). Label a cell \(A\), \(B\),
or \(0\) according as its source mass, target mass, or neither is
positive. Both masses cannot be positive in one cell. These symbols
are measurable functions of the conditional kernels, so there are
finitely many Borel word types. Delete zero cells, group consecutive
surviving cells of the same color, and take the closed interval hull
of each group. These are fixed containers on a word type.
For adjacent groups their last and first occupied cells contain
opposite-color support points at least \(\delta_0\) apart; padding
by their cell endpoints loses at most \(2h\). Thus the container
gap is at least \(\delta_0/2\). Selecting one support point per group
gives the same bound on their number. Containers may include empty
regions; they are not asserted to be smooth geometric chart domains.
\end{proof}

\begin{proposition}[Forward feasibility and zero-flux blocks]
\label{prop:ray-forward-criterion}
For compactly supported probability measures \(\alpha,\beta\) on
\(\R\), put \(H(t)=\alpha(({-\infty},t])-\beta(({-\infty},t])\).
A coupling supported on \(a\le b\) exists if and only if \(H\ge0\).
Every such coupling satisfies
\begin{equation}
\label{eq:ray-cut-flux}
                         H(t)=\kappa\{a\le t<b\}.
\end{equation}
For the separated absolutely continuous conditionals above, the set
\(\{H>0\}\) has at most
\(\lfloor(1+\lfloor D/\delta_0\rfloor)/2\rfloor\) nonempty
interval components \(I_j\). Both measures vanish on \(\{H=0\}\),
and \(\alpha(I_j)=\beta(I_j)=w_j>0\). Their forward-coupling set
consists exactly of
\[
 \kappa=\sum_j w_j\kappa_j,\qquad
 \kappa_j\in\Pi(\alpha|_{I_j}/w_j,\beta|_{I_j}/w_j),\qquad
 \kappa_j\{a<b\}=1.
\]
These block decompositions can be made measurable over the quotient.
\end{proposition}
\begin{proof}
If \(a\le b\), subtracting the probabilities of \(b\le t\)
and \(a\le t\) proves \eqref{eq:ray-cut-flux}. Conversely \(H\ge0\)
implies \(Q_\alpha(u)\le Q_\beta(u)\) for generalized quantiles;
their common-uniform coupling is forward. Separated supports make
this inequality strict almost surely.

For absolutely continuous separated measures, \(H\) is absolutely
continuous and \(H'=f_\alpha-f_\beta\). Its derivative vanishes
almost everywhere on its zero set. The two densities are mutually
singular, so both vanish there. Each component \((l_j,r_j)\) of
\(\{H>0\}\) is bounded and has zero endpoint values. Hence its
two masses agree and are positive. Equation~\eqref{eq:ray-cut-flux}
at a countable dense subset of the closed zero set forbids every
cross-component edge, outside one null set. This proves the necessary
block decomposition; the converse follows by adding forward block
couplings with the displayed weights.

The positive-mass color word starts with \(A\) and ends with \(B\),
since otherwise a cut at its first or last band would have negative
\(H\). Within an \(A\)-band \(H\) is nondecreasing, within a
\(B\)-band it is nonincreasing, and in an empty gap it is constant.
Each positive component needs a new \(A\)-band and a \(B\)-band;
Proposition~\ref{prop:ray-finite-bands} gives the bound.
On each of its finite word types, choose a fixed cut in every
inter-container gap. The conditions \(H_z(t)=0\) at these finitely
many cuts are measurable. They identify exactly the nonzero-mass
blocks: a monochromatic band cannot meet two positive components.
Restriction and normalization therefore give measurable block kernels.
\end{proof}

\begin{corollary}[Continuous nonuniqueness]
\label{cor:general-continuum-nonunique}
For compactly supported atomless \(\alpha,\beta\) with \(H\ge0\),
the forward-coupling set is a singleton exactly when
\(\alpha=\beta\), equivalently \(H\equiv0\).
Under Proposition~\ref{prop:general-optimal-face}, the population
problem has a continuum of distinct optimal plans.
\end{corollary}
\begin{proof}
If \(H=0\), equation~\eqref{eq:ray-cut-flux} at rational cuts
forces \(a=b\), giving the unique diagonal coupling. Otherwise any
forward coupling has positive mass on \(a<b\). Cover that set by
bounded open rational rectangles \(U\times V\) with
\(\sup U<\inf V\). Some restriction \(\eta\) has mass \(p>0\).
Its normalized marginals \(\alpha',\beta'\) are atomless. Replacing
\(\eta\) either by \(p\alpha'\otimes\beta'\) or by
\(p(Q_{\alpha'},Q_{\beta'})_\#du\) preserves the marginals and
forward support. These replacements differ: the second is supported
on a graph and the first gives that graph zero mass.

To make the construction global, fix one optimal kernel
\(\kappa_z^0\), enumerate these rational rectangles \(R_j\), and
on \(\{z:H_z\not\equiv0\}\) choose the least \(j\) with
\(\kappa_z^0(R_j)>0\). The index, restricted measures and normalized
marginals are measurable. Generalized quantiles of probability kernels
are jointly measurable, as their sublevel sets can be tested at
rational endpoints. Thus both replacements define measurable forward
kernels, unchanged on the complementary labels.
They differ at every selected label. Uniqueness of disintegration and
the injective ray coordinates imply that their physical reconstructions
differ whenever this label set has positive \(\Lambda\)-measure.
Here it has full measure: the two conditional probabilities have
disjoint supports, so cannot agree. Convex mixtures of the resulting
distinct optimal plans give the asserted continuum.
\end{proof}

The quotient and its labels may depend on the chosen optimal potential;
the physical optimal face in \eqref{eq:general-optimal-face} does not.
Distinct plans can have identical raywise cut fluxes and total length.
No additional homology, simple-connectivity or curvature restriction
was used. On a compact manifold a global smooth unit-eikonal real
function is impossible at its maximum, but the Lipschitz potential and
measurable local ray coordinates used here remain available, including
on a torus. These are exact criteria for the continuous optimal face
and conditional feasibility; they are not criteria for deterministic
selection by independent empirical sampling.

\begin{corollary}[Intrinsic occupation measure and current]
\label{cor:general-occupation-current}
Under Proposition~\ref{prop:general-optimal-face}, let
\(\gamma_{xy}:[0,d_F(x,y)]\to M\) be the unique unit-speed minimizing
segment for almost every pair of an optimal plan \(\pi\). Define its
arclength occupation measure and oriented current by
\begin{equation}
\label{eq:general-intrinsic-occupation}
\begin{aligned}
 \tau_\pi(E)&=\int\!\int_0^{d_F(x,y)}
                   \mathbf1_E(\gamma_{xy}(s))\,ds\,\pi(dx,dy),\\
 J_\pi(\omega)&=\int\!\int_0^{d_F(x,y)}
       \omega_{\gamma_{xy}(s)}(\dot\gamma_{xy}(s))\,ds\,\pi(dx,dy),
\end{aligned}
\end{equation}
for Borel \(E\) and smooth compactly supported one-forms \(\omega\).
Every optimal plan induces the same finite measure \(\tau\) and normal
one-current \(J\), independent also of the optimal potential used to
describe the rays. Moreover,
\[
 \tau\ll m,\qquad \tau(M)=W_1(\mu,\nu),\qquad
 \partial J=\nu-\mu,
 \quad\text{where }\partial J(f):=J(df).
\]
For any fixed optimal potential, put
\(H_z(t)=\alpha_z(({-\infty},t])-\beta_z(({-\infty},t])\).
Then, for every nonnegative Borel \(f\),
\begin{equation}
\label{eq:general-occupation-cut-flux}
\begin{aligned}
 \int f\,d\tau
   &=\int\!\int H_z(t)f(\gamma_z(t))\,dt\,\Lambda(dz),\\
 J(\omega)
   &=\int\!\int H_z(t)
       \omega_{\gamma_z(t)}(\dot\gamma_z(t))\,dt\,\Lambda(dz).
\end{aligned}
\end{equation}
The inner integrals are over the ray interiors.
\end{corollary}
\begin{proof}
On the open regular-pair locus of
Proposition~\ref{prop:general-regular-pairs}, the local inverse exponential
map and geodesic flow make the segment and its unit tangent jointly
Borel in \((x,y,s/d_F(x,y))\). Extending the integrands by zero elsewhere
therefore defines \eqref{eq:general-intrinsic-occupation}; that extension
is null for every optimal plan. Every segment lies in the closed
\(D\)-neighborhood of \(A\cup B\), which is compact by completeness.
Tonelli gives a finite Radon measure with
\(\tau_\pi(M)=\int d_F\,d\pi=W_1(\mu,\nu)\le D\).
Its mass bounds the current against the supremum dual Finsler norm of
\(\omega\). Norm equivalence on this compact set also gives finite
current mass for any auxiliary smooth Riemannian metric; no ambient
orientation is needed.

Disintegrate \(\pi\) as in \eqref{eq:general-optimal-face}.
Its segment is the ray interval from \(a\) to \(b\), and
\eqref{eq:ray-cut-flux} gives
\(\kappa_z\{a\le t<b\}=H_z(t)\).
Tonelli proves the first identity in
\eqref{eq:general-occupation-cut-flux}, and signed Fubini proves the
second. The kernels make \(H_z(t)\) jointly measurable. The ray tangent
is jointly Borel by coordinate difference quotients of the Borel ray
map, since each individual ray is a smooth geodesic. The displayed
identities contain only the fixed marginal kernels, so are independent
of \(\pi\). For two different optimal potentials, apply both identities
to one fixed optimal plan: each equals its intrinsic definition
\eqref{eq:general-intrinsic-occupation}. This also proves potential
independence, without identifying the two ray partitions.

As in Proposition~\ref{prop:general-optimal-face}, choose an equivalent
smooth positive finite volume \(m_0\), and write its interior-ray
disintegration as
\[
 m_0|_N=\int(\gamma_z)_\#(h_z(t)\,dt)\,v(dz),\qquad
 \Lambda(dz)=a(z)\,v(dz).
\]
By \cite[Proposition~6.4(i),(iii)]{OhtaNeedles}, \(h_z>0\) on
almost every ray interior. The occupation formula consequently gives
the density \(a(z)H_z(t)/h_z(t)\) relative to \(m_0\) there,
and zero outside \(N\). Thus \(\tau\ll m_0\ll m\), with finite
integral already proved. Finally, the fundamental theorem of calculus
along every segment gives
\[
 J(df)=\int(f(y)-f(x))\,\pi(dx,dy)=\int f\,d(\nu-\mu).
\]
The boundary has finite mass, so \(J\) is normal.
\end{proof}

This is the ray cut-flux mechanism behind transport-density uniqueness
in norm geometry; see \cite[Section~4 and Theorem~4.1]{FeldmanMcCannDensity}.
It fixes spatial occupation and oriented traffic while leaving the
endpoint coupling free: the distinct plans in
Corollary~\ref{cor:general-continuum-nonunique} induce exactly this same
\(\tau,J\). No uniqueness among arbitrary currents with the same
boundary, or empirical selector, follows from this corollary.

\begin{proposition}[A finite conditional-entropy competitor]
\label{prop:ray-finite-entropy}
In Proposition~\ref{prop:general-optimal-face}, set
\[
\pi_0=\int(\gamma_z\times\gamma_z)_\#
             (\alpha_z\otimes\beta_z)\,\Lambda(dz),
\qquad K=1+\lfloor D/\delta_0\rfloor.
\]
There is an optimal plan \(\bar\pi\) with
\(\operatorname{KL}(\bar\pi\mid\pi_0)\le\log K\).
The reference \(\pi_0\) has marginals \(\mu,\nu\), but need not be
optimal when the colors interleave. In contrast, for \(d\ge2\),
every optimal plan satisfies
\(\operatorname{KL}(\pi\mid\mu\otimes\nu)=+\infty\).
\end{proposition}

\begin{proof}
Fix an optimal kernel \(\kappa_z\). Partition the two conditional
supports into their alternating monochromatic order runs, as in
Proposition~\ref{prop:ray-finite-bands}. Write \(A_i\) and \(B_j\)
for the source and target runs and put
\[
a_i=\alpha_z(A_i),\qquad b_j=\beta_z(B_j),\qquad
m_{ij}=\kappa_z(A_i\times B_j).
\]
Omitting zero-mass runs, define
\[
\bar\kappa_z=
\sum_{i,j}m_{ij}
  \frac{\alpha_z|_{A_i}}{a_i}\otimes
  \frac{\beta_z|_{B_j}}{b_j}.
\]
The row and column sums of \(m\) are \(a\) and \(b\).
Each positive entry joins a source run wholly preceding its target
run. Thus the new kernel has the prescribed marginals and is forward.
Moreover,
\[
\operatorname{KL}(\bar\kappa_z\mid\alpha_z\otimes\beta_z)
=\sum_{i,j}m_{ij}\log\frac{m_{ij}}{a_i b_j}
\le \min\{H(a),H(b)\}\le\log K,
\]
where \(H\) denotes finite Shannon entropy and \(0\log0=0\).
The middle quantity is the mutual information of the two run labels.
The measurable grid construction in
Proposition~\ref{prop:ray-finite-bands} makes this a measurable
probability kernel. Reconstruction gives an optimal \(\bar\pi\).
Since the ray label is recovered from either endpoint, disintegration
of relative entropy gives
\[
\operatorname{KL}(\bar\pi\mid\pi_0)
=\int\operatorname{KL}(\bar\kappa_z\mid
                 \alpha_z\otimes\beta_z)\,\Lambda(dz)
\le\log K.
\]
For the final assertion, a fixed ray is a countable union of smooth
arcs and hence has zero \(d\)-dimensional volume for \(d\ge2\).
Thus \(\Lambda\) is atomless. The same-label relation has zero
\(\mu\otimes\nu\)-measure, since its label pushforward is the
diagonal in \(\Lambda\otimes\Lambda\), whereas every optimal plan
gives that relation full mass.
\end{proof}

This proposition supplies a finite-entropy competitor relative to a
derived same-ray reference. It does not identify either the
zero-temperature entropic limit or the iid empirical selector.
The former additionally depends on transverse asymptotics; the latter
requires the empirical matching and recovery arguments. The two
reference measures in the proposition must not be interchanged.

These are structural consequences of transport-ray disintegration.
The one-dimensional origin of nonuniqueness for smooth uniformly
convex norm costs is classical; see
\cite[pp.~81--83]{FeldmanMcCannDensity}.
The finite block-product construction has a Euclidean entropy
counterpart in \cite[Theorem~5.2(a), equations~(5.4)--(5.6)]{NutzZhong}.
Here the qualitative Finsler ray theorem supplies the measurable
conditional laws, while compact support separation supplies the
uniform finite-run bound.

\begin{proposition}[Localization to measurable whole-ray labels]
\label{prop:general-ray-mass-localization}
In the setting of Proposition~\ref{prop:general-regular-pairs}, let
\(N=T_\phi^\circ\) be the common full-marginal-measure interior-ray
set. Let \(\ell:N\to Z\) be a fixed Borel map into a Polish space
with bounded compatible metric \(d_Z\le1\), constant on each whole
calibrated ray. Then
\(\Lambda:=\ell_\#\mu=\ell_\#\nu\), and, almost surely,
\begin{gather}
\sup_{\pi\in\mathcal O_n^F}
\pi\{d_Z(\ell(x),\ell(y))\ge\delta\}\longrightarrow0
\quad(\delta>0),
\label{eq:general-label-localization}\\
\sup_{\pi\in\mathcal O_n^F}
d_{\rm BL}\bigl((\ell,\ell)_\#\pi,
                  (\operatorname{diag})_\#\Lambda\bigr)
\longrightarrow0,
\label{eq:general-label-diagonal}
\end{gather}
where \((\ell,\ell)(x,y)=(\ell(x),\ell(y))\) and
\(\operatorname{diag}(z)=(z,z)\).
For every fixed countable partition into whole-ray clusters
\(A_j,B_j\), this implies
\[
\sup_{\pi\in\mathcal O_n^F}\sum_{j\ne k}\pi(A_j\times B_k)
\longrightarrow0.
\]
For each fixed cluster, the original restricted plans
\(\pi|_{A_j\times B_j}\) approach the population optimal set for
the balanced submarginals \(\mu|_{A_j},\nu|_{B_j}\).
\end{proposition}

\begin{proof}
Every population optimum couples points of \(N\) on the same ray,
so its label marginals agree. Extend \(\ell\) arbitrarily outside
\(N\); almost surely no sample point uses that extension.
Fix \(\epsilon,\delta>0\). Inner regularity and Lusin's theorem
give compact \(K_A\subset A\cap N\), \(K_B\subset B\cap N\),
losing at most \(\epsilon\) marginal mass each, on which \(\ell\)
is continuous. The compact set
\[
D_{\epsilon,\delta}
=\{(x,y)\in K_A\times K_B:
                   d_Z(\ell(x),\ell(y))\ge\delta\}
\]
is disjoint from \(\Gamma_\phi\). Thus the calibration slack
\(s\) has a positive minimum \(\kappa_{\epsilon,\delta}\)
there; choose any positive constant if the set is empty. Since
\(\int s\,d\pi=a_n\), every empirical optimum satisfies
\begin{equation}
\label{eq:general-label-slack}
\pi\{d_Z(\ell(x),\ell(y))\ge\delta\}
\le\mu_n(A\setminus K_A)+\nu_n(B\setminus K_B)
   +\frac{a_n}{\kappa_{\epsilon,\delta}}.
\end{equation}
The fixed-complement empirical laws and \(a_n\to0\), followed by
countable sequences \(\epsilon,\delta\downarrow0\), prove
\eqref{eq:general-label-localization}. Put \(\Lambda_n=\ell_\#\mu_n\).
The empirical law on the Polish label space converges weakly almost
surely. With the sum metric on \(Z\times Z\),
\[
d_{\rm BL}\bigl((\ell,\ell)_\#\pi,
                   (\operatorname{diag})_\#\Lambda\bigr)
\le \delta+\pi\{d_Z(\ell(x),\ell(y))\ge\delta\}
       +2d_{\rm BL}(\Lambda_n,\Lambda),
\]
which proves \eqref{eq:general-label-diagonal}.
For triangular iid samples these empirical convergences follow from
bounded-test Hoeffding bounds and Borel--Cantelli, without requiring
independence between sample sizes.

For the cluster conclusion use the countable discrete label space
with distance one between distinct indices. Write
\(\pi_n^j=\pi|_{A_j\times B_j}\). Its marginal deficits are
positive measures with exact total variation masses
\[
\|\mu_n|_{A_j}-(\operatorname{pr}_1)_\#\pi_n^j\|_{\rm TV}
  =\pi(A_j\times B_j^c),\qquad
\|\nu_n|_{B_j}-(\operatorname{pr}_2)_\#\pi_n^j\|_{\rm TV}
  =\pi(A_j^c\times B_j).
\]
The sum of either family is the cross-cluster mass, which tends
to zero. The empirical restrictions converge weakly to their fixed
submarginals by applying the empirical law to tests
\(f\mathbf1_{A_j}\) and \(g\mathbf1_{B_j}\). Any weak limit of
\(\pi_n^j\) therefore has those marginals and zero slack, because
\(0\le\int s\,d\pi_n^j\le a_n\). It is optimal for the
submarginals. Compactness makes the convergence to that optimal set
uniform over the original empirical optimizers.
\end{proof}

The rough atlas of Section~\ref{sec:rough-rays} permits the label
\(\ell(x)=(j,z_j(x))\), using distance one between chart indices
and \(\min(1,|z-w|)\) within a chart. Thus the result also
localizes the paired transverse coordinates without a smooth global
foliation. It supplies no rate at the microscopic scale. For example,
with uniform independent samples on
\([0,1]\times[0,1]^q\) and \([2,3]\times[0,1]^q\), partition the
transverse cube into \(n=k^q\) equal cubes. The fraction of target-empty
cubes tends in \(L^2\) to \(e^{-1}\): its expectation is
\((1-1/n)^n\), and distinct emptiness indicators have nonpositive
covariance. Conditional on the targets, the fraction of source points
in those cubes has that empty-cube fraction as its mean and variance
at most \(1/(4n)\); it therefore also tends to \(e^{-1}\).
Every matching must send all those
sources across cube boundaries. Consequently fixed-chart localization
does not imply negligible crossings for shrinking microscopic cells
or justify independently reoptimizing such cells.


\subsection{Intrinsic active ray components}
\label{subsec:intrinsic-active-components}

The maximal ray quotient may depend on the optimal potential. Its
positive-flux components have a stronger invariance. Retain the
assumptions of Proposition~\ref{prop:general-optimal-face}. For a fixed
optimal potential \(\phi\), write its ray data as
\(\gamma_z^\phi,\Lambda_\phi,\alpha_z^\phi,\beta_z^\phi\), and set
\[
 H_z^\phi(r)=F_{\alpha_z^\phi}(r)-F_{\beta_z^\phi}(r).
\]
Choose a Borel full set of whole-ray labels on which these data satisfy
Propositions~\ref{prop:general-optimal-face} and
\ref{prop:ray-forward-criterion}. A component \(I\) of
\(\{H_z^\phi>0\}\) defines the \emph{active block}
\(C_\phi(z,I)=\gamma_z^\phi(I)\). This is an oriented open geodesic
arc of length at most \(D\), not necessarily an ambient open set.
Its two conditional marginal masses agree and are positive.

Here is an explicit Borel block space. Enumerate the rationals as
\((r_j)_{j\ge1}\). For \(H_z^\phi(r)>0\), let \(j_\phi(z,r)\)
be the first index whose rational belongs to the same component as
\(r\). The test for an index is
\begin{equation}
\label{eq:intrinsic-active-marker}
 \inf_{u\in\mathbb Q\cap[0,1]}
 H_z^\phi((1-u)r+ur_j)>0.
\end{equation}
Continuity identifies this infimum with the minimum on the compact
segment, so the test and the least-index function are Borel. Put
\[
 \mathcal B_\phi
 =\{(z,j):H_z^\phi(r_j)>0,\ j_\phi(z,r_j)=j\},\qquad
 b_\phi(\gamma_z^\phi(r))=(z,j_\phi(z,r)).
\]
The injective Borel ray coordinates make the active physical set
\(E_\phi\) and its block map \(b_\phi\) Borel.
Proposition~\ref{prop:ray-forward-criterion} gives
\begin{equation}
\label{eq:intrinsic-active-quotient}
 \mu(E_\phi)=\nu(E_\phi)=1,\qquad
 \eta_\phi=(b_\phi)_\#\mu=(b_\phi)_\#\nu.
\end{equation}
Writing \(w_\phi(z,j)\) for the conditional block mass, one has
\(\eta_\phi(dz,\{j\})=w_\phi(z,j)\Lambda_\phi(dz)\).
Let \(\mu_c^\phi,\nu_c^\phi\) be the conditional physical
probability laws over this block quotient.

\begin{theorem}[Potential independence of active blocks]
\label{thm:intrinsic-active-invariance}
For any two fixed optimal potentials \(\phi,\psi\), there are Borel
conull block-label sets \(\mathcal B_\phi^0,\mathcal B_\psi^0\)
and a Borel bijection \(T:\mathcal B_\phi^0\to\mathcal B_\psi^0\)
such that
\[
 C_\phi(c)=C_\psi(Tc),\qquad T_\#\eta_\phi=\eta_\psi.
\]
The first equality holds for every retained label as equality of whole
oriented open arcs. For \(\eta_\phi\)-almost every \(c\),
\[
 \mu_c^\phi=\mu_{Tc}^\psi,\qquad
 \nu_c^\phi=\nu_{Tc}^\psi.
\]
On every identified arc, \(\psi-\phi\) is constant. The omitted
physical sets are null for \(\mu,\nu\) and the occupation measure
\(\tau\) of Corollary~\ref{cor:general-occupation-current}.
\end{theorem}
\begin{proof}
First, let \(\kappa\) be any forward coupling of atomless scalar
laws with continuous surplus \(H\), and let \(G\) be a single
full-\(\kappa\)-measure set whose pairs have endpoints in the same
positive component. Then
\begin{equation}
\label{eq:intrinsic-active-cover}
             \{H>0\}=\bigcup_{(a,b)\in G}(a,b).
\end{equation}
Indeed, each interval stays in its endpoint component. Conversely,
\eqref{eq:ray-cut-flux} and atomlessness give
\(\kappa\{a<r<b\}=H(r)\) for every \(r\). A positive crossing
event meets the same fixed full set \(G\). This proves the identity
pointwise without intersecting uncountably many full-measure sets.
If intersecting covering intervals have the same label, their unions
by label form a relatively open partition of each component.
Connectedness forces one label on that component.

Fix one optimal plan \(\pi\). Take a Borel full-\(\pi\)-measure
set \(R\) on which both potentials calibrate the unique minimizing
segment, both endpoints belong to \(E_\phi\cap E_\psi\), and
the endpoints have equal block labels for each potential.
Such a set exists by Propositions~\ref{prop:general-regular-pairs}
and~\ref{prop:ray-forward-criterion}. Its entire closed transported
segment lies in each of its two active blocks: the endpoints are
inside each positive coordinate interval, and the physical minimizing
segment is unique.

Disintegrate \(\pi\) over each whole-ray quotient. Let
\(Z_\phi^0,Z_\psi^0\) be Borel full sets on which the conditional
pair kernels give \(R\) mass one. For every ray in \(Z_\phi^0\),
\eqref{eq:intrinsic-active-cover} covers each phi block by transported
intervals lying in psi blocks. Overlapping intervals lie in the same
psi block, since distinct physical psi blocks are disjoint.
The connected-cover observation gives containment of each entire
phi block in one psi block. Reversing the roles gives the opposite
containment for every ray in \(Z_\psi^0\).

To ensure that this reverse argument applies to the containing block,
set
\[
 R_1=R\cap\{\ell_\phi(x)\in Z_\phi^0,
                  \ell_\psi(x)\in Z_\psi^0\}.
\]
It is still full for \(\pi\). Let \(Z_\phi^1\subset Z_\phi^0\)
be Borel and full, with conditional probability one for \(R_1\).
The cover by \(R_1\) gives, for every block \(C\) over this set,
\(C\subset D\), where the containing psi block lies on a ray in
\(Z_\psi^0\). The reverse cover by \(R\) gives
\(D\subset C'\) for a phi block \(C'\). Disjointness of distinct
phi blocks and \(C\subset C'\) imply \(C'=C\), hence \(C=D\).

Let \(\mathcal B_\phi^0\) consist of the blocks over
\(Z_\phi^1\), and define
\[
            T(z,j)=b_\psi(\gamma_z^\phi(r_j)).
\]
Exact arc equality makes this expression defined for every retained
label. It is Borel and injective. The injective-image theorem for
standard Borel spaces makes its image \(\mathcal B_\psi^0\)
Borel and its inverse Borel. Moreover,
\(b_\psi=T\circ b_\phi\) almost everywhere for both marginals.
This proves the quotient pushforward and that the image is conull;
uniqueness of disintegration gives the conditional-law identities.
Both potentials increase by arclength on each covering interval.
Overlaps and connectedness make their difference constant on the
whole block and identify its orientation.

Finally, an omitted quotient-null family loses no endpoint mass.
Disintegrating \(\pi\) over its blocks, every transported segment
stays inside its block and has length at most \(D\). The occupation
of an omitted family is therefore at most \(D\) times its quotient
mass. This proves the remaining null-set assertion.
\end{proof}

Thus the measured partition into active arcs, with its oriented
physical endpoint laws, is intrinsic to \((F,\mu,\nu)\).
Its normalized surplus is \(H_z^\phi/w_\phi(z,j)\) on a block,
since the surplus vanishes at the component endpoints; it is canonical
after the additive coordinate shift. The whole-ray conditional surplus
and its separate normalization factors need not be canonical.
The theorem compares any two fixed potentials, without asserting one
exceptional set for an uncountable family. It identifies neither
zero-flux continuations nor ambient connected components of the
occupation support, and it selects no endpoint coupling.


\section{Local geometry of calibrated Finsler ray tubes}
\label{sec:raytube}

We record a local second-variation formula used in the measurable atlas. Let
\((M,F)\) be a complete smooth reversible strongly convex Finsler
manifold of dimension \(q+1\), and let
\(X:J\times B\to M\) be a smooth coordinate diffeomorphism,
where \(B\subset\R^q\) is open. Assume a globally 1-Lipschitz
function \(\varphi:M\to\R\) satisfies
\[
\varphi(X(r,z))=r,\qquad \widetilde F(r,z;1,0)=1,
\]
where \(\widetilde F\) is the pulled-back Finsler norm. This
\emph{calibrated tube} assumption is a local geometric hypothesis;
we do not assume that arbitrary marginals already supply such tubes.
Define the positive definite transverse Hessian and its integrated
inverse by
\begin{equation}
H(r,z)=D^2_{vv}\widetilde F(r,z;1,0),\qquad
S(a,b,z)=\int_a^b H(r,z)^{-1}\,dr.
\label{eq:raytube-compliance}
\end{equation}
Here \(v\in\R^q\) denotes transverse velocity. The matrix \(S\)
records the transverse compliance of the whole ray segment.

\begin{lemma}[The Hessian of the actual length distance]
\label{lem:raytube-distance}
Uniformly on compact sets of \((a,b,z)\) whose segments lie in
the tube and satisfy \(b-a\ge\ell>0\), the global distance has
a smooth branch near the corresponding endpoint pairs and
\[
d_F(X(a,z),X(b,z+\delta))
=b-a+\tfrac12\delta^T S(a,b,z)^{-1}\delta+O(|\delta|^3).
\]
Moving the two endpoints by \(\alpha,\beta\) gives quadratic term
\(\tfrac12(\beta-\alpha)^TS^{-1}(\beta-\alpha)\) and remainder
\(O((|\alpha|+|\beta|)^3)\).
\end{lemma}

\begin{proof}
Calibration gives \(\widetilde F(r,z;v_r,v_z)\ge v_r\).
The radial segment has length \(b-a\), so it attains the global
lower bound. For any other unit-speed global minimizer \(c\),
the Lipschitz inequalities from its two endpoints give
\(\varphi(c(t))=a+t\). Near its initial point \(\varphi\) is
smooth. Strict convexity therefore forces the unique calibrated
initial vector \((1,0)\), and uniqueness of the geodesic equation
identifies the entire curve. This uses no smoothness of
\(\varphi\) outside the tube.

For completeness, the energy \(E=\widetilde F^2/2\) satisfies,
along the radial velocity, \(E_{xx}=E_{xv}=0\) and
\(E_{vv}=\operatorname{diag}(1,H)\): both
\(E(x;(1,0))=1/2\) and \(E_v(x;(1,0))=(1,0)\) are
independent of position. Its endpoint-fixed second variation is
\[
\int_a^b\bigl((s')^2+(v')^THv'\bigr)\,dr.
\]
It is positive for every nonzero variation. There is no
endpoint-fixed Jacobi field and the exponential endpoint map is
nonsingular; see also \cite[Proposition~2.16]{BhowmickPrasad}.

Completeness gives compactness of bounded-length minimizing
segments. If perturbed endpoint minimizers failed to lie on the
smooth local branch, a subsequence of their initial unit vectors
and lengths would converge to a minimizing baseline geodesic.
The uniqueness just proved identifies its initial vector; the
endpoint inverse function theorem then puts all nearby minimizers
on the same branch, a contradiction. Compactness of the parameter
set makes the neighborhood uniform. In particular these actual
global minimizers stay in the tube and have positive radial
velocity.

After using \(r\) as parameter, the transverse second variation
is \(\int_a^b(v')^THv'\,dr\). Among paths with
\(v(a)=0,v(b)=\delta\), completing the square gives the unique
minimizer
\(v'(r)=H(r,z)^{-1}S(a,b,z)^{-1}\delta\), of energy
\(\delta^TS^{-1}\delta\). Smooth dependence of the minimizing
branch and Taylor's theorem prove the expansion. For two moving
endpoints, replace the base point \(z\) by \(z+\alpha\);
the resulting coefficient change contributes only a cubic error.
\end{proof}

\section{Almost-everywhere ray geometry and local sampling}
\label{sec:rough-rays}

The smooth-tube assumption can be weakened for the local quadratic
calculation and the input point-process limit. Throughout this section
\(d=q+1\ge2\), \((M,F)\) is complete, smooth, reversible and strongly
convex, and \(u\) is a globally 1-Lipschitz potential. The ray-cluster
results of \cite[Theorem~4.11, Proposition~4.18 and Lemma~4.21]{Assouline}
apply to the auxiliary Lagrangian \(L(v)=(F(v)^2+1)/2\).
This Lagrangian is \(C^1\) on \(TM\) and smooth off the zero
section, as allowed by \cite[Definition~2.1 and Example~2.6]{Assouline}.
Its free-time action is exactly the original distance: a curve of
length \(\ell\), traversed in time \(T\), satisfies
\[
\int_0^T L(\dot\gamma)\,dt
\ge \frac{\ell^2}{2T}+\frac T2\ge\ell,
\]
with equality on a unit-speed minimizing geodesic when the endpoints
are distinct. Constant curves with \(T\downarrow0\) cover coincident
endpoints. Completeness gives minimizing geodesics and compact
containment for compact endpoint sets, verifying the remaining
geometric hypotheses of that reference. No curvature or topological
restriction is used here.

\subsection{The quadratic defect on a typical ray}

By \cite[Lemma~4.23]{Assouline}, the moving ray set is covered, up to
volume zero, by countably many disjoint ray clusters. Each has an
injective locally Lipschitz parametrization on a Borel domain with
compact transverse base. Reparametrize longitudinally by \(u\), and
write it as \(X(r,z)\), so that \(u(X(r,z))=r\). The endpoints of the
longitudinal interval depend continuously on the allowed label \(z\).
At a full-measure set of transverse density points \(z_0\), the
differential of \(X\) exists and is nonsingular along the whole
interior ray. On each trimmed segment, a \(C^{1,1}\) extension \(u_k\)
agrees with both \(u\) and \(du\); its differential is differentiable
at every point of this ray, with symmetric derivative propagated by
the Hamiltonian flow. These are the conclusions of the cited
ray-cluster package, rather than an assumed smooth open foliation.

\begin{proposition}[Quadratic defect at almost every ray]
\label{prop:rough-ray-quadratic}
Fix such a label \(z_0\), put \(\gamma(r)=X(r,z_0)\), and define
\[
J(r)=D_zX(r,z_0),\qquad
H(r)=J(r)^\top D^2_{vv}F(\gamma(r),\dot\gamma(r))J(r),\qquad
S(a,b)=\int_a^b H(r)^{-1}\,dr.
\]
The columns of \(J\) span \(\ker du\); thus \(H\) and \(S(a,b)\)
are positive definite. For every fixed interior pair \(a<b\), as
\(\alpha,\beta\to0\) through allowed transverse labels,
\begin{align}
&d_F\bigl(X(a,z_0+\alpha),X(b,z_0+\beta)\bigr)-(b-a)
\notag\\
&\hspace{12mm}=
\frac12(\beta-\alpha)^\top S(a,b)^{-1}(\beta-\alpha)
+o(|\alpha|^2+|\beta|^2).
\label{eq:rough-ray-quadratic}
\end{align}
No common interface separating the source and target populations is
required. The remainder is pointwise in \(z_0,a,b\).
\end{proposition}

\begin{proof}
The change from the original unit-speed time coordinate to \(r=u\)
is the shear \(r=t+w(z)\), where \(w\) is the restriction of \(u_k\)
to the smooth seed. Differentiating it shows that
\([\dot\gamma,J]\) is nonsingular and \(du\,J=0\).
Choose a level \(c\) inside a slightly enlarged compact ray segment,
and one trim index \(k\) covering this segment and nearby allowed
rays. Such a common trim follows from continuity of the ray endpoints.

We construct a smooth auxiliary family matching the original first
variation. In ambient coordinates at \(x_0=\gamma(c)\), let
\(p_0=du_k(x_0)\) and \(B_0=D(du_k)(x_0)\). Take a smooth function
\(f\) with the same value, differential and Hessian there, and set
\(\Sigma=\{f=c\}\). For the kinetic Hamiltonian
\(K(x,p)=F_x^*(p)^2/2\),
\begin{equation}
D_xK(x_0,p_0)[v]+D_pK(x_0,p_0)[B_0v]=0
\quad\text{for every }v.
\label{eq:rough-eikonal-jet}
\end{equation}
Indeed, \(K(x,du_k(x))\) is differentiable at \(x_0\) and equals
\(1/2\) on a strain set having density one there. A differentiable
function constant on a density-one set has zero differential.
Consequently the normalized conormal
\(n=df/F^*(df)\) on \(\Sigma\) has, on tangent vectors at \(x_0\),
the same bundle derivative as \(du_k\).

Parametrize \(\Sigma\) smoothly by \(\sigma(z)\), with
\(\sigma(z_0)=x_0\) and \(D\sigma(z_0)=J(c)\). If \(\Phi_s\) is the Hamiltonian flow of
\(K\), define
\[
Y(r,z)=\pi\Phi_{r-c}\bigl(\sigma(z),n(\sigma(z))\bigr).
\]
The initial cotangent data and their first transverse derivatives
agree with those of the original ray family. Smooth dependence of
the flow gives
\(Y(r,z_0)=\gamma(r)\) and \(D_zY(r,z_0)=J(r)\) along the whole
compact segment. In particular \(J\) is smooth in \(r\).
The differential \([\dot\gamma,J]\) is nonsingular there, so a
common small transverse neighborhood makes \(Y\) a local
diffeomorphism. It can be made injective on the enlarged compact
segment: otherwise colliding parameter pairs at widths tending to
zero have limits on the baseline; injectivity of \(\gamma\), followed
by local injectivity at the common limit, is a contradiction.

On this smooth tube define \(\psi(Y(r,z))=r\). The canonical
one-form on the unit conormal flow evaluates to zero on the
transverse directions and to one on the longitudinal direction;
hence \(d\psi\) is the transported covector and \(F^*(d\psi)=1\).
This is a local eikonal solution; global Lipschitz extension of
\(\psi\) is not needed.

The baseline segment is the unique global minimizer: calibration by
the original \(u\) forces any minimizing initial unit velocity to
satisfy \(du(v)=1\), and strong convexity determines that velocity
uniquely. In the \(Y\) coordinates the energy \(E=F^2/2\) satisfies,
at radial velocity, \(E_{xx}=E_{xv}=0\) and
\(E_{vv}=\operatorname{diag}(1,H)\). Its endpoint-fixed index form is
\[
\int_a^b\bigl((s')^2+(v')^\top Hv'\bigr)\,dr,
\]
which is strictly positive for nonzero endpoint-fixed variations.
Thus the endpoint exponential map is nonsingular. Compactness of
bounded-length minimizing geodesics and uniqueness of the limiting
segment show that all sufficiently close global minimizers belong
to this smooth branch and stay in the auxiliary tube, as in the
proof of Lemma~\ref{lem:raytube-distance}.

The same second-variation calculation as in
Lemma~\ref{lem:raytube-distance} gives
\(\tfrac12(\beta-\alpha)^\top S^{-1}(\beta-\alpha)\) for the
quadratic part of
\(e_\psi(x,y)=d_F(x,y)-\psi(y)+\psi(x)\): completing the square
in \(\tfrac12\int_a^b v'^\top Hv'\) with endpoint values
\(\alpha,\beta\) proves the formula without requiring commuting
matrices.

To transfer this expansion, let \(\Gamma_t(\zeta)\) be the geodesic
with initial phase \(\zeta=(x,v)\), put \(L(\zeta)=F_x(v)\), and
write \(t_a=a-c\), \(t_b=b-c\). On nearby initial phases consider
\[
 A(\zeta_1,\zeta_2)
 =d_F(\Gamma_{t_a}(\zeta_1),\Gamma_{t_b}(\zeta_2))
  -f(x_2)+f(x_1)-t_bL(\zeta_2)+t_aL(\zeta_1).
\]
This is \(C^2\) near the baseline pair and stationary there. Writing
\(p_t\) for the central length covector and \(\pi(x,v)=x\),
first variation gives \(p_tD\Gamma_t=p_0D\pi+t\,dL\), also for
negative \(t\); \(df(x_0)=p_0\) cancels the remaining terms.
The original level-\(c\) phases \(\zeta_X(z)\) and the smooth
conormal phases \(\zeta_Y(z)\) have the same first jet, as established
above. Stationary Taylor expansion therefore changes \(A\) by
\(o(|\alpha|^2+|\beta|^2)\) when one pair of phases replaces the other.

Both seed families have unit speed. On the original allowed seeds,
\(u(X(c,z))=u_k(X(c,z))=c\), so calibration gives
\(u(X(r,z))=u_k(X(c,z))+(r-c)=r\). Set \(R=f-u_k\) and
\(e_u(x,y)=d_F(x,y)-u(y)+u(x)\). Exactly,
\[
\begin{split}
 A(\zeta_X(z_1),\zeta_X(z_2))
 &=e_u(X(a,z_1),X(b,z_2))
   +R(X(c,z_1))-R(X(c,z_2)),\\
 A(\zeta_Y(z_1),\zeta_Y(z_2))
 &=e_\psi(Y(a,z_1),Y(b,z_2)).
\end{split}
\]
The second identity uses \(f(\sigma(z))=c\). Initial jet matching
makes \(R(x)=o(|x-x_0|^2)\), while
\(X(c,z)-x_0=O(|z-z_0|)\). Thus the correction is also
\(o(|\alpha|^2+|\beta|^2)\), proving
\eqref{eq:rough-ray-quadratic}.
\end{proof}

\subsection{Marked sampling at a fixed typical root}

Let the marginal densities relative to a smooth volume \(m\) be
\(\rho_-,\rho_+\), with compact support, and suppose now that \(u\)
is a Kantorovich potential for these marginals. In a whole-ray
cluster with coordinate domain \(D\), write its volume Jacobian as
\(\mathcal J(z,r)\), and fix the zero-extended densities
\[
f_\pm(z,r)=\mathbf1_D(z,r)\rho_\pm(X(r,z))\mathcal J(z,r).
\]
Take a bounded interval \(I\) containing \(u\) on both marginal
supports. Change of variables makes \(f_\pm\) integrable on
\(\R^q\times I\), after extension by zero outside the transverse
base. Each cloud consists of \(n\) iid samples from its marginal;
the source and target families are independent, and only samples
inside the chart are retained. Whole-ray balance, \cite[Theorem~5.9]{Assouline}, gives
\[
\lambda(z):=\int_I f_-(z,r)\,dr=\int_I f_+(z,r)\,dr
\quad\text{for almost every }z.
\]

\begin{proposition}[Fixed-root Poisson input law]
\label{prop:rough-ray-poisson}
At almost every transverse root \(z_0\), the chart-restricted
source and target clouds, with transverse coordinate
\(n^{1/q}(z-z_0)\) and unchanged longitudinal mark \(r\), converge
jointly to independent Poisson random measures with intensities
\[
dx\,f_-(z_0,r)\,dr,\qquad dx\,f_+(z_0,r)\,dr.
\]
In each bounded transverse window the convergence of the finite
coordinate configurations is in total variation. At roots with
\(0<\lambda(z_0)<\infty\), the two spatial intensities agree and
their mark laws have densities \(f_\pm(z_0,r)/\lambda(z_0)\).
No transverse continuity of these densities is required.
\end{proposition}

\begin{proof}
Lebesgue differentiation of the Bochner-integrable maps
\(z\mapsto f_\pm(z,\cdot)\in L^1(I)\) gives, for almost every
\(z_0\) and every finite \(R\),
\begin{equation}
\int_{B_R}\int_I
|f_\pm(z_0+\varepsilon x,r)-f_\pm(z_0,r)|\,dr\,dx\longrightarrow0.
\label{eq:ray-density-differentiation}
\end{equation}
For completeness, approximate the maps in \(L^1_z(L^1_r)\) by
continuous compactly supported maps and apply the scalar maximal
inequality to the \(L^1_r\) norm of the error. The usual
differentiation argument gives the assertion, simultaneously for
integer \(R\) and both signs.

Put \(\varepsilon=n^{-1/q}\). For a nonnegative bounded test \(h\)
supported in a bounded transverse window, independence within one
sample family gives the Laplace functional
\[
\left(1-\frac{A_n}{n}\right)^n,
\qquad A_n=\int(1-e^{-h(x,r)})f(z_0+\varepsilon x,r)\,dx\,dr.
\]
Equation~\eqref{eq:ray-density-differentiation} sends this to the
claimed Poisson functional. Independence of the sample families
gives the joint assertion. More strongly, in a bounded Borel window
\(W\), the count is binomial with parameter
\(n^{-1}\Lambda_n\), where
\(\Lambda_n=\int_{W\times I}f(z_0+\varepsilon x,r)\,dx\,dr\)
converges to \(\Lambda=|W|\int_I f(z_0,r)\,dr\).
Conditional on the count, enumeration in original sample-index
order gives iid coordinates with density
\(f(z_0+\varepsilon x,r)/\Lambda_n\).
For \(\Lambda>0\), these densities converge in \(L^1\), and the
counts converge in total variation to a Poisson count. This proves
total-variation convergence of the finite configurations.
For \(\Lambda=0\), the count vanishes in probability.
\end{proof}

\begin{corollary}[Finite local cost arrays, including interleaving]
\label{cor:rough-ray-cost-array}
Choose a root satisfying both preceding propositions and a bounded
transverse window. Decorate each source--target pair in its finite
clouds by the scaled calibration defect
\[
n^{2/q}\bigl(d_F(X_i,Y_j)-u(Y_j)+u(X_i)\bigr).
\]
The finite configurations and all their decorated costs converge
in distribution to the marked Poisson clouds with extended cost
\begin{equation}
K((x,a),(y,b))=
\begin{cases}
\frac12(y-x)^\top S(a,b)^{-1}(y-x),&a<b,\\
+\infty,&a>b.
\end{cases}
\label{eq:rough-ray-extended-cost}
\end{equation}
The event \(a=b\) is null. Costs take values in the compact space
\([0,\infty]\). No separating section for the mark populations
is assumed.
\end{corollary}

\begin{proof}
Condition first on finite source and target counts. By the
total-variation statement, their coordinate parameters can be
replaced by limiting independent parameters with error tending
to zero, while retaining the \(n\)-dependent cost function.
Chart-domain failures under these limiting parameters also vanish:
if \(q_n(x,r)=f(z_0+n^{-1/q}x,r)\), \(q_0(x,r)=f(z_0,r)\), and
\(D_n=\{(x,r):(z_0+n^{-1/q}x,r)\in D\}\), then
\[
\int_{D_n^c}q_0\le\|q_n-q_0\|_{L^1(W\times I)}\longrightarrow0.
\]
The limiting marks are almost surely interior and distinct across
the two families. For fixed admissible parameters with \(a<b\),
Proposition~\ref{prop:rough-ray-quadratic} gives the cost limit
along every admissible sequence. For \(a>b\), apply it with the
endpoints reversed: reversibility gives the unscaled defect
\(2(a-b)+O(n^{-2/q})\), so the scaled cost diverges.
Compactify costs by \(t\mapsto t/(1+t)\) and set their comparison
error to zero off admissible domains. Pointwise convergence and
bounded convergence prove joint convergence of the finitely many
costs. Tightness of the counts removes the conditioning. Thus no
uniform Taylor remainder in varying longitudinal marks is needed.
\end{proof}

These results concern a fixed typical root, or an independent root
whose law charges only the good-root set (in particular one with
an absolutely continuous transverse law). In the latter case the
limiting clouds are independent conditional on the root; their
mixture need not be independent. These results identify the local
input and finite pair costs in every dimension \(d\ge2\).
They do not justify restricting an empirical optimizer to a cluster,
discarding exterior assignments, or replacing its sample-dependent
tilt by a deterministic root. Matched-edge tightness, compensated
sampling estimates and uniqueness of the limiting paired-mark law
remain necessary for the full selection theorem. The next section
proves the optimizer-limit and reconstruction steps under additional
regular separated-tube hypotheses in \(d\ge4\).

\subsection{A finite upper atlas and global energy comparison}
\label{subsec:general-upper-atlas}

Keep the general geometric and marginal assumptions above, write
\(q=d-1\), \(D=\operatorname{diam}_{d_F}(A\cup B)\), and use the
potential \(\phi\) and nonnegative slack \(s\) of
Proposition~\ref{prop:general-regular-pairs}. The geometric construction
below holds for \(d\ge2\). Its empirical energy consequence holds for
\(q>4\), hence \(d\ge6\), with arbitrary interleaving and merely
measurable conditional radial laws. No curvature sign or topology
condition is added.

\begin{proposition}[Finite whole-ray upper atlas]
\label{prop:general-upper-atlas}
Up to a common marginal null set, the transporting rays admit a finite
Borel partition into whole-ray sets \(T_1,\ldots,T_J\). For each set
there is an injective Borel ray label \(z_j\) in a bounded convex cube
\(Q_j\subset\R^q\) and a Lipschitz map
\[
 X_j:I\times Q_j\longrightarrow M,
 \qquad I=[\min_{A\cup B}\phi,\max_{A\cup B}\phi],
\]
such that every actual endpoint \(x\in(A\cup B)\cap T_j\) equals
\(X_j(\phi(x),z_j(x))\). There is a finite constant \(C\) such that,
for actual endpoints \(x=X_j(a,z)\in A\cap T_j\) and
\(y=X_j(b,w)\in B\cap T_j\),
\begin{equation}
\label{eq:general-upper-quadratic}
 0\le d_F(x,y)-|b-a|\le C|z-w|^2.
\end{equation}
The unnormalized label measures
\[
 \Lambda_j=(z_j)_\#(\mu|_{T_j})
           =(z_j)_\#(\nu|_{T_j})
\]
have bounded densities with respect to Lebesgue measure on \(Q_j\).
The auxiliary maps \(X_j\) need not be injective or calibrated outside
the selected rays. In particular this is not a global smooth foliation
assumption.
\end{proposition}

\begin{proof}
Choose \(h=\delta_0/32\) and \(k\) with \(2^{-k}<h\). Let
\(\mathcal S_k\) be the closed strain set consisting of points with a
calibrated extension of length \(2^{-k}\) on both sides.
By \cite[Lemma~4.10 and Theorem~4.11]{Assouline}, there exists a
\(C^{1,1}\) function \(\phi_k\) with
\[
 \phi_k=\phi,\qquad d\phi_k=d\phi\quad\hbox{on }\mathcal S_k.
\]
The result applies directly to \(L=(F^2+1)/2\):
\cite[Definition~2.1 and Example~2.6]{Assouline} allow \(C^1\)
regularity on the tangent bundle and \(C^2\) regularity off its zero
section. The inequality \((F^2+1)/2\ge F\), with equality at unit
speed, identifies its free-time action with \(d_F\). Completeness and
properness provide the required existence and compactness assumptions.

Use finitely many levels \(c_\ell\) from an \(h\)-spaced grid in
\(I\). Every transporting ray contains a calibrated pair from
\(A\times B\), of length at least \(\delta_0\), and therefore meets
one such level with two-sided calibrated margin at least
\(\delta_0/3-h>h\). These seed points lie in the compact closed
\(D\)-neighborhood \(K\) of \(A\). Each set
\[
 K\cap\mathcal S_k\cap\{\phi_k=c_\ell\}
\]
is compact, and \(F^*(d\phi_k)=1\) there. The \(C^{1,1}\)
implicit-function theorem and a finite covering give level-seed
parametrizations
\(\sigma_j:Q_j\to\{\phi_k=c_{\ell(j)}\}\), with closures inside
larger coordinate cubes and uniformly controlled first derivatives.
Normalize their conormals by
\[
 p_j(z)=\frac{d\phi_k(\sigma_j(z))}
              {F^*(d\phi_k(\sigma_j(z)))}.
\]
They are Lipschitz, satisfy \(p_jD\sigma_j=0\), and have unit dual
norm. On the actual strain seeds they equal \(d\phi\).

Let \(\mathcal H=(F^*)^2/2\) and let \(\Phi_t\) be its Hamiltonian
flow. Define
\begin{equation}
\label{eq:general-upper-characteristics}
 (X_j(r,z),P_j(r,z))
 =\Phi_{r-c_{\ell(j)}}(\sigma_j(z),p_j(z)).
\end{equation}
Completeness defines this flow for all \(r\in I\). The initial unit
covectors form a compact family and the times have a common finite
bound. Smooth flow dependence gives uniform Lipschitz bounds in
\((r,z)\), on a compact part of the cotangent bundle away from zero.
For a true strain seed, geodesic uniqueness identifies
\eqref{eq:general-upper-characteristics} with its calibrated ray at
all of that ray's actual endpoints in \(A\cup B\).
Each ray meets any fixed calibrated level at most once. Assign it to
the first eligible patch in the finite list. This is a Borel whole-ray
assignment: level hitting and strain membership are Borel in the
existing regular-ray quotient. It defines \(T_j\) and \(z_j\).

We prove the upper estimate even if the auxiliary flow has folds.
Write \(\vartheta\) for the canonical cotangent one-form.
Homogeneity gives
\begin{equation}
\label{eq:general-upper-contact}
 \Phi_t^*\vartheta=\vartheta+t\,d\mathcal H.
\end{equation}
Indeed \(\vartheta(X_{\mathcal H})=2\mathcal H\), so Cartan's
formula gives \(\mathcal L_{X_{\mathcal H}}\vartheta=d\mathcal H\).
For any absolutely continuous label path \(z(v)\), its initial
cotangent path \(G(v)=(\sigma_j(z(v)),p_j(z(v)))\) is absolutely
continuous and satisfies, almost everywhere,
\[
 \vartheta(G')=0,\qquad d\mathcal H(G')=0.
\]
The first equality uses the exact level-seed identity, and the second
uses \(\mathcal H(G)=1/2\). Consequently, along an absolutely
continuous time function \(t(v)\),
\begin{equation}
\label{eq:general-upper-ac-cancellation}
 P(v)\frac{d}{dv}\operatorname{pr}_M\Phi_{t(v)}(G(v))=t'(v).
\end{equation}
This is a one-dimensional chain rule along every such path; it does
not presume differentiability of a Lipschitz label map at every point
of a prescribed segment.

For \(a<b\), put \(\ell=b-a\) and use the competitor
\[
 \gamma(r)=X_j\left(r,z+\frac{r-a}{\ell}(w-z)\right),
 \qquad a\le r\le b.
\]
Its velocity is the unit characteristic velocity \(v\) plus a
perturbation \(\xi\) satisfying
\( |\xi|\le L|w-z|/\ell\) and \(P\xi=0\), by
\eqref{eq:general-upper-ac-cancellation} and the flow bounds.
Here an auxiliary smooth Riemannian norm is fixed on the compact
region. Uniformly for the unit velocities in question,
\[
 F(v+\xi)\le 1+D_vF(v)[\xi]+C_0|\xi|^2.
\]
Taylor's theorem proves this for small \(\xi\); uniform norm
comparison and homogeneity prove it for the remaining \(\xi\).
Since \(D_vF(v)=P\), integration yields
\begin{equation}
\label{eq:general-upper-path}
 d_F(X_j(a,z),X_j(b,w))
 \le |b-a|+C_1\frac{|z-w|^2}{|b-a|}\qquad(a\ne b).
\end{equation}
Reversibility gives the same argument for \(b<a\), using velocity
\(-v\) and supporting covector \(-P\).

Let \(L_1\ge1\) bound the spatial Lipschitz constants of the maps.
For actual source and target endpoints, the unit-speed path and
triangle inequality give
\[
 \delta_0\le |b-a|+L_1|z-w|.
\]
Thus \(|z-w|\le\delta_0/(2L_1)\) forces
\(|b-a|\ge\delta_0/2\), and
\eqref{eq:general-upper-path} proves the required quadratic bound.
For larger label separation use \(d_F(x,y)\le D\) and enlarge the
constant. The lower bound follows from the global Lipschitz potential.

Population optimality gives the equality of the two label measures.
For any Borel \(E\subset Q_j\), corresponding source points lie in
\(X_j(I\times E)\). The area bound for Lipschitz maps, smooth-volume
comparison on the compact image, and boundedness of \(\rho_A\) give
\[
 \Lambda_j(E)\le \|\rho_A\|_\infty m^*(X_j(I\times E))
                   \le C_j|I|\,\operatorname{Leb}_q(E).
\]
Here outer measure avoids any image-measurability issue. This proves
the bounded density assertion without a lower Jacobian bound.
\end{proof}

\begin{lemma}[Distribution-free squared matching above dimension four]
\label{lem:general-upper-dyadic}
If \(q>4\), two independent samples of size \(m\) from any probability
law supported in a cube of side \(L\) satisfy
\[
 \E\min_\sigma\sum_{i=1}^m|Z_i-W_{\sigma(i)}|^2
       \le C_qL^2m^{1-2/q}.
\]
\end{lemma}

\begin{proof}
Scale to \(L=1\). Match inside fine dyadic cells and pass each color
surplus to the parent. At level \(k\), the expected sum of absolute
color discrepancies is at most
\(\min(2m,\sqrt{2m}\,2^{kq/2})\), by the binomial variance and
Cauchy--Schwarz over the cells. Each parent repair costs at most a
constant times \(2^{-2k}\) per pair. Thus the expected total cost is
bounded by
\[
 C_q\sum_{k\ge0}2^{-2k}
             \min(m,\sqrt m\,2^{kq/2})
       \le C_qm^{1-2/q}.
\]
Split the sum at \(2^{kq}\) comparable to \(m\); both parts are
geometric because \(q>4\). A finite stopping depth adds at most
\(qm2^{-2k}\), which vanishes as that depth increases.
\end{proof}

\begin{proposition}[General interleaving energy and small population pieces]
\label{prop:general-upper-energy}
Suppose \(q=d-1>4\), and let
\[
 \mathcal E_n=\min_\sigma\sum_{i=1}^n s(X_i,Y_{\sigma(i)})
\]
for the original two independent iid samples. There are finite
constants, independent of \(n\), such that
\begin{equation}
\label{eq:general-upper-global-energy}
 \E\mathcal E_n\le Cn^{1-2/q}+C'\sqrt n.
\end{equation}
In particular every exact empirical optimum, including a fractional
optimum, has uniformly bounded expected slack at the microscopic
normalization \(n^{2/q}\int s\,d\pi\).

More precisely, fix any continuum optimum \(\pi_0\), and let
\(\eta\le\pi_0\) be a deterministic subplan of mass \(\varepsilon\).
There is an independently thinned empirical competitor, constructed
jointly for \(\eta\) and \(\pi_0-\eta\), whose normalized cost
attributed to the \(\eta\) part, including its share of count repair,
has expectation at most
\begin{equation}
\label{eq:general-upper-small-mass}
 C\varepsilon^{1-2/q}
       +C'n^{2/q-1/2}\sqrt\varepsilon.
\end{equation}
The constants are uniform over \(\eta\). The assertion concerns a
competitor constructed from deterministic population pieces, not
repair of points adaptively removed from an empirical optimizer.
\end{proposition}

\begin{proof}
Partition \(I\) into finitely many half-open bins of length at most
\(\delta_0/4\). Split \(\pi_0\) into finitely many types
\(\pi_t\), indexed by whole-ray chart, source bin and target bin.
Every positive type has its entire source bin below its target bin,
with gap at least \(\delta_0/2\): one actual calibrated pair in that
type has radial separation at least \(\delta_0\), and moving to the
two adverse bin endpoints loses at most twice the bin length.
The two label submarginals of \(\pi_t\) agree because each pair lies
on one ray. On full-measure sets of actual endpoints in this assigned
Borel chart and type, \eqref{eq:general-upper-quadratic} gives
\begin{equation}
\label{eq:general-upper-type-cost}
 s(x,y)\le C|z_j(x)-z_j(y)|^2.
\end{equation}
No topological closure of the endpoint sets is taken; such a closure
could meet a different assigned chart.

Assign each source point an auxiliary type with conditional
probabilities \(d(\pi_t)_1/d\mu\), independently over points.
Assign target types independently with probabilities
\(d(\pi_t)_2/d\nu\). This keeps the original samples unchanged.
The two augmented samples are independent, their common type masses
are \(p_t=\pi_t(M\times M)\), and the two conditional label laws
coincide within each type. Conditional on the type counts \(N_t,M_t\),
the first \(m_t=\min(N_t,M_t)\) points in original sample-index
order on each side have independent iid labels with that common law.
Writing \(\gamma_q=1-2/q\),
Lemma~\ref{lem:general-upper-dyadic} and
\eqref{eq:general-upper-type-cost} bound their expected matching
cost by
\[
 C_t\E m_t^{\gamma_q}\le C_t(np_t)^{\gamma_q}.
\]
Only finitely many fixed label cubes occur, so the constants \(C_t\)
have a common finite bound.

There remain \(\tfrac12\sum_t|N_t-M_t|\) points on each side.
Pair them arbitrarily; \(0\le s\le2D\) on \(A\times B\), and
\(\E|N_t-M_t|\le\sqrt{2np_t}\). Summing over finitely many types
proves \eqref{eq:general-upper-global-energy}. Auxiliary randomness
only constructs a feasible assignment, so its expectation bounds
the original minimum. Since the calibration endpoint sum is fixed,
all genuine-distance optima have this same total slack; Birkhoff's
theorem includes fractional optima.

For the last assertion refine the same types by membership in
\(\eta\) or \(\pi_0-\eta\), using their marginal
Radon--Nikodym densities for the independent thinning. If
\(p_t^\eta\) are the \(\eta\)-type masses and \(T\) is the fixed
number of unrefined types, then
\[
 \sum_t p_t^\eta=\varepsilon,\qquad
 \sum_t(p_t^\eta)^{\gamma_q}
       \le T^{1-\gamma_q}\varepsilon^{\gamma_q},\qquad
 \sum_t\sqrt{p_t^\eta}\le\sqrt{T\varepsilon}.
\]
Let \(\mathcal C_{n,\eta}\) be the internal cost of these types and
\(H_{n,\eta}=\sum_t|N_t^\eta-M_t^\eta|\) their total unmatched
endpoint count, counting both colors. The preceding estimates give
\[
 \E\mathcal C_{n,\eta}\le Cn^{\gamma_q}
                         \varepsilon^{\gamma_q},\qquad
 \E H_{n,\eta}\le C\sqrt{n\varepsilon}.
\]
After the complementary types are treated, the residuals have equal
global color counts. Arbitrary completion costs at most
\(D(H_{n,\eta}+H_{n,\pi_0-\eta})\). Attribute
\(D H_{n,\eta}\) to the first part and divide by
\(n^{\gamma_q}\), obtaining \eqref{eq:general-upper-small-mass}.
This attribution does not assert that the two thinned \(\eta\)
clouds have equal empirical counts or can be completed separately.
\end{proof}

The same energy bound holds for minimizers of
\(s+n^{-2/q}t\Psi\), where \(\Psi\) is bounded measurable and
\(|t|\) is bounded: comparison with the unperturbed competitor adds
at most \(2|t|\|\Psi\|_\infty\) after normalization.
In \eqref{eq:general-upper-small-mass}, first let \(n\to\infty\)
and then \(\varepsilon\downarrow0\). This gives negligible comparison
cost for deterministic discarded whole-ray sets, and also for
pair-space trimmings of a fixed population plan. It avoids any
prescribed rate for a large-measure compact-core exhaustion.
It does not retain an already chosen empirical submatching.

These results supply global upper energy and deterministic small-mass
comparison in \(d\ge6\). The finite seeds are only \(C^{1,1}\);
no ambient twice-differentiability assertion on the chosen levels is
used. Relative coefficient freezing and local lower bounds still use
the separate compact rough-ray cores. These interfaces are completed
below by Theorem~\ref{thm:general-finsler-selector} for \(d\ge6\). In \(d=4,5\), the count-repair term above is not
negligible at the proposed normalization, so this argument does not
supply the corresponding global estimate.

\subsection{Conditional laws on compact whole-ray cores}
\label{subsec:compact-conditional-laws}

Retain the general assumptions of
Proposition~\ref{prop:general-optimal-face}, with \(q=d-1\).
The regularity below holds for every \(d\ge2\). Continuity is relative
to a compact set of actual whole-ray labels; it does not imply that
the density obtained by extending that set by zero is continuous in
an ambient transverse cube. Write
\(\|\xi\|_{\mathrm{TV}}=|\xi|(\R)\) for the total variation norm
of a signed measure.

Total variation continuity alone follows directly from the measurable
conditional densities. Let \(H_C(z,r)\), \(C\in\{A,B\}\), be jointly
measurable normalized densities of the original scalar conditionals
against \(dr\), under the finite common label law \(\Lambda\).
The maps \(z\mapsto H_C(z,\cdot)\) are Borel measurable into
\(L^1(\R)\) on their normalization set: distance to any fixed
\(L^1\) function is a measurable parameter integral, and
\(L^1(\R)\) is separable. Extend the maps by zero elsewhere.
In a Polish label space, remove a measurable null set containing
all failures of normalization and of the original conditional-law
identities, then apply Lusin's theorem to the pair of maps on its
complement. This gives a common compact continuity set with
arbitrarily small loss of \(\Lambda\), retaining the original
conditionals at every point, since
\[
 \|H_C(z,\cdot)\,dr-H_C(z',\cdot)\,dr\|_{\mathrm{TV}}
   =\|H_C(z,\cdot)-H_C(z',\cdot)\|_1.
\]
For a Borel seed-label subset, apply the same argument after zero
extension to its Polish ambient space, retaining membership in the
assigned subset. The next proposition supplies the additional
geometric regularity and explicit Jacobian formulas used below.

\begin{proposition}[Whole-ray conditional continuity]
\label{prop:compact-conditional-whole}
In a bounded rough-ray cluster, let \(\Lambda\) denote the common
population quotient measure. For every \(\varepsilon>0\) there is
a compact set \(K\) of its actual assigned labels, losing less than
\(\varepsilon\) of \(\Lambda\), with the following properties.
Let \(J\) be a fixed bounded closed interval containing all potential
values of marginal-support points, with a positive buffer. There is
a \(C^1\) characteristic map \(Y\) on a neighborhood of \(J\times K\)
which agrees, in value and full first derivative, with the actual ray
map \(X\) throughout its interior domain. The maximal endpoint levels
\(\ell_-(z),\ell_+(z)\), clipped to \(J\), are continuous on \(K\).

In local coordinates with \(m=\omega(x)|dx^1\cdots dx^d|\), put
\begin{equation}
\label{eq:compact-conditional-jacobian}
\begin{aligned}
 \mathcal J_m(r,z)
   &=\omega(Y(r,z))\left|\det D_{(r,z)}Y(r,z)\right|,\\
 f_C(r,z)&=\rho_C(Y(r,z))\mathcal J_m(r,z),
 \qquad C\in\{A,B\}.
\end{aligned}
\end{equation}
The factors \(f_A,f_B\) are jointly continuous on \(J\times K\).
The entire physical coordinate marginals on the retained rays have
densities
\begin{equation}
\label{eq:compact-conditional-whole-densities}
 g_C(r,z)=f_C(r,z)
       \mathbf 1_{\{\ell_-(z)<r<\ell_+(z)\}},
 \qquad r\in J,\quad C\in\{A,B\},
\end{equation}
and \(z\mapsto g_C(\cdot,z)\) is continuous into \(L^1(J)\).
Their common quotient density
\begin{equation}
\label{eq:compact-conditional-normalization}
 \theta(z)=\int_Jg_A(r,z)\,dr=\int_Jg_B(r,z)\,dr
\end{equation}
is continuous on \(K\). On \(\{\theta>0\}\), the original
whole-ray conditional laws admit the total-variation-continuous versions
\[
 \alpha_z(dr)=\theta(z)^{-1}g_A(r,z)\,dr,
 \qquad
 \beta_z(dr)=\theta(z)^{-1}g_B(r,z)\,dr.
\]
In particular the continuity is uniform on \(K\cap\{\theta\ge c\}\)
for \(c>0\); removing \(\{\theta<c\}\) costs at most
\(c\operatorname{Leb}_q(K)\) of common quotient mass.
No common continuous latent representation of the conditionals is assumed.
\end{proposition}

\begin{proof}
Use the rough-ray coordinates of Section~\ref{sec:rough-rays}.
Their initial position, covector and potential offset form a locally
Lipschitz map of the seed label on a neighborhood of its bounded base.
The quotient is absolutely continuous with respect to label Lebesgue
measure, by the ray change-of-variables formula; compare also
Proposition~\ref{prop:general-upper-atlas}. Retain only density points
of the assigned label base at which the initial map is differentiable,
the actual full ray Jacobian is
nonsingular at every interior time, and whole-ray balance and forward
feasibility hold. These are full-measure conditions by
\cite[Proposition~4.18 and Lemma~4.21]{Assouline} and
Proposition~\ref{prop:general-optimal-face}.

Here only a \(C^1\) Whitney restriction is needed. In finitely many
seed and cotangent coordinates, Lusin and Egorov give a compact set
on which the initial covector's derivatives are continuous and its
first-order differentiability remainders are uniform. Its values and
derivatives therefore satisfy the order-one Whitney conditions.
In local bundle trivializations extend these jets to a \(C^1\)
covector field \(p(z)\) over the original smooth seed \(\sigma(z)\);
a partition of unity preserves the matching jets on overlaps.
The original seed potential offset \(f_0\) is already \(C^{1,1}\).
Normalize the extended covector to unit dual norm. This preserves
the specified first jet: on the actual seed set
\(F^*(\sigma,p)=1\), and its derivative is zero at every retained
density label. Shrinking the neighborhood keeps the extended covectors
nonzero. With \(\mathcal H=(F^*)^2/2\), define
\begin{equation}
\label{eq:compact-conditional-flow}
 Y(r,z)=\operatorname{pr}_M
          \Phi_{r-f_0(z)}(\sigma(z),p(z)).
\end{equation}
Completeness defines this unit-energy flow for the entire bounded
time range. Compact initial data and smooth flow dependence give a
\(C^1\) map on a neighborhood of \(J\times K\). Matching initial
first jets identifies both \(Y=X\) and \(DY=DX\) on the entire
actual ray interior. Off that domain \(Y\) may fold; no injectivity
or eikonal identity there is needed.

In Assouline's cluster coordinates the extended-real endpoint times
are continuous by \cite[Definition~4.16(iii) and
Proposition~4.18(i)]{Assouline}. Adding the continuous seed potential
gives the endpoint levels used here. For a different Borel whole-ray
assignment, one can instead include continuity of their clipped Borel
versions in the same Lusin restriction. All restrictions cost the
prescribed total quotient mass. Choose \(K\) inside the good assigned
label set, rather than taking the closure of an arbitrary Borel class.

The full volume density in
\eqref{eq:compact-conditional-jacobian} is intrinsic: it evaluates
\(m\) on \(\partial_rY,\partial_{z_1}Y,\ldots,\partial_{z_q}Y\).
It is not a determinant of the rectangular matrix \(D_zY\), and
there is no additional coarea factor. It is continuous and bounded
on the compact cylinder, and positive on its actual interior domain.
The injective ray change-of-variables formula
\cite[Proposition~4.18(iv), equation~(4.13)]{Assouline} gives, for
bounded Borel \(h\),
\[
 \int_{\text{retained rays}}h(x)\,d\mu_C(x)
   =\int_K\int_J h(X(r,z))g_C(r,z)\,dr\,dz,
 \qquad \mu_A=\mu,\quad\mu_B=\nu.
\]
The integrand is evaluated only in the actual ray domain.
Finite ray endpoints have zero marginal mass, and the interval \(J\)
contains every supported population coordinate.

If \(z_n\to z\) in \(K\), the continuous factors \(f_C(\cdot,z_n)\)
converge uniformly on \(J\). The domain indicators converge except
at the at most two limiting endpoint levels. Bounded convergence
proves \(\|g_C(\cdot,z_n)-g_C(\cdot,z)\|_1\to0\).
The equality in \eqref{eq:compact-conditional-normalization} is
whole-ray balance on the chosen good labels. Integration and division
by the positive continuous mass prove the remaining assertions.
The densities \(g_C\) themselves may jump at endpoint graphs;
their asserted continuity is into \(L^1\), not jointly in \((r,z)\).
\end{proof}

On a fixed closed interval \(I=[a,b]\) lying strictly inside every
retained ray, the indicator in
\eqref{eq:compact-conditional-whole-densities} is absent. Hence
\(f_A,f_B\) are the jointly continuous densities of the actual
restrictions to \(X(I\times K)\). Their two integrals over \(I\)
need not agree. The next construction restores balance while retaining
all forward coupling choices for the resulting submarginals.

\begin{lemma}[Continuous balanced trimming with a forward complement]
\label{lem:compact-conditional-trimming}
Define
\begin{align*}
 A_z(t)&=\int_a^t f_A(r,z)\,dr,&
 B_z(t)&=\int_a^t f_B(r,z)\,dr,\\
 U_z&=A_z(b),&V_z&=B_z(b),\\
 k_z&=-\min_{t\in I}\bigl(A_z(t)-B_z(t)\bigr),&
 M_z&=V_z-k_z.
\end{align*}
Then \(M_z\) is continuous and \(0\le M_z\le\min(U_z,V_z)\).
The cumulative functions
\begin{equation}
\label{eq:compact-conditional-trimmed-cdfs}
 A_z^{\rm ret}(t)=\min(A_z(t),M_z),
 \qquad
 B_z^{\rm ret}(t)=(B_z(t)-k_z)_+
\end{equation}
define source and target submeasures of equal mass \(M_z\),
continuous in total variation with respect to \(z\), and admitting
a forward coupling. Their normalized laws are TV-continuous on
\(\{M_z>0\}\).

Their complements in the original whole-ray marginals are also
balanced and forward feasible. In particular,
\begin{equation}
\label{eq:compact-conditional-trimming-loss}
 \theta(z)-M_z
       \le (\theta(z)-U_z)+(\theta(z)-V_z).
\end{equation}
Every measurable choice of forward couplings of the retained
submarginals extends, using a fixed forward coupling of the complement,
to a population optimum. It is not required to be a restriction of an
arbitrarily prescribed optimal plan.
\end{lemma}

\begin{proof}
The cumulative difference starts at zero and ends at \(U_z-V_z\),
so
\[
 \max(0,V_z-U_z)\le k_z\le V_z.
\]
Uniform continuity of the cumulatives on \(I\times K\) gives
continuity of their minimum and of \(M_z\). Formula
\eqref{eq:compact-conditional-trimmed-cdfs} retains the earliest
\(M_z\) source mass and the latest \(M_z\) target mass. Both
\(A_z(t)\) and \(M_z\) are at least \(B_z(t)-k_z\);
their minimum is therefore at least \((B_z(t)-k_z)_+\).
Proposition~\ref{prop:ray-forward-criterion} gives forward feasibility.

For atomless finite measures, prefix truncation satisfies
\[
 \|\operatorname{Prefix}_m(\lambda)
          -\operatorname{Prefix}_n(\eta)\|_{\rm TV}
       \le 2\|\lambda-\eta\|_{\rm TV}+|m-n|
\]
whenever the retained masses are admissible. To verify this, first
compare equal masses \(\min(m,n)\). On the common prefix the
variation is at most \(\|\lambda-\eta\|_{\rm TV}\); the
additional interval between the two cutoffs has mass at most the
same quantity, since the prefix totals agree. Changing the mass
costs \(|m-n|\). The suffix estimate is identical. Apply these
estimates to \(f_C(\cdot,z)\,dr\), which vary continuously in
\(L^1(I)\). This proves the asserted total variation continuity,
including when a cutoff moves across an empty radial gap.

Let \(H_z\ge0\) be the unnormalized whole-ray source-minus-target
cumulative difference, and put \(h_z=H_z(a)\). On \(I\),
\(H_z=h_z+A_z-B_z\), whence \(k_z\le h_z\). The retained surplus
satisfies
\begin{equation}
\label{eq:compact-conditional-surplus-domination}
 0\le H_z^{\rm ret}(t)
   \le A_z(t)-B_z(t)+k_z\le H_z(t),\qquad t\in I.
\end{equation}
Outside \(I\) the retained surplus is zero. Thus the complementary
surplus is everywhere nonnegative. Also \(h_z\) is at most the
source mass before \(a\), hence at most \(\theta(z)-U_z\).
The identity \(\theta-M=\theta-V+k\) proves
\eqref{eq:compact-conditional-trimming-loss}.

Generalized inverses of the Borel cumulative functions give measurable
forward quantile couplings for the retained and complementary laws.
These are existence witnesses. Any other measurable retained forward
coupling can be combined with the same complement, and reconstructed
by Proposition~\ref{prop:general-optimal-face}. The resulting full
plan is calibrated and optimal. At every zero of the original surplus,
\eqref{eq:compact-conditional-surplus-domination} forces the retained
surplus to vanish as well. Both parts therefore respect the intrinsic
active blocks of Theorem~\ref{thm:intrinsic-active-invariance}.
\end{proof}

These constructions retain arbitrarily large balanced population
pieces in fixed interior collars. Indeed, compact interior intervals
exhaust both conditional marginals on almost every whole ray. Choose
the first rational interval, with a slightly larger interior collar,
capturing all but a prescribed conditional mass of each color. This
gives a countable Borel partition by whole-ray labels. Retain finitely
many pieces and choose compact cores inside them, charging their
omitted quotient mass. Equation~\eqref{eq:compact-conditional-trimming-loss}
bounds the additional balancing loss by the two omitted endpoint
masses. Removing \(\{M<c\}\) then costs at most
\(c\operatorname{Leb}_q(K)\) of retained mass. Every remainder
belongs to the forward complement; no endpoint mass is silently
discarded. In particular a deterministic retained subplan and its
complement are eligible for the comparison in
Proposition~\ref{prop:general-upper-energy}.

Finally, the mark classes can be fixed on sufficiently small compact
label pieces. On the bounded flow range, for actual supported endpoints,
\[
 \delta_0\le d_F(X(r,z),X(s,w))
                   \le |r-s|+C|z-w|.
\]
The upper bound uses the unit-speed characteristic also outside its
maximal calibrated domain, followed by the label Lipschitz bound.
Thus \(C\operatorname{diam}(K)<\delta_0/2\) separates the source
and target radial projections by at least \(\delta_0/2\). Their
closures in the fixed bounded radial range remain separated and are
compact. They have finitely many alternating monochromatic runs, as
in Proposition~\ref{prop:ray-finite-bands}. On a closed interior
cylinder the supported preimages are already compact. This supplies
the fixed compact color classes of
Proposition~\ref{prop:forward-weak-continuity} without imposing global
source-before-target ordering.

The balanced trimmed densities can jump at their scalar cutoff graphs,
just as the whole-ray densities can jump at endpoint graphs; their
conditional laws nevertheless have the continuity proved above.
Arbitrary restrictions of a preselected pair kernel, or a selection of
individual active blocks at opening or closing zero-surplus cuts, need
not have that continuity. The compact transverse support may still
have holes, and the constants of different Whitney cores need not be
uniform. These population regularity statements supply neither an
empirical repair theorem on that support nor the global selector.


\section{From empirical optimizers to local matching laws}
\label{sec:highd-optimizer}

The fixed-root input limit can be connected to actual optimal
matchings on a regular separated tube when \(q=d-1\ge3\).
The tube and separation hypotheses in this section are explicit
intermediate assumptions; a general rough ray cluster need not
satisfy them.

Use the smooth calibrated coordinates of Section~\ref{sec:raytube}.
Let \(D=[0,1]^q\) be compactly contained in their transverse domain,
and let the source and target intervals \(I_A,I_B\) be compactly
contained in the radial domain, with their joining interval also
compactly contained there and \(b-a\ge\delta>0\).
Let \(\theta\) be a Lipschitz probability density on \(D\), bounded
above and below by positive constants. Write the coordinate marginals as
\[
\mu(dz,da)=\theta(z)\,dz\,\alpha_z(da),\qquad
\nu(dz,db)=\theta(z)\,dz\,\beta_z(db),
\]
where the conditional mark laws are probability kernels supported in
\(I_A,I_B\). Their dependence on \(z\) is allowed.
The physical marginals are their pushforwards under \(X\).

\begin{lemma}[Uniform transverse coercivity]
\label{lem:highd-coercivity}
There are constants \(0<c\le C<\infty\) such that, on these compact
endpoint sets,
\[
c|z-w|^2\le
 e(a,z;b,w):=d_F(X(a,z),X(b,w))-(b-a)
 \le C|z-w|^2.
\]
In particular, the continuum optimal couplings are exactly those
preserving the transverse label \(z\).
\end{lemma}

\begin{proof}
If the nonnegative defect vanishes, every unit-speed minimizing
geodesic between the endpoints is calibrated by the global potential.
At its source the potential is smooth, because the endpoints lie in
an open calibrated collar. Strong convexity fixes its initial unit
tangent to the radial vector. Uniqueness of the geodesic equation then
forces the target to have the same label \(z\).
Lemma~\ref{lem:raytube-distance}, uniformly over compact \((a,b,z)\),
gives a positive quadratic lower bound near \(z=w\). The continuous
ratio \(e/|z-w|^2\) has a positive minimum on the remaining compact
set, proving the global lower bound.

For the upper bound follow
\(r\mapsto X(r,z+(r-a)(w-z)/(b-a))\).
Convexity of \(D\) keeps it in the tube. Calibration gives
\(\widetilde F(r,z;1,0)=1\) and
\(D_v\widetilde F(r,z;1,0)=0\). Boundedness of the transverse
velocity Hessian on the relevant compact set yields
\(\widetilde F(r,z;1,v)\le1+K|v|^2\).
Its length is at most \(b-a+(K/\delta)|w-z|^2\).
Finally, the product kernel \(\alpha_z\otimes\beta_z\) gives a
coupling preserving \(z\), attaining the calibration lower bound.
Coercivity forces every optimal coupling to preserve \(z\) as well.
\end{proof}

\begin{proposition}[Uniform energy and no loss of matched endpoints]
\label{prop:highd-energy}
Assume \(q\ge3\), let \(\varepsilon_n=n^{-1/q}\), and sample the
two marginals independently. For the set \(\mathcal S_n\) of all
optimal permutations for the genuine length distance,
\begin{equation}
\E\!\left[\sup_{\sigma\in\mathcal S_n}
 \frac1n\sum_i
 \left|\frac{W_{\sigma(i)}-Z_i}{\varepsilon_n}\right|^2\right]
 \le K.
\label{eq:highd-energy}
\end{equation}
Choose an optimizer from the samples and any auxiliary randomness,
and subsequently choose \(U\) independently and uniformly in \(D\).
For every bounded Borel window \(B\subset\R^q\), the expected number
of its source endpoints whose partners are farther than \(R\) in
rescaled transverse coordinates is at most \(K|B|/R^2\).
The same bound holds for target endpoints. The estimates are uniform
over the initial choice of optimizer.
\end{proposition}

\begin{proof}
The subtracted radial sum is independent of the permutation.
Lemma~\ref{lem:highd-coercivity} therefore gives, simultaneously
for every optimal \(\sigma\),
\[
\sum_i|Z_i-W_{\sigma(i)}|^2
 \le\frac Cc\min_\tau\sum_i|Z_i-W_{\tau(i)}|^2.
\]
For independent uniform samples on \(D\),
\cite[Theorem~1.1]{GoldmanTrevisan2021} bounds the expected minimum
on the right by a constant times \(n^{1-2/q}\), for every \(q\ge3\).
This upper bound also holds for \(\theta\). Indeed its coordinate
marginals remain positive and Lipschitz. Their conditional densities
are uniformly bounded below and Lipschitz. The corresponding conditional
CDFs have Lipschitz inverses, jointly in their probability argument and
the preceding coordinates: a uniform positive slope controls the
inverse in its scalar argument, and the Lipschitz bound controls its
other arguments. Recursively inverting these CDFs gives a Lipschitz
map \(T:D\to D\) pushing uniform volume to \(\theta(z)dz\).
Applying \(T\) to a uniform comparison assignment multiplies its
quadratic cost by at most \(\Lip(T)^2\). This proves
\eqref{eq:highd-energy}.

Put \(d_i=(W_{\sigma(i)}-Z_i)/\varepsilon_n\).
Integrating first in the independent root gives
\[
\E\sum_i\mathbf1_{\{(Z_i-U)/\varepsilon_n\in B,\ |d_i|>R\}}
 \le |B|\varepsilon_n^q\E\sum_i\mathbf1_{\{|d_i|>R\}}
 \le \frac{K|B|}{R^2}.
\]
The identical argument anchors the window at the target endpoint.
Conditional on \(U\), each ordinary endpoint count in \(B\) is
binomial with one-point probability at most
\(\|\theta\|_\infty|B|/n\). Its second moment is uniformly bounded.
Thus local endpoint counts are uniformly integrable, and the two tail
bounds prevent either endpoint projection from losing mass to partners
at infinity in a subsequential local limit. The optimizer is fixed
before the independent root is introduced; no root-dependent choice
is used in this calculation.
\end{proof}

The energy estimate also implies
\begin{equation}
\E\!\left[\sup_{\sigma\in\mathcal S_n}
 \frac1n\sum_i|Z_i-W_{\sigma(i)}|\right]
 \le\sqrt K\,n^{-1/q}.
\label{eq:highd-macro-collapse}
\end{equation}
Consequently a Lipschitz test of physical endpoint pairs may replace
the target label \(W_{\sigma(i)}\) by the source label \(Z_i\),
uniformly over optimizers, with expected error tending to zero.
A uniform second-moment bound does not assert uniform integrability
of the squared edge lengths or equality of limiting energies.

\subsection{Changes in a common sample count}
\label{subsec:highd-balanced-count}

The same squared-matching input gives a useful comparison between
different sample sizes. Both colors retain the common transverse
density \(\theta\); their conditional mark laws may differ and
may depend on position.

\begin{lemma}[Expected assignment cost under equal-count changes]
\label{lem:highd-count-continuity}
Let \(q\ge3\), and suppose every source--target pair is allowed
for a fixed measurable cost satisfying
\[
0\le c((z,a),(w,b))\le\Lambda|w-z|^2.
\]
Write \(\mathcal C_m\) for the minimum total assignment cost of
two independent \(m\)-point samples, with \(\mathcal C_0=0\).
Then a constant \(K\), depending only on \(q,\theta\), satisfies
\[
\big|\E\mathcal C_m-\E\mathcal C_n\big|
 \le K\Lambda|m-n|^{1-2/q}
\qquad(m,n\ge0).
\]
This applies in particular to the fixed-potential defect in
Lemma~\ref{lem:highd-coercivity}.
\end{lemma}

\begin{proof}
Put \(s=1-2/q\). If \(\eta_h\) is the empirical probability
measure of \(h\) iid transverse points and \(\eta_h'\) is an
independent copy, Proposition~\ref{prop:highd-energy}'s unmarked
comparison gives
\(\E W_2^2(\eta_h,\eta_h')\le K_0h^{-2/q}\).
Averaging an optimal coupling conditional on \(\eta_h\) gives
a coupling to \(\theta(z)dz\), hence
\[
\E W_2^2(\eta_h,\theta(z)dz)\le K_0h^{-2/q}.
\]
Consequently, if \(\eta_h,\zeta_h\) have these two marginal
empirical laws but are possibly dependent, the triangle inequality
through \(\theta(z)dz\) yields
\begin{equation}
\E\big[hW_2^2(\eta_h,\zeta_h)\big]\le4K_0h^s.
\label{eq:highd-dependent-clouds}
\end{equation}
The value for \(h=0\) is taken as zero.

Let \(m=n+k\). Keeping an optimal \(n\)-point matching and
separately matching the \(k\) new iid points gives
\(\E\mathcal C_m\le\E\mathcal C_n+K_0\Lambda k^s\).
For the reverse inequality choose any measurable optimal permutation
\(\sigma\) for the \(m\)-point data. Independently delete uniform
\(k\)-subsets \(I,J\) of source and target indices. Retain edges
with neither endpoint deleted. The unmatched surviving index sets are
\[
A=\sigma^{-1}(J)\setminus I,\qquad
B=\sigma(I)\setminus J,
\qquad |A|=|B|=:h\le k.
\]
Conditional on the full data and \(\sigma\), the sets
\(I,T:=\sigma^{-1}(J)\) are independent uniform \(k\)-subsets.
The law of \(h=k-|I\cap T|\) is therefore independent of the
data. Given \(h\) and the data, \(T\setminus I\) is a uniform
\(h\)-subset of source indices, and \(\sigma(I\setminus T)\)
is a uniform \(h\)-subset of target indices. After averaging over
the original samples, each unmatched cloud conditional on \(h\)
has the iid transverse empirical law. Their mutual dependence is
allowed in \eqref{eq:highd-dependent-clouds}.

Repairing these unmatched points thus costs at most
\(4K_0\Lambda\E[h^s]\le4K_0\Lambda k^s\) in expectation.
Retained edges cost at most \(\mathcal C_m\) by nonnegativity.
The entire surviving clouds have, unconditionally, the two independent
\(n\)-point input laws. Their optimal expected cost is therefore
at most \(\E\mathcal C_m+4K_0\Lambda k^s\), proving the claim.
No uniqueness of the finite optimum was used.
\end{proof}

If \(N\) is Poisson with mean \(n\), independent of the input
arrays, use the same count \(N\) for both colors on the same cube.
The lemma and \(\E(N-n)^2=n\) imply
\[
\big|\E\mathcal C_N-\E\mathcal C_n\big|
 \le K\Lambda n^{(1-2/q)/2}.
\]
After normalization by \(n^{1-2/q}\), this error is
\(O(n^{-1/2+1/q})\), which vanishes for every \(q\ge3\).
Equivalently the bound on a cube of side \(L\) has an additional
factor \(L^2\); at \(L=n^{1/q}\) the same rate holds after
division by volume. Both costs in the comparison use the same cube.

Sharing a random count differs from taking two independent Poisson
restrictions, whose counts generally differ. The lemma alone does not
control their adaptively unmatched vertices; the next proposition
supplies that additional argument for uniform spatial marginals.
Spatial-boundary repair is a further question.

\subsection{Independent Poisson sample counts}
\label{subsec:highd-independent-count}

The common-count comparison extends to two independent Poisson
counts by repairing adaptively unmatched vertices. The following
statement uses uniform spatial marginals; marks may depend on their
own positions.

\begin{proposition}[Independent-count correction]
\label{prop:highd-independent-count}
Let \(q\ge3\), \(Q_n=[0,L_n]^q\), and \(L_n=n^{1/q}\), so
\(|Q_n|=n\). For each \(n\), take two independent iid marked
sample families with uniform spatial marginals on \(Q_n\), and
suppose every pair is allowed for a measurable finite cost satisfying
\[
0\le c_n((x,a),(y,b))\le\Lambda|y-x|^2,
\]
with \(\Lambda\) independent of \(n\). The per-color laws and
cost may depend on \(n\), but are fixed within each comparison
on \(Q_n\). Let \(F_n\) be the expected optimal total cost for
\(n\) points of each color. Let \(P_n\) be the expected minimum
cost of a maximum-cardinality matching when the two counts are
independent \(\operatorname{Poisson}(n)\) variables, independent
of the input arrays. Then
\[
\frac{|F_n-P_n|}{n}\longrightarrow0.
\]
\end{proposition}

\begin{proof}
Put \(s=1-2/q>0\) and \(t=1+2/q<2\).
For a deterministic \(m\), write \(f_m^{(n)}\) for the expected
\(m\)-by-\(m\) optimum on this same cube and for these same laws
and cost. Scaling Lemma~\ref{lem:highd-count-continuity} gives
\[
|f_m^{(n)}-F_n|\le C_q\Lambda L_n^2|m-n|^s.
\]
If \(M=\min(N_S,N_T)\), matching the first \(M\) points of
each color is a maximum-cardinality Poisson competitor. Therefore
\[
P_n\le\E f_M^{(n)}
 \le F_n+C_q\Lambda n^{1/2+1/q},
\]
using \(\E|M-n|^s\le C_q n^{s/2}\).

For the other direction use two infinite iid arrays, independent of
the counts. Their first \(N_S,N_T\) entries form the Poisson
clouds and their first \(n\) entries form the final clouds. Choose
a measurable optimal Poisson maximum matching and retain its edges
whose endpoints both occur in the final clouds. If \(k_0\) vertices
of each final color are unpaired, then
\[
k_0\le|N_S-n|+|N_T-n|,\qquad
\E k_0\le2\sqrt n,\qquad
\E[k_0^t]\le(4n)^{t/2}.
\]
These unpaired vertices need not have iid locations.

Fix \(0<\rho<1\). Conditional on all data and the retained
matching, give every retained edge an independent Bernoulli\((\rho)\)
coin shared by its endpoints. Give each unpaired final vertex its own
independent such coin. Keep only retained edges with coin zero.
The residual clouds are
\[
\mu=R_S+A,\qquad \nu=R_T+B,
\]
where \(R_S,R_T\) contain the coin-one vertices and \(A,B\)
the unpaired coin-zero vertices. Write their counts as
\(h_S,h_T,k_S,k_T\), respectively. Then
\[
h_S+k_S=h_T+k_T,\qquad k_S,k_T\le k_0.
\]
Within each color, distinct vertices use distinct independent coins.
Thus its full reserve-indicator vector has the Bernoulli product law
conditional on all data, with no data dependence. Each reserve cloud
separately is an independent thinning of an iid final cloud;
\(h_S,h_T\) separately have law \(\operatorname{Binomial}(n,\rho)\),
and conditional on its own count each reserve cloud is iid uniform.
The two reserve clouds need not be independent, and this assertion
does not condition on both counts.

Here \(W_2^2\) denotes total quadratic transport cost for equal
finite masses, without probability normalization. Set
\[
D_S=W_2^2\bigl(R_S,(h_S/n)\,dx\bigr),\qquad
D_T=W_2^2\bigl(R_T,(h_T/n)\,dx\bigr)
\]
on \(Q_n\), with value zero for an empty reserve.
Averaging an independent comparison cloud in
\cite[Theorem~1.1]{GoldmanTrevisan2021} bounds the expected transport
of \(h\) iid points to uniform mass \(h\) by \(C_qL_n^2h^s\).
The separate reserve laws and concavity therefore give
\[
\E(D_S+D_T)\le C_qL_n^2\E(h_S^s+h_T^s)
 \le C_q\rho^s n.
\]
This expectation is taken before imposing any further good event.

Partition \(Q_n\) into cubes of side
\(L_n/\lceil L_n\rceil\in[1/2,1]\). Smear every atom of
\(A,B\) uniformly in its own cell, producing densities \(a,b\)
at transport cost at most \(C_qk_0\).
For any fixed \(A_0>1+2/q\), put
\(B_n=C_{A_0}\log(en)\). The full final clouds have binomial
cell counts of mean at most one and there are at most \(2^qn\)
cells. A union bound gives failure probability \(O(n^{-A_0})\)
for the event that all these counts are at most \(B_n\).
On this event the adaptive subsets satisfy deterministically
\[
0\le a,b\le2^qB_n,\qquad
\int|a-b|\le k_S+k_T\le2k_0.
\]
Also
\(\Pp(\min(h_S,h_T)<\rho n/2)\le2e^{-\rho n/8}\).
Let \(G_n\) be the intersection of these occupancy and count
events; no independence between them is needed.

On \(G_n\), the densities
\[
f=h_S/n+a,\qquad g=h_T/n+b
\]
have equal total mass and are both at least \(\rho/2\).
Solve the zero-mean weak Neumann equation
\(-\Delta u=f-g\), \(\partial_\nu u=0\), on \(Q_n\).
The linear density interpolation with flux \(-\nabla u\) is a
continuity-equation competitor. The Neumann/Benamou--Brenier estimate
in \cite[proof of Lemma~3.4]{GoldmanTrevisan2021} yields
\[
W_2^2(f\,dx,g\,dx)\le\frac2\rho\int|\nabla u|^2.
\]
This also holds for the stated finite-mass convention by normalization
and multiplication by the common mass. For
\(p=2q/(q+2)\), critical Sobolev--Poincare gives
\(\|u\|_{2q/(q-2)}\le C_q\|\nabla u\|_2\).
Its constant is independent of \(L_n\): both norms scale as
\(L_n^{(q-2)/2}\). The constant density difference disappears
against the mean-zero test \(u\), so
\[
\int|\nabla u|^2=\int(a-b)u
 \le C_q\|a-b\|_p\|\nabla u\|_2.
\]
Using ordinary Lebesgue norms, the interpolation inequality
\(\|a-b\|_p^2\le
\|a-b\|_\infty^{2-2/p}\|a-b\|_1^{2/p}\) now gives
\[
W_2^2(f\,dx,g\,dx)
 \le\frac{C_q}{\rho}B_n^s k_0^t.
\]
This critical Sobolev refinement is the additional step beyond the
cited transport estimate; the unsmoothed atomic ports are never used
as an \(L^p\) density.

Combine reserve transport and port spreading to couple \(\mu\)
to \(f\,dx\), and likewise \(\nu\) to \(g\,dx\).
All intermediate measures have the same residual mass. The squared
triangle inequality with three legs gives, on \(G_n\),
\[
W_2^2(\mu,\nu)
 \le C_q\bigl(D_S+D_T+k_0+\rho^{-1}B_n^sk_0^t\bigr).
\]
Equal integer endpoint masses admit a minimizing permutation.
Its marked cost is at most \(\Lambda\) times this bound.
Adding the kept edges costs no more than the original Poisson optimum
plus the repair, since the original costs are nonnegative.
On \(G_n^c\), instead use any complete matching of the final
clouds, of cost at most \(\Lambda qnL_n^2\). The failure estimates
make its expected contribution \(o_\rho(n)\).
Taking expectations, using the unconditional nonnegative reserve
bounds, proves
\begin{equation}
\begin{split}
F_n\le P_n+C_q\Lambda\rho^s n
 &+C_q\Lambda\rho^{-1}(\log(en))^s n^{t/2}\\
 &+C_q\Lambda\sqrt n+o_\rho(n).
\end{split}
\label{eq:highd-count-repair}
\end{equation}
For fixed \(\rho\), all terms after the first error are \(o(n)\).
Divide by \(n\), send \(n\) to infinity, and then send
\(\rho\) to zero. Together with the reverse comparison, this
proves the proposition.
\end{proof}

This closes the independent-count correction for expected values.
The same bound need not absorb the larger number of endpoints exposed
by restricting an infinite stationary matching to a spatial cube.
That spatial-boundary repair, and uniqueness of paired-mark laws,
are not established here. The proposition also requires finite
all-pair costs; it does not apply directly to an infinite backward
cost or a sample-dependent radial gauge.

\subsection{Concentration of the optimized cost}
\label{subsec:highd-value-concentration}

For this subsection assume product mark laws
\(\alpha_z=\alpha\), \(\beta_z=\beta\) in the regular
separated tube. The common transverse density remains \(\theta\).
The marks themselves need no density or regularity assumption.

\begin{proposition}[Concentration of the genuine length defect]
\label{prop:highd-value-concentration}
Let \(q\ge3\), let \(g\) be a bounded measurable function of
the source and target marks, and put \(G=\|g\|_\infty\).
For \(t\in\R\) define
\[
C_n(t)=\min_\sigma\sum_{i=1}^n
 \bigl[n^{2/q}e(A_i,Z_i;B_{\sigma(i)},W_{\sigma(i)})
       +t g(A_i,B_{\sigma(i)})\bigr],
\]
where \(e\) is the actual distance defect in
Lemma~\ref{lem:highd-coercivity}. For every \(T<\infty\),
\begin{equation}
\sup_{|t|\le T}\E\left|
 \frac{C_n(t)}n-\frac{\E C_n(t)}n\right|
 \le K_T n^{-1/2+1/q}.
\label{eq:highd-value-concentration}
\end{equation}
Moreover, the supremum of the centered absolute difference over
\(|t|\le T\) tends to zero in probability.
The parameter \(t\) is a proof perturbation; \(t=0\) retains
the original distance optimizers.
\end{proposition}

\begin{proof}
We first verify a weak spatial derivative estimate for the genuine
defect. The distance triangle inequality and the smooth coordinate
maps give a uniform Lipschitz bound in \((z,w)\) on the compact
endpoint sets. Lemma~\ref{lem:raytube-distance} gives a uniform
smooth neighborhood of the diagonal pairs with bounded spatial
Hessian. Nonnegativity and \(e(a,z;b,z)=0\) force both endpoint
gradients to vanish there. Taylor's theorem consequently gives
\(|\nabla_z e|+|\nabla_w e|\le K|w-z|\) near the diagonal.
Away from it the uniform Lipschitz bound gives the same estimate
almost everywhere, after enlarging \(K\). Coercivity yields
\begin{equation}
|\nabla_z e|^2+|\nabla_w e|^2\le K e
\quad\text{almost everywhere, uniformly in }(a,b).
\label{eq:highd-defect-gradient}
\end{equation}
No smoothness of the distance at other cut points is required.

Fix \(t\), put \(L=n^{1/q}\), \(b=1+|t|G\), and write
\[
d_{ij}=L^2e(A_i,Z_i;B_j,W_j)+t g(A_i,B_j)+b,
\qquad V=C_n(t)+nb,\qquad Y=\sqrt V.
\]
Every \(d_{ij}\) is at least \(L^2e_{ij}+1\), so \(V\ge n\).
For fixed marks and a fixed permutation, the squared norm of the
full positional gradient of its total cost is at most
\(KL^2\sum_i L^2e_{i,\sigma(i)}\). Thus its square root has
squared gradient bounded by \(KL^2/4\). The minimum of the
finitely many square-root branches has the same almost-everywhere
bound. This uses no uniqueness of the minimizing permutation.

The single-point law \(\theta(z)dz\) has a Poincare inequality:
compare its variance with the squared deviation from the uniform
mean, then compare the two Dirichlet integrals using the upper and
lower bounds on \(\theta\). Tensorizing this single-point
inequality over the \(2n\) independent positions gives
\begin{equation}
\E\operatorname{Var}(Y\mid\text{all marks})\le K L^2.
\label{eq:highd-positional-variance}
\end{equation}

It remains to control mark randomness without a Poincare inequality
for the mark laws. Replacing one source mark changes only its row;
replacing one target mark changes only its column. By coercivity,
every affected edge satisfies \(d_{ij}'\le R d_{ij}\), where
\[
R=\max\{C/c,1+2|t|G\}
\]
and \(c,C\) are the constants in
Lemma~\ref{lem:highd-coercivity}. Choose a measurable old optimal
permutation. If its incident edge in the replaced row or column
has cost \(d_i\), then, on \(V'>V\), using that permutation
as a competitor gives
\[
V'-V\le(R-1)d_i,
\qquad (Y'-Y)_+^2\le(R-1)^2d_i^2/V.
\]
Let \(H=\E[Y\mid\text{all marks}]\). Efron--Stein for the
\(2n\) independent marks, Jensen using shared positions, and
exchangeability of each mark with its independent replacement give
\[
\begin{split}
\operatorname{Var}H
&\le\frac12\sum_{\text{rows, columns}}\E(H'-H)^2\\
&\le\sum_{\text{rows, columns}}\E(Y'-Y)_+^2
 \le2(R-1)^2\E\frac{\sum_i d_i^2}{V}.
\end{split}
\]
Each chosen edge is counted once by its source and once by its
target. Since \(\sum_i d_i=V\), the last ratio is at most
\(\max_i d_i\le CqL^2+1+2|t|G\).
Together with \eqref{eq:highd-positional-variance}, conditional
variance decomposition proves
\[
\operatorname{Var}\sqrt V\le K_T L^2
\qquad(|t|\le T).
\]

The unmarked comparison from Proposition~\ref{prop:highd-energy}
gives \(\E V\le K_T n\). For a nonnegative random variable
\(Y\), with mean \(m\) and variance \(v\), Cauchy--Schwarz gives
\[
\E|Y^2-\E Y^2|
 \le\E[|Y-m|(Y+m)]+v
 \le2\sqrt{v\E Y^2}+v.
\]
Apply this to \(Y^2=V\); the deterministic shift \(nb\)
cancels. Division by \(n\) bounds the centered cost by
\(K_T(n^{-1/2+1/q}+n^{-1+2/q})\), implying
\eqref{eq:highd-value-concentration}. Finally both
\(C_n(t)/n\) and its expectation are \(G\)-Lipschitz in
\(t\). A finite mesh on \([-T,T]\), followed by decreasing
its spacing, gives uniform convergence in probability.
\end{proof}

The same proof applies to the homogeneous matrix kinetic cost
\(k_{ab}(y-x)=\tfrac12(y-x)^TS(a,b)^{-1}(y-x)\),
with uniform positive upper and lower compliance bounds, using
\(L^2k_{ab}(w-z)\) on the unit cube. Its spatial derivative
estimate is immediate. The mark laws remain independent of the
positions in both applications.

The deterministic mean functions \(\E C_n(t)/n\) are concave,
locally bounded, and \(G\)-Lipschitz, so they have locally uniform
subsequential limits. Concentration does not show that those limits
agree or that they are differentiable at \(t=0\). For example,
every unperturbed optimal permutation satisfies, for \(t>0\),
\[
\frac{C_n(t)-C_n(0)}{nt}
\le\frac1n\sum_i g(A_i,B_{\sigma(i)})
\le\frac{C_n(0)-C_n(-t)}{nt}.
\]
A unique mean limit differentiable at zero would make these two
bounds agree and identify the paired test for all optimizers.
Neither conclusion follows from scalar concentration alone.

\subsection{A finite-volume limit for homogeneous matrix costs}
\label{subsec:highd-finite-pressure}

We now specialize to uniform transverse sampling and product mark
laws. Fix compact metric mark spaces \(K_-,K_+\), probability laws
\(\alpha,\beta\), and a measurable symmetric positive definite compliance
matrix \(S(a,b)\) such that
\[
k_{ab}(v)=\tfrac12v^TS(a,b)^{-1}v,
\qquad c_-|v|^2\le k_{ab}(v)\le c_+|v|^2
\]
for fixed \(0<c_-\le c_+<\infty\). Every mark pair is allowed.
The matrix may depend on both marks but not on spatial position.

\begin{proposition}[Homogeneous finite-volume pressure]
\label{prop:highd-finite-pressure}
Let \(q\ge3\), let \(Z_i,W_j\) be independent uniform points
in \((0,1)^q\), and let their independent marks \(A_i,B_j\)
have laws \(\alpha,\beta\), independently of all positions.
For bounded measurable \(g:K_-\times K_+\to\R\), put
\[
\mathcal K_n(t)=\min_\sigma\sum_{i=1}^n
 \bigl[n^{2/q}k_{A_iB_{\sigma(i)}}(W_{\sigma(i)}-Z_i)
       +t g(A_i,B_{\sigma(i)})\bigr].
\]
There is a finite concave \(\|g\|_\infty\)-Lipschitz function
\(p_g:\R\to\R\) such that, for every \(T<\infty\),
\[
\sup_{|t|\le T}\left|\frac{\E\mathcal K_n(t)}n-p_g(t)\right|
 \longrightarrow0,
\qquad
\sup_{|t|\le T}\left|\frac{\mathcal K_n(t)}n-p_g(t)\right|
 \xrightarrow{\Pp}0.
\]
\end{proposition}

\begin{proof}
Write \(G=\|g\|_\infty\), \(s=1-2/q>0\), and
\(Q_R=(0,R)^q\). Let \(\mathcal C_t(R)\) be the minimum
\(k+t g\) cost of a maximum-cardinality matching between two
independent unit-intensity marked Poisson processes restricted to
\(Q_R\), and set \(f_t(R)=R^{-q}\E\mathcal C_t(R)\).
Exactly the smaller input count is matched, even for signed costs.
The unmarked quadratic comparison gives, for \(R\ge1\),
\[
-|t|G\le f_t(R)\le Cc_++|t|G,
\qquad |f_t(R)-f_u(R)|\le G|t-u|.
\]
Each \(f_t(R)\) is concave in \(t\). For fixed \(t\),
it is continuous in \(R>0\): under a common Poisson realization,
the point sets are locally constant at each deterministic cube side
almost surely, and a bounded-cube point count dominates the cost
integrably.

Fix \(\rho\in(0,1/2)\) and \(\lambda=1-\rho\). Split each
input into independent main and reserve processes of intensities
\(\lambda\) and \(\rho\). In each of the \(m^q\) side-\(L\)
blocks of \(Q_{mL}\), choose an optimal maximum-cardinality main
matching by the same measurable translated rule. If \(U_i,V_i\)
are its unmatched sets and \(N_i^-,N_i^+\) its input counts, then
\[
|U_i|=(N_i^--N_i^+)_+,
\qquad |V_i|=(N_i^+-N_i^-)_+,
\qquad
\E(|U_i|^r+|V_i|^r)\le C_rL^{qr/2}\quad(r\ge1).
\]
The last bound is the centered Poisson moment estimate, uniform
for \(\lambda\in[1/2,1]\) and \(L\ge1\).
The pairs \((U_i,V_i)\) are independent translated copies and
are independent of the reserves; their within-block positions may
be arbitrary functions of the marked samples.

These are precisely the hypotheses of
\cite[Proposition~6.5]{GoldmanTrevisan2023}, with dimension \(q\)
and power \(2<q\). Writing \(M_2\) for the minimum squared
cost of a maximum-cardinality matching of the residual clouds,
that proposition gives, uniformly in \(m\),
\[
\frac{\E M_2}{(mL)^q}\le C\rho^s+C(\rho)L^{-\gamma}
\]
for some \(\gamma>0\) and all sufficiently large \(L\).
The leading \(C\) is independent of \(\rho\). The expected
number \(H\) of residual edges satisfies
\(\E H/(mL)^q\le\rho+CL^{-q/2}\).
Use this Euclidean repair, whose marked cost is at most
\(c_+M_2+|t|GH\). Its union with the block edges has the required
global cardinality: if \(J\) edges were kept, its total is
\(J+\min(N^--J,N^+-J)=\min(N^-,N^+)\).

Scaling a main block by \(\lambda^{1/q}\) makes its intensity
one and gives the exact cost identity
\[
k_{ab}(y-x)+t g(a,b)
 =\lambda^{-2/q}\bigl[k_{ab}(y'-x')
                   +t\lambda^{2/q}g(a,b)\bigr].
\]
Consequently
\begin{equation}
\begin{split}
f_t(mL)&\le\lambda^s
 f_{t\lambda^{2/q}}(\lambda^{1/q}L)
 +Cc_+\rho^s+c_+C(\rho)L^{-\gamma}\\
&\hspace{18mm}+|t|G(\rho+CL^{-q/2}).
\end{split}
\label{eq:highd-pressure-block}
\end{equation}
This comparison uses no nonnegativity of the signed block costs.
For arbitrary \(R\to\infty\), take
\(m_R=\lfloor R/L\rfloor\) and \(\ell_R=R/m_R\to L\).
Apply \eqref{eq:highd-pressure-block} with side \(\ell_R\).
Continuity at fixed side shows that its right side with \(L\)
also bounds \(\limsup_{R\to\infty}f_t(R)\).

Put \(b_t=\liminf_{R\to\infty}f_t(R)\), choose
\(R_j\to\infty\) with \(f_t(R_j)\to b_t\), and use
\(L_j=\lambda^{-1/q}R_j\) in that bound. The parameter
Lipschitz estimate gives
\[
\limsup_{R\to\infty}f_t(R)
 \le\lambda^s b_t
 +\lambda^s|t|G(1-\lambda^{2/q})+Cc_+\rho^s+|t|G\rho.
\]
Sending \(\rho\downarrow0\) proves existence of
\(p_g(t)=\lim_{R\to\infty}f_t(R)\).
Concavity and the common Lipschitz bound pass to the limit and
imply locally uniform convergence in \(t\).

We check the bounded additive term when passing to equal counts;
Proposition~\ref{prop:highd-independent-count} cannot be invoked
for signed costs directly. Put \(B=|t|G\) and
\(d_t=k+t g+B\), so \(0\le d_t\le c_+|y-x|^2+2B\).
On the cube of volume \(n\), the proof of
Lemma~\ref{lem:highd-count-continuity} now gives
\[
|\E D_m-\E D_n|
 \le Cc_+ n^{2/q}|m-n|^s+2B|m-n|,
\]
where \(D_m\) is the optimal \(m\)-by-\(m\) \(d_t\)-cost
on that same cube. Indeed at most \(|m-n|\) repair edges each
contribute the added bound \(2B\).
In the independent-count proof, after thinning with reserve
fraction \(\eta\), the residual count has expectation at most
\(\eta n+2\sqrt n\). Its repair therefore gains at most
\(2B\eta n+4B\sqrt n\); the bad-event fallback gains at most
\(2Bn\) on an event of vanishing probability. All retained
costs remain nonnegative. The reverse comparison gains only
\(O(B\sqrt n)\). First \(n\to\infty\), then
\(\eta\downarrow0\), proves the same normalized count
comparison for \(d_t\). Subtracting its shift changes the two
expected values by amounts differing by \(O(B\sqrt n)\), since
\(\E|\min(N^-,N^+)-n|\le2\sqrt n\).
Thus \(\E\mathcal K_n(t)/n\to p_g(t)\).
A finite mesh and the common Lipschitz bound make this convergence
locally uniform. The kinetic version of
Proposition~\ref{prop:highd-value-concentration} gives the asserted
locally uniform convergence in probability.
\end{proof}

The finite-volume argument alone gives neither a stationary variational
identity nor differentiability at zero. For common-cumulative
compliance, Proposition~\ref{prop:highd-stationary-pressure} below
proves the variational identity. For continuous compliance,
Theorem~\ref{thm:highd-pressure-differentiability} proves differentiability.
Proposition~\ref{prop:highd-variable-pressure} extends the pressure limit
to regular ordered tubes with varying coefficients and the positive
conditional densities specified there.

\begin{corollary}[Paired-mark selection from pressure]
\label{cor:highd-pressure-selection}
Suppose \(p_g\) is differentiable at zero for every \(g\) in a
countable uniformly dense subset of \(C(K_-\times K_+)\).
This hypothesis holds whenever \(S\) is continuous, by
Theorem~\ref{thm:highd-pressure-differentiability}. Then a deterministic \(\kappa\in\Pi(\alpha,\beta)\) exists
such that the paired-mark empirical measures of all unperturbed
kinetic optimizers converge weakly to \(\kappa\), uniformly over
optimizers in probability. The same holds for optimal fractional
couplings with their paired-mark pushforwards.
\end{corollary}

\begin{proof}
For every optimal permutation and \(t>0\), its paired test lies
between \((\mathcal K_n(t)-\mathcal K_n(0))/(nt)\) and
\((\mathcal K_n(0)-\mathcal K_n(-t))/(nt)\).
The proposition, followed by \(t\downarrow0\), makes both bounds
converge to \(p_g'(0)\), uniformly over optimizers. Compactness,
the dense tests, and convergence of the two empirical mark
marginals identify a unique deterministic \(\kappa\) and give
the stated uniform weak convergence. Birkhoff decomposition into
optimal permutations extends it to every optimal fractional
coupling. The cited theorem verifies the differentiability hypothesis
for continuous compliance.
\end{proof}

The corollary concerns the global paired-mark pushforward. The
spatially resolved reconstruction in
Corollary~\ref{cor:highd-uniqueness-criterion} still requires
identification of the conditional paired laws at almost every
transverse label.

\subsection{Dual rigidity and transfer to the actual distance}
\label{subsec:highd-dual-transfer}

Return to the regular ordered tube, with \(q\ge3\), product
mark laws \(\alpha,\beta\), and the common positive Lipschitz
density \(\theta\) on \(D=[0,1]^q\). Set
\(K_A=\operatorname{supp}\alpha\),
\(K_B=\operatorname{supp}\beta\), and write the coordinate
supports as \(\mathsf X=K_A\times D\),
\(\mathsf Y=K_B\times D\). We first consider any continuous
cost \(c\) on their product satisfying
\begin{equation}
c_0|z-w|^2\le c((a,z),(b,w))\le C_0|z-w|^2,
\qquad 0<c_0\le C_0<\infty.
\label{eq:highd-dual-cost}
\end{equation}
Fix an anchor \(x_0\in\mathsf X\). A normalized continuous
dual pair satisfies \(f(x_0)=0\) and \(f(x)+h(y)\le c(x,y)\).

\begin{lemma}[Rigidity of the continuum dual]
\label{lem:highd-dual-rigidity}
The unique normalized continuous optimal dual pair is \((0,0)\).
\end{lemma}

\begin{proof}
The same-label coupling
\(\theta(z)\,dz\,\alpha(da)\,\beta(db)\) has zero cost and
full support on the diagonal relation
\(\{((a,z),(b,z)):a\in K_A,b\in K_B,z\in D\}\).
For an optimal dual, the nonnegative continuous slack integrates
to zero against this coupling. Hence
\(f(a,z)+h(b,z)=0\) at every such triple. It follows that
\(f(a,z)=u(z)\), \(h(b,z)=-u(z)\). Feasibility in both
orders of \(z,w\) gives
\(|u(z)-u(w)|\le C_0|z-w|^2\).
Subdividing the segment from \(z\) to \(w\) into \(m\)
equal pieces bounds this difference by \(C_0|z-w|^2/m\).
Thus \(u\) is constant, and normalization makes it zero.
\end{proof}

This dual uniqueness does not make the same-label primal coupling
unique. Its quantitative consequence below concerns optimal edges.

\begin{proposition}[Exponential localization of all optimal edges]
\label{prop:highd-exponential-localization}
Fix \(g\in C(K_A\times K_B)\) and \(T\ge0\). For
\(|t|\le T\), let
\[
d_{n,t}(x,y)=c(x,y)+n^{-2/q}t g(a,b).
\]
Let \(R_n^c(T)\) be the largest transverse distance of an edge
in the support of any optimal empirical coupling for any parameter
\(|t|\le T\). This includes optimal fractional couplings.
For each \(\varepsilon>0\), there are positive constants
\(A_{\varepsilon,T},a_{\varepsilon,T}\) such that, for all
sufficiently large \(n\),
\[
\Pp\{R_n^c(T)>\varepsilon\}
 \le A_{\varepsilon,T}e^{-a_{\varepsilon,T}n}.
\]
In particular \(R_n^c(T)\to0\) almost surely under a common
iid sample realization.
\end{proposition}

\begin{proof}
Write \(G=\|g\|_\infty\), \(\delta_n=n^{-2/q}TG\),
and denote the sampled coordinate points by \(x_i,y_j\).
Take any optimal finite dual vectors \((u_i,v_j)\) for
\(d_{n,t}\) and extend them by
\[
F(x)=\min_j\{d_{n,t}(x,y_j)-v_j\},\qquad
H(y)=\min_{x\in\mathsf X}\{d_{n,t}(x,y)-F(x)\}.
\]
Every row and column has a tight edge in an optimal permutation.
Feasibility and these tight edges give
\(F(x_i)=u_i\), \(H(y_j)=v_j\). Thus the extension is
globally feasible and preserves the empirical dual objective.
Normalize it as \(f=F-F(x_0)\), \(h=H+F(x_0)\), and
put \(\widehat h=h-\delta_n\). This pair is feasible for
the base cost \(c\). Writing \(J,J_n\) for its population
and empirical dual objectives, respectively, gives
\begin{equation}
J_n(f,\widehat h)
 =\min_\sigma\frac1n\sum_i d_{n,t}(x_i,y_{\sigma(i)})
   -\delta_n\ge-2\delta_n.
\label{eq:highd-shifted-dual}
\end{equation}

These pairs lie in one deterministic compact class. Indeed put
\(M=C_0q\), \(B=M+TG\), and \(R=4B\). All costs
\(d_{n,t}\) have absolute value at most \(B\) and common
source and target continuity moduli, bounded by those of \(c\)
plus \(T\) times those of \(g\). Taking minima preserves
these moduli. The oscillation of \(F\) is at most \(2B\),
so \(\|f\|_\infty\le2B\). The second transform and
\(f(x_0)=0\) give \(-3B\le h\le B\), hence
\(\|\widehat h\|_\infty\le4B\).
Let \(\mathcal D\) consist of all normalized \(c\)-feasible
pairs with both norms bounded by \(R\) and these common
moduli. It is compact in the maximum uniform norm. Its population
objective is continuous, at most zero, and, by
Lemma~\ref{lem:highd-dual-rigidity}, maximized only at \((0,0)\).

For each \(\eta>0\), compactness therefore supplies
\(\Delta_\eta>0\) such that
\[
J(f,h)\le-\Delta_\eta
\quad\text{if }(f,h)\in\mathcal D,
\quad\max\{\|f\|_\infty,\|h\|_\infty\}\ge\eta.
\]
If that subset is empty, the required norm bound is automatic.
For \(s>0\), choose a finite \(s/8\)-net of
\(\mathcal D\), with \(N(s)\) elements. Replacing a pair
by its representative changes \(J_n-J\) by at most \(s/2\).
For each representative the iid variables
\(f(x_i)+h(y_i)\) lie in \([-2R,2R]\). Hoeffding's
inequality and the union bound yield
\begin{equation}
\Pp\!\left\{\sup_{\mathcal D}|J_n-J|>s\right\}
 \le2N(s)\exp\!\left(-\frac{ns^2}{32R^2}\right).
\label{eq:highd-dual-uniform-lln}
\end{equation}
The supremum is measurable by compactness and continuity in the
dual pair.

Take \(\eta=c_0\varepsilon^2/8\). On the event
\(\sup_{\mathcal D}|J_n-J|\le\Delta_\eta/4\), once
\(\delta_n\le\Delta_\eta/8\),
\eqref{eq:highd-shifted-dual} excludes the separated subset.
Thus both shifted dual norms are below \(\eta\), simultaneously
for all parameters and all finite optimal dual choices. Every edge
in any optimal empirical coupling is tight for these duals.
Consequently, once also \(\delta_n\le c_0\varepsilon^2/8\),
\[
c_0|z-w|^2\le c(x,y)
 =f(x)+\widehat h(y)+\delta_n-n^{-2/q}t g(a,b)
 \le2\eta+2\delta_n\le\tfrac12c_0\varepsilon^2.
\]
Equation~\eqref{eq:highd-dual-uniform-lln} proves the claimed
tail bound. Complementary slackness applies also to every positive
entry of an optimal fractional coupling; equivalently, use its
Birkhoff decomposition into optimal permutations. The event
involving all parameters is measurable: for each of finitely many
permutations, its cost excess above the minimum is continuous in
\(t\), and its minimum on \([-T,T]\) is measurable.
Finally apply Borel--Cantelli to \(\varepsilon=1/m\).
\end{proof}

Define the source-frozen quadratic defect by
\[
k((a,z),(b,w))
 =\tfrac12(w-z)^TS(a,b,z)^{-1}(w-z),
\]
and, for \(c=e\) or \(k\), define its normalized value
\[
V_n^c(t)=\frac1n\min_\sigma\sum_i
 \bigl[n^{2/q}c(x_i,y_{\sigma(i)})
       +t g(A_i,B_{\sigma(i)})\bigr].
\]
Both values use the same samples. The matrix in \(k\) may vary
with the source transverse coordinate \(z\).

\begin{proposition}[Uniform transfer of the optimized value]
\label{prop:highd-quadratic-transfer}
For every \(g\in C(K_A\times K_B)\) and \(T<\infty\),
\[
\E\sup_{|t|\le T}|V_n^e(t)-V_n^k(t)|\longrightarrow0.
\]
\end{proposition}

\begin{proof}
The two continuous costs obey \eqref{eq:highd-dual-cost} with
common constants. Lemma~\ref{lem:raytube-distance}, uniformly on
the compact endpoint sets, gives \(A,r_0>0\) such that
\[
|e((a,z),(b,w))-k((a,z),(b,w))|
 \le A|z-w|^3\qquad(|z-w|\le r_0).
\]
Let \(Q_n\) be the minimum total unmarked squared transverse
distance, and put \(\delta_n=n^{-2/q}T\|g\|_\infty\).
Comparing any optimizer for either perturbed cost with
an unmarked optimum gives, simultaneously for \(|t|\le T\),
\begin{equation}
\sum_i|Z_i-W_{\sigma(i)}|^2
 \le c_0^{-1}(C_0Q_n+2n\delta_n).
\label{eq:highd-transfer-energy}
\end{equation}
The bound used in Proposition~\ref{prop:highd-energy} gives
\(\E Q_n\le K n^{1-2/q}\), including for the stated
density \(\theta\).

Fix \(0<\varepsilon\le r_0\). On the event that both
\(R_n^e(T)\) and \(R_n^k(T)\) are at most \(\varepsilon\),
use an optimizer of either cost as a competitor for the other.
The perturbations cancel on each chosen edge, and the cubic bound
and \eqref{eq:highd-transfer-energy} yield
\[
\sup_{|t|\le T}|V_n^e(t)-V_n^k(t)|
 \le A\varepsilon n^{2/q-1}c_0^{-1}
       (C_0Q_n+2n\delta_n).
\]
The expectation of its right side is at most \(K_T\varepsilon\).
On the complementary event the deterministic comparison
\(\sup_{|t|\le T}|V_n^e(t)-V_n^k(t)|
 \le n^{2/q}\|e-k\|_\infty\)
still holds. Proposition~\ref{prop:highd-exponential-localization}
makes this latter contribution tend to zero. First let
\(n\to\infty\), then \(\varepsilon\downarrow0\).
Only bounded expected energy and the exponential bad-event bound
were used; uniform integrability of the microscopic edge energy
was not assumed.
\end{proof}

\begin{corollary}[Pressure for genuine homogeneous-compliance tubes]
\label{cor:highd-length-pressure}
Suppose \(\theta=1\) and \(S(a,b,z)=S(a,b)\) is independent
of \(z\). For every continuous paired-mark test \(g\), the
pressure \(p_g\) from Proposition~\ref{prop:highd-finite-pressure}
also satisfies
\[
\sup_{|t|\le T}|\E V_n^e(t)-p_g(t)|\longrightarrow0,
\qquad
\sup_{|t|\le T}|V_n^e(t)-p_g(t)|\xrightarrow{\Pp}0.
\]
\end{corollary}

\begin{proof}
Here \(nV_n^k=\mathcal K_n\) in
Proposition~\ref{prop:highd-finite-pressure}. Apply that proposition
and the uniform \(L^1\) transfer just proved.
\end{proof}


The value is the normalized distance defect; subtracting the radial
endpoint sum preserves every unperturbed genuine-distance optimizer.
Theorem~\ref{thm:highd-pressure-differentiability} verifies the
differentiability required by Corollary~\ref{cor:highd-pressure-selection}.
The finite-difference argument therefore identifies the global paired-mark pushforwards of
all these optimizers, including fractional ones. Tests \(g(a,b)\)
alone do not identify their spatially conditional paired laws.
Corollary~\ref{cor:highd-spatial-pressure} below supplies the spatial
pressure needed for full-plan selection in this homogeneous-compliance class.
For spatially varying compliance or nonuniform \(\theta\), the
value comparison above remains valid. Proposition~\ref{prop:highd-variable-pressure}
below proves the corresponding pressure limit on regular ordered tubes,
and also allows the positive continuous conditional mark densities
specified there. The present value comparison retains its product-mark
assumptions.

\subsection{Finite-energy screening for ordered kinetic matchings}
\label{subsec:highd-kinetic-screening}

Let \(q\ge3\), and let the two independent marked Poisson inputs have
common spatial intensity \(0<\lambda\le1\) and product mark laws
\(\alpha,\beta\) on compact intervals \(I_A,I_B\), with
\(\max I_A+\delta\le\min I_B\). Atoms in either mark law are allowed.
Let \(T\) be a \(C^1\) symmetric-matrix-valued function on the convex
hull of \(I_A\cup I_B\), with \(hI\le T'(r)\le HI\), where
\(0<h\le H<\infty\).
For an edge \(e=((x,a),(y,b))\), write
\[
 S_e=T(b)-T(a),\qquad v_e=S_e^{-1}(y-x),\qquad
 X_e(r)=x+[T(r)-T(a)]v_e\quad(a\le r\le b).
\]
The cost is \(c(e)=\tfrac12(y-x)^TS_e^{-1}(y-x)\).
All pairs are admissible, and this cost is uniformly comparable to
\(|y-x|^2\). The scalar model has \(T(r)=rI\).
Throughout, \(M\) denotes a stationary perfect matching of these inputs;
\(\E^0\) denotes its source-Palm expectation. Suppose
\(\E^0|y-x|^2<\infty\), equivalently \(\E^0|v_e|^2<\infty\).

Fix a nonnegative radial \(\eta\in C_c^\infty(B_1)\), positive on
\(B_{1/2}\), with integral one. Define the smoothed current and its
nonnegative traffic majorant by
\begin{equation}
\label{eq:highd-kinetic-current}
 J(z)=\sum_{e\in M}\int_a^b T'(r)v_e\eta(z-X_e(r))\,dr,
 \qquad
 U(z)=\sum_{e\in M}\int_a^b |T'(r)v_e|\eta(z-X_e(r))\,dr.
\end{equation}
Each path has length and diameter at most \(C|y-x|\). Campbell's
formula gives local integrability and
\begin{equation}
\label{eq:highd-current-divergence}
 \operatorname{div}J=\eta*(\mu_S-\mu_T),\qquad
 \E J(0)=\lambda\E^0(y-x),
\end{equation}
where \(\mu_S,\mu_T\) are the unmarked input counting measures.
Indeed, the contribution of one edge to the divergence is
\(\eta(\,\cdot-x)-\eta(\,\cdot-y)\), and
\(\int_a^b T'(r)v_e\,dr=y-x\).

\begin{lemma}[Interface velocity bound]
\label{lem:highd-interface-velocity}
Suppose two-cycle optimality holds for \(M\). Consider edges active at
a common time \(r_0\), and put
\[
 A_i=T(r_0)-T(a_i),\quad B_i=T(b_i)-T(r_0),\quad
 z_i=X_i(r_0).
\]
If, for \(0<t\le t_*\),
\(c_0tI\le A_i\le C_0tI\) and \(c_0I\le B_i\le C_0I\), then
\begin{equation}
\label{eq:highd-interface-velocity}
 |v_i-v_j|\le C|z_i-z_j|/t.
\end{equation}
The analogous assertion holds with \(A\) and \(B\) exchanged.
\end{lemma}
\begin{proof}
Set \(F_i(z)=\tfrac12(z-x_i)^TA_i^{-1}(z-x_i)\) and
\(G_i(z)=\tfrac12(y_i-z)^TB_i^{-1}(y_i-z)\).
Their sum has minimum \(c_{ii}\) at \(z_i\), with individual gradients
\(v_i,-v_i\). At \(\bar z=(z_i+z_j)/2\),
\[
 F_i(\bar z)+G_j(\bar z)+F_j(\bar z)+G_i(\bar z)
 \le c_{ii}+c_{jj}+C|z_i-z_j|^2/t.
\]
The gradient of \(F_i+G_j\) there equals
\(v_i-v_j+\tfrac12(A_i^{-1}-B_j^{-1})(z_j-z_i)\).
Its Hessian is at most \(CI/t\), so minimization saves at least
\(ct\) times the squared norm of this gradient. Minimizing the other
swapped objective cannot increase it. Their minima are the two crossed
costs, since their compliances are \(A_i+B_j\) and \(A_j+B_i\).
Two-cycle optimality bounds the saving by \(C|z_i-z_j|^2/t\), proving
\eqref{eq:highd-interface-velocity}. No commutation of the matrices is
used.
\end{proof}

\begin{theorem}[Square-integrable kinetic traffic]
\label{thm:highd-kinetic-traffic}
If \(M\) is two-cycle optimal, then
\begin{equation}
\label{eq:highd-traffic-ltwo}
 \E U(0)^2\le C\bigl(1+\lambda\E^0|v_e|^2\bigr)<\infty.
\end{equation}
The constant depends only on the dimension, coefficient bounds, mark
intervals, separation, and smoothing kernel.
\end{theorem}
\begin{proof}
We first prove a bound for a short time window \(I\), centered at
\(r_0\), of length at most \(2\varepsilon t\). Fix a deterministic
source-mark set \(A\), and any collection of edges with \(a\in A\),
active throughout \(I\), satisfying the source-side hypotheses of
Lemma~\ref{lem:highd-interface-velocity} at \(r_0\). Also suppose
\(\|[T(r_0)-T(a)]/t\|\le C_0\) for every \(a\in A\).
Write \(U_{A,I}\) for this collection's traffic at zero integrated over
\(I\). For \(K\ge K_0\) we claim
\begin{equation}
\label{eq:highd-short-window-traffic}
 \E U_{A,I}^2
 \le CKt^2\lambda\E^0[|v_e|^2\mathbf1_{\{a\in A\}}]
       +Ce^{-cK}.
\end{equation}
All constants are uniform in \(t,A,I\) and the selected collection.

For an empty collection the claim is immediate. Otherwise the
velocity map \(z_i\mapsto v_i\) has a Lipschitz extension \(V\)
to \(\R^q\) with constant \(C/t\). Choose \(\varepsilon\) sufficiently
small. For \(C_r=T(r)-T(r_0)\), the map
\(F_r(z)=z+C_rV(z)\) is a bijection with inverse Lipschitz constant
at most two. Set
\(\gamma(r)=F_r^{-1}(0)\), \(w(r)=tV(\gamma(r))\), and
\(W_0=t|V(0)|\). The contraction identity gives
\[
 |\gamma(r)|\le C\varepsilon W_0,\qquad
 |w(r)-tV(0)|\le C\varepsilon W_0,\qquad
 \operatorname{length}(\gamma,w)\le CW_0.
\]
For the last bound, subtraction at \(r,s\) gives
\(|\gamma(r)-\gamma(s)|\le C(W_0/t)|r-s|\); the map \(tV\)
has uniformly bounded Lipschitz constant. A contributing edge satisfies,
at some \(r\in I\),
\[
 |z_i-\gamma(r)|\le2,\qquad |tv_i-w(r)|\le C.
\]
Thus its phase point \((z_i,tv_i)\) lies in a fixed-radius tube about
a curve of length \(CW_0\) in \(\R^{2q}\), contained in a ball of
radius \(C(1+W_0)\). When \(W_0\) is large, all contributing
\(t|v_i|\) are comparable to \(W_0\). When \(W_0\) is bounded,
their source positions are bounded as well, since
\(x_i=z_i-C_{a_i}tv_i\), where \(C_a=[T(r_0)-T(a)]/t\).

We control every such curve using the original Poisson source input.
A unit phase cell with integer corner \((k,l)\) has source preimage
contained in the deterministic marked set
\[
 \{(x,a):a\in A,\ |x-(k-C_al)|\le C\}.
\]
Its intensity is at most \(C\lambda\alpha(A)\le C\). The union
of the preimages of \(m\) cells therefore has intensity at most \(Cm\),
even if the sets overlap. In a lattice with fixed finite-degree
adjacency, the number of connected sets of \(m\) cells containing a
specified root is at most \(C^{m}\). A fixed-radius tube around a curve
of length \(CW\) is covered by a connected set of at most \(C'(1+W)\)
cells. For \(W\ge1\), pad this set to exactly \(\lceil C''W\rceil\)
cells. Its root has \(O(W^{2q})\) possibilities in the phase ball.
The Poisson upper tail and this union bound yield a data-only event
\(E_{W,K}\) such that
\begin{equation}
\label{eq:highd-phase-animal-tail}
 \Pp(E_{W,K})\le Ce^{-cKW},
\end{equation}
and outside this event every such cell set contains at most \(CKW\)
distinct sources in its preimage. Indeed, the tail at threshold
\(CKW\) for a Poisson variable of mean \(O(W)\) is
\(\exp[-cK\log(K)W]\), which absorbs the cell-set count and root
polynomial for \(K\ge K_0\). No Poisson law for selected phase points
is being asserted.

Let \(u_i\) be a single edge's contribution. Projection of its path
onto \(v_i/|v_i|\) has derivative at least \(h|v_i|\), so it spends
at most \(2/(h|v_i|)\) time in \(B_1\). Consequently
\(u_i\le C\min(1,t|v_i|)\). On a good large scale the preceding
count bound gives
\(U_{A,I}^2\le CK\sum_i t|v_i|u_i\).
For bounded \(W_0\), the same bound is replaced by
\(CKt^2\sum_{x_i\in B_C,\ a_i\in A}|v_i|^2\), outside an event
\(E_{0,K}\) on which the fixed-ball source count exceeds \(CK\).

The following pointwise formulation avoids any measurable-choice issue
for the extension \(V\). Let \(N_W\) be the full source count in
\(B_{CW}\), increasing the fixed constants if necessary. The case
analysis just proved implies
\begin{align}
\label{eq:highd-measurable-traffic-majorant}
 U_{A,I}^2\le{}&CK\left(\sum_i t|v_i|u_i
       +t^2\!\sum_{x_i\in B_C,\ a_i\in A}|v_i|^2\right)
       +CN_1^2\mathbf1_{E_{0,K}}\notag\\
 &+C\sum_{W\in\{W_*2^j:j\ge0\}}
                  N_W^2\mathbf1_{E_{W,K}},
\end{align}
where \(W_*\ge1\) is fixed. Every quantity on the right is measurable
and independent of the choice of extension. On an exceptional scale,
the spatial confinement bounds all contributors by \(N_W\); otherwise
one of the two energy terms applies. Since
\(\E N_W^4\le CW^{4q}\), Cauchy--Schwarz and
\eqref{eq:highd-phase-animal-tail} bound the expected exception terms
by \(Ce^{-c'K}\) after summing dyadic \(W\). The bounded-scale term
has the same estimate.

Campbell's formula controls the second energy term in
\eqref{eq:highd-measurable-traffic-majorant}. For the first, dropping
the collection restriction and integrating the source anchor gives
\[
 \E\sum_i t|v_i|u_i
 \le\lambda t\int_I H\E^0[|v_e|^2\mathbf1_{\{a\in A\}}]\,dr
 \le C\lambda t^2\E^0[|v_e|^2\mathbf1_{\{a\in A\}}].
\]
The translated smoothing kernel integrates to one even though the
edge displacement depends on the entire input. This proves
\eqref{eq:highd-short-window-traffic}.

Choose \(0<t_*<\delta/8\). Split each path into its source layer
\(0<r-a<t_*\), target layer \(0<b-r<t_*\), and middle portion.
For the source layer put \(t_k=t_*2^{-k-1}\); its level \(k\) has
\(t_k<r-a\le2t_k\). Partition \(I_A\) into half-open mark bands
\(A_{k,j}\) of width at most \(t_k/8\), assigning each boundary atom
to one band. There are \(J_k\le C/t_k\) bands. For a band starting
at \(u\), the relevant times lie in
\([u+t_k,u+(2+1/8)t_k]\). Cover this interval by a fixed number of
windows of length at most \(2\varepsilon t_k\), enlarged by at most
\(t_k/8\). On each window, for every mark in the band,
\(ct_k\le r-a\le Ct_k\) and \(b-r\ge\delta/2\).
All band edges are active; extending their contribution to these full
windows only increases traffic. Set
\(e_{k,j}=\E^0[|v_e|^2\mathbf1_{\{a\in A_{k,j}\}}]\).
Minkowski and \eqref{eq:highd-short-window-traffic} give
\begin{equation}
\label{eq:highd-endpoint-traffic-sum}
 \|U_k\|_2
 \le C\sqrt{K_kt_k}\sqrt{\lambda\E^0|v_e|^2}
       +Ct_k^{-1}e^{-cK_k/2}.
\end{equation}
Here we used \(\sum_j e_{k,j}=\E^0|v_e|^2\) and
\(\sum_j\sqrt{e_{k,j}}\le\sqrt{J_k\sum_j e_{k,j}}\).
Taking \(K_k=K_0+C_1(k+1)\), with \(C_1\) sufficiently large,
makes both series in \eqref{eq:highd-endpoint-traffic-sum} summable.

Reverse time, replace \(T(r)\) by \(-T(-r)\), and exchange the
colors for the target layer. The new derivative remains positive
definite; the velocity is \(-v_e\), with the same Palm second moment
by mass transport. Finally cover the middle portion by finitely many
windows of size \(2\varepsilon t_*\). At center \(r_0\), select
edges with \(a\le r_0-(1-\varepsilon)t_*\) and
\(b\ge r_0+(1-\varepsilon)t_*\). They are active throughout the
window, both compliances are bounded below, and every middle path
contribution occurs in one such window. The fixed-\(t_*\) version of
\eqref{eq:highd-short-window-traffic}, with the entire source mark
interval, applies. Combining these bounds proves
\eqref{eq:highd-traffic-ltwo}. The argument uses no mark-density bound
or displacement moment above two.
\end{proof}

\begin{lemma}[Subvolume Neumann boundary energy]
\label{lem:highd-current-boundary}
Let \(J\) be a stationary smooth square-integrable vector field on
\(\R^q\), \(q\ge3\), whose invariant conditional mean is zero and
whose divergence is \(\eta*(\mu_S-\mu_T)\) for independent equal-intensity
Poisson processes. For \(Q_L=[-L/2,L/2]^q\), define
\[
 \mathcal B_L(J)=
 \sup_{\substack{\phi\in H^1(Q_L),\ \int_{Q_L}\phi=0\\
                         \int_{Q_L}|\nabla\phi|^2\le1}}
 \left|\int_{\partial Q_L}(J\cdot n)\phi\right|^2.
\]
Then \(\E\mathcal B_L(J)=o(L^q)\). No mixing assumption is required.
\end{lemma}
\begin{proof}
The squared norm of the sum of the \(2q\) face functionals is at most
\(2q\) times their squared norms. On a face normal to \(e_i\), let
\(\sigma_i\) be the scalar spectral measure of \(J_i\), including any
atom at zero. Trace and Poincar\'e inequalities give
\[
 \|\phi\|_{L^2(\partial Q_L)}
       \le C L^{1/2}\|\nabla\phi\|_{L^2(Q_L)}
       \quad\text{when }\int_{Q_L}\phi=0.
\]
The squared face-functional norm of the datum \(e^{ik\cdot x}\) is
therefore at most \(CL^q\). After scaling to \(Q_1\), its norm divided
by \(L^{q/2}\) tends to zero whenever \(k\) has a nonzero tangential
component: \(e^{iLk_{\rm tan}\cdot x}\) converges weakly to zero in
\(L^2\) of the unit face, and the trace map from mean-zero \(H^1(Q_1)\),
with gradient norm, to this \(L^2\) space is compact.

The exceptional normal axis has zero \(\sigma_i\)-mass. To see this,
let \(\Sigma\) be the matrix spectral measure of \(J\), with convention
\(\E[J(x)J(0)^T]=\int e^{ik\cdot x}\Sigma(dk)\) and
\(\widehat\eta(k)=\int e^{-ik\cdot x}\eta(x)\,dx\).
The exact Poisson divergence implies
\begin{equation}
\label{eq:highd-current-spectrum}
 \sum_{i,j}k_i k_j\Sigma_{ij}(dk)
       =\frac{2\lambda}{(2\pi)^q}|\widehat\eta(k)|^2\,dk.
\end{equation}
On \(\R e_i\setminus\{0\}\), the left side is \(k_i^2\sigma_i\),
whereas the right side charges no such axis. At zero the atom is the
squared \(L^2\) norm of the invariant conditional mean of \(J_i\),
which is zero. Spectral isometry for the Hilbert-valued face functional
now expresses its expected squared norm as the integral of the
plane-wave squared norms against \(\sigma_i\). This identity follows
also by expansion in an orthonormal basis of the mean-zero \(H^1\)
test space and Tonelli. The uniform \(CL^q\) bound and dominated
convergence prove the assertion for each face and then for their sum.
\end{proof}

\begin{proposition}[Signed density of the remaining endpoints]
\label{prop:highd-signed-ports}
Assume the hypotheses of Theorem~\ref{thm:highd-kinetic-traffic} and
suppose the source-Palm mean displacement is zero in each stationary
ergodic component. Choose deterministic \(\ell_L\to\infty\) with
\(\ell_L/L\to0\), and retain precisely the matching edges whose entire
kinetic paths lie in \(Q_{L-\ell_L}\). Among input points in \(Q_L\),
call the remaining unmatched endpoints the main ports.
There are nonnegative mass-preserving smoothings \(a_L,b_L\) of the
source and target ports, costing at most one unit of quadratic transport
per endpoint, such that
\begin{equation}
\label{eq:highd-port-neumann}
 \E\|a_L-b_L\|_{H_N^{-1}(Q_L)}^2=o(L^q).
\end{equation}
The expected number of ports is \(o(L^q)\). Here \(H_N^{-1}\) is the
dual norm of mean-zero \(H^1\) functions with gradient norm at most one;
its argument need not have zero total mass.
\end{proposition}
\begin{proof}
Theorem~\ref{thm:highd-kinetic-traffic} supplies \(\E U(0)^2<\infty\).
In particular \(J\) is square-integrable; convolution gives a smooth
version. The latter can also be checked by Campbell domination of all
kernel derivatives on compact sets. The independent Poisson input law
is ergodic, so its law and intensity remain unchanged in almost every
matching component. Componentwise Campbell identifies the invariant
conditional mean of \(J\) with the componentwise displacement mean
times this intensity. Thus Lemma~\ref{lem:highd-current-boundary} applies.

Let \(K_L\) be the retained edges' contribution to the current, and
let \(U_{>R}\) be the traffic of edges with \(|y-x|>R\).
The support of \(K_L\) is strictly inside \(Q_L\) when \(\ell_L>2\).
At points whose distance to the boundary exceeds
\(\ell_L/2+CR+1\), every contributing edge of displacement at most
\(R\) has its entire path in the retention cube. Hence
\begin{equation}
\label{eq:highd-bulk-current-error}
 \frac1{L^q}\E\int_{Q_L}|J-K_L|^2
 \le C\frac{\ell_L+R+1}{L}\E U(0)^2+\E U_{>R}(0)^2.
\end{equation}
For fixed \(R\), send \(L\) to infinity and then \(R\) to infinity.
The last term tends to zero by domination by \(U^2\); each edge has
finite displacement. Thus the left side tends to zero. The same path
diameter bound and Campbell imply that the expected port count divided
by \(L^q\) is at most
\(C[\ell_L/L+\E^0\min(1,C|y-x|/L)]\), which tends to zero.
The target estimate follows from mass transport.

Let \(\eta_L^N(z,x)\) be the coordinate-reflected kernel on \(Q_L\),
formed by summing \(\eta\) over the reflection images of \(x\).
It is symmetric, nonnegative, preserves mass and constant density, and
is supported where \(|z-x|\le1\). For endpoints farther than one
unit from the boundary it agrees with \(\eta(z-x)\). Use this kernel
to define \(a_L,b_L\). Define the full-input boundary discrepancy
\[
 Z_L=\int_{Q_L}\eta_L^N(\,\cdot,x)(\mu_S-\mu_T)(dx)
       -\bigl[\eta*(\mu_S-\mu_T)\bigr]\big|_{Q_L}.
\]
It satisfies
\begin{equation}
\label{eq:highd-reflected-input-error}
 \E\|Z_L\|_{H_N^{-1}(Q_L)}^2\le C\lambda L^{q-1}.
\end{equation}
Indeed, write \(Z_L\) as the Poisson integral of
\(\mathbf1_{Q_L}(x)\eta_L^N(\,\cdot,x)-\eta(\,\cdot-x)\).
Only inputs in the unit boundary layer, inside or outside, contribute.
Each such kernel difference has bounded \(L^1\) and \(L^\infty\)
norms and therefore bounded \(L^{2q/(q+2)}\) norm. The scale-independent
Sobolev--Poincar\'e inequality on the cube bounds its \(H_N^{-1}\) norm
uniformly. The Hilbert-space Poisson variance identity proves
\eqref{eq:highd-reflected-input-error}; the two equal-intensity colors
cancel its mean. This calculation concerns the full independent inputs.

Retained endpoints are in the interior, so for \(p_L=a_L-b_L\) the
exact identity on \(Q_L\) is
\(p_L=\operatorname{div}(J-K_L)+Z_L\).
For a mean-zero test \(\phi\), integration by parts gives
\[
 \int_{Q_L}p_L\phi
 =-\int_{Q_L}(J-K_L)\cdot\nabla\phi
       +\int_{\partial Q_L}(J\cdot n)\phi+\int_{Q_L}Z_L\phi.
\]
The squared triangle inequality, Lemma~\ref{lem:highd-current-boundary},
\eqref{eq:highd-bulk-current-error}, and
\eqref{eq:highd-reflected-input-error} prove
\eqref{eq:highd-port-neumann}. No independent-input assertion about
the selected ports is used.
\end{proof}

\begin{corollary}[Independent reserves screen the boundary]
\label{cor:highd-kinetic-screening}
In Proposition~\ref{prop:highd-signed-ports}, let \(\lambda=1-\rho\),
where \(0<\rho<1\), and add independent source and target reserve
Poisson processes of intensity \(\rho\), with the same respective mark
laws. There is a maximum-cardinality
matching of the full inputs restricted to \(Q_L\) that preserves every
retained main edge and whose additional squared-distance cost satisfies
\begin{equation}
\label{eq:highd-kinetic-screening}
 \E[\text{additional quadratic cost}]
       \le C_q\rho^{1-2/q}L^q+o_{\rho,M}(L^q).
\end{equation}
The leading constant is independent of \(\rho\) and \(M\). Any
additional all-pair marked cost bounded by \(\Lambda|y-x|^2\) is bounded
by \(\Lambda\) times the right side.
\end{corollary}
\begin{proof}
Put \(V=L^q\). Let \(H_S,H_T\) be the reserve counts, \(N_S,N_T\)
the full counts, and \(\Delta=N_S-N_T\). On the event
\[
 G_L=\{\min(H_S,H_T)\ge\rho V/2,\ |\Delta|\le\rho V/4\},
\]
discard \(|\Delta|\) reserve vertices of the excess color, using only
auxiliary iid sample indices. The remaining reserve counts \(h_S,h_T\)
are at least \(\rho V/4\), and the two residual atomic clouds have
equal mass. Conditional on the main data and all relevant counts, each
remaining reserve cloud separately consists of iid uniform points;
independence between the two selected clouds is unnecessary. Poisson
Chernoff bounds give
\(\Pp(G_L^c)\le C\exp(-c_\rho V)\).

Transport the reserve clouds to their constant densities
\(h_S/V,h_T/V\). The sharp uniform matching estimate
\cite[Theorem~1.1]{GoldmanTrevisan2021}, followed by Jensen against an
independent comparison cloud, bounds each conditional unnormalized
cost by \(C_qL^2h^{1-2/q}\). Concavity and \(h_S\le H_S,h_T\le H_T\)
therefore give
\begin{equation}
\label{eq:highd-reserve-cost}
 \E[\mathbf1_{G_L}(D_S+D_T)]\le C_q\rho^{1-2/q}V.
\end{equation}
The smoothed residual densities
\(f=h_S/V+a_L\) and \(g=h_T/V+b_L\) have equal mass and are at
least \(\rho/4\). Solving the mean-zero Neumann problem and using
the linear interpolation between these densities in the
Benamou--Brenier formula yields
\[
 W_2^2(f\,dz,g\,dz)
 \le\frac4\rho\|p_L\|_{H_N^{-1}(Q_L)}^2.
\]
Here \(W_2^2\) is total quadratic cost for equal finite masses.
Indeed, the constant difference \((h_S-h_T)/V\) vanishes on mean-zero
tests, the Neumann field has squared norm \(\|p_L\|_{H_N^{-1}}^2\),
and the interpolated density stays above \(\rho/4\); this is the
construction in \cite[proof of Lemma~3.4]{GoldmanTrevisan2021}.
Spreading the ports costs at most their number. The three-leg squared
triangle inequality, \eqref{eq:highd-port-neumann}, and
\eqref{eq:highd-reserve-cost} bound the expected repair cost on \(G_L\)
by the right side of \eqref{eq:highd-kinetic-screening}. Equal integer
atomic masses admit a minimizing permutation, so this transport
construction gives an actual residual matching.

On \(G_L^c\), keep the same retained edges and match residual points
arbitrarily up to maximal cardinality. The cost is at most
\(qL^2(N_S+N_T)\). Its expectation on this event is \(o(V)\) by
Cauchy--Schwarz, the Poisson second moment, and the exponential tail.
The construction leaves precisely the global excess color count
unmatched and proves the corollary.
\end{proof}

The result uses the ordered common-cumulative-matrix structure to
derive square-integrable absolute traffic from two-cycle optimality.
It does not assert this for arbitrary pairwise elliptic matrices or
interleaving supports with forbidden pairs. It supplies a boundary
comparison for centered stationary kinetic laws; paired-mark uniqueness
is a separate question.

\subsection{Spatial probes and an upper comparison}
\label{subsec:highd-localized-pressure}

Keep uniform product sampling and the homogeneous kinetic cost of
Proposition~\ref{prop:highd-finite-pressure}. For a continuous
\(\Psi:D\times K_A\times K_B\to\R\), put
\[
F_n^\Psi(t)=\frac1n\min_\sigma\sum_i
 \bigl[n^{2/q}k_{A_iB_{\sigma(i)}}(W_{\sigma(i)}-Z_i)
       +t\Psi(Z_i,A_i,B_{\sigma(i)})\bigr].
\]
The source position in the probe retains spatial information about
the paired marks. Write \(p_0\) for the common value \(p_g(0)\).

\begin{proposition}[Localized upper pressure and concentration]
\label{prop:highd-localized-upper}
For every continuous \(\Psi\) and every fixed \(t\),
\begin{equation}
\limsup_n\E F_n^\Psi(t)
 \le\int_D p_{\Psi(z,\cdot,\cdot)}(t)\,dz.
\label{eq:highd-localized-upper}
\end{equation}
If \(\Psi(z,a,b)=\phi(z)g(a,b)\), with \(\phi\) Lipschitz
and \(g\) continuous, the centered values converge to zero in
probability locally uniformly in \(t\), with the pointwise bound
\[
\sup_{|t|\le T}\E|F_n^\Psi(t)-\E F_n^\Psi(t)|
 \le C_{T,\phi,g}n^{-1/2+1/q}.
\]
\end{proposition}

\begin{proof}
First freeze the spatial probe to \(g_j\) on each cell \(D_j\)
of a fixed finite cubical partition. Use independent Poisson inputs
of mean \(n\), split into main intensity \((1-\rho)n\) and
independent reserve intensity \(\rho n\). Optimize the main
inputs separately in each cell. Their residual total counts have
every \(r\)-th moment bounded by \(C_r n^{r/2}\): each cell
leaves precisely the absolute difference of its two Poisson counts.
The residual locations may depend on the cell's optimizer.

Apply \cite[Proposition~6.3]{GoldmanTrevisan2023} to these residuals
and the independent reserves, with \(h=\rho n\), \(H=\sqrt n\)
and power two. Its proof gives exponents \(1+2/q<2\) and one,
so their squared repair cost \(M_2\) obeys
\[
n^{2/q-1}\E M_2
 \le C\rho^{1-2/q}+C\rho^{-2/q}n^{-1/2+1/q}.
\]
Constants may depend on the fixed partition. This proposition
allows different cell probe laws; identical residual laws are not
required. Ellipticity bounds the marked repair. Its probe cost is
at most \(|t|\|\Psi\|_\infty\) times its number of edges,
whose expectation is at most \(\rho n+C\sqrt n\).
The main contribution divided by \(n\) tends to
\[
\sum_j |D_j|(1-\rho)^{1-2/q}
 p_{g_j}\bigl(t(1-\rho)^{2/q}\bigr).
\]
First take \(n\to\infty\), then \(\rho\downarrow0\).
The bounded-additive count comparison used in
Proposition~\ref{prop:highd-finite-pressure} applies to position-
dependent probes as well: all kernels remain fixed on the same cube
during the comparison. It gives the fixed-count conclusion.
Uniform approximation by cellwise spatial probes finishes the proof
of \eqref{eq:highd-localized-upper}, since both the finite values
and the limiting pressures change by at most
\(|t|\|\Psi-\widetilde\Psi\|_\infty\).

For concentration use the positive shift and square-root proof of
Proposition~\ref{prop:highd-value-concentration}. The extra spatial
gradient of the probe contributes at most \(C_{T,\phi,g}n\)
to the squared gradient of each permutation's total cost. This is
absorbed in its previous \(CL^2V\) bound, where \(L=n^{1/q}\)
and the shifted optimum satisfies \(V\ge n\). The one-row or
one-column mark replacement estimate is unchanged. Thus
\(\operatorname{Var}\sqrt V\le C_{T,\phi,g}L^2\), giving
the asserted rate and, by a parameter mesh, local uniform convergence
in probability.
\end{proof}

The compact dual class and two-cost comparison in
Subsection~\ref{subsec:highd-dual-transfer} also allow a continuous
\(\Psi(z,a,b)\): the perturbation is bounded, has a uniform
continuity modulus, and cancels in the comparison. Consequently
their actual-length transfer remains valid for these spatial probes
in the regular product tube.

For an arbitrary pairwise elliptic compliance, the upper estimate
alone does not give a spatial lower bound. For the common-cumulative
compliance supplied by the tube geometry,
Subsection~\ref{subsec:highd-stationary-pressure} will prove the
matching lower estimate through a stationary variational formula.
The concentration and actual-length transfer just established then
retain the spatial paired-mark information required for a full plan.

\subsection{Stationary variational pressure and full spatial localization}
\label{subsec:highd-stationary-pressure}

The screening argument also applies to bounded continuous mark
perturbations. This extension identifies the pressure with a stationary
variational problem and supplies the spatial lower comparison.
Throughout this subsection \(q\ge3\), transverse sampling is uniform
on \(D=[0,1]^q\), and the independent product mark laws are supported
on compact ordered separated intervals. The compliance has the common-cumulative
form \(S(a,b)=T(b)-T(a)\) used in the preceding screening result.

\begin{lemma}[Traffic under bounded mark perturbations]
\label{lem:highd-perturbed-traffic}
For a stationary perfect matching that is two-cycle optimal for
\(k+s g(a,b)\), where \(g\) is bounded and \(|s|\le S_0\),
finite Palm squared displacement implies square-integrable absolute
kinetic-path traffic. Its second moment is bounded by
\(C_{S_0,g}(1+\lambda\E^{\rm Palm}|v|^2)\).
Consequently the reserve-screening estimate holds if the invariant
conditional mean displacement is zero.
\end{lemma}

\begin{proof}
A two-edge swap now bounds the excess kinetic cost by
\(4S_0\|g\|_\infty\). The interface midpoint and minimization
argument, with source compliance comparable to \(t\), gives
\[
|v_i-v_j|\le \frac C t|z_i-z_j|
                  +C\sqrt{S_0\|g\|_\infty/t}.
\]
Put \(L=C/t\), \(B=C\sqrt{S_0\|g\|_\infty/t}\), and,
for each coordinate \(\ell\), define the measurable function
\[
V_\ell(z)=\inf_i\{v_{i,\ell}+L|z-z_i|\}.
\]
For a nonempty edge collection the displayed coarse Lipschitz bound
makes this infimum finite everywhere. Each coordinate is
\(L\)-Lipschitz, and
\(v_{i,\ell}-B\le V_\ell(z_i)\le v_{i,\ell}\).
For an empty collection set \(V=0\). Thus
\(\operatorname{Lip}(V)\le\sqrt q L\) and
\(|v_i-V(z_i)|\le\sqrt q B\), even when interface points coincide.

On a window \(|r-r_0|\le\epsilon t\), choose \(\epsilon\)
so that \(F_r(z)=z+[T(r)-T(r_0)]V(z)\) has inverse Lipschitz
constant at most two. The inverse trajectory
\(\gamma(r)=F_r^{-1}(0)\) and its phase curve
\((\gamma(r),tV(\gamma(r)))\) obey the previous length and
position bounds in terms of \(W_0=t|V(0)|\). The actual path has
\[
|X_i(r)-F_r(z_i)|\le C\epsilon\sqrt{S_0\|g\|_\infty t}.
\]
Every path meeting the unit ball therefore has
\(|z_i-\gamma(r)|+|tv_i-tV(\gamma(r))|\le C_{S_0,g}\).
This is a fixed-radius phase tube, uniformly over the endpoint
scales \(0<t\le t_*\). Large \(W_0\) still makes
\(t|v_i|\) comparable to \(W_0\); bounded \(W_0\) confines
all contributing sources to a fixed ball.

The source preimage of a phase cell remains
\(x=z-[T(r_0)-T(a)]v\), with the same bounded marked intensity.
The animal count, bounded contribution of one trajectory, and
Campbell energy estimate consequently give
\[
\E U_{A,I}^2\le C_{S_0,g}Kt^2\lambda
             \E^{\rm Palm}[|v|^2\mathbf1_{\{a\in A\}}]
                       +C_{S_0,g}e^{-c_{S_0,g}K}.
\]
The dyadic endpoint summation is unchanged, with a sufficiently
large linear choice of \(K_k\). The paths, divergence, and
retention rule are identical, so the centered screening conclusion
follows as before.
\end{proof}

Let \(\mathcal M\) consist of stationary perfect integer matching
laws of the prescribed independent unit-intensity marked Poisson
inputs, with finite kinetic energy. Define
\[
I_g(s)=\inf_{P\in\mathcal M}
               \E_P^{\rm Palm}[k+s g(a,b)].
\]
Here \(g\in C(K_A\times K_B)\). No paired-mark law is fixed
in this infimum.

\begin{proposition}[Stationary variational formula]
\label{prop:highd-stationary-pressure}
The infimum \(I_g(s)\) is attained and equals the homogeneous
pressure \(p_g(s)\) of Proposition~\ref{prop:highd-finite-pressure}.
Every subsequential random conditional scene law of finite
\(k+s g\)-optimizers is almost surely mean-minimizing at almost
every macroscopic root. The same assertion holds for actual costs \(n^{2/q}e+s g\) under the
homogeneous-compliance hypotheses
of Corollary~\ref{cor:highd-length-pressure}.
\end{proposition}

\begin{proof}
The unmarked quadratic bound in \(q\ge3\), followed by the
two-endpoint local compactness argument of
Proposition~\ref{prop:highd-local-limit}, gives a stationary perfect
matching with finite second moment. Independent marks and ellipticity
make \(\mathcal M\) nonempty. A minimizing sequence for \(I_g(s)\)
has uniformly bounded kinetic energy since \(g\) is bounded.
Local Poisson count bounds make its joint input and integer edge
measures vaguely tight. Both source- and target-anchored escape
probabilities are \(O(R^{-2})\), the latter by mass transport.
As in Proposition~\ref{prop:highd-local-limit}, compact truncation
and uniform integrability of local counts retain both full endpoint
projections. The limiting law is stationary and perfect, with the
prescribed input law. Lower semicontinuity of the nonnegative cost
\(k+s g+|s|\|g\|_\infty\) proves attainment.

A minimizer is two-cycle optimal. Otherwise some bounded box has,
with positive probability, two complete edges admitting a swap
with gain at least \(\varepsilon>0\). Tile space by sufficiently
large boxes with an independent uniform translation. In each box
perform one measurably selected such improvement when available.
The swaps are disjoint, preserve stationarity and finite energy,
and strictly lower the mean cost, a contradiction.

For a bounded measurable vector function \(h(a)\), translate each
source to \(x+\varepsilon h(a)\) and keep its assigned target.
This preserves the prescribed product marked Poisson law and
stationarity. Differentiating the exact quadratic competitor cost
at zero gives \(\E^{\rm Palm}[v\cdot h(a)]=0\).
The mark perturbation is unchanged. Target translations give the
analogous identity conditioned on \(b\); mass transport converts
between the two Palm anchors. Hence
\[
\E^{\rm Palm}[v\mid a]=\E^{\rm Palm}[v\mid b]=0,
\qquad
\E^{\rm Palm}(y-x)
 =\E^{\rm Palm}[(T(b)-T(a))v]=0.
\]
The input Poisson law is ergodic, so every ergodic component of a
stationary minimizer retains that input law. Each component must
also minimize mean cost. Applying these translations within each
component gives invariant conditional zero displacement.

Fix \(\rho\in(0,1/2)\) and \(\lambda=1-\rho\).
Scale an intensity-one minimizer for
\(k+s\lambda^{2/q}g\) by \(\lambda^{-1/q}\).
It minimizes \(k+s g\) at intensity \(\lambda\), with cost
density
\[
\lambda^{1-2/q}I_g(s\lambda^{2/q}).
\]
It satisfies Lemma~\ref{lem:highd-perturbed-traffic} and the
invariant centering just proved. Retain its paths in an inner
cube and complete with independent reserves of intensity \(\rho\).
The kinetic repair costs at most
\(C_{s,g}\rho^{1-2/q}L^q+o(L^q)\).
Only \(\rho L^q+o(L^q)\) repair edges occur in expectation, so
their mark contribution is at most
\(|s|\|g\|_\infty(\rho L^q+o(L^q))\).
Dropping \(o(L^q)\) main edges increases the retained signed cost
by at most \(|s|\|g\|_\infty o(L^q)\); this follows by adding
that constant to each edge before comparing retained and full costs.
For the independent-Poisson optimum \(C_L^P(s)\), therefore,
\[
\limsup_{L\to\infty}\frac{\E C_L^P(s)}{L^q}
 \le\lambda^{1-2/q}I_g(s\lambda^{2/q})
          +C_{s,g}\rho^{1-2/q}+|s|\|g\|_\infty\rho.
\]
Use the bounded-perturbation independent-count comparison and let
\(\rho\downarrow0\). Since \(I_g\) is
\(\|g\|_\infty\)-Lipschitz, this gives \(p_g(s)\le I_g(s)\).

Conversely, any subsequential random scene limit \(\Gamma\) of
finite optimizers has conditional laws \(P_z^\Gamma\) in
\(\mathcal M\) almost everywhere. Nonnegative compact truncations
of \(k+s g+|s|\|g\|_\infty\) give
\[
I_g(s)\le
\E\int_D\E_{P_z^\Gamma}^{\rm Palm}[k+s g]\,dz
\le p_g(s).
\]
Equality proves the identity and forces each conditional energy to
equal \(I_g(s)\) almost surely for almost every \(z\).
For actual length optimizers, the uniform Taylor expansion gives
the same truncated lower bound, while
Proposition~\ref{prop:highd-quadratic-transfer} identifies their
expected normalized values with \(p_g(s)\). The identical equality
argument applies. No uniform integrability of squared edge lengths
is assumed.
\end{proof}

\begin{corollary}[Spatial pressure and full-plan selection]
\label{cor:highd-spatial-pressure}
For every continuous \(\Psi:D\times K_A\times K_B\to\R\),
\[
\lim_n\E F_n^\Psi(s)
 =\int_D p_{\Psi(z,\cdot,\cdot)}(s)\,dz.
\]
For \(\Psi(z,a,b)=\phi(z)g(a,b)\), with \(\phi\) Lipschitz,
this convergence also holds in probability, locally uniformly in
\(s\). Every unperturbed empirical optimal coupling
converges in probability to one deterministic coordinate plan
\(dz\,\kappa(da,db)\), with \(\kappa(g)=p_g'(0)\).
This implication also holds for the genuine homogeneous-compliance
length model. The required pressure differentiability is
Theorem~\ref{thm:highd-pressure-differentiability}.
\end{corollary}

\begin{proof}
The upper bound is Proposition~\ref{prop:highd-localized-upper}.
For the lower bound, take a subsequence attaining the lower limit
and its random scene limit. Boundedness of the probe retains the
same energy and endpoint bounds. On each fixed microscopic window,
the probe's source position tends to the retained macroscopic root
\(z\). Continuity and nonnegative compact truncation give
\[
\begin{split}
\liminf_n\E F_n^\Psi(s)
&\ge\E\int_D\E_{P_z^\Gamma}^{\rm Palm}
                        [k+s\Psi(z,a,b)]\,dz\\
&\ge\int_D I_{\Psi(z,\cdot,\cdot)}(s)\,dz
 =\int_D p_{\Psi(z,\cdot,\cdot)}(s)\,dz.
\end{split}
\]
These deterministic integrands are measurable by uniform-norm
Lipschitz continuity of the pressure in its probe. For product
probes, concentration and parameter meshes give the claimed
probabilistic convergence.

By Theorem~\ref{thm:highd-pressure-differentiability}, all \(p_g\)
are differentiable at zero. Since
\(p_{\phi(z)g}(s)=p_g(s\phi(z))\), bounded slopes allow
differentiation under the integral, giving derivative
\(p_g'(0)\int_D\phi\). The two finite-difference inequalities
bound the paired test of every unperturbed optimizer between the
same random pressure secants. They converge to this derivative.
A countable dense family of products \(\phi g\) determines the
coordinate plan. Compactness supplies its probability law
\(\kappa\); a worst-test optimizer and convexity give uniformity
over all empirical optimal couplings. The spatial-probe form of
Proposition~\ref{prop:highd-quadratic-transfer} transfers the pressure
and the same sandwich to the actual length defect.
\end{proof}

The diffuse common-environment comparison in
Theorem~\ref{thm:highd-pressure-differentiability} completes this
homogeneous selection argument. Stationary mean minimality supplies
the pressure localization; the separate finite Dirichlet identity and
local fractional swaps supply differentiability.

\subsection{Finite empirical stability for the actual curved distance}
\label{subsec:highd-curved-empirical-stability}

We give a finite-sample comparison for the actual distance, allowing
varying geometric coefficients and conditional mark densities. It controls
unrestricted finite reoptimization and will support the cavity argument.

Use transverse coordinates \(D=[0,1]^q\), after an affine change on
a sufficiently thin convex calibrated tube, with \(q\ge3\) and
compact ordered radial intervals \(I_A,I_B\). Write
\[
 e(a,z;b,w)=d_F(X(a,z),X(b,w))-(b-a).
\]
Assume this defect is \(C^2\) in its transverse arguments on all
cross pairs, with continuous derivatives on the compact parameter
set, and constants \(\gamma,\Gamma,K,m>0\) such that
\begin{gather}
 \gamma|w-z|^2\le e(a,z;b,w)\le\Gamma|w-z|^2,\qquad
 \partial_z e(a,z;b,z)=0, \notag\\
 \|\partial_{zz}^2 e\|,\ \|\partial_{zw}^2 e\|\le K,\qquad
 \operatorname{sym}(-\partial_{zw}^2 e)\succeq mI.
 \label{eq:highd-curved-stability-hypotheses}
\end{gather}
Here \(\operatorname{sym}(A)=(A+A^\top)/2\).
Lemma~\ref{lem:raytube-distance} supplies these conditions after
shrinking the transverse cube; their geometric verification is
given below.

Let the two independent iid sampling laws be
\[
 \rho_A(a,z)\,\alpha_0(da)\,dz,\qquad
 \rho_B(b,w)\,\beta_0(db)\,dw,
\]
where \(\alpha_0,\beta_0\) are fixed probability laws on the radial
intervals, the densities are positive and jointly continuous, and
both transverse marginals equal a positive Lipschitz density
\(\theta\). Product radial laws, including atomic ones, are allowed.
Write \(L_n=n^{1/q}\), \(X_i=L_nZ_i\), \(Y_j=L_nW_j\), and
\[
 c_n((x,a),(y,b))=L_n^2e(a,x/L_n;b,y/L_n).
\]
Minimizing this cost gives exactly the original genuine-distance
optimizers: the subtracted radial sum is independent of the
permutation, and \(L_n^2>0\).

\begin{proposition}[A finite empirical margin for genuine length]
\label{prop:highd-curved-empirical-stability}
The finite optimal permutation \(M_n\) is almost surely unique.
Choose a source index independently and uniformly, write \(x,a\)
for its microscopic location and mark, and let \(j_0=M_n(x)\).
For every target \(j\), let \(B_{n,j}\) be the minimum excess total
cost among permutations assigning this source to \(j\). Put
\[
 v_{n,j}=-\partial_x c_n((x,a),(Y_j,B_j)),\qquad
 b_n=\min_{j\ne j_0}
       \frac{B_{n,j}}{1+|v_{n,j}-v_{n,j_0}|}.
\]
For \(n=1\), set this empty minimum to \(+\infty\).
Then
\begin{equation}
 \lim_{\varepsilon\downarrow0}\limsup_{n\to\infty}
       \Pp\{b_n\le\varepsilon\}=0.
 \label{eq:highd-curved-empirical-margin}
\end{equation}
For every deterministic \(r_n\downarrow0\), with probability tending
to one, moving this one source by any microscopic vector
\(|h|\le r_n\) that keeps it in \([0,L_n]^q\), with its mark fixed,
leaves the entire original permutation uniquely optimal. Thus the
physical displacement is \(o(n^{-1/q})\). The expected fraction of
source indices failing this individual stability property tends
to zero.
\end{proposition}

\begin{proof}
Comparison with the unmarked quadratic assignment, using the
common transverse marginal and the \(q\ge3\) matching bound from
Proposition~\ref{prop:highd-energy}, gives expected total cost
at most \(Cn\). Coercivity and uniform random rooting therefore
bound the selected squared displacement uniformly. By integrating
\(\partial_{zw}^2 e\) from the diagonal,
\(|v_{n,j}|\le K|Y_j-x|\), so
\begin{equation}
 \sup_n\E|v_{n,j_0}|^2=:C_E<\infty.
 \label{eq:highd-curved-root-gradient-moment}
\end{equation}

Conditional on \(x,a,b\), the map
\[
 y\longmapsto-L_n\partial_z e(a,x/L_n;b,y/L_n)
\]
has derivative \(-\partial_{zw}^2e\). Its symmetric part is at
least \(mI\), making it injective on the convex cube and giving
absolute Jacobian determinant at least \(m^q\). The bounded
target density and change of variables thus bound the conditional
velocity density, after averaging over \(b\), by \(C_v/n\).
Conditional on \(x,a\), all target velocities are independent.
If \(s_n=\min_{j\ne j_0}|v_{n,j}-v_{n,j_0}|\), counting all
ordered target pairs gives
\begin{equation}
 \Pp\{s_n\le\delta\}
 \le C_E/V^2+C_v^2|B_1(0)|^2 V^q\delta^q .
 \label{eq:highd-curved-gradient-separation}
\end{equation}
In particular this probability tends to zero uniformly in \(n\)
as first \(\delta\downarrow0\), then \(V\uparrow\infty\).
No conditioning on the chosen partner was used in the pair count.

We also obtain finite uniqueness. Compare two permutations that
give some source different targets, conditioning on every other
source. Their cost difference is a \(C^1\) function of this
source's position. On its zero level set its gradient vanishes
at Lebesgue-almost every point of that set. Equality of the two
candidate gradients has probability zero by the preceding target
density bound and Fubini. The cost tie therefore has probability
zero. A finite union over permutation pairs proves uniqueness.

For the comparison, fix \(\sigma>0\) and take \(n\) large enough
that \(\sigma<L_n/2\). Let \(U\) be uniform on
\([-1/\sqrt q,1/\sqrt q]^q\), independently of the samples.
Propose \(x+\sigma U\), reflecting each coordinate at the endpoints
of \([0,L_n]\), and call the proposal \(\widetilde x\). Accept it
with probability
\[
 \min\left(1,\frac{\rho_A(a,\widetilde x/L_n)}
                       {\rho_A(a,x/L_n)}\right);
\]
otherwise retain \(x\). The reflected proposal kernel is reversible
for Lebesgue measure: in one coordinate each fixed step magnitude
uses its two signs equally, and reflection pairs their forward and
reverse moves; the product kernel has the same property.
The displayed acceptance rule consequently preserves the
conditional spatial law given \(a\). The complete iid input law
is unchanged, including its conditional marks.

Write \(\Delta\) for the accepted displacement and
\(G_n(\Delta)\ge0\) for the difference between retaining the original
permutation and globally reoptimizing after the edit. Equality of
the input laws implies that its expectation equals the expected
retained root-row increment. The microscopic root Hessian is
\(\partial_{zz}^2e\), hence is bounded by \(K\). Define
\[
 \omega(t)=\frac{1}{\inf\rho_A}
   \sup_{\substack{a\in I_A\\|z-z'|\le t}}
       |\rho_A(a,z)-\rho_A(a,z')|\longrightarrow0.
\]
The rejection probability is at most \(\omega(\sigma/L_n)\).
Outside the width-\(\sigma\) boundary layer, the unreflected
proposal is centered, so the conditional mean accepted displacement
has norm at most \(\sigma\omega(\sigma/L_n)\).
The boundary-layer probability is at most
\(2q\|\theta\|_\infty\sigma/L_n\), and always
\(|\Delta|\le\sigma\). Taylor's theorem and Cauchy--Schwarz yield
\begin{equation}
 \E G_n(\Delta)\le
 \frac K2\sigma^2+\sqrt{C_E}\,\sigma
 \left[\omega(\sigma/L_n)
       +\sqrt{2q\|\theta\|_\infty\sigma/L_n}\right].
 \label{eq:highd-curved-global-repair}
\end{equation}
This controls unrestricted finite reoptimization, without bounding
the number of changed partners.

Moving the root changes every permutation assigning it to \(j\)
by the same root-row increment. Thus its best possible gain has
the envelope
\[
 G_n(h)=\max_j
   [h\cdot(v_{n,j}-v_{n,j_0})-B_{n,j}+R_{n,j}(h)]_+,
 \qquad |R_{n,j}(h)|\le K|h|^2 .
\]
On \(\{b_n\le\varepsilon,\ s_n\ge\delta\}\), select one
qualifying target before sampling \(U\), and put
\(d=v_{n,j}-v_{n,j_0}\).
Then \(B_{n,j}\le\varepsilon(1+\delta^{-1})|d|\).
Let \(c_q=1/(2\sqrt q)\), \(\beta_q=4^{-q}\), and set
\(\sigma=4\varepsilon(1+\delta^{-1})/c_q\).
With probability at least \(\beta_q\),
\(U\cdot d\ge c_q|d|\): choose each coordinate's sign to agree
with \(d\) and its magnitude in
\([1/(2\sqrt q),1/\sqrt q]\).
For \(K\sigma\le c_q\delta/4\), off the boundary layer and
upon acceptance this gives gain at least
\(\sigma c_q\delta/2\). Acceptance has conditional probability
at least \(1-\omega(\sigma/L_n)\).
Combining with \eqref{eq:highd-curved-global-repair}, and taking
\(n\to\infty\) at fixed \(\varepsilon,\delta\), gives
\[
 \limsup_n\Pp\{b_n\le\varepsilon,\ s_n\ge\delta\}
 \le \frac{4K\varepsilon(1+\delta^{-1})}
            {\beta_q c_q^2\delta}.
\]
First let \(\varepsilon\downarrow0\), then
\(\delta\downarrow0\), using
\eqref{eq:highd-curved-gradient-separation}. This proves
\eqref{eq:highd-curved-empirical-margin}.

If \(b_n>2r_n+2Kr_n^2\), every root-changing competitor has
strictly positive excess after every valid \(|h|\le r_n\), by
the same envelope bound. Root-preserving competitors retain
their original positive excess. Equation
\eqref{eq:highd-curved-empirical-margin} proves the asserted
stability. Uniform random rooting identifies its failure
probability with the expected fraction of failing source indices.
This last assertion tests one source at a time.
\end{proof}

To verify the geometric hypotheses, at coincident transverse
labels Lemma~\ref{lem:raytube-distance} gives
\[
 \partial_{ww}^2 e=S(a,b,z)^{-1},\qquad
 \partial_{zw}^2 e=-S(a,b,z)^{-1},\qquad \partial_z e=0.
\]
Compactness of the separated radial intervals and continuity of
these derivatives give
\eqref{eq:highd-curved-stability-hypotheses} on every sufficiently
small transverse cube, with quadratic coercivity obtained by
integrating \(\partial_{ww}^2e\) from \(w=z\).
These are properties of the smooth branch of the \emph{global}
distance established in that lemma.

The reflected acceptance step is solely an input-preserving
comparison in the proof; it does not modify the empirical objective
or prescribe a new selection rule. The result does not bound the
size of a changing cycle under a nonvanishing displacement or an
independent sample replacement.


\subsection{A finite cavity tree and paired response}
\label{subsec:highd-cavity-tree}

Deleting one source gives a common finite structure for every possible
target vacancy. The assignment tree is classical; see
\cite[proof of Theorem~5.3 and Corollary~5.4]{RichterGebertSturmfelsTheobald}
and \cite[Corollary~3.6 and Theorem~3.8]{CaraccioloSicuroSportiello}.
We record its consequence for the paired observables. The statement
applies to actual distance matrices without a quadratic approximation.

\begin{proposition}[The one-hole target tree]
\label{prop:highd-cavity-tree}
Let \(C=(c_{ij})\) be a finite real \((n-1)\times n\) matrix,
\(n\ge2\). Assume that, for every target \(j\), the assignment
leaving exactly \(j\) vacant has a unique optimizer \(M^j\), of
cost \(R_j\). The potentials
\[
 \beta_j=-R_j,\qquad \alpha_i=\min_j(c_{ij}+R_j)
\]
satisfy \(\alpha_i+\beta_j\le c_{ij}\) and calibrate every
\(M^j\). Their entire bipartite tight graph is a tree, with every
source vertex of degree two. Contracting its source vertices gives
a tree on the targets, whose edges are labeled by the sources.
Orienting this tree away from target \(j\), and assigning each
edge's source to its child target, gives exactly \(M^j\).
\end{proposition}

\begin{proof}
Shift column \(j\) by \(R_j\), setting \(d_{ij}=c_{ij}+R_j\).
Every vacancy optimizer then has cost \(\sum_jR_j\), so all are
optimal for the rectangular assignment allowing any vacant target.
Its column dual multipliers \(\gamma_j\) satisfy \(\gamma_j\le0\).
Complementary slackness with an optimizer leaving \(j\) vacant
forces \(\gamma_j=0\). Thus all these multipliers vanish, the
optimal row potentials are \(\alpha_i=\min_jd_{ij}\), and
\(\sum_i\alpha_i=\sum_jR_j\). Returning to \(C\) gives the stated
potentials and the exact lower bound \(R_j\) for each fixed vacancy.
Every \(M^j\) consequently uses only tight edges.

Fix a reference vacancy \(r\). For each nonmatching tight edge
\((i,u)\), put an arc \(u\to v\), where \(v=M^r(i)\).
A directed cycle would flip a zero-cost alternating cycle of
\(M^r\), contradicting its uniqueness. Every target \(j\) is
reachable from \(r\) along the vacancy path in the symmetric
difference of \(M^r\) and \(M^j\). Every directed \(r\)-to-\(j\)
path gives a tight vacancy-\(j\) assignment, so uniqueness gives
exactly one such path. A reachable acyclic directed graph with
unique root paths is an arborescence. Each \(v\ne r\) therefore
has exactly one incoming arc, and its matched source has exactly
one nonmatching tight neighbor. Adding the reference matched edges
gives the asserted subdivided target tree. Rerooting it constructs
the tight assignment for any other vacancy.
\end{proof}

The finite uniqueness argument of
Proposition~\ref{prop:highd-curved-empirical-stability} applies to
each deleted-source, deleted-target problem, so the tree exists
almost surely in that genuine iid setting. Write \(U_i=(x_i,a_i)\)
and \(V_j=(y_j,b_j)\) for complete endpoint data. For a fixed bounded
pair test \(g(U,V)\), put \(G=\|g\|_\infty\) and define
\[
 Q_j=\sum_{i=1}^{n-1}g(U_i,V_{M^j(i)}).
\]
If targets \(u,v\) are adjacent through source \(i\), changing
the vacancy from \(u\) to \(v\) changes only that source's partner.
Thus \(Q_v-Q_u=g(U_i,V_u)-g(U_i,V_v)\); all cavity-score
differences are sums along the unique target-tree path. Equivalently,
\(Q_j=R_j'(0)\) when the finite remainder cost is perturbed by
\(t g\), since its optimizer is unique at zero.

Let \(W\) denote the punctured input and insert an independent
source \(s=(x,a)\). Its target is
\(J(s)=\arg\min_j[c(s,V_j)+R_j]\), almost surely unique in the
iid setting; for arbitrary insertion laws fix the smallest-index
tie rule. The full paired average satisfies
\begin{equation}
 \begin{aligned}
 nT_n(s,W)&=g(s,V_{J(s)})+Q_{J(s)},\\
 n\{T_n(s',W)-T_n(s,W)\}
   &=g(s',V_{J(s')})-g(s,V_{J(s)})+Q_{J(s')}-Q_{J(s)}.
 \end{aligned}
 \label{eq:highd-cavity-tree-response}
\end{equation}
Both insertions use the same tree. Their remainder matchings differ
exactly along the path between \(J(s)\) and \(J(s')\).
The test is fixed under resampling; in particular this includes fixed
spatial probes of the full plan, without root-dependent recentering.

Condition on \(W\), and put \(p_j=\Pp\{J(s)=j\mid W\}\).
For a reference tree root, write \(C_e\) for an edge's descendant
target set and \(\Delta_e=g(U_i,V_u)-g(U_i,V_v)\) for its increment
from parent \(u\) to child \(v\). Then
\[
 \operatorname{Var}_p Q=
 \sum_{e,f}\Delta_e\Delta_f
   [p(C_e\cap C_f)-p(C_e)p(C_f)].
\]
For independent insertions \(s,s'\), set
\(H=n\{T_n(s',W)-T_n(s,W)\}\). Conditional independence and the
triangle inequality for centered \(L^2\) variables give
\begin{equation}
 \begin{gathered}
 \E[H^2\mid W]
 =2\operatorname{Var}_s[g(s,V_{J(s)})+Q_{J(s)}\mid W],\\
 \left|\sqrt{\E[H^2\mid W]/2}
        -\sqrt{\operatorname{Var}_p Q}\right|\le G.
 \end{gathered}
 \label{eq:highd-cavity-tree-variance}
\end{equation}
Consequently \(\E H^2=o(n)\) is equivalent to
\(\E\operatorname{Var}_p Q=o(n)\). Proving this estimate for both
source and target deletion would give \(\operatorname{Var}(T_n)\to0\)
by Efron--Stein. This influence estimate remains unproved and is
not needed by the pressure proof below.
Proposition~\ref{prop:iid-cavity-coarea} supplies conditional
independence of edge source positions given the prices and tree; it
does not center their score increments or bound long tree paths.
The separate exact mean-value identity and diffuse comparison prove
pressure differentiability instead.

\subsection{Exact iid cavity disintegration}
\label{subsec:iid-cavity-disintegration}

The finite cavity tree has an exact distributional consequence. Fix
\(n\ge2\) target labels and let an iid source row \(X\sim\mu\) have
finite measurable costs \(c_1(X),\ldots,c_n(X)\). Assume that
\begin{equation}
\label{eq:iid-cavity-atomless}
  c_u(X)-c_v(X)\quad\hbox{has an atomless law whenever }u\ne v.
\end{equation}
For \(n-1\) retained rows let \(R_j\) be the optimal cost with target
\(j\) vacant, and normalize their common prices by
\(\beta_j=R_n-R_j\). An independent source chooses
\[
 J_\beta(X)=\mathop{\rm argmin}_{1\le j\le n}
                  \{c_j(X)-\beta_j\},\qquad
 p_j(\beta)=\mu\{J_\beta=j\}.
\]
For each fixed price vector the minimizing label is unique
\(\mu\)-almost everywhere by \eqref{eq:iid-cavity-atomless}.

\begin{theorem}[Distribution-free cavity cell masses]
\label{thm:iid-cavity-dirichlet}
Under \eqref{eq:iid-cavity-atomless}, the random vector
\(p(\beta)\) has law \(\operatorname{Dirichlet}(1,\ldots,1)\).
In particular its coordinates are positive almost surely. This holds
for arbitrary source support and does not require a density for the
cost vector or for the prices.
\end{theorem}

\begin{proof}
Conditioning on all but one row on which two assignment maps differ
shows that their costs tie with probability zero. Thus every finite
assignment problem with specified integer target counts has an almost
surely unique optimizer.

We first extend Proposition~\ref{prop:highd-cavity-tree} to positive
integer target capacities \(r_1,\ldots,r_n\) and
\(N=\sum_j r_j-1\) rows. Write \(M^j,R_j\) for the optimum and its
cost with target counts \(r-e_j\). For shifted costs
\(d_{ij}=c_{ij}+R_j\), all these optima have the same value
\(D=\sum_j r_jR_j\). Every assignment obeying the upper capacities
has exactly one vacancy, so its shifted cost is at least \(D\).
The column variables \(\gamma_j\) in the capacitated assignment dual
satisfy \(\gamma_j\le0\). Complementary slackness with \(M^j\)
forces \(\gamma_j=0\). Applying this to every vacancy gives the
common dual
\[
 \beta_j=R_n-R_j,\qquad
 \alpha_i=\min_j(c_{ij}-\beta_j).
\]
Every edge used by any \(M^j\) is tight for this dual.

Fix a reference vacancy \(r\). For every nonmatching tight incidence
\((i,u)\), draw an arc \(u\to M^r(i)\) labeled by row \(i\).
A directed simple cycle has distinct heads and hence distinct row
labels. Flipping those rows would preserve the target counts and cost,
contradicting uniqueness. Comparing \(M^j\) with \(M^r\) gives a
flow with one unit leaving \(r\), one entering \(j\), and balance
elsewhere; thus every target is reachable from \(r\). Each directed
\(r\)-to-\(j\) path flips distinct rows and gives a tight assignment
with vacancy \(j\). Two different paths give different assignments,
so uniqueness implies a unique such path. The directed tight graph is
therefore an arborescence. Its \(n-1\) arcs have distinct row labels,
since arcs carrying one label have the same head. Exactly \(n-1\)
rows have two tight targets and form a subdivided target tree; all
other rows have one tight target. The tree rows in \(M^r\) occupy
each target other than \(r\) once, so exactly \(r_j-1\) of the
remaining, strict rows choose target \(j\). This tree-row subset is
determined by the common dual, independently of the reference vacancy.

Now fix nonnegative integers \(k_1,\ldots,k_n\), put
\(K=\sum_j k_j\), and sample \(N=n-1+K\) rows. For every subset
\(I\) of \(n-1\) rows, form its ordinary cavity prices. Call \(I\)
successful when its complementary rows have strict cell counts \(k\).
All complementary rows avoid cell ties almost surely: condition on
\(I\), use \eqref{eq:iid-cavity-atomless}, and take a finite union.

There is exactly one successful subset. Indeed, apply the capacitated
construction with \(r_j=k_j+1\). Its canonical tree rows succeed,
because the strict rows have counts \(k\); removing these same strict
rows from every quota optimum changes every vacancy cost by the same
constant and leaves the ordinary cavity prices unchanged. Conversely,
if \(I\) succeeds, append its strict complementary choices to each
ordinary vacancy matching. The resulting quota assignments are
calibrated by the same prices and hence are optimal. Their costs have
the form \(C-\beta_j^I\), so their normalized quota prices equal
\(\beta^I\). The rows of \(I\) are precisely the two-tight-target
rows, while its complement is strict. Thus \(I\) is the canonical
quota tree-row subset.

Exchangeability and the conditional multinomial law of the complement
now give
\[
 \binom{n-1+K}{n-1}\frac{K!}{\prod_j k_j!}
       \E\prod_j p_j^{k_j}=1.
\]
Consequently
\begin{equation}
\label{eq:iid-cavity-moments}
 \E\prod_{j=1}^n p_j^{k_j}
   =\frac{(n-1)!\prod_{j=1}^n k_j!}{(n-1+\sum_j k_j)!}.
\end{equation}
These are the Dirichlet moments. Polynomial approximation on the
compact probability simplex determines the law.
\end{proof}

The atomlessness hypothesis applies to genuine global length costs,
including cross-pairs for which distance is not smooth.

\begin{lemma}[Fixed-target distance differences]
\label{lem:distance-difference-nonatomic}
Let \((M,F)\) be a complete smooth reversible strongly convex Finsler
manifold of dimension \(d\ge2\), and let \(\mu\) be absolutely
continuous with respect to smooth volume. For any distinct fixed
targets \(y,z\), the law of
\(d_F(X,y)-d_F(X,z)\), \(X\sim\mu\), is atomless.
\end{lemma}

\begin{proof}
Both distance functions are locally Lipschitz and differentiable almost
everywhere. Where their differentials agree, strong convexity forces
their minimizing geodesics from \(x\) to have the same initial unit
velocity. Thus \(x,y,z\) lie on one oriented geodesic, with \(y,z\)
reached at different times. If such an \(x\) exists, the minimizing
route between \(y\) and \(z\) is unique: splicing any alternative
route to the segment from \(x\) to the nearer target would give a
minimizing curve, which cannot have a corner; uniqueness for the
geodesic equation then identifies the routes. Hence every point of
equal differentials lies on the two-sided continuation of this fixed
\(y\)-to-\(z\) geodesic. Its image is a countable union of smooth
finite arcs and has zero \(d\)-dimensional volume.

The locally Lipschitz difference therefore has nonzero differential
almost everywhere. A locally Lipschitz function has zero differential
at almost every density point of each of its level sets. Each level
set must consequently have volume zero, and hence \(\mu\)-measure
zero.
\end{proof}

In the full geometric setting of
Proposition~\ref{prop:general-finite-uniqueness}, the iid targets are
distinct almost surely. Conditional on each such target cloud,
Theorem~\ref{thm:iid-cavity-dirichlet} therefore applies without a
global smooth ray chart.

\begin{proposition}[Conditional mean plans and values]
\label{prop:iid-cavity-mean}
Fix target points \(V=(V_1,\ldots,V_n)\), put
\(c_j(x)=c(x,V_j)\), and assume
\eqref{eq:iid-cavity-atomless} and
\(\int\max_j|c_j|\,d\mu<\infty\). Let \(\pi_n\) be the
optimal coupling of \(n\) iid sources to these targets, all atoms
having mass \(1/n\). For Dirichlet-almost every \(p\), let
\(\gamma_p^V\) denote the unique semidiscrete optimum from \(\mu\)
to \(\sum_jp_j\delta_{V_j}\). Then, as measures,
\begin{equation}
\label{eq:iid-cavity-mean-plan}
 \E[\pi_n\mid V]
   =\E_{p\sim\operatorname{Dirichlet}(1,\ldots,1)}\gamma_p^V.
\end{equation}
For a bounded measurable full-endpoint test \(g\), let \(V_n(t)\)
be the normalized empirical minimum for \(c+tg\), and let
\(V_{\rm sd}(t,p)\) be the semidiscrete minimum with weights \(p\).
For every \(t\in\R\),
\begin{equation}
\label{eq:iid-cavity-pressure}
 \E[V_n(t)\mid V]
   =\E_{p\sim\operatorname{Dirichlet}(1,\ldots,1)}
                   V_{\rm sd}(t,p).
\end{equation}
\end{proposition}

\begin{proof}
For fixed prices \(\beta\), allocate \(x\) to \(J_\beta(x)\).
The inequality
\[
 c_j(x)\ge \min_k(c_k(x)-\beta_k)+\beta_j
\]
calibrates this allocation for its own target masses \(p(\beta)\).
Any other optimum with those masses must use the same minimizing
cells. Their ties are \(\mu\)-null, so this optimum is unique and
depends only on \(p(\beta)\), even if its prices are not unique.
Theorem~\ref{thm:iid-cavity-dirichlet} supplies existence and uniqueness
for Dirichlet-almost every \(p\); conditioning the cell allocation on
its mass vector gives a measurable version.

Distinguish the last source and condition on the other \(n-1\).
Its partner is
\(\mathop{\rm argmin}_j\{c_j(X_n)+R_j\}=J_\beta(X_n)\).
Its conditional source-partner law is the cell allocation just
described. Source exchangeability identifies the expectation of the
empirical average with this distinguished-row law. Averaging over the
cavity proves \eqref{eq:iid-cavity-mean-plan}, including its formulation
against every bounded measurable paired test.

For the value identity first choose \(t\) such that the perturbed pair
differences are atomless. The same mean-plan identity for \(c+tg\),
integrated against its cost, gives
\eqref{eq:iid-cavity-pressure}. All but countably many \(t\) have this
property. To see this, fix a target pair and write
\(h=c_u-c_v\), \(k=g(\cdot,V_u)-g(\cdot,V_v)\). For independent
\(X,X'\), let \(\Delta h=h(X)-h(X')\), and likewise define
\(\Delta k\). Atomlessness gives \(\Pp(\Delta h=0)=0\).
The law of \(h(X)+tk(X)\) has an atom exactly when
\(\Pp(\Delta h+t\Delta k=0)>0\). Such a \(t\) must be an atom
of the random ratio \(-\Delta h/\Delta k\) on
\(\{\Delta k\ne0\}\), so there are at most countably many.
Take a finite union over target pairs. Finally, both sides of
\eqref{eq:iid-cavity-pressure} are \(\|g\|_\infty\)-Lipschitz in
\(t\); equality on the dense good set extends to every \(t\).
The integrability hypothesis justifies integration of the base cost.
\end{proof}

One can realize the weights as \(p_j=E_j/\sum_kE_k\), where the
\(E_j\) are iid mean-one exponentials independent of the targets.
Thus \eqref{eq:iid-cavity-mean-plan} replaces the source sampling in
the conditional mean by random target weights. Equation
\eqref{eq:iid-cavity-pressure} remains exact after averaging over iid
targets and applying a common deterministic normalization. At exceptional
perturbation parameters the value identity does not assert uniqueness
or a mean-plan identity for an arbitrary optimizer selection. These
finite identities do not establish concentration of the empirical plan.

Under regularity of the row costs there is also a density-level
description of the cavity tree itself.

\begin{proposition}[Coarea law of the cavity tree]
\label{prop:iid-cavity-coarea}
Let the source law have density \(\rho\) on a bounded open
\(\Omega\subset\R^m\). Assume that each \(c_j\) is \(C^1\) and
Lipschitz on a neighborhood of \(\Omega\), and that
\(\nabla(c_u-c_v)\ne0\) almost everywhere for every \(u\ne v\).
For \(\beta_n=0\), define the regular strict facet
\[
 \Sigma_{uv}(\beta)=\left\{x\in\Omega:
 \begin{array}{l}
 c_u(x)-\beta_u=c_v(x)-\beta_v
       <c_j(x)-\beta_j\quad(j\notin\{u,v\}),\\
 \nabla(c_u-c_v)(x)\ne0
 \end{array}\right\},
\]
and put
\[
 w_{uv}(\beta)=\int_{\Sigma_{uv}(\beta)}
      \frac{\rho(x)}{|\nabla(c_u-c_v)(x)|}\,
          d\mathcal H^{m-1}(x),\qquad
 \tau(w)=\sum_{T\text{ spanning tree}}\prod_{uv\in T}w_{uv}.
\]
The normalized cavity prices have density
\((n-1)!\tau(w(\beta))\) with respect to
\(d\beta_1\cdots d\beta_{n-1}\). Conditional on \(\beta\), the
target tree has law
\[
 \Pp(T\mid\beta)=\frac{\prod_{uv\in T}w_{uv}(\beta)}
                           {\tau(w(\beta))}.
\]
The source-to-edge bijection is conditionally uniform. Given the prices,
tree, and bijection, its source rows are independent on their respective
facets, with probability density
\(\rho/(w_{uv}|\nabla(c_u-c_v)|)\) relative to
\(\mathcal H^{m-1}\).
\end{proposition}

\begin{proof}
The coarea formula first implies
\eqref{eq:iid-cavity-atomless}. Fix a target tree, orient its edges,
and assign its \(N=n-1\) edges to the labeled source rows. Its reduced
incidence matrix \(B\) is square with determinant \(\pm1\).
Writing \(d_i(x_i)=c_{u_i}(x_i)-c_{v_i}(x_i)\), the prices required
by these edge equalities are
\(\beta=B^{-1}(d_i(x_i))_{i=1}^N\). On the domain where every
assigned edge is a strict regular active facet, these are exactly the
cavity tree and prices: all its vacancy matchings are calibrated.
Conversely Proposition~\ref{prop:highd-cavity-tree} and the assumed
gradient condition partition the iid inputs into these domains up to
a null set.

The derivative before multiplication by \(B^{-1}\) has disjoint
source-coordinate blocks. Its Gram matrix is
\(\operatorname{diag}(|\nabla d_i|^2)\); therefore the coarea
Jacobian of the price map is exactly \(\prod_i|\nabla d_i|\).
At fixed prices the fiber is the Cartesian product of the assigned
facets, in orthogonal coordinate blocks. Coarea thus gives the joint
price density \(\prod_{uv\in T}w_{uv}\) for this labeled tree and
the asserted conditional product facet law. Summing over the \(N!\)
edge labelings and all target trees proves the proposition, including
normalization. All conditional formulas are understood at almost every
price under this probability law, where the total tree weight is
positive and finite; facet products can equivalently be defined as
product measures before normalizing.
\end{proof}


\subsection{Differentiability from a common diffuse representation}
\label{subsec:highd-diffuse-pressure}

The exact finite identity in Proposition~\ref{prop:iid-cavity-mean}
allows opposite perturbations to be compared using a common diffuse
source and the same random target masses. Their average is nearly
optimal for the unperturbed cost. Local fractional exchanges then
force its limiting allocation to be a map, identifying the two pressure
slopes without a new thermodynamic comparison for diffuse transport.

\begin{theorem}[Differentiability of homogeneous paired pressures]
\label{thm:highd-pressure-differentiability}
Let \(q\ge3\), let \(K_-,K_+\) be compact metric mark spaces with
arbitrary probability laws \(\alpha,\beta\), and let
\(S:K_-\times K_+\to\R^{q\times q}\) be continuous, symmetric,
and uniformly positive definite. Every mark pair is allowed. For
\[
 k((x,a),(y,b))=\tfrac12(y-x)^TS(a,b)^{-1}(y-x)
\]
and every \(g\in C(K_-\times K_+)\), the pressure \(p_g\) of
Proposition~\ref{prop:highd-finite-pressure} belongs to \(C^1(\R)\).
In particular, the paired-mark measures of all unperturbed empirical
kinetic optimizers converge in probability to one deterministic
\(\kappa\in\Pi(\alpha,\beta)\), uniformly over permutation and
fractional optimizers, with \(\kappa(g)=p_g'(0)\).
\end{theorem}

\begin{proof}
Write \(G=\|g\|_\infty\), \(L_n=n^{1/q}\),
\(D_n=(0,L_n)^q\), and
\(p_n(t)=n^{-1}\E\mathcal K_n(t)\), using the microscopic
rescaling in Proposition~\ref{prop:highd-finite-pressure}.
That proposition gives locally uniform convergence
\(p_n\to p_g\) on \(\R\) and \(\sup_n p_n(0)<\infty\).

First fix distinct spatial targets \((y_j,b_j)\). For every fixed
\(a\) and \(j\ne l\), the difference
\[
 d_{jl}(x,a)=k((x,a),(y_j,b_j))-k((x,a),(y_l,b_l))
\]
is a nonconstant polynomial in \(x\). Indeed, if
\(S(a,b_j)^{-1}\ne S(a,b_l)^{-1}\), its quadratic part is
nonzero. Otherwise its linear part is
\(S(a,b_j)^{-1}(y_l-y_j)\ne0\). Every level set of a
nonconstant polynomial has Lebesgue measure zero, by induction on
dimension and the one-dimensional finite-root property. Fubini gives
\begin{equation}
 (dx\,\alpha(da))\{(x,a):d_{jl}(x,a)=h\}=0
 \qquad(h\in\R).
 \label{eq:highd-diffuse-levelsets}
\end{equation}
Adding \(t[g(a,b_j)-g(a,b_l)]\) changes only the constant term
in \(x\). Thus the atomless pair-difference hypothesis of
Proposition~\ref{prop:iid-cavity-mean} holds for every real \(t\),
including atomic mark laws. This argument applies to each fixed
countable target configuration with distinct spatial locations.

We will also use its following consequence. A coupling from
\(dx\,\alpha(da)\), or a measure absolutely continuous with respect
to it, to such a countable target set is induced by a map whenever
it has a two-cyclically monotone full-measure support. To see this,
let \(f_j\) be the density of the source mass sent to target \(j\).
If \(f_j,f_l>0\) on a positive-measure set \(A\), remove the null
exceptions so both target choices lie in the monotone support at
every \(r\in A\). For \(r,r'\in A\), the two-cycle inequality
gives \(d_{jl}(r)\le d_{jl}(r')\); interchanging them gives
equality. This contradicts \eqref{eq:highd-diffuse-levelsets}.
Countability of the targets proves the assertion. Consequently, if
two couplings have the same diffuse first marginal and their average
is two-cyclically monotone, the two couplings coincide: the average
conditional target law is a Dirac mass, so both summands have that
same conditional law.

Sample the \(n\) iid target endpoints \((Y_j,B_j)\) in
\(D_n\times K_+\). Independently take
\(P\sim\mathrm{Dirichlet}(1,\ldots,1)\), and set
\[
 \lambda_n=\mathbf1_{D_n}(x)\,dx\,\alpha(da),
 \qquad \nu_n=\sum_{j=1}^n nP_j\delta_{(Y_j,B_j)}.
\]
Both measures have mass \(n\). Define
\[
 F_n(t)=\frac1n\min_{\Gamma\in\Pi(\lambda_n,\nu_n)}
                \int(k+t g)\,d\Gamma.
\]
The compact finite problem has measurable optimal plans, and the
exact identity \eqref{eq:iid-cavity-pressure}, after rescaling, says
\begin{equation}
 \E F_n(t)=p_n(t)\qquad(t\in\R).
 \label{eq:highd-diffuse-value}
\end{equation}
In particular the same target-and-weight environment can be used for
every \(t\); no joining with an original atomic optimizer is needed.

Choose \(t_n\downarrow0\), sufficiently slowly as specified below,
and on this common environment let \(\Gamma_n^+,\Gamma_n^-\)
minimize the costs \(k+t_ng,k-t_ng\), respectively. Let
\(\Gamma_n^0\) minimize \(k\), and put
\(\overline\Gamma_n=(\Gamma_n^++\Gamma_n^-)/2\).
The optimality comparisons with \(\Gamma_n^0\) give
\begin{equation}
 \begin{split}
 0\le\Delta_n
   &:=\int k\,d\overline\Gamma_n-\int k\,d\Gamma_n^0\\
   &\le\frac{t_n}{2}\int g\,d(\Gamma_n^--\Gamma_n^+)
    \le t_nGn,
 \end{split}
 \label{eq:highd-diffuse-excess}
\end{equation}
and, for either sign,
\[
 \int k\,d\Gamma_n^\pm
    \le\int k\,d\Gamma_n^0+2t_nGn.
\]
By \eqref{eq:highd-diffuse-value}, their expected kinetic energy
per unit volume is uniformly bounded. Uniform ellipticity gives the
same bound for squared spatial displacement.

Choose an independent uniform root \(U_n\in D_n\), and write
\(\tau_u\) for translation of both spatial endpoints by \(-u\).
Retain the common target process, including its weights, in the
joint rooted scene of \(\Gamma_n^+,\Gamma_n^-\). For every bounded
spatial set \(B\), root averaging gives
\begin{equation}
 \begin{split}
 \E\int\mathbf1_{\{x-U_n\in B,\ |y-x|>R\}}\,d\Gamma_n^\pm
   &\le C|B|R^{-2},\\
 \E\int\mathbf1_{\{y-U_n\in B,\ |y-x|>R\}}\,d\Gamma_n^\pm
   &\le C|B|R^{-2}.
 \end{split}
 \label{eq:highd-diffuse-tails}
\end{equation}
For example, integration of \(\mathbf1_{\{x-U_n\in B\}}\)
over the root contributes at most \(|B|/n\), after which the
energy bound applies. The source measures converge locally to
\(\lambda=dx\,\alpha(da)\).
Writing \(P_j=E_j/\sum_iE_i\) with iid mean-one exponential
variables, the marked binomial local limit and
\(n/\sum_iE_i\to1\) in probability give the common target limit
\[
 \nu=\sum_j E_j\delta_{(Y_j,B_j)},
\]
where the spatial process is a unit-intensity PPP and its marks and
exponential weights are independent. In particular its spatial
locations are distinct and locally finite. Local target masses are
uniformly integrable, since the Dirichlet moments give
\[
 \E[\nu_n(U_n+B)^2]\le2|B|+|B|^2.
\]
Local mass bounds and \eqref{eq:highd-diffuse-tails} therefore give
joint subsequential vague limits \((\Gamma^+,\Gamma^-,\nu)\).
The two tail bounds retain both marginals, so both plans belong to
\(\Pi(\lambda,\nu)\) almost surely. The joint law is stationary:
translating the root integral by a fixed vector changes it only on
the symmetric difference of the cube and its translate, whose
relative volume tends to zero.

We next transfer \eqref{eq:highd-diffuse-excess} to local
two-cycle optimality of the limiting average. Fix bounded open source
rectangles \(A_1,A_2\), disjoint in the full source space, and
bounded disjoint target rectangles \(B_1,B_2\), such that for some
\(\varepsilon>0\),
\begin{equation}
 k(r_1,s_1)+k(r_2,s_2)-k(r_1,s_2)-k(r_2,s_1)
       \ge\varepsilon
 \quad(r_i\in A_i, s_i\in B_i).
 \label{eq:highd-diffuse-certificate}
\end{equation}
Rational spatial rectangles, countable bases of the compact mark
spaces, and rational \(\varepsilon>0\) give a countable family
detecting every strict support violation. The functional
\[
 W(\Gamma)=\min\{1,\Gamma(A_1\times B_1),
                        \Gamma(A_2\times B_2)\}
\]
is bounded and lower semicontinuous for vague convergence. If
\(m=W(\Gamma)>0\), scale the two indicated restrictions to mass
\(m\), calling them \(\gamma_1,\gamma_2\), with source marginals
\(\lambda_i\) and target marginals \(\nu_i\). Replacing their
sum by
\[
       m^{-1}(\lambda_1\otimes\nu_2+\lambda_2\otimes\nu_1)
\]
preserves every marginal and decreases cost by at least
\(\varepsilon m\), by integration of
\eqref{eq:highd-diffuse-certificate}.

Take a coarse cube of side \(K\) containing both source spatial
projections in its interior. Translate this correction to centers
\(z\) of a randomly shifted \(K\mathbb Z^q\) grid, keeping centers
in \(D_n\). The removed source subplans lie in disjoint cubes,
so all corrections can be performed simultaneously. Overlapping
target neighborhoods are harmless because each correction preserves
each target marginal. Source rectangles crossing the boundary of
\(D_n\) carry zero mass outside the domain. The resulting plan is
still feasible for the finite problem, and its total improvement is
at most \(\Delta_n\). Averaging the grid phase gives exactly
\begin{equation}
 \E W(\tau_{U_n}\overline\Gamma_n)
   \le\frac{K^q\E\Delta_n}{\varepsilon n}
   \le\frac{K^qG}{\varepsilon}t_n\longrightarrow0.
 \label{eq:highd-diffuse-packing}
\end{equation}
Indeed the phase-averaged sum is
\(K^{-q}\int_{D_n}W(\tau_z\overline\Gamma_n)\,dz\).
This argument needs no energy uniform integrability or boundary
repair estimate. Lower semicontinuity and Portmanteau imply that
every such witness vanishes on \((\Gamma^++\Gamma^-)/2\)
almost surely. Its support is consequently two-cyclically monotone.
The map observation following
\eqref{eq:highd-diffuse-levelsets} now yields
\begin{equation}
                 \Gamma^+=\Gamma^-\quad\text{almost surely}.
 \label{eq:highd-diffuse-equality}
\end{equation}

To read out the pressure derivatives, define
\(q_n^\pm=n^{-1}\E\int g\,d\Gamma_n^\pm\).
By \eqref{eq:highd-diffuse-value}, these are supergradients of
\(p_n\) at \(\pm t_n\). Choose \(t_n\downarrow0\) so that
the uniform error \(p_n-p_g\) on a fixed neighborhood of zero
is \(o(t_n)\). The concave secant bounds are
\[
 \frac{p_n(2t_n)-p_n(t_n)}{t_n}
 \le q_n^+\le\frac{p_n(t_n)-p_n(0)}{t_n},
\]
\[
 \frac{p_n(0)-p_n(-t_n)}{t_n}
 \le q_n^-\le\frac{p_n(-t_n)-p_n(-2t_n)}{t_n}.
\]
They imply \(q_n^+\to p_g'(0+)\) and
\(q_n^-\to p_g'(0-)\).
For a nonnegative \(\eta\in C_c(\R^q)\) with
\(\int\eta=1\), the deterministic diffuse source marginal gives
\[
 \left|q_n^\pm-
   \E\int\eta(x)g(a,b)\,d(\tau_{U_n}\Gamma_n^\pm)\right|
       \le C G/L_n.
\]
The discrepancy is supported in a fixed-width source boundary
layer. Truncate the target endpoint to a bounded spatial set; then
the integral is a continuous compactly supported plan test, bounded
by the fixed diffuse source mass. Removing the cutoff costs at most
\(CG/R^2\) by \eqref{eq:highd-diffuse-tails}. Its expectation
therefore converges along the joint subsequence to
\(\E\int\eta(x)g(a,b)\,d\Gamma^\pm\).
Equation~\eqref{eq:highd-diffuse-equality} identifies the two
limits, proving \(p_g'(0+)=p_g'(0-)\).

For an arbitrary fixed \(t_0\), repeat the argument with base cost
\(k+t_0g\). It remains coercive up to a bounded shift, and the
energy comparison gains only \(2|t_0|Gn\). Its pair differences
have the same nonconstant spatial polynomial, so the map argument
and the local certificates apply unchanged. Thus \(p_g\) is
differentiable everywhere. A differentiable concave function on
\(\R\) has continuous derivative. Finally apply
Corollary~\ref{cor:highd-pressure-selection} to a countable uniformly
dense family of paired tests; its finite-difference argument gives
the asserted uniform conclusion for all optimizers.
\end{proof}


\subsection{Pressure in regular tubes with varying data}
\label{subsec:highd-variable-pressure}

Let \(q\ge3\), and retain the smooth ordered tube and compact collars
of Lemma~\ref{lem:highd-coercivity}. The common transverse probability
density \(\theta\) on \(D=[0,1]^q\) is positive, bounded, and Lipschitz.
The radial intervals \(I_A,I_B\) are fixed and separated. Assume the
conditional mark laws have jointly measurable representations
\begin{equation}
\label{eq:highd-conditional-representations}
 \begin{gathered}
 Q_A:D\times(0,1)\longrightarrow I_A,\qquad
 Q_B:D\times(0,1)\longrightarrow I_B,\\
 \alpha_z=Q_A(z,\cdot)_\#du,\qquad
 \beta_z=Q_B(z,\cdot)_\#du,
 \end{gathered}
\end{equation}
uniformly continuous in \(z\), uniformly over the latent variable:
\[
 \sup_{\substack{|z-z'|\le h\\0<u<1}}
       |Q_A(z,u)-Q_A(z',u)|\longrightarrow0
 \quad(h\downarrow0),
\]
and similarly for \(B\). No continuity or monotonicity in the latent
variable is required. This includes fixed product laws, with atoms.
It also includes conditional densities \(f_A(z,a),f_B(z,b)\) jointly
continuous and bounded positively on the fixed compact rectangles.
Indeed, their inverse CDFs satisfy
\[
 |Q_A(z,u)-Q_A(z',u)|
 \le(\inf f_A)^{-1}\sup_a|F_A(z,a)-F_A(z',a)|,
\]
and similarly for \(B\). The two
sample families are independent and iid within each color, with laws
\(\theta(z)\,dz\,\alpha_z(da)\) and
\(\theta(w)\,dw\,\beta_w(db)\).

Write \(e(a,z;b,w)=d_F(X(a,z),X(b,w))-(b-a)\) and
\[
 \widehat e(a,z;b,w)=\tfrac12(w-z)^TS(a,b,z)^{-1}(w-z),
 \qquad S(a,b,z)=\int_a^b H(r,z)^{-1}\,dr.
\]
For either \(c=e\) or \(c=\widehat e\), and
\(\Psi\in C(D\times I_A\times I_B)\), set
\begin{equation}
\label{eq:highd-variable-finite-pressure}
 V_n^\Psi(t)=\frac1n\min_\sigma\sum_{i=1}^n
 \left[n^{2/q}c(A_i,Z_i;B_{\sigma(i)},W_{\sigma(i)})
             +t\Psi(Z_i,A_i,B_{\sigma(i)})\right].
\end{equation}
For fixed \(z\), let \(p_{z,g}(s)\) be the homogeneous unit-intensity
pressure with compliance \(S(\,\cdot,\,\cdot,z)\), mark laws
\(\alpha_z,\beta_z\), and paired-mark probe \(g\). By
Proposition~\ref{prop:highd-stationary-pressure}, it is the corresponding
stationary Palm mean-cost infimum.

\begin{proposition}[Variable pressure and local mean minimality]
\label{prop:highd-variable-pressure}
Define
\begin{equation}
\label{eq:highd-variable-pressure}
 P_\Psi(t)=\int_D\theta(z)^{1-2/q}
       p_{z,\Psi(z,\cdot,\cdot)}(t\theta(z)^{2/q})\,dz.
\end{equation}
For every \(T<\infty\),
\[
 \E\sup_{|t|\le T}|V_n^\Psi(t)-P_\Psi(t)|\longrightarrow0.
\]
For fixed \(t\), every subsequential random conditional scene law of
optimizers of \eqref{eq:highd-variable-finite-pressure} is almost surely,
at almost every macro label \(z\), a stationary mean minimizer for
\(k_z+t\Psi(z,\cdot,\cdot)\) with the prescribed intensity-\(\theta(z)\)
inputs. Here \(k_z((x,a),(y,b))=\tfrac12(y-x)^TS(a,b,z)^{-1}(y-x)\).
\end{proposition}
\begin{proof}
First record continuity of the frozen values. Couple neighboring frozen
models using the same independent uniform latent marks in
\eqref{eq:highd-conditional-representations}. Its uniform transverse
modulus makes the corresponding physical marks uniformly close.
Their kinetic costs differ by \(o(1)|y-x|^2\), uniformly, and their
probes differ uniformly by \(o(1)\). The unmarked comparator gives
uniform expected kinetic-energy bounds for their finite optimizers on
bounded parameter intervals. Testing each optimizer in the other model
therefore proves continuity of
\((z,s)\mapsto p_{z,\Psi(z,\cdot,\cdot)}(s)\).
The stationary infimum per unit volume at intensity \(\lambda\) is
\begin{equation}
\label{eq:highd-variable-intensity-scaling}
 \lambda^{1-2/q}p_{z,g}(t\lambda^{2/q}),
\end{equation}
by the spatial scaling \(x\mapsto\lambda^{1/q}x\).

For the upper bound, partition \(D\) into finitely many cubes \(D_j\)
of side \(\ell\), choose \(z_j\in\overline D_j\), and put
\(\bar\theta_j=|D_j|^{-1}\int_{D_j}\theta\).
There is a map \(\Phi_j\) from uniform probability on \(D_j\) to
its normalized \(\theta\)-density with
\begin{equation}
\label{eq:highd-local-density-map}
 \operatorname{Lip}(\Phi_j-\operatorname{Id})\le C\ell.
\end{equation}
Indeed, after scaling to the unit cube the density is \(1+O(\ell)\)
with Lipschitz constant \(O(\ell)\). Successive conditional densities
in the inverse Rosenblatt construction have the same properties.
Their inverse CDFs \(Q_k(p,u)\), where \(p\) consists of earlier
output coordinates, satisfy
\[
 \begin{split}
 |Q_k(p,u)-Q_k(p,v)-(u-v)|&\le C\ell|u-v|,\\
 |Q_k(p,u)-Q_k(p',u)|&\le C\ell|p-p'|.
 \end{split}
\]
Recursive composition proves \eqref{eq:highd-local-density-map}, with
constants uniform in the cube.

In each cube write \(z=\Phi_j(u)\), \(w=\Phi_j(v)\), and set
\(a=Q_A(z,\xi),b=Q_B(w,\zeta)\),
\(a_0=Q_A(z_j,\xi),b_0=Q_B(z_j,\zeta)\).
For product laws these maps can be constant in \(z\), so
\(a=a_0,b=b_0\). Uniform local Taylor expansion,
\eqref{eq:highd-local-density-map}, and uniform continuity give
\begin{align}
\label{eq:highd-block-freezing}
 \left|c(a,z;b,w)-\tfrac12(v-u)^TS(a_0,b_0,z_j)^{-1}(v-u)\right|
      &\le\varepsilon(\ell)|v-u|^2,\notag\\
 |\Psi(z,a,b)-\Psi(z_j,a_0,b_0)|&\le\varepsilon(\ell),
 \qquad \varepsilon(\ell)\longrightarrow0.
\end{align}
All within-cube edges are short enough for this uniform expansion as
\(\ell\downarrow0\); no microscopic rate for the freezing error is
required.

Start with independent Poisson samples of intensity \(n\theta\), and
split them into main intensity \(\lambda n\theta\) and reserve intensity
\(\rho n\theta\), where \(\lambda=1-\rho\) and \(0<\rho<1/2\).
Pull main positions back by \(\Phi_j\), freeze marks and coefficients,
and optimize a maximum-cardinality matching in each cube. The frozen
positions have intensity \(\lambda n\bar\theta_j\), so their expected
contribution divided by \(n\) tends to
\[
 |D_j|(\lambda\bar\theta_j)^{1-2/q}
       p_{z_j,\Psi(z_j,\cdot,\cdot)}
                         (t(\lambda\bar\theta_j)^{2/q}).
\]
The unmarked comparator bounds the sum of their normalized kinetic
energies by \(C_T\) in the limit, uniformly in the partition and
\(\rho<1/2\). Pushing back to the actual data therefore costs at most
\(C_T\varepsilon(\ell)\) in \eqref{eq:highd-block-freezing}.

The number of unused main points in each cube is the absolute difference
of its Poisson color counts. The total residual count \(R_n\) obeys
\(\E R_n^r\le C_{\mathcal D,r}n^{r/2}\) for every \(r>0\), with
the partition \(\mathcal D\) fixed. These adaptive residual clouds
are independent of the reserves. Pull both back by the global inverse
of the Lipschitz Rosenblatt map sending uniform volume to \(\theta\).
The reserves become uniform; the residual count bounds are unchanged.
The reserve estimate \cite[Proposition~6.3]{GoldmanTrevisan2023}, as
used in Proposition~\ref{prop:highd-localized-upper}, and pushforward
by that Lipschitz map bound the normalized quadratic repair cost by
\[
 C_{\mathcal D,\theta}\rho^{1-2/q}
       +C_{\mathcal D,\theta}\rho^{-2/q}n^{-1/2+1/q}.
\]
The expected number of repair edges is at most
\(\rho n+O_{\mathcal D}(\sqrt n)\), so their probe costs at most
\(|t|\|\Psi\|_\infty\rho+o(1)\) after normalization.
This gives a maximum-cardinality matching of the full Poisson inputs.
Send \(n\to\infty\), then \(\rho\downarrow0\) with the partition
fixed, and finally \(\ell\downarrow0\). The frozen sums converge to
\eqref{eq:highd-variable-pressure}. This order allows the reserve
constants to depend on the partition.

The bounded-additive version of the independent-count comparison
from Proposition~\ref{prop:highd-independent-count} and the proof of
Proposition~\ref{prop:highd-finite-pressure} transfers this upper bound
to exactly \(n\) points per color. To check its hypotheses, use the
global uniform spatial coordinates and independent uniform latent
marks. At each fixed \(n\), the resulting
kernel on the cube of side \(n^{1/q}\) is bounded by \(C|x-y|^2\)
plus a bounded probe. It may depend on both latent coordinates and
on \(n\); the comparison keeps this kernel and cube fixed while
varying counts. No selected residual cloud is assumed iid.

For the lower bound, the unmarked comparator gives the same uniform
energy estimate for all optimizers, including the bounded probe.
The random-scene argument of Proposition~\ref{prop:highd-local-limit}
applies to these weakly continuous conditional mark laws. To check
the input statement directly, for a nonnegative continuous compactly
supported marked-window test \(f\), the single-color binomial Laplace
functional is \((1-a_{n,z})^n\), where
\[
 n a_{n,z}=\int \theta(z+n^{-1/q}x)
   \int\bigl(1-e^{-f(x,Q_A(z+n^{-1/q}x,u))}\bigr)\,du\,dx.
\]
At interior roots this converges to the marked Poisson expression
by the uniform transverse modulus and dominated convergence.
For two disjoint bounded spatial windows the covariance error is
\(O(n^{-1})\), since both one-point window probabilities are \(O(n^{-1})\).
Pairs of roots with overlapping windows have total volume \(O(n^{-1})\).
Thus integrated input tests concentrate; marked counts are bounded
by their spatial counts, giving the needed moment bounds. The two
colors are independent. Consequently the conditional inputs at a
macro root \(z\) have intensity \(\theta(z)\) and mark laws
\(\alpha_z,\beta_z\), almost surely, including atomic kernels.
Applying nonnegative compact
truncations to \(k_z+t\Psi(z,a,b)+|t|\|\Psi\|_\infty\) yields
\begin{equation}
\label{eq:highd-variable-lower}
 \liminf_n\E V_n^\Psi(t)
 \ge\E\int_D\theta(z)\E_{P_z^\Gamma}^{\rm Palm}
                 [k_z+t\Psi(z,a,b)]\,dz
 \ge P_\Psi(t).
\end{equation}
The first inequality uses the actual local Taylor expansion and
two-endpoint retention; the second uses the stationary identity and
\eqref{eq:highd-variable-intensity-scaling}. The same statement holds
for each realized random scene measure, whose input marginal is the
prescribed one almost surely. Equality of the upper and lower values
forces the nonnegative conditional excess above the stationary infimum
to vanish almost everywhere, proving local mean minimality. No uniform
integrability of squared displacements is inferred.

It remains to prove concentration with conditional marks. The actual
defect has a uniform relative mark modulus
\begin{equation}
\label{eq:highd-relative-mark-modulus}
 |e(a,z;b,w)-e(a',z;b',w)|
 \le\omega(|a-a'|+|b-b'|)|w-z|^2,
 \qquad \omega(r)\longrightarrow0.
\end{equation}
Near the diagonal, the value and first transverse derivative vanish
for every pair of marks, and the integral second-order Taylor formula
uses a Hessian uniformly continuous in those marks. Away from that
neighborhood, radial Lipschitz continuity of the distance gives the
bound after division by a fixed positive lower bound on \(|w-z|^2\).
The corresponding assertion for \(\widehat e\) follows directly from
continuity of \(S^{-1}\).

Approximate the conditional representations uniformly by maps \(Q_A^m,Q_B^m\)
that are Lipschitz in \(z\), uniformly in the latent mark: take finite
Lipschitz partitions of unity in \(z\) and convex combinations of
representations at their centers. These maps are jointly measurable
and stay in the same radial intervals; no regularity in the latent
variable is used.
Also approximate \(\Psi\) uniformly by Lipschitz probes. Let \(V_n^m\)
be the resulting values on the shared latent sample coupling. The
two-competitor comparison, \eqref{eq:highd-relative-mark-modulus}, and
the uniform optimizer energy bound give
\begin{equation}
\label{eq:highd-quantile-value-approximation}
 \sup_n\E\sup_{|t|\le T}|V_n^\Psi(t)-V_n^m(t)|
       \le\delta_m,\qquad \delta_m\longrightarrow0.
\end{equation}
For fixed product marks, only the probe approximation is needed.

For each fixed \(m\), positions and latent marks are independent, so
the proof of Proposition~\ref{prop:highd-value-concentration} applies
to the transformed cost. Its required derivative bound follows as
follows. Near the transverse diagonal the full defect is smooth in
all endpoint coordinates. Its spatial and mark derivatives vanish
on the diagonal; bounded mixed derivatives make them \(O(|w-z|)\).
The Lipschitz chain rule for \(Q_A^m,Q_B^m\), followed by coercivity,
bounds the squared spatial derivative of the composed cost by
\(C_m\) times that cost. Away from a fixed diagonal neighborhood,
the defect is uniformly Lipschitz in all endpoint coordinates, so
its composition is Lipschitz with constant \(C_m\). The lower bound
on \(|w-z|\) and coercivity give the same estimate there. This uses
the Lipschitz bound directly, without pulling back an ambient
almost-everywhere differentiability set through a conditional graph.
For \(\widehat e\) the same argument is immediate. The composed Lipschitz
probe adds only \(C_{m,T}n\) to the squared gradient of a permutation's
total cost. After the per-edge shift \(1+T\|\Psi^m\|_\infty\),
the total value is at least \(n\); its square root still has squared
spatial gradient at most \(C_{m,T}n^{2/q}\). Resampling one latent
mark changes one row or column, with the same multiplicative shifted
edge-cost bound supplied by coercivity. Poincar\'e and Efron--Stein
therefore give
\[
 \sup_{|t|\le T}\E|V_n^m(t)-\E V_n^m(t)|
       \le C_{m,T}n^{-1/2+1/q}.
\]
Finite parameter meshes upgrade this to expected locally uniform
concentration. Equation~\eqref{eq:highd-quantile-value-approximation}
then gives it for \(V_n^\Psi\) by sending \(m\to\infty\).
The expectation limit already proved is locally uniform by the common
\(\|\Psi\|_\infty\)-Lipschitz bound in \(t\), completing the proof.
\end{proof}

These regularized estimates will be used on compact pieces of the
measurable atlas in the general pressure argument.

\subsection{Random empirical scene measures}

For the rest of this section assume that \(\alpha_z,\beta_z\)
depend weakly continuously on \(z\). Let \(\pi_n\) be any measurable
choice of an optimal empirical coupling. Write its matrix as
\(r_{ij}/n\), where every row and column of \(r\) sums to one.
A finite Birkhoff decomposition
\(r=\sum_\sigma w_n(\sigma)P_\sigma\) can be chosen measurably,
using a greedy decomposition with a fixed lexicographic rule.
Every positive-weight permutation is optimal, since their average
cost is the minimum. This represents the given coupling exactly.

Let \(\omega_{n,z,\sigma}\) contain the two marked clouds centered
at \(z\) and scaled by \(\varepsilon_n^{-1}\), together with the
integer measure of their matched endpoint pairs. The scene space
\(\Omega\) is the product of the spaces of locally finite integer
Radon measures for these three objects, with vague topology.
Define the random probability measure
\begin{equation}
\Gamma_n=\int_D\sum_\sigma w_n(\sigma)
 \delta_{(z,\omega_{n,z,\sigma})}\,dz.
\label{eq:highd-scene}
\end{equation}
The weights are fixed from the data before the root is introduced.
Keeping \(\Gamma_n\) random retains possible randomness of the
macroscopic empirical plan.

At a fixed label \(z\), let \(\mathcal L_z\) consist of the
stationary perfect-matching laws whose input clouds are independent
marked Poisson processes of intensities
\[
dx\,\theta(z)\alpha_z(da),\qquad dy\,\theta(z)\beta_z(db),
\]
whose source-Palm squared displacement is integrable, and whose
matched support is cyclically monotone for
\begin{equation}
k_z((x,a),(y,b))=
 \tfrac12(y-x)^T S(a,b,z)^{-1}(y-x).
\label{eq:highd-local-cost}
\end{equation}
Stationarity means invariance under simultaneous translations of
both transverse clouds and the matching. For \(P\in\mathcal L_z\),
write \(\kappa(P)\) for the longitudinal mark pair law at a typical
source. Equivalently, for any nonnegative \(\eta\in C_c(\R^q)\)
with integral one and any continuous mark test \(g\),
\[
\E_P\sum_{((x,a),(y,b))\text{ matched}}\eta(x)g(a,b)
 =\theta(z)\int g\,d\kappa(P).
\]
Its marginals are \(\alpha_z,\beta_z\), by perfection and mass
transport. The definition of \(\mathcal L_z\) imposes finite local
optimality; it does not impose minimality of the expected cost
among stationary laws.

\begin{proposition}[Local optimizer limits and reconstruction]
\label{prop:highd-local-limit}
The laws of \(\Gamma_n\) are tight. Every subsequential limit
\(\Gamma\) has, almost surely, a disintegration
\[
\Gamma(dz,d\omega)=dz\,P_z^\Gamma(d\omega),
\qquad P_z^\Gamma\in\mathcal L_z
\quad\text{for almost every }z.
\]
If \(\pi_n\) and \(\Gamma_n\) converge jointly in distribution
along a subsequence, the limiting physical plan is the pushforward of
\begin{equation}
\theta(z)\,dz\,\kappa(P_z^\Gamma)(da,db)
\label{eq:highd-reconstruction}
\end{equation}
under \((z,a,b)\mapsto(X(a,z),X(b,z))\).
The conditional matching laws and their paired-mark laws may still
be random.
\end{proposition}

\begin{proof}
Local endpoint counts have bounded second moments. In a compact
product of endpoint windows, the edge count is at most the source
count in its source projection. These bounds give tightness of the
averaged scene laws. They also give tightness of \(\Gamma_n\) as
random probability measures: choose compact scene sets with successively
small expected outside mass and use Markov's inequality with summable
error levels to obtain a tight set of probability measures.

Proposition~\ref{prop:highd-energy} controls both endpoint projections.
A local vague edge limit has degree at most one at every simple input
vertex. A missing source incidence in a bounded window must come from
partners escaping every bounded endpoint window. The expected number
of such possible losses is bounded by \(K|B|/R^2\) for each \(R\),
and hence vanishes as \(R\to\infty\). The target bound gives the
same conclusion there. Count uniform integrability permits passage of
the endpoint projections. Thus limiting scenes are perfect matchings.
Their simplicity and their input law follow from the next argument.

Let \(\Pi_z\) denote the independent Poisson input law appearing
in the definition of \(\mathcal L_z\). The input-only marginal of
\(\Gamma_n\) converges in probability to the deterministic measure
\(dz\,\Pi_z\). To see this, test it against a continuous macro
function \(f\) and a bounded local Laplace cylinder function \(G\)
of the two clouds. The binomial Laplace calculation and weak continuity
of the mark kernels give convergence of its expectation. Its variance
is an integral of covariances over two roots \(z,z'\).
For every fixed distinct pair of interior roots, the microscopic
windows are eventually disjoint. The multinomial rare-event formula
gives independent Poisson limits there, so that covariance tends to
zero. The diagonal has zero product volume and the covariances are
bounded. Dominated convergence gives vanishing variance.
A countable determining family and tightness establish the asserted
probability limit. Consequently almost every realization of a limiting
\(\Gamma\) has input marginal \(dz\,\Pi_z\). Its conditional laws
have the prescribed independent input law for almost every \(z\);
this assertion concerns the random measure itself, not just its mean.

Translate a scene by \(-h\) in both transverse coordinates, writing
this operation as \(\tau_h\). The exact finite identity is
\[
\tau_h\omega_{n,z,\sigma}
 =\omega_{n,z+\varepsilon_nh,\sigma}.
\]
Changing root variables in \eqref{eq:highd-scene} shows that
\(\Gamma_n(f(z)[G(\tau_h\omega)-G(\omega)])\) tends to zero
for bounded continuous \(G\) and continuous \(f\): the error is
controlled by the modulus of continuity of \(f\) and the volume
of \(D\mathbin\triangle(D+\varepsilon_nh)\).
The bound is deterministic and uniform in the Birkhoff weights.
A countable family of tests and rational \(h\), followed by
continuity of translations and disintegration, proves stationarity
of \(P_z^\Gamma\).

The smooth distance expansion is uniform on the compact radial
intervals and bounded rescaled endpoint windows. Hence
\(\varepsilon_n^{-2}e(a,z+\varepsilon_nx;
b,z+\varepsilon_ny)\) converges locally uniformly to
\eqref{eq:highd-local-cost}. Every finite permutation in the mixture
is optimal for \(e\). A strict violation of a finite cyclic inequality
in a limiting scene would persist on finitely many small endpoint
neighborhoods, with a positive cost margin. Vague convergence of the
integer edge measures and uniform cost convergence would then violate
the corresponding finite-sample inequality. Exhausting a countable
base of bounded neighborhoods proves local cyclic optimality.
Moreover, for nonnegative compactly supported \(\eta\), the energy
bound and root integration give
\[
\E\int\Gamma_n(dz,d\omega)
 \sum_{\text{edges}}\eta(x)|y-x|^2\le K\int\eta.
\]
Lower semicontinuity on bounded endpoint windows, then monotone
convergence, preserves this bound in the limit. Campbell's formula
gives a finite Palm second moment for almost every conditional law.
This proves membership in \(\mathcal L_z\).

It remains to connect the random scene measure to the empirical plan.
Put \(\rho_n=n^{-1}\sum_{ij}r_{ij}\delta_{(Z_i,A_i,B_j)}\).
Fix a continuous \(f(z,a,b)\) and nonnegative \(\eta\in C_c(\R^q)\)
with integral one, and set
\(F_f(z,\omega)=\sum_{\text{edges}}\eta(x)f(z,a,b)\).
The exact finite representation gives
\begin{align*}
\Gamma_n(F_f)=\frac1n\sum_{ij}r_{ij}\int\eta(v)
 &\mathbf1_D(Z_i-\varepsilon_nv)\\
 &\cdot f(Z_i-\varepsilon_nv,A_i,B_j)\,dv.
\end{align*}
If \(\eta\) is supported in \(B_R\), its difference from
\(\rho_n(f)\) is bounded by the modulus of continuity of \(f\)
at \(\varepsilon_nR\) and a constant times the empirical source
mass within \(\varepsilon_nR\) of \(\partial D\). The latter
tends to zero in expectation because \(\theta\) is bounded.

To pass \(F_f\) through the random-measure limit, first use a continuous cutoff of
edges of displacement greater than \(L\). The mean error is
\(O(L^{-2})\). The remaining edge test has compact support in both
endpoints and is continuous for vague convergence. Its absolute value
is bounded by a constant times a local source count, whose second
moment is uniformly bounded. Bounded clipping, followed by removal
of the clip and then of the displacement cutoff, therefore passes
its integral to the limit. For a countable determining family of
\(f\), this proves almost surely
\(\rho(f)=\Gamma(F_f)\) in every joint distributional limit.
Campbell's formula identifies \(\rho\) as
\eqref{eq:highd-reconstruction}. Finally
\eqref{eq:highd-macro-collapse} replaces the target transverse label
by the source label in the physical plan, completing reconstruction.
\end{proof}

\begin{corollary}[Selection from cyclic-law uniqueness]
\label{cor:highd-uniqueness-criterion}
Suppose that, for almost every \(z\), all laws in \(\mathcal L_z\)
have the same paired-mark law \(\kappa_z\). Then independent
empirical sampling in the present regular tube satisfies
\[
\sup_{\pi\in\mathcal O_n^F}d_{\rm BL}(\pi,\pi_*)
 \xrightarrow{\Pp}0,
\]
where \(\pi_*\) is the pushforward of
\(\theta(z)dz\,\kappa_z(da,db)\) under
\((z,a,b)\mapsto(X(a,z),X(b,z))\).
\end{corollary}

\begin{proof}
A measurable version of \(\kappa_z\) is obtained by taking the
mean of any subsequential \(\Gamma\), disintegrating it, and using
the Campbell formula with one fixed \(\eta\). Its conditional laws
belong to \(\mathcal L_z\) for almost every \(z\): stationarity,
the prescribed input law, local optimality and the integrated finite
moment bound are preserved by this averaging. The uniqueness
hypothesis makes every random reconstructed limit equal to this
same deterministic \(\pi_*\). Tightness therefore gives convergence
in probability for every measurable choice of empirical optimizer.

Finite Birkhoff decomposition and convexity of \(d_{\rm BL}\) in
its first argument give
\[
\sup_{\pi\in\mathcal O_n^F}d_{\rm BL}(\pi,\pi_*)
 =\max_{\sigma\in\mathcal S_n}d_{\rm BL}(\pi_\sigma,\pi_*).
\]
Select a maximizing permutation measurably by a fixed finite
enumeration, before introducing the independent root. Failure of the
displayed maximum to converge to zero in probability would contradict
the convergence just proved for that selected sequence.
\end{proof}

Uniqueness over the whole cyclic class \(\mathcal L_z\) is a
stronger hypothesis than the proved regular-tube selection needs.
Propositions~\ref{prop:highd-stationary-pressure}
and~\ref{prop:highd-variable-pressure} place the relevant limits
in the stationary mean-minimizing class, where
Theorem~\ref{thm:highd-pressure-differentiability} identifies one
paired law. This proves the regular ordered results above. Stationary
mean minimality is not asserted here for arbitrary weakly continuous
conditional kernels or for all cyclic laws.

\section{Measurable conditional laws and forward pressures}\label{sec:measurable-forward-pressure}

\subsection{Weak mark stability and measurable conditional populations}
\label{subsec:highd-measurable-pressure}

The common-latent continuity assumption in
\eqref{eq:highd-conditional-representations} can be removed in transverse
dimension \(q>4\). The proof uses fixed microscopic blocks and a direct
repair of their count imbalances. We first record a stability statement
for the homogeneous pressure, valid already for \(q\ge3\).

\begin{lemma}[Weak continuity of the homogeneous ordered pressure]
\label{lem:highd-weak-mark-pressure}
Let \(K_A,K_B\) be compact metric mark spaces, and let
\(\alpha_j\Rightarrow\alpha\), \(\beta_j\Rightarrow\beta\).
Suppose the continuous compliance matrices \(S_j\) converge uniformly
to \(S\), with common positive ellipticity bounds, and the continuous
paired probes \(g_j\) converge uniformly to \(g\).
Assume every model, including the limit, satisfies the common-cumulative
ordered hypotheses of Proposition~\ref{prop:highd-stationary-pressure};
in particular every mark pair is admissible. For \(q\ge3\) and
\(t_j\to t\), their unit-intensity pressures satisfy
\[
                    p_j(t_j)\longrightarrow p(t).
\]
For fixed varying model data, convergence is locally uniform in the
real pressure parameter. The mark laws may have atoms.
\end{lemma}
\begin{proof}
Choose a stationary mean minimizer \(P\) for the limiting model.
Couple \(\alpha_j\) to \(\alpha\), and \(\beta_j\) to \(\beta\),
by Skorokhod representations with almost sure convergence of the
marks. Disintegrate the joint sequences over their limiting marks,
and independently attach the corresponding conditional sequences
to every input point. The independent marking construction preserves
stationarity and gives, for each \(j\), the required independent
marked Poisson inputs. Keep every spatial edge of \(P\).
All mark pairs remain admissible. Its costs converge almost surely
and are bounded in absolute value by \(C(1+|y-x|^2)\), integrable
under the fixed Palm law \(P\). The stationary variational formula
and dominated convergence therefore give
\[
                  \limsup_j p_j(t_j)\le p(t).
\]

For the reverse inequality, stationary minimizers \(P_j\) have a
common Palm second-moment bound: compare with one unmarked
finite-energy stationary matching and use ellipticity and bounded
probes. The marked inputs converge locally. The two-endpoint
compactness argument of
Proposition~\ref{prop:highd-stationary-pressure} retains both endpoint
projections; source- and target-anchored escape probabilities are
\(O(R^{-2})\). Every limiting law is consequently a stationary
perfect matching of the limiting inputs. Compact truncations of
\(k_j+t_jg_j+C\ge0\), uniform coefficient convergence, and lower
semicontinuity yield
\[
                 p(t)\le\liminf_j p_j(t_j).
\]
The common probe bound gives equi-Lipschitz dependence on \(t\),
which upgrades convergence to locally uniform convergence.
No uniform integrability of the displacements of varying optimizers
is needed for the upper comparison: that comparison used one fixed
stationary matching.
\end{proof}

Here is the extension with measurable conditional populations.
Let \(D=[0,1]^q\), let \(\theta\) be a bounded nonnegative measurable
probability density on \(D\), and let \(\alpha_z,\beta_z\) be Borel
probability kernels on fixed compact mark spaces \(K_A,K_B\).
The support of \(\theta\) may have an irregular boundary and
\(\theta\) need not be bounded below. Every mark pair is admissible.
Let \(S(a,b,z)\) be continuous and uniformly positive definite,
with each frozen model satisfying the common-cumulative ordered
hypotheses of Proposition~\ref{prop:highd-stationary-pressure}.
For example, this includes the compact ordered radial mark intervals
and \(S(a,b,z)=\int_a^b H(r,z)^{-1}\,dr\) used above.
Assume the nonnegative actual cost \(c\) satisfies, on its supported
endpoint pairs,
\begin{align}
\label{eq:highd-measurable-comparison}
 c_0|w-z|^2&\le c(a,z;b,w)\le C_0|w-z|^2,\\
\label{eq:highd-measurable-expansion}
 \left|c(a,z;b,w)-\tfrac12(w-z)^TS(a,b,z)^{-1}(w-z)\right|
       &\le\omega(|w-z|)|w-z|^2,
       \qquad \omega(r)\longrightarrow0.
\end{align}
The constants are uniform in marks and labels. In a genuine ordered
ray chart, \(c=e\) is the calibrated length defect; its removed
endpoint sum does not change any optimizer. The kinetic cost itself
also satisfies these assumptions.

For independent iid color samples with laws
\(\theta(z)\,dz\,\alpha_z(da)\) and
\(\theta(w)\,dw\,\beta_w(db)\), use
\(V_n^\Psi(t)\) from \eqref{eq:highd-variable-finite-pressure}, with
\(\Psi\in C(D\times K_A\times K_B)\). Let \(p_{z,g}\) denote
its frozen unit-intensity homogeneous pressure.

\begin{proposition}[Ordered pressure with measurable conditional laws]
\label{prop:highd-measurable-pressure}
Under the preceding assumptions, if \(q>4\), then
\[
 \E\sup_{|t|\le T}|V_n^\Psi(t)-P_\Psi(t)|\longrightarrow0
 \qquad(T<\infty),
\]
where
\begin{equation}
\label{eq:highd-measurable-pressure}
 P_\Psi(t)=\int_D\theta(z)^{1-2/q}
       p_{z,\Psi(z,\cdot,\cdot)}(t\theta(z)^{2/q})\,dz,
\end{equation}
and the integrand is zero when \(\theta(z)=0\).
For each fixed \(t\), every subsequential random conditional scene
law of optimizers is almost surely, at almost every positive-intensity
root, a stationary mean minimizer for its frozen intensity-\(\theta(z)\)
model. No continuity or common-latent representation of the conditional
mark laws is assumed.
\end{proposition}
\begin{proof}
Fix \(L\ge1\), put
\[
 r_n=\lceil n^{1/q}/L\rceil,\qquad
 \ell_n=r_n^{-1},\qquad L_n=n^{1/q}\ell_n\longrightarrow L,
\]
and partition \(D\) into \(r_n^q\) cubes. Retain a
maximum-cardinality optimum within each leaf cube. At successive
dyadic parent scales, match opposite-colored residuals arbitrarily.
Embed the grid into the next larger dyadic grid if necessary, with
empty padding. At the root the two total counts equal \(n\), so this
constructs a perfect matching.

For any cell \(C\), write \(\Delta(C)=N_A(C)-N_B(C)\).
After all its descendants have been repaired, exactly
\(|\Delta(C)|\) points remain, all of one color. Equal transverse
marginals and independent colors imply
\begin{equation}
\label{eq:highd-measurable-count}
 \E|\Delta(C)|
 \le\left(2n\int_C\theta\right)^{1/2}
 \le\bigl(2n\|\theta\|_\infty|C\cap D|\bigr)^{1/2}.
\end{equation}
At parent scale \(h=2^k\ell_n\), the expected number of repair
edges is therefore at most
\(C_q\sqrt{n\|\theta\|_\infty}\,h^{-q/2}\).
Every such edge has squared length at most \(qh^2\). By
\eqref{eq:highd-measurable-comparison}, the total expected normalized
kinetic repair cost is at most
\begin{equation}
\label{eq:highd-measurable-repair}
 C\sum_{k\ge1}(2^kL_n)^{2-q/2}
             \le C_q L_n^{2-q/2},
\end{equation}
whereas its absolute probe contribution is at most
\(C_TL_n^{-q/2}\). The sum in
\eqref{eq:highd-measurable-repair} converges because \(q>4\).
These estimates use exact counts only; the adaptively retained
residual positions are not claimed to be iid. All mark pairs must
be admissible for this repair. Starting at microscopic unit scale
also supplies a uniform quadratic comparator and hence a uniform
expected second-moment bound for every optimizer, on bounded
\(t\)-intervals.

Let \(b_L(z,t)\) be \(L^{-q}\) times the expected
maximum-cardinality optimum in \([0,L]^q\) for two independent
Poisson clouds of intensity \(\theta(z)\), with the frozen marks,
kinetic coefficient, and probe at \(z\). This is measurable in
\(z\): conditional on counts, a finite assignment value is the
minimum of finitely many continuous functions, integrated against
Borel kernels.

At almost every \(z\), Lebesgue differentiation holds simultaneously
for \(\theta\) and for
\[
 w\longmapsto\theta(w)\int f\,d\alpha_w,
 \qquad
 w\longmapsto\theta(w)\int g\,d\beta_w,
\]
for countable dense families in \(C(K_A)\) and \(C(K_B)\).
Let \(C_n(z)\) be its leaf cube. These containing cubes form a
bounded-eccentricity differentiation basis. After rescaling the cube
to side \(L_n\), its input converges to the frozen independent
marked Poisson input. To see the product position/mark law, test by
a continuous position function times a mark function: the error is
bounded by the cell average of the absolute error in the displayed
differentiated functions. The binomial rare-event limit gives the
counts, and the two colors remain independent. At zero-density
roots the limiting input is empty.

On a bounded-count configuration, continuity of \(S\) and \(\Psi\)
and \eqref{eq:highd-measurable-expansion} give uniform convergence of
the rescaled finite cost array to the frozen one. Finite assignment
values are continuous even with atomic mark laws. The absolute leaf
value is bounded by \(C_{L,T}(N_A(C_n)+N_B(C_n))\), and these counts
have uniformly bounded second moments for fixed \(L\).
Consequently
\[
 \frac{\E[\text{leaf optimum in }C_n(z)]}{n|C_n(z)|}
                         \longrightarrow b_L(z,t)
 \quad\text{for almost every }z.
\]
The bound is uniform in \(z\), since \(\theta\) is bounded.
Integrating over roots, and then using the repair estimate, gives
\begin{equation}
\label{eq:highd-measurable-upper-block}
 \limsup_n\E V_n^\Psi(t)
 \le\int_D b_L(z,t)\,dz+C L^{2-q/2}+C_TL^{-q/2}.
\end{equation}

Homogeneous independent-count pressure convergence and intensity
scaling give
\[
 b_L(z,t)\longrightarrow
 \theta(z)^{1-2/q}
 p_{z,\Psi(z,\cdot,\cdot)}(t\theta(z)^{2/q}).
\]
The same dyadic comparator, now for homogeneous Poisson inputs,
bounds \(b_L\) uniformly for \(L\ge1\) and bounded \(t\), since
\(\theta\) is bounded. Its lower bound is
\(-|t|\|\Psi\|_\infty\theta(z)\). Dominated convergence in
\eqref{eq:highd-measurable-upper-block}, followed by \(L\to\infty\),
therefore gives
\begin{equation}
\label{eq:highd-measurable-upper}
                    \limsup_n\E V_n^\Psi(t)\le P_\Psi(t).
\end{equation}
This also proves measurability of the integrand in
\eqref{eq:highd-measurable-pressure}.

For the lower bound, use the uniform second moment and the
input-only measurable-kernel Poisson argument to extract joint
subsequential limits \((V,\Gamma)\) of values and random scenes.
The two-endpoint compactness argument of
Proposition~\ref{prop:highd-local-limit} gives
\(\Gamma(dz,d\omega)=dz\,P_z^\Gamma(d\omega)\), whose
positive-intensity conditional laws are stationary perfect matchings
of the prescribed frozen inputs, almost surely.
The measurable-kernel input identification uses only the preceding
Lebesgue-point argument, as in
Proposition~\ref{prop:rough-ray-poisson}.

Here the lower cost comparison holds for each realized limiting
scene. To make this precise, choose a nonnegative compactly supported
continuous anchor \(\eta\) with integral one. Put
\(B=T\|\Psi\|_\infty\), use increasing compact endpoint cutoffs
\(\chi_R\), and clip the nonnegative local sum
\[
 \sum_{\text{edges }(x,a;y,b)}
  \eta(x)\chi_R(x,y)[k_z((x,a),(y,b))+t\Psi(z,a,b)+B]
\]
at height \(N\). Denote this local functional by \(F_{R,N}\).
For the corresponding actual finite-sample sum, integration over
macro roots is at most \(V_n^\Psi(t)+B\): each edge has anchor
integral at most one, and cutoffs and clipping can only decrease the
sum. On fixed cutoffs, the relative expansion and continuity of
coefficients and probes, together with local count bounds, give
\[
       \Gamma_n(F_{R,N})\le V_n^\Psi(t)+B+o_{\mathbb P}(1).
\]
Pass jointly to the limit, and then exhaust the cutoffs and clips.
The Campbell identity and \(\int_D\theta=1\) imply, almost surely,
\begin{equation}
\label{eq:highd-measurable-pathwise-lower}
 V\ge\int_D\theta(z)\,
       \E_{P_z^\Gamma}^{\rm Palm}[k_z+t\Psi(z,a,b)]\,dz
   \ge P_\Psi(t).
\end{equation}
The last inequality is the stationary variational formula with
intensity scaling. No convergence of untruncated energies is assumed.

The values have a deterministic lower bound \(-B\) and are tight
above by their expectation bound. Thus
\eqref{eq:highd-measurable-pathwise-lower} for every joint limit gives
\(\liminf_n V_n^\Psi(t)\ge P_\Psi(t)\) in probability.
Together with \eqref{eq:highd-measurable-upper} and the lower bound,
this implies \(\E V_n^\Psi(t)\to P_\Psi(t)\). Indeed, the negative
part relative to \(P_\Psi(t)\) is bounded and vanishes in probability;
hence
\[
 \E|V_n^\Psi(t)-P_\Psi(t)|
 =\E V_n^\Psi(t)-P_\Psi(t)
       +2\E[(P_\Psi(t)-V_n^\Psi(t))_+]\longrightarrow0.
\]
Finite meshes and the common \(\|\Psi\|_\infty\)-Lipschitz bound
in \(t\) give expected locally uniform convergence. Equality in
\eqref{eq:highd-measurable-pathwise-lower} then forces the
nonnegative conditional excess over the stationary infimum to vanish
almost everywhere, proving local mean minimality.

For an optimal fractional plan, use a data-measurable Birkhoff mixture
supported on optimal permutations and draw the permutation before
introducing the independent macro root. The scene argument applies
to that integer matching realization; it does not regard a fractional
edge measure as an integer matching.
\end{proof}

\begin{corollary}[Ordered selector with measurable conditional laws]
\label{cor:highd-measurable-selector}
Under Proposition~\ref{prop:highd-measurable-pressure}, the entire
unperturbed empirical optimal-plan set converges in probability to the
deterministic coordinate plan
\[
           \theta(z)\,dz\,\kappa_z(da,db),
             \qquad \kappa_z(g)=p_{z,g}'(0).
\]
The conclusion is uniform over permutation and fractional optimizers.
For a genuine ordered length defect it gives the corresponding
true-distance conclusion after the chart pushforward.
\end{corollary}
\begin{proof}
Theorem~\ref{thm:highd-pressure-differentiability} applies to every
frozen compact mark law. For a spatial probe
\(\Psi(z,a,b)=\phi(z)g(a,b)\), bounded slopes permit differentiation
of \eqref{eq:highd-measurable-pressure} at zero, giving
\[
                P_\Psi'(0)
        =\int_D\theta(z)\phi(z)p_{z,g}'(0)\,dz.
\]
The finite-difference sandwich fixes this observable simultaneously
for all unperturbed optimizers. A countable dense family of spatial
and paired-mark tests, the second-moment no-escape bound, and the
existing reconstruction determine the full same-label plan.
A deterministic measurable version of \(\kappa_z\) is obtained from
an annealed limiting scene; differentiation fixes its tests in every
other limiting phase. This is convergence in probability of the full
optimal set, including fractional couplings by Birkhoff decomposition.
\end{proof}

The dimension restriction here is \(q>4\), or \(d\ge6\); it comes
from \eqref{eq:highd-measurable-repair}, and is not asserted necessary
for selection. The result permits measurable conditional laws and
irregular supported label sets within its ordered core. Applying it
to a rough ray core requires the two-sided comparison and uniform
relative expansion stated above. Arbitrary pairings in the dyadic
repair are not admissible for general interleaving populations, and
this subsection supplies neither their sharp pressure recovery nor
a global countable-atlas selector.

\subsection{A sharp local feasibility and energy criterion}
\label{subsec:forward-poisson-criterion}

The finite color-run reduction has a direct geometric-probability
consequence. It uses the existing homogeneous quadratic tangent model,
without changing the original distance objective.

\begin{proposition}[Forward Poisson matching: existence and finite energy]
\label{prop:forward-poisson-criterion}
Let \(q\ge1\), and let \(\alpha,\beta\) be probability laws on compact
disjoint subsets of \(\R\) with positive separation. Take independent
unit-intensity Poisson source and target processes on \(\R^q\), with
independent marks of laws \(\alpha,\beta\). A stationary perfect matching
whose paired marks satisfy \(a<b\) exists if and only if
\[
 F_\alpha(t)\ge F_\beta(t)\quad\text{for all }t.
\]
Suppose the quadratic edge cost is
\(\frac12(w-z)^T A(a,b)(w-z)\), with measurable symmetric coefficients
\(\lambda I\le A(a,b)\le\Lambda I\) on allowed pairs, for fixed
\(0<\lambda\le\Lambda<\infty\). There exists such a matching with finite mean
quadratic cost per unit volume if and only if the same Hall condition
holds and \(q\ge3\).
Equivalently, the infimum \(I_+(t)\) of its source-Palm mean cost plus
\(t g(a,b)\), for any bounded measurable \(g\) and fixed real \(t\),
is finite exactly under these conditions. In the sufficient case there
is a comparison matching with
\(\mathbb P^0\{|w-z|>r\}\le e^{-c r^q}\) for some \(c>0\).
\end{proposition}
\begin{proof}
Separate the finitely many alternating source and target color runs,
omitting zero-mass runs. Denote their masses by \(a_i\), \(1\le i\le k\),
and \(b_j\), \(1\le j\le l\). Independent Poisson marking gives colors
with intensity vector \(\xi=(a_1,\ldots,a_k,b_1,\ldots,b_l)\).
Allow the family vector \(e_i+e_{k+j}\) exactly when the whole source
run \(i\) precedes target run \(j\); call the finite set of these
vectors \(V\). Membership \(\xi\in\operatorname{cone}(V)\) is exactly
the existence of a nonnegative allowed traffic matrix with row masses
\(a_i\) and column masses \(b_j\). A forward quantile coupling gives
such a matrix when the Hall condition holds. Conversely, replacing
each matrix entry by a product of its normalized run laws gives a
forward coupling, whose cut-flux identity implies that condition.

For any stationary perfect matching, mass transport supplies the same
traffic matrix of edge intensities, proving necessity. Conversely,
\cite[Theorem~1]{AmirAngelHolroyd} gives a stationary allowed matching
whenever \(\xi\) belongs to this cone. Restore the finer marks
independently with their conditional run laws. This preserves the
independent marked Poisson input laws and all forward constraints.

Every vector in \(V\) satisfies total-source count equal to total-target
count. Its cone therefore lies in a proper linear hyperplane: every
feasible intensity is on its ambient boundary. The critical cases of
\cite[Theorem~1(ii),(iii)]{AmirAngelHolroyd} give an exponential
\(r^q\)-tail when \(q\ge3\), whereas for \(q=1,2\) every stationary
matching has infinite \(q/2\)-th displacement moment. In a bipartite
pair matching the source- and target-Palm displacement laws agree by
mass transport. Thus those typical-family-diameter assertions apply
to the source-Palm displacement here. The upper ellipticity bound
gives finite quadratic cost in the first case. The lower bound forces
infinite quadratic cost in the second, since an infinite moment of
order at most one forces an infinite second moment. A bounded probe
changes neither conclusion and bounds the infimum below by
\(-|t|\|g\|_\infty\).
\end{proof}

Strict Hall slack at internal cuts does not turn this into an interior
cone problem: the total bipartite balance remains exact. On each
regular compact ray core the uniformly elliptic second variation has
\(A=S^{-1}\), so the proposition supplies a finite-energy stationary
forward competitor for \(d=q+1\ge4\). It does not give a sharp
finite-box recovery sequence or tail bounds for minimizing matchings.
The criterion concerns this local stationary problem, not existence
of a deterministic empirical selector on the manifold.

\begin{proposition}[Disconnected populations and uncentered pressure]
\label{prop:disconnected-pressure-obstruction}
For \(d=q+1\), \(q=3,4\), there are smooth compactly supported
Lipschitz densities on Euclidean space, and on a flat torus, with
disjoint supports and globally ordered product ray laws, for which
the uncentered microscopic calibrated pressure fails to converge to a
finite deterministic constant. For \(q=3\) it diverges to infinity in
probability. This conclusion alone does not preclude deterministic
convergence of the empirical plans.
\end{proposition}
\begin{proof}
Take \(\mu=\alpha\otimes\theta\), \(\nu=\beta\otimes\theta\), with
smooth probability densities \(\alpha,\beta\) in separated compact
intervals \(I_A<I_B\), and
\(\theta=p\theta_1+(1-p)\theta_2\), \(0<p<1\). The two smooth
transverse bumps have compact supports \(C_1,C_2\) at distance
\(h>0\). The potential \(\phi(r,z)=r\) calibrates the population
problem. For independent samples of size \(n\), their counts \(K_n,L_n\)
in \(C_1\) are independent \(\operatorname{Bin}(n,p)\) variables.
Every assignment has at least \(|K_n-L_n|\) cross-component edges.
For \(\Delta=b-a>0\), its defect is
\[
 e((a,z),(b,w))=
 \frac{|w-z|^2}{\sqrt{\Delta^2+|w-z|^2}+\Delta}.
\]
Compactness bounds the denominator by \(C_D<\infty\), so every such
edge costs at least \(\kappa=h^2/C_D>0\). Thus the minimum total
defect \(E_n\), including fractional assignments by mass balance,
satisfies
\[
 E_n\ge\kappa|K_n-L_n|,\qquad
 \frac{K_n-L_n}{\sqrt n}\ \Rightarrow\
 Z\sim N(0,2p(1-p)).
\]
Write \(P_n=n^{2/q-1}E_n\). When \(q=3\), for each fixed \(M\) and
each \(\eta>0\), eventually
\(\mathbb P\{P_n\le M\}\le
\mathbb P\{|K_n-L_n|/\sqrt n\le\eta\}\).
The limiting upper bound tends to zero with \(\eta\), proving
divergence in probability. When \(q=4\), for every \(M>0\),
\[
 \liminf_n\mathbb P\{P_n>M\}
 \ge\mathbb P\{\kappa|Z|>M\}>0.
\]
Choosing \(M\) above any proposed finite constant rules out convergence
to that constant. Adding a fixed bounded probe cannot remove either
obstruction. The same construction fits in a flat-torus coordinate
box with span less than half the period in each coordinate. Every
relevant endpoint distance is then Euclidean. The local coordinate
\(r\) is 1-Lipschitz on the endpoint set and has a global McShane
extension, providing the calibration without a global smooth
eikonal function.
\end{proof}

The forced cross-component fraction is only of order \(n^{-1/2}\).
It can disappear from a weak plan limit while dominating the
uncentered pressure. A centered pressure or a pressure-difference
argument would require a separate cancellation proof. Thus the example
does not justify a topological restriction on the selector; it exposes
a population-dependent obstruction to this particular energy method.

\subsection{Weak stability of the forward stationary value}
\label{subsec:forward-weak-stability}

Fix \(q\ge3\) and a finite ordered family of nonempty compact
classes \(K_1,\ldots,K_N\subset\R\), with
\(\max K_h<\min K_{h+1}\), each assigned source or target color.
Let \(K_A,K_B\) be their respective unions, and put
\[
 \mathcal E=\{(h,k):h<k,\ K_h\text{ source},\ K_k\text{ target}\},
 \qquad \mathcal A=\bigcup_{(h,k)\in\mathcal E}K_h\times K_k.
\]
Zero-probability classes remain in this fixed partition. A pair of
probability laws \((\alpha,\beta)\) on \(K_A,K_B\) is called
feasible when it has a coupling supported on \(\mathcal A\).
Equivalently, the source-minus-target mass is nonnegative at every
gap between consecutive classes, by
Proposition~\ref{prop:ray-forward-criterion}.

For a continuous symmetric field \(A\) on \(\mathcal A\), with
\(cI\le A\le CI\), and \(g\in C(\mathcal A)\), define
\begin{equation}
\label{eq:forward-weak-value}
 I_+(t;\alpha,\beta,A,g)
 =\inf_P\E_P^{\rm Palm}
       \left[\tfrac12 v^TA(a,b)v+t g(a,b)\right].
\end{equation}
Here \(P\) ranges over stationary forward perfect matchings of
independent unit-intensity marked Poisson inputs, \(v=y-x\), and
Palm expectation is source-anchored. Auxiliary randomness is allowed.

\begin{proposition}[Weak continuity across tight Hall cuts]
\label{prop:forward-weak-continuity}
Let \(\alpha_j\Rightarrow\alpha\), \(\beta_j\Rightarrow\beta\),
with every pair feasible. Suppose \(A_j\to A\) and \(g_j\to g\)
uniformly on \(\mathcal A\), with all fields continuous and all
matrices satisfying the same bounds \(cI\le A_j,A\le CI\).
The infima in \eqref{eq:forward-weak-value} are finite and attained
by finite-second-moment matchings. If \(t_j\to t\), then
\[
 I_+(t_j;\alpha_j,\beta_j,A_j,g_j)
          \longrightarrow I_+(t;\alpha,\beta,A,g).
\]
With a common variable \(t\), convergence is locally uniform.
Atoms, zero limiting class masses, and tight Hall cuts are permitted.
\end{proposition}
\begin{proof}
Write \(s=1-2/q>0\). We first give a uniform reserve estimate.
For feasible class intensities \((r_h)\), with total mass \(\rho\)
in each color, choose a nonnegative traffic matrix \((m_{hk})\)
on \(\mathcal E\) with the prescribed row and column sums.
For each positive entry use an independent copy of a fixed stationary
two-color unit-intensity matching, scaled by \(m_{hk}^{-1/q}\).
Such a scheme has finite Palm second moment \(C_q\), by
Proposition~\ref{prop:forward-poisson-criterion}, or
\cite[Theorem~1(iii)]{AmirAngelHolroyd} in the two-color case.
Attach independent marks with the prescribed conditional class laws.
Independent superposition produces exactly the required independent
marked Poisson inputs and a legal perfect matching. Its squared
displacement cost per volume is at most
\begin{equation}
\label{eq:forward-weak-reserve}
 C_q\sum_{(h,k)\in\mathcal E}m_{hk}^{s}
       \le C_q|\mathcal E|^{2/q}\rho^s.
\end{equation}
Indeed, intensity contributes \(m_{hk}\), squared distance contributes
\(m_{hk}^{-2/q}\), and concavity gives the final bound.
Zero entries are omitted. Thus the reserve's perturbed cost is at most
\((C C_q/2)|\mathcal E|^{2/q}\rho^s+|t|\|g\|_\infty\rho\),
uniformly even at the boundary of the feasible traffic cone.

Taking \(\rho=1\) gives a uniform finite competitor. Ellipticity
and the bounded probe bound the Palm second moment of every
bounded-cost minimizing sequence. The two-endpoint stationary
compactness argument in Proposition~\ref{prop:highd-stationary-pressure}
therefore applies: both anchored escape probabilities are
\(O(R^{-2})\), the target bound following by mass transport;
local Poisson count bounds retain the full integer endpoint
projections. The limit is stationary and perfect. The allowed set
\(\mathcal A\) is closed, so legality survives. Adding
\(|t|\|g\|_\infty\) makes the cost nonnegative, and compact
truncation gives lower semicontinuity and attainment. This compactness
step uses no common-cumulative representation of \(A\).

For the upper continuity bound fix a limiting minimizer \(P\).
Let \(u_h,u_{j,h}\) be the limiting and new class masses. Each
class is clopen in its same-color compact space, so \(u_{j,h}\to u_h\).
For a class cut \(\ell<N\), write
\[
 H_\ell(u)=\sum_{\substack{h\le\ell\\K_h\text{ source}}}u_h
           -\sum_{\substack{h\le\ell\\K_h\text{ target}}}u_h.
\]
For sufficiently large \(j\), set
\begin{equation}
\label{eq:forward-weak-residual}
\begin{aligned}
 \rho_j=\max\bigg\{&\frac1j,\quad
    \max_{h:u_h>0}\frac{(u_h-u_{j,h})_+}{u_h},\\
 &\max_{\ell:H_\ell(u)>0}
       \frac{(H_\ell(u)-H_\ell(u_j))_+}{H_\ell(u)}\bigg\},\\
 \lambda_j&=1-\rho_j,\qquad r_{j,h}=u_{j,h}-\lambda_j u_h.
\end{aligned}
\end{equation}
Empty maxima are zero. Then \(\rho_j\to0\), eventually
\(0<\rho_j<1\), and \(r_{j,h}\ge0\). Each residual color has
total mass \(\rho_j\). At a positive limiting Hall cut, its defining
ratio gives
\[
 H_\ell(r_j)=H_\ell(u_j)-H_\ell(u)+\rho_jH_\ell(u)\ge0.
\]
At a tight limiting cut this is simply \(H_\ell(u_j)\ge0\).
Thus the residual is feasible, including new mass in a previously
empty class or new flow across a previously tight cut.

Within every class of positive limiting mass, the new normalized
conditional laws converge weakly to the limiting conditional law.
Take Skorokhod representations of these sequences, disintegrate
over the limiting mark, and attach the resulting sequence kernels
independently at every input point of \(P\), separately for the two
colors. This is stationary independent marking. The new inputs have
the old class masses and the new conditional class laws. Keep every
old edge: its class rectangle remains legal. Almost sure convergence
holds at every point of the countable input, hence also at its
matching-selected partner.

Dilate the re-marked main process by \(\lambda_j^{-1/q}\).
Its class intensities become \(\lambda_j u_h\), and its new cost
per volume is
\begin{equation}
\label{eq:forward-weak-main-cost}
 \lambda_j^s\E_P^{\rm Palm}
       [\tfrac12 v^TA_j(a_j,b_j)v]
       +\lambda_j t_j\E_P^{\rm Palm}[g_j(a_j,b_j)].
\end{equation}
Continuity and uniform coefficient convergence give pointwise
convergence of these integrands. The kinetic part is dominated by
\(C|v|^2/2\) under the one fixed finite-energy law \(P\), and
the probes are uniformly bounded. Dominated convergence therefore
makes \eqref{eq:forward-weak-main-cost} tend to the limiting value.
Add an independent reserve matching with intensities \(r_{j,h}\)
and the new conditional class laws. Its cost tends to zero by
\eqref{eq:forward-weak-reserve}. In each class the superposed mark
intensity is exactly \((\lambda_j u_h+r_{j,h})\) times that
conditional law, namely the required new intensity measure.
Previously empty classes are supplied wholly by the reserve.
The two colors remain independent PPPs, and the union of their
matchings is stationary and legal. This proves the limsup inequality.

For the liminf choose new-model minimizers. The uniform competitor
and coercivity bound their Palm second moments uniformly. The same
two-endpoint compactness argument gives a stationary perfect legal
matching of the limiting independent marked inputs; input convergence
follows also directly from the Poisson Laplace functionals.
Choose \(B\) bounding all \(|t_j|\|g_j\|_\infty\) and
\(|t|\|g\|_\infty\). The nonnegative shifted edge costs converge
uniformly on bounded displacements. Compact endpoint cutoffs and
input-count clipping, followed by monotone exhaustion, yield
\[
 \E_{P_*}^{\rm Palm}[\tfrac12v^TA(a,b)v+t g(a,b)]
       \le\liminf_j I_+(t_j;\alpha_j,\beta_j,A_j,g_j).
\]
The shift \(B\) cancels because all source intensities are one.
The limiting infimum bounds the left side below, completing the proof
of convergence. Finally all values are Lipschitz in \(t\), with
the common bound on \(\|g_j\|_\infty\); a finite parameter net
gives local uniformity.
\end{proof}

For fixed \(A,g\), the feasible mark-law set is weakly compact:
the finitely many Hall inequalities are closed on this fixed clopen
class partition. Proposition~\ref{prop:forward-weak-continuity}
therefore makes \(I_+\) continuous, hence Borel, on the entire feasible
set. Measurable feasible conditional laws consequently define a
measurable frozen stationary value, including at tight cuts.
This conclusion concerns stationary values. It supplies neither
finite-sample forward repair nor a finite-volume pressure formula for
arbitrary coefficient fields; any such identification requires its
own recovery argument.

% Complete conversion of accepted sources:
% interleaving-pressure-strict-slack.md
% SHA256 258ede3b79e629204f4d97ef1caea787825e8b54589dbd5b7f17b946ce5e274c
% interleaving-recovery.md
% SHA256 cc2cfcc79c5cb30b5a00d62b7814f0ebeebabb1087bc8216db2854b8c85fc173
\subsection{Forward stationary pressure and differentiability}
\label{forward-full-section}

Let \(q\ge3\). The source and target mark laws \(\alpha,\beta\)
have compact supports in \(\R\), separated by \(\delta>0\), with
finitely many ordered monochromatic classes. Assume
\(F_\alpha\ge F_\beta\), so a forward population coupling exists;
tight Hall cuts are permitted. On the scalar span let \(T\) be a
\(C^1\) symmetric matrix function satisfying
\(hI\le T'(r)\le HI\), where \(0<h\le H<\infty\).
For allowed pairs \(\mathcal F=\{a<b\}\), put
\[
 S(a,b)=T(b)-T(a),\qquad
 k((x,a),(y,b))=\tfrac12(y-x)^TS(a,b)^{-1}(y-x).
\]
The allowed and forbidden sets are clopen in the actual compact mark
product. Choose a continuous uniformly positive definite extension
\(S_{\rm e}\) to that product, agreeing with \(S\) on
\(\mathcal F\); for example use a constant positive matrix on the
forbidden part. Denote its quadratic cost by \(k_{\rm e}\), and set
\[
 \chi(a,b)=\mathbf1_{\{a>b\}},\qquad
 c_{M,t}=k_{\rm e}+M\chi+t g,
 \qquad M\ge0,\quad g\in C(\operatorname{supp}\alpha
                                  \times\operatorname{supp}\beta).
\]
Write \(G=\|g\|_\infty\) and \(s=1-2/q\). There are constants
\(0<c_-\le c_+<\infty\) with
\(c_-|y-x|^2\le k_{\rm e}\le c_+|y-x|^2\).

In the cube \(D_n=(0,n^{1/q})^q\), take two independent lists of
\(n\) uniform spatial points, with independent product marks, and put
\[
 f_{n,M}(t)=\frac1n\E\min_\sigma\sum_{i=1}^n
                    c_{M,t}((X_i,a_i),(Y_{\sigma(i)},b_{\sigma(i)})),
 \qquad p_M(t)=\lim_{n\to\infty}f_{n,M}(t).
\]
For fixed \(M\), existence and locally uniform convergence follow from
Proposition~\ref{prop:highd-finite-pressure}: its bounded mark addition
can be \(M\chi+t g\). Its proof also identifies \(p_M\) with the
expected cost per volume of maximum-cardinality matchings of independent
Poisson inputs in growing cubes. No stationary variational formula for
the unrestricted extension \(S_{\rm e}\) is used.

\begin{lemma}[Uniform penalty bound]
\label{forward-full-penalty-bound}
There is \(C<\infty\), independent of \(M,t\), such that
\[
 -|t|G\le p_M(t)\le C+|t|G.
\]
Consequently \(p_\infty(t)=\sup_{M\ge0}p_M(t)\) is finite,
concave and \(G\)-Lipschitz, and \(p_M\uparrow p_\infty\)
locally uniformly in \(t\).
\end{lemma}
\begin{proof}
Condition on the marks and pair their order statistics. This minimizes
their empirical scalar absolute-distance cost. Divide these pairs by their
source-class and target-class labels, obtaining at most \(J\) groups.
Within a group of size \(m\), replace the spatial pairing by an optimal
squared-distance matching. Group membership used only marks, so the two
spatial lists remain independent uniform lists conditional on these groups.
The unmarked quadratic bound used in
Proposition~\ref{prop:highd-energy}, scaled to \(D_n\), gives total
expected kinetic cost at most
\[
 C n^{2/q}\sum_{\rm groups}m^{1-2/q}
          \le C J^{2/q}n.
\]
Each class rectangle has one orientation. Its rematching therefore preserves
both its backward count and its total scalar absolute-distance cost. If
\(B_n\) is the resulting backward count, then
\[
 \frac{2\delta B_n}{n}
 \le W_1(\alpha_n,\beta_n)
       -\int b\,d\beta_n+\int a\,d\alpha_n.
\]
Compactly supported empirical measures converge in \(W_1\), and forward
population feasibility makes the limiting right side zero. It is bounded,
so \(\E B_n/n\to0\). At each fixed \(M\), the comparison proves
the upper bound; nonnegative kinetic and penalty costs prove the lower
bound. This argument takes the volume limit before the penalty limit.
The pressures are increasing in \(M\). Their pointwise limit inherits
concavity and the common Lipschitz bound. A finite parameter net then gives
locally uniform convergence.
\end{proof}

\begin{theorem}[Differentiability of the limiting penalty pressure]
\label{forward-full-pressure-c1}
The function \(p_\infty\) belongs to \(C^1(\R)\).
\end{theorem}
\begin{proof}
Fix \(t_0\) and \(h_j\downarrow0\). First choose \(M_j\to\infty\),
then \(n_j\to\infty\), so that for \(\ell=-2,-1,0,1,2\),
\begin{equation}
\label{forward-full-joint-diagonal}
 \begin{split}
 |p_{M_j}(t_0+\ell h_j)-p_\infty(t_0+\ell h_j)|&\le h_j^2,\\
 |f_{n_j,M_j}(t_0+\ell h_j)-p_{M_j}(t_0+\ell h_j)|&\le h_j^2.
 \end{split}
\end{equation}
These choices use the two successive value limits, not convergence of
fixed-penalty derivatives.

For each \(j\), use one target cloud \(V_l=(Y_l,b_l)\) and one
independent \(P\sim\operatorname{Dirichlet}(1,\ldots,1)\). With
\(n=n_j\), define the mass-\(n\) measures
\[
 \lambda_n=\mathbf1_{D_n}(x)\,dx\,\alpha(da),\qquad
 \nu_n=\sum_{l=1}^n nP_l\delta_{V_l}.
\]
The exact value identity in Proposition~\ref{prop:iid-cavity-mean}
identifies their expected normalized optimum for \(c_{M_j,t}\) with
\(f_{n_j,M_j}(t)\). To check its hypothesis at every \(M,t\), fix
distinct spatial targets and a source mark \(a\). Their row-cost
difference is a nonconstant polynomial in \(x\): either the inverse
matrices differ, giving a nonzero quadratic part, or the common inverse
matrix multiplies \(Y_l-Y_m\ne0\), giving a nonzero linear part.
Both mark additions change only its constant term. Polynomial level sets
are Lebesgue null, and Fubini handles arbitrary, including atomic, mark
laws. The same argument holds for a fixed countable target configuration.

On this common environment choose diffuse optima
\(\Gamma_j^+,\Gamma_j^-,\Gamma_j^0\) at
\(t_0+h_j,t_0-h_j,t_0\). Let
\(C_j=k_{\rm e}+M_j\chi+t_0g\) and
\(\overline\Gamma_j=(\Gamma_j^++\Gamma_j^-)/2\).
Finite optimality gives
\begin{equation}
\label{forward-full-diffuse-excess}
 \begin{split}
 0\le\Delta_j
  &:=\int C_j\,d\overline\Gamma_j-\int C_j\,d\Gamma_j^0
       \le h_jGn,\\
 \int(k_{\rm e}+M_j\chi)\,d\Gamma_j^\pm
  &\le\int C_j\,d\Gamma_j^0+(|t_0|+2h_j)Gn.
 \end{split}
\end{equation}
The value identity, Lemma~\ref{forward-full-penalty-bound} and
\eqref{forward-full-joint-diagonal} thus bound the expected normalized
kinetic energies uniformly and the forbidden masses by \(C/M_j\).

Root uniformly in \(D_n\) and translate both plans and their shared
target measure by the negative root. For a bounded anchor window \(B\),
the source- and target-anchored masses of edges longer than \(R\) have
expectation at most \(C|B|/R^2\). Indeed root integration contributes
at most \(|B|/n\), after which the kinetic bound applies. The common
local marginal limit is
\[
 \lambda=dx\,\alpha(da),\qquad
 \nu=\sum_l E_l\delta_{(Y_l,b_l)},
\]
where the targets form a unit-intensity marked Poisson process and their
weights \(E_l\) are independent mean-one exponentials. This follows
from \(nP_l=nE_l/\sum_iE_i\), the binomial local limit and the law of
large numbers. The Dirichlet second moments give
\(\E[\nu_n(U_n+B)^2]\le2|B|+|B|^2\).
These local mass bounds and both endpoint tails yield a joint stationary
subsequential limit \((\Gamma^+,\Gamma^-,\nu)\), with both full
marginals retained. This is the compactness argument in the proof of
Theorem~\ref{thm:highd-pressure-differentiability}; it needs only the
uniform energy bounds just proved. Nonnegative compact tests times the
continuous indicator \(\chi\) have expectation \(O(M_j^{-1})\),
so both limiting plans are supported on \(\mathcal F\).

Their average has no strict legal two-cycle improvement for
\(k_{\rm e}+t_0g\). Here are the required finite corrections. A strict
violation has a certificate of bounded open source rectangles \(A_1,A_2\)
and target rectangles \(B_1,B_2\), disjoint within each color, such that
all four mark products are in \(\mathcal F\) and the swapped cost saves
at least \(\varepsilon>0\). Continuity, openness of \(\mathcal F\)
and countable bases give a countable family of such certificates. Set
\[
 W(\Gamma)=\min\{1,\Gamma(A_1\times B_1),\Gamma(A_2\times B_2)\}.
\]
Scale the two restrictions to their common mass \(m=W\), and cross
their source and target marginals using normalized product couplings.
Every marginal is preserved and the cost saving is at least
\(\varepsilon m\). The penalty is zero on all old and new edges.

Place these corrections on a uniformly translated coarse grid of side
\(K\), with the two source projections inside one grid cube. Removed
source pieces in different cubes are disjoint. Overlapping target
neighborhoods cause no problem, since each correction preserves every
target mass. Their total saving is at most \(\Delta_j\). Grid-phase
averaging, including boundary cells where absent source mass is zero,
therefore gives
\[
 \E W(\tau_{U_n}\overline\Gamma_j)
       \le\frac{K^q\E\Delta_j}{\varepsilon n}
       \le\frac{K^qG}{\varepsilon}h_j\longrightarrow0.
\]
The witness is bounded and lower semicontinuous for vague convergence.
Portmanteau and the countable certificate family show that the limiting
average is two-cyclically monotone whenever its crossed pairs are legal.

This constrained monotonicity still forces the average to be a map. If
positive source density were sent to two distinct countable target labels
on a set \(A\) of positive \(\lambda\)-measure, both choices would
be forward at every point of \(A\), after removing null exceptions.
The crossed assignments for any two points of \(A\) would then be
legal. The two-cycle inequalities in both orders force the difference of
the two row costs to be constant on \(A\), contrary to the polynomial
level-set argument. Countability of target pairs rules out every split.
The average conditional target law is therefore a Dirac mass, which forces
\begin{equation}
\label{forward-full-diffuse-agreement}
                     \Gamma^+=\Gamma^-\quad\text{almost surely}.
\end{equation}

Put \(q_j^\pm=n_j^{-1}\E\int g\,d\Gamma_j^\pm\).
They are supergradients of \(f_{n_j,M_j}\) at \(t_0\pm h_j\).
For the positive sign, the adjacent secants give
\[
 \frac{f_{n_j,M_j}(t_0+2h_j)-f_{n_j,M_j}(t_0+h_j)}{h_j}
 \le q_j^+\le
 \frac{f_{n_j,M_j}(t_0+h_j)-f_{n_j,M_j}(t_0)}{h_j}.
\]
The analogous negative-side bounds use \(t_0-2h_j,t_0-h_j,t_0\).
By \eqref{forward-full-joint-diagonal}, they imply
\(q_j^+\to p_\infty'(t_0+)\) and
\(q_j^-\to p_\infty'(t_0-)\).
To pass these expectations through the joint local limit, test with
\(\eta(x)g(a,b)\), where \(\eta\ge0\) is continuous, compactly
supported and has integral one. Root averaging differs from \(q_j^\pm\)
by \(O(Gn_j^{-1/q})\), using the deterministic diffuse source marginal
and its fixed-width boundary layer. A compact target cutoff makes the
integral a continuous local plan test bounded by fixed source mass;
removing it costs \(O(G/R^2)\) by the endpoint tail estimate.
Equation~\eqref{forward-full-diffuse-agreement} identifies the two
limits. Thus \(p_\infty\) is differentiable at every \(t_0\).
A differentiable concave function on \(\R\) has continuous derivative.
\end{proof}

Let \(\mathcal M_+\) be the stationary perfect integer matching laws
of independent unit-intensity product-marked Poisson inputs, supported on
\(\mathcal F\) and having finite source-Palm squared displacement.
Auxiliary randomization is allowed. Define
\[
 I_+(t)=\inf_{P\in\mathcal M_+}\E_P^0[k+t g].
\]

\begin{lemma}[Attainment and the lower stationary comparison]
\label{forward-full-stationary-existence}
The class \(\mathcal M_+\) is nonempty, \(I_+(t)\) is finite and
attained, and \(I_+(t)\le p_\infty(t)\).
\end{lemma}
\begin{proof}
Proposition~\ref{prop:forward-poisson-criterion} gives a forward
finite-second-moment competitor under exactly the stated Hall condition
and \(q\ge3\). A minimizing sequence has a uniform second-moment
bound because the probe is bounded and \(k\) is uniformly elliptic on
allowed pairs. Local Poisson counts make the input and integer edge
measures tight. Source- and target-anchored escape bounds are \(O(R^{-2})\),
the latter by mass transport. Compact endpoint truncations and uniform
integrability of local counts retain both complete endpoint projections.
The limit is a stationary perfect integer matching of the prescribed
inputs. Forward support is closed on the separated mark product; the
nonnegative cost \(k+t g+|t|G\) is lower semicontinuous there. This
proves attainment, by the same two-endpoint argument as
Proposition~\ref{prop:highd-local-limit}.

For the comparison choose \(M_j\to\infty\) and then \(n_j\to\infty\)
with \(f_{n_j,M_j}(t)\to p_\infty(t)\), and use the original atomic
finite optimizers. The penalty bound gives uniformly bounded expected
kinetic energy and forbidden fraction \(O(M_j^{-1})\). Their rooted
limits are stationary perfect integer matchings with the independent
marked Poisson input law; both endpoint tails again retain perfection.
They are forward, and lower semicontinuity of the cost shifted by \(|t|G\)
gives \(I_+(t)\le p_\infty(t)\). This uses finite all-pairs values,
not a stationary identity for their extended kernel.
\end{proof}

\begin{lemma}[Forward minimizing traffic and centering]
\label{forward-full-traffic-centering}
Every stationary forward mean minimizer for \(k+t g\) has
square-integrable smoothed absolute kinetic traffic and source-Palm mean
displacement zero in every translation-ergodic component. Its paths have
the total current divergence of the full independent inputs. Thus the
independent-reserve screening conclusion of
Corollary~\ref{cor:highd-kinetic-screening} applies, with the resulting
repair allowed to use all mark pairs. The same assertions hold at any
equal input intensity \(0<\lambda\le1\).
\end{lemma}
\begin{proof}
Let \(\lambda\) be the common input intensity. A minimizer admits no
strict legal two-edge improvement. Otherwise a
bounded complete-edge certificate with a fixed positive gain occurs with
positive probability. On a randomly translated sufficiently coarse grid,
perform one measurable such swap in every successful cube. The swaps are
disjoint, preserve forward perfection and finite energy, and strictly
decrease the stationary mean cost.

Partition the matching edges by their finitely many legal source-class and
target-class rectangles. Within one rectangle, all source marks precede
all target marks, so every two-edge crossing is legal. Its selected edges
therefore satisfy exactly the perturbed two-cycle inequality used in
Lemma~\ref{lem:highd-perturbed-traffic}. For clarity the extension to a
selected collection uses its proof, not a Poisson assumption for that
collection. The interface estimate is
\[
 |v_i-v_j|\le C|z_i-z_j|/u+C\sqrt{|t|G/u},
 \qquad v_e=[T(b_e)-T(a_e)]^{-1}(y_e-x_e),
\]
on a source-side window of scale \(u\). The Lipschitz approximate
velocity extension and its phase tube are unchanged. Every phase-cell
source preimage is a deterministic marked set with intensity bounded by
that of the \emph{full} source input. The Poisson-animal exceptions
therefore hold simultaneously for the selected edges. The nonexceptional
Campbell terms are bounded by their contribution to the full matching's
second moment, as in \eqref{eq:highd-measurable-traffic-majorant}.
The target-layer proof uses the full target input in the same way.
For the fixed nonnegative compact smoothing kernel \(\eta\), write
\[
 U_{ij}(z)=\sum_{e\in(i,j)}\int_{a_e}^{b_e}
 |T'(r)v_e|\eta(z-x_e-[T(r)-T(a_e)]v_e)\,dr.
\]
Dyadic endpoint summation proves, for each allowed rectangle,
\[
 \E U_{ij}(0)^2\le C_{ij,t,g}(1+\lambda\E^0|v_e|^2)<\infty.
\]
Finite summation makes \(U=\sum_{ij}U_{ij}\) square integrable.

For every edge use its actual forward kinetic path
\[
 X_e(r)=x_e+[T(r)-T(a_e)]v_e\quad(a_e\le r\le b_e).
\]
Its length and diameter are at most \(C|y_e-x_e|\). Summing its
oriented current gives
\[
 \operatorname{div}J=\eta*(\mu_S-\mu_T),\qquad
 \E J(0)=\lambda\E^0(y-x).
\]
These identities concern the total current and full independent inputs;
no such independent-input identity is assigned to a selected rectangle.

For any bounded measurable vector function \(h(a)\), translate each
source by \(\varepsilon h(a)\), retaining its assignment and marks.
This preserves its product Poisson law, independence from targets,
stationarity and every forward constraint. The exact quadratic cost
difference has linear part \(-\varepsilon\E^0[v\cdot h(a)]\)
and a finite nonnegative quadratic part; the probe is unchanged.
Minimality for both signs yields \(\E^0[v\mid a]=0\).
Target translations and mass transport give \(\E^0[v\mid b]=0\).
Since \(T\) is bounded on the scalar span,
\[
 \E^0(y-x)=\E^0[(T(b)-T(a))v]=0.
\]
The independent marked Poisson input law is ergodic, so every ergodic
component retains that input law. Its mean cost is at least the stationary
minimum at that fixed intensity, and their average equals that number.
Almost every component is therefore
itself minimizing, and the preceding translations give zero mean
displacement in each component.

This verifies every analytic hypothesis used in the proofs of
Lemma~\ref{lem:highd-current-boundary},
Proposition~\ref{prop:highd-signed-ports} and
Corollary~\ref{cor:highd-kinetic-screening}: total square-integrable
traffic, the full-input divergence, invariant conditional centering, and
the kinetic path-diameter bound. Their boundary-energy, smoothing and
independent-reserve arguments use no ordering beyond these facts. They
therefore give the same screening conclusion here. The repair is an
unrestricted matching, which will only be used at a fixed penalty.
\end{proof}

\begin{theorem}[Forward stationary variational identity]
\label{forward-full-stationary-identity}
For every real \(t\),
\[
                         p_\infty(t)=I_+(t).
\]
In particular \(I_+\in C^1(\R)\), and the limiting penalty pressure
does not depend on the positive definite extension on forbidden pairs.
Every forward stationary mean minimizer has paired expectation
\(\E^0g=I_+'(t)\).
\end{theorem}
\begin{proof}
Fix \(M,t\) and \(0<\rho<1/2\), and put \(\lambda=1-\rho\).
Take an intensity-one forward minimizer for
\(k+t\lambda^{2/q}g\) and dilate space by \(\lambda^{-1/q}\).
Its inputs now have intensity \(\lambda\), and its \(k+t g\) cost
density is exactly
\[
                  \lambda^s I_+(t\lambda^{2/q}).
\]
Scaling gives a bijection of the finite-energy forward matching classes,
so this law minimizes the mean cost at intensity \(\lambda\).
Add independent source and target reserves of intensity \(\rho\) with
the original mark laws. Each full color input is exactly the prescribed
unit-intensity product-marked Poisson process, and the colors are
independent.

Lemma~\ref{forward-full-traffic-centering} permits retention of the main
paths in an inner cube and maximum-cardinality completion in \(Q_L\).
The completion preserves the retained edges and has
\[
 \E R_2\le C\rho^sL^q+o_{\rho,t,P}(L^q),\qquad
 \E N_{\rm repair}\le\rho L^q+o_{\rho,t,P}(L^q).
\]
The second bound follows because the main port count is \(o(L^q)\)
and the reserve count has mean \(\rho L^q\). The construction includes
its exponentially unlikely reserve-count failures, as in the cited
screening proof. Its repair need not be forward. Its full penalized cost
is bounded above by \(c_+R_2+(M+|t|G)N_{\rm repair}\).
All retained main penalties are zero. Omitting \(o(L^q)\) main endpoints
raises an upper bound for their signed cost by at most \(|t|G\) times
that count, because the kinetic part is nonnegative.

Let \(C_L^P(M,t)\) denote the independent-Poisson finite optimum. Then
\[
 \limsup_{L\to\infty}\frac{\E C_L^P(M,t)}{L^q}
 \le\lambda^s I_+(t\lambda^{2/q})
             +C c_+\rho^s+(M+|t|G)\rho.
\]
The left side is \(p_M(t)\) by the independent-count part of the
finite-pressure proof. At fixed \(M,t\), let \(\rho\downarrow0\).
The stationary infimum is \(G\)-Lipschitz in its probe parameter,
so \(p_M(t)\le I_+(t)\). Taking the supremum over \(M\) and using
Lemma~\ref{forward-full-stationary-existence} proves equality.
The order is volume first, reserve intensity second, and penalty last;
no uniform screening rate along varying reserves is needed.

Theorem~\ref{forward-full-pressure-c1} now gives differentiability of
\(I_+\). Its definition involves only forward costs, giving extension
independence. A minimizer's paired expectation is a supergradient of
\(I_+\) at \(t\), so differentiability identifies it with \(I_+'(t)\).
At \(t=0\), a countable dense family of continuous paired tests therefore
identifies one common paired-mark mean measure for all stationary forward
mean minimizers. Tight Hall cuts were permitted throughout.
\end{proof}

\begin{corollary}[Uniform integrability on value-recovery diagonals]
\label{forward-full-diagonal-ui}
Fix \(t\), and let \(M_j\to\infty\), \(n_j\to\infty\) satisfy
\(f_{n_j,M_j}(t)\to I_+(t)\). For any measurable choices of exact
permutation optimizers, the kinetic-plus-penalty cost at a uniformly selected
source is uniformly integrable. In particular, uniformly over choices of
exact permutation or fractional optima,
\[
 \E\left[\frac1{n_j}\sum_{\rm backward}k_{\rm e}\right]\to0,
 \qquad
 \E\left[\frac1{n_j}\sum M_j\chi\right]\to0,
 \qquad
 \E\left[\frac1{n_j}\sum_{\rm backward}|y-x|^2\right]\to0.
\]
This assertion requires the stated value convergence; it is not a claim
for arbitrary diverging penalties.
\end{corollary}
\begin{proof}
Along any joint rooted subsequence the uniform second moment and vanishing
forbidden fraction give a forward stationary perfect limit. Its mean
\(k+t g\) cost is at least \(I_+(t)\), while lower semicontinuity and
the assumed value convergence give the reverse inequality. Thus every
such limit is mean minimizing.

Anchor at a uniformly chosen source and write \(X_j=k_{\rm e}\),
\(Z_j=M_j\chi\), and \(H_j=g\) for its chosen edge. The
root-to-Palm passage uses compact edge tests and uniform integrability
of the local Poisson counts, as in the preceding atomic compactness
argument. Along a further subsequence, \(X_j\) converges in law to the
limiting forward kinetic cost \(X\), \(\E H_j\to\E H\), and
\(\mathbb P(Z_j\ne0)\to0\). Mean minimality gives
\[
 \E(X_j+Z_j)=f_{n_j,M_j}(t)-t\E H_j
         \longrightarrow I_+(t)-t\E H=\E X.
\]
Also \(X_j+Z_j\) converges in law to \(X\). Nonnegative weak
convergence together with convergence of means implies uniform
integrability. Hence \(\E Z_j\to0\) and
\(\E[X_j\mathbf1_{\{\chi=1\}}]\to0\); uniform ellipticity gives
the squared-displacement conclusion. Applying this reasoning to a
contrary subsequence proves the assertion for the full chosen sequence.

Uniformity over permutation optima follows by choosing a measurable
maximizer of a contrary bad-edge observable within the finite optimal set.
A fractional finite optimum is a Birkhoff mixture supported on optimal
permutations, so the same uniform bounds hold for it. The argument supplies
no spatial independence assertion for its selected backward endpoints.
\end{proof}


% Full conversion of strict-hall-genuine-recovery.md
% source SHA256 43e23d5d48c85e040620a8539e6e241e6534558ee7bc57ad68011d27777bb730.
% This supplies the proof behind measurable-forward-strict-input.
% Sections 2--5 and 8 of the accepted source are reproduced below.
% The variable-tube theorem of source Section 6 is not asserted in q=3,4.
\subsection{Strict-Hall screening and boundary-port injection}
\label{strict-screen-section}

Fix compact disjoint scalar mark supports \(K_A,K_B\), partitioned into
finitely many ordered monochromatic classes of positive respective
probability. Let \(\alpha,\beta\) be their laws, and suppose that for every
proper gap \(c\) between consecutive classes,
\begin{equation}
\label{strict-screen-hall}
 H(c):=\alpha((-\infty,c])-\beta((-\infty,c])\ge\eta>0.
\end{equation}
The first class is then a source class and the last a target class.
A source-target pair is legal exactly when \(a<b\), so each class
rectangle is either entirely legal or entirely illegal. Let \(T\) be
\(C^1\) on the compact scalar span, with
\[
 h\,\mathrm{Id}\le T'(r)\le H\,\mathrm{Id},\qquad 0<h\le H<\infty,
\]
and use the legal kinetic cost
\[
 k((x,a),(y,b))
   =\tfrac12(y-x)^T[T(b)-T(a)]^{-1}(y-x).
\]
Fix a continuous uniformly elliptic quadratic extension \(k_e\) to illegal
rectangles, and a continuous paired probe \(g\), with \(G=\|g\|_\infty\).
Write \(I_+(t)\) for the unit-intensity stationary forward mean-cost
infimum of \(k+tg\) with finite source-Palm displacement second moment.
We work in \(q\ge3\), put \(s=1-2/q\), and specify where \(q>4\) is needed.

Finite-energy stationary minimizers and the legal-rectangle traffic and
centering argument are given by
Proposition~\ref{prop:forward-weak-continuity} and
Lemma~\ref{forward-full-traffic-centering}. In particular a minimizing
forward law has square-integrable smoothed absolute kinetic traffic and
zero invariant conditional mean displacement. Its kinetic paths have
diameter at most \(C|D|\), where \(D=y-x\). The signed-port and unmarked
screening results, Proposition~\ref{prop:highd-signed-ports} and
Corollary~\ref{cor:highd-kinetic-screening}, therefore apply to its total
current. Those results concern the full independent input, not
independently thinned selected edge subcollections. We now prove the
additional forward repair and its exceptional-energy estimate.

\subsubsection{Boundary depth of the actual stationary ports}

Put \(Q_L=(-L/2,L/2)^q\), \(V=L^q\), and
\(d_L(x)=\operatorname{dist}(x,\partial Q_L)\).
For a stationary finite-second-moment matching of independent Poisson
inputs with common intensity \(\lambda\le1\), retain exactly the edges
whose entire chosen path lies in \(Q_{L-\ell_L}\), with
\(\ell_L=L^{1/4}\). Call the unretained source endpoints in \(Q_L\) the
source ports \(B_L\).

\begin{lemma}[Vanishing boundary-depth moment]
\label{strict-screen-depth}
The source ports satisfy
\begin{equation}
\label{strict-screen-depth-moment}
 \E|B_L|=o(V),\qquad
 \E S_L=o(V),\qquad S_L=\sum_{x\in B_L}d_L(x)^2.
\end{equation}
The analogous target-port count is \(o(V)\).
\end{lemma}

\begin{proof}
An omitted endpoint satisfies
\(d_L(x)\le C\ell_L+C|D|\), after adjusting a harmless collar constant.
Campbell's formula and the cube collar-volume bound show that its
count divided by \(V\) is at most
\[
 C\ell_L/L+C\E^{\rm S}\min(1,C|D|/L)=o(1).
\]
For the weighted count, first integrate the layer \(d_L\le C\ell_L\).
Its contribution per volume is at most \(C\ell_L^3/L\). Outside a
slightly enlarged layer, omission implies \(d_L\le C|D|\).
Integrating the possible endpoint locations with Campbell gives
\begin{equation}
\label{strict-screen-depth-estimate}
 \frac{\E S_L}{V}
 \le C\frac{\ell_L^3}{L}
       +C\E^{\rm S}\!\left[
          |D|^2\min(1,C|D|/L)\right]\longrightarrow0.
\end{equation}
Here \(\ell_L^3/L=L^{-1/4}\), and dominated convergence uses only
\(\E^{\rm S}|D|^2<\infty\). A third moment or a quantitative tail
rate is unnecessary. The target count follows by the same argument
and mass transport. This collar is admissible in the existing
screening proof because \(\ell_L\to\infty\) and \(\ell_L/L\to0\).
\end{proof}

\subsubsection{Injection into an independent reserve}

\begin{lemma}[Boundary-sensitive reserve injection]
\label{strict-screen-injection}
Fix \(r>0\). Let \(R_L\) be either an independent Poisson process of
intensity \(r\) in \(Q_L\), or an independent uniform cloud of \(m\)
points with \(rV/2\le m\le2rV\). For the source ports above there is,
on an event \(\mathcal G_L\), an injection \(\iota_L:B_L\to R_L\)
such that
\begin{equation}
\label{strict-screen-injection-cost}
 \E\!\left[\mathbf1_{\mathcal G_L}
        \sum_{x\in B_L}|x-\iota_L(x)|^2\right]=o_r(V).
\end{equation}
If \(N_L\) is the total full-input count, with
\(\E N_L^2=O(V^2)\) and exponentially small upper tail above \(CV\),
the failure event can be charged at whole-cube cost:
\begin{equation}
\label{strict-screen-injection-exception}
 \E\!\left[\mathbf1_{\mathcal G_L^c}L^2N_L\right]=o_r(V).
\end{equation}
No independence or Poisson law is required of the selected port subset.
Its domination by the main source process and
\eqref{strict-screen-depth-moment} are the inputs.
\end{lemma}

\begin{proof}
Partition \(Q_L\setminus Q_{L/2}\) into its \(2q\) face pyramids,
assigning null ties arbitrarily. In a pyramid normal to coordinate \(i\),
write
\[
 x_i=\pm(L/2-d),\qquad
 x_j=(1-2d/L)u_j\quad(j\ne i),\qquad
 0\le d\le L/4,
\]
where \(u\) lies in a \((q-1)\)-cube of side \(L\). Its Jacobian is
\((1-2d/L)^{q-1}\), between \(2^{-(q-1)}\) and one, and the coordinate
map and inverse have Lipschitz constants independent of \(L\).

Choose \(C_r\ge8\,2^{q-1}/r\), set \(h_L=L/(4C_r)\), and expand the
thin collar by
\[
 \Phi_L(u,d)=(u,C_rd),\qquad d\le h_L.
\]
The displacement is at most \(C(q,r)d\). The image of the
\emph{entire} main source Poisson process in that collar is an
inhomogeneous Poisson process of intensity at most
\begin{equation}
\label{strict-screen-expanded-intensity}
 \frac{\lambda}{C_r}
       \left(\frac{1-2d/L}{1-2C_rd/L}\right)^{q-1}
                 \le r/8.
\end{equation}
It is independent of the reserve. This reference process is formed
before restricting to any selected port subset.

Pull the reference red cloud and the blue reserve in each face pyramid
back to its rectangular coordinate domain. Use a dyadic hierarchy whose
smallest cells have side lengths between fixed positive constants.
Within each cell match as many remaining red and blue points as possible,
then merge to the parent, up to the full face domain. A red point remains
unmatched after its cell is exhausted only if that cell's total red
count exceeds its total blue count. At scale \(a\) its volume is comparable
to \(a^q\) and diameter is \(O(a)\). By
\eqref{strict-screen-expanded-intensity}, Poisson Chernoff bounds give,
also after insertion of one specified red point,
\begin{equation}
\label{strict-screen-reference-tail}
 \Pp(\text{red count exceeds blue count in its cell})
                    \le C_r e^{-c_r a^q}.
\end{equation}
For the fixed uniform reserve the blue count in a physical cell is
binomial, with mean at least \(r/2\) times its volume; the red mean
is at most \(r/8\) times that volume. Independent-cloud Chernoff bounds
give the same estimate uniformly for \(rV/2\le m\le2rV\).
Counts of different cells need not be independent.

Summing the scale tails proves that the assigned displacement \(Z\) of
every reference red point has every fixed finite moment bounded by
\(C_{p,q,r}\), uniformly in its location and \(L\). Set \(Z=0\) if the
root-face injection fails. The probability that any face fails to match
all its red points is at most \(C_r e^{-c_rV}\).
The coordinate Jacobian and Lipschitz bounds justify both the volume
and diameter estimates for these dyadic rectangles.

For a near-boundary port \(x\), use the reference target assigned to
\(\Phi_L(x)\). Distinct ports get distinct targets. The deterministic
part of the squared displacement sums to at most \(C_r S_L\). For the
reference part, Cauchy--Schwarz on the joint sample-and-point measure gives
\begin{align}
 \E\sum_{\substack{x\in B_L\\d_L(x)\le h_L}}Z(x)^2
 &\le(\E|B_L|)^{1/2}
       \left(\E\sum_{\substack{\text{all main sources}\\d_L(x)\le h_L}}
                         Z(x)^4\right)^{1/2}\notag\\
 &\le C_r(\E|B_L|)^{1/2}V^{1/2}=o_r(V).
\label{strict-screen-adaptive-restriction}
\end{align}
Campbell and the uniform reference fourth moment give the second factor.
This remains valid for an arbitrary dependence of \(B_L\) on the full
main matching. It does not condition selected ports to be iid.

The central reserve cube \(Q_{L/2}\) has not been used. Except with
exponentially small probability it has at least \(c_rV\) reserve points.
If there are at most \(c_rV\) deep ports, inject them arbitrarily into
distinct central points. Since their depth exceeds \(h_L\), their cost is
\[
 qL^2\#\{\text{deep ports}\}\le C_r S_L.
\]
If instead there are more than \(c_rV\) deep ports, then
\begin{equation}
\label{strict-screen-deep-failure}
 S_L\ge c'_rL^2V.
\end{equation}
On \(\{N_L\le CV\}\), this pays the entire fallback \(L^2N_L\) by
\(C_r S_L\). Its expectation is \(o_r(V)\). On the remaining input-count,
face, or central-reserve failures, exponential tails and count moments
pay the polynomial fallback. This proves both conclusions.

For a balanced reserve with random size
\(m=\min(N^S,N^T)\), where the parent counts are independent Poisson
of mean \(rV\), first restrict to \(rV/2\le m\le2rV\). Its complement
is exponentially unlikely. Auxiliary-index selection leaves independent
uniform positions conditional on \(m\), so the preceding version applies.
\end{proof}

\subsubsection{Three independent pools and universal bridge swaps}

\begin{proposition}[Strict-Hall forward screening with charged exceptions]
\label{strict-screen-three-pools}
Suppose \(q>4\) and \eqref{strict-screen-hall} holds. Fix
\(0<\rho<1/2\), set \(\lambda=1-\rho\), and start from a stationary
unit-intensity forward minimizer at parameter \(t\lambda^{2/q}\),
spatially dilated to intensity \(\lambda\). Its cost density for
\(k+tg\) is
\begin{equation}
\label{strict-screen-main-density}
 \lambda^s I_+(t\lambda^{2/q}).
\end{equation}
Independent product-marked reserves of intensity \(\rho\) per color
complete its retained interior paths to a maximum-cardinality matching
of the original independent unit-intensity Poisson inputs in \(Q_L\).
The quadratic repair cost, apart from the retained main edges, has
expectation at most
\begin{equation}
\label{strict-screen-repair-energy}
 C\rho^sV+o_\rho(V).
\end{equation}
All pool-count vertices put aside have universal earliest-source or
latest-target marks and expected number \(O_\rho(\sqrt V)\). They can be
paired forward up to the total color excess. The exceptional event
\(\mathcal E_{L,\rho}\) satisfies
\begin{equation}
\label{strict-screen-exceptional-energy}
 \E[L^2N_{\rm tot}\mathbf1_{\mathcal E_{L,\rho}}]=o_\rho(V).
\end{equation}
The construction preserves every retained main edge and never asserts
an iid law for the main ports.
\end{proposition}

\begin{proof}
Choose
\[
 0<\kappa<
 \frac18\min\{\eta,\alpha(\text{first source class}),
                    \beta(\text{last target class}),1\},
 \qquad r=\kappa\rho.
\]
Independent Poisson marking extracts two universal reserve pools
\(R_0,R_1\), each with source and target intensity \(r\), using only
the first source and last target classes. Let \(R_2\) be the remainder.
The three pools are independent. The remaining proper-cut slack is
\begin{equation}
\label{strict-screen-residual-slack}
 \rho H(c)-2r\ge\rho(\eta-2\kappa)>0,
\end{equation}
and its extreme-class intensities are positive. This is where strict
population slack is used.

Apply Corollary~\ref{cor:highd-kinetic-screening} with \(R_0\) to the
main paths retained in \(Q_{L-\ell_L}\). Its proof deletes the full
color excess from reserve vertices by auxiliary indices before spatial
repair. Outside its exponentially unlikely count failure, every main
port is matched and only that reserve excess is put aside. Its repair
squared cost obeys
\begin{equation}
\label{strict-screen-first-pool}
 \E C_0\le Cr^sV+o_r(V).
\end{equation}
Every edge using a universal reserve endpoint has the required forward
orientation. Thus only repair edges with two main endpoints can be illegal.

In the independent pool \(R_1\), remove the color-count excess by
auxiliary indices, and optimally match the remaining
\(m=\min(N_1^S,N_1^T)\) uniform clouds for squared distance. Its edges are
all forward and the sharp iid quadratic estimate
\cite[Theorem~1.1]{GoldmanTrevisan2021} gives
\begin{equation}
\label{strict-screen-bridge-pool}
 \E C_1\le Cr^sV.
\end{equation}
Conditional on counts, its bridge-source positions are an independent
uniform cloud, independent of the full main input. Apply
Lemma~\ref{strict-screen-injection} to \emph{all} main source ports.
Only afterwards restrict its injection to source ports belonging to
illegal screened repair edges. The injection cost is \(o_r(V)\) and
its failure is paid by \eqref{strict-screen-injection-exception}.

For an illegal screened edge \(x\to y\) and its distinct assigned
bridge edge \(u\to v\), replace the two edges by
\[
 x\to v,\qquad u\to y.
\]
Both are forward by the universal endpoint marks. Injectivity makes the
used bridge sources, targets, and old repair edges distinct, so all
matching constraints are preserved. The squared cost satisfies
\begin{equation}
\label{strict-screen-swap-cost}
 |x-v|^2+|u-y|^2
       \le4|x-u|^2+2|u-v|^2+2|x-y|^2.
\end{equation}
Charge the selected old illegal-edge term to the \emph{whole} screened
repair cost, and the selected bridge term to the whole bridge cost.
Uniform ellipticity and
\eqref{strict-screen-first-pool}--\eqref{strict-screen-bridge-pool}
give total kinetic repair \(Cr^sV+o_r(V)\). No uniform integrability of
an optimizer-selected subset has been presumed. The bridge pool is
separate from \(R_0\), so its matching has no competing use in the swap.

For \(R_2\), delete its total color excess from the earliest source class
if sources are in excess, or from the latest target class otherwise.
Use marks and auxiliary indices only. Except with exponentially small
probability, there are enough extreme points and every remaining prefix
Hall inequality holds: the means in
\eqref{strict-screen-residual-slack} are of order \(\rho V\), and the
total excess is smaller than half the retained margins. Finitely many
Poisson Chernoff bounds establish these simultaneous events. Sorted
class matching supplies legal integer quotas. Conditional on marks and
counts, each quota group contains two independent uniform spatial
samples. The finite number of groups and concavity give
\begin{equation}
\label{strict-screen-ordinary-pool}
 \E C_2\le C L^2(\rho V)^s=C\rho^sV.
\end{equation}

The remaining put-aside counts are bounded by
\[
 U_L=
 |N_{\rm main+0}^S-N_{\rm main+0}^T|
       +|N_1^S-N_1^T|+|N_2^S-N_2^T|.
\]
Each difference is centered with variance \(O(V)\), so
\(\E U_L=O_\rho(\sqrt V)\). Every put-aside source has the earliest
source class, and every put-aside target the latest target class.
Consequently all opposite-color leftover pairs are forward; pair them
up to the overall color excess. This is a maximum-cardinality matching
of the full input. A crude bound for their quadratic cost is
\(CL^2U_L=o(V)\) when \(q>4\). Even allowing a penalty
\(M_L\le CL^2\), the bound
\[
 C(L^2+M_L+|t|)U_L=o(V)
\]
is sufficient in that dimension.

On any preceding exception, use an arbitrary full maximum-cardinality
matching, at cost at most \(C(L^2+M_L+|t|)N_{\rm tot}\).
Exponential count failures are negligible; the remaining injection
failure has exactly the energy estimate
\eqref{strict-screen-injection-exception}. Since \(M_L\le CL^2\),
these estimates also pay a growing penalty. They prove
\eqref{strict-screen-exceptional-energy} and
\eqref{strict-screen-repair-energy}.
\end{proof}

Deleting omitted main edges increases the comparison value only by the
negative part of their probe, at most \(|t|G\) per omitted source.
Lemma~\ref{strict-screen-depth} makes this \(o(V)\). Repair probes cost
\(C|t|\rho V+o_\rho(V)\). Hence the independent-count comparison, for
\(M_L\le CL^2\), is
\begin{equation}
\label{strict-screen-poisson-upper}
 \limsup_{L\to\infty}\frac{\E[\text{finite cost}]}V
 \le\lambda^s I_+(t\lambda^{2/q})+C\rho^s+C|t|\rho.
\end{equation}
There is no \(M_L\rho\) term. The volume limit precedes
\(\rho\downarrow0\). Since \(I_+\) is \(G\)-Lipschitz in \(t\), the latter
limit gives the sharp value \(I_+(t)\).

Proposition~\ref{strict-screen-three-pools}, including the universal
put-aside classes and \eqref{strict-screen-exceptional-energy}, proves the
supporting input stated at \ref{measurable-forward-strict-input}.
Its use in Lemma~\ref{measurable-forward-homogeneous} requires no iid
claim for the adaptive ports. The final leftover pairs there can be
chosen forward precisely because both put-aside classes are universal.

\subsection{Measurable forward recovery by finite-type augmentation}
\label{measurable-forward-section}

We extend the measurable conditional-law argument to finitely interleaving
radial classes. The new repair changes a bounded number of edges along
each alternating path and estimates its frequency by the original input
counts. Its dimension range is \(q>4\), or \(d=q+1\ge6\).

Let \(D=[0,1]^q\), let \(\theta\) be a bounded nonnegative measurable
probability density on \(D\), and put \(s=1-2/q\). Let
\(K_1,\ldots,K_N\subset\R\) be compact nonempty classes satisfying
\(\max K_h<\min K_{h+1}\), each assigned source or target color.
Write \(K_A,K_B\) for the corresponding unions. The allowed class graph
has an edge \(h\to k\) precisely when \(h<k\), \(K_h\) is a source
class, and \(K_k\) a target class. Every allowed rectangle is complete
bipartite. The laws \(\alpha_z,\beta_z\) are Borel probability kernels
on \(K_A,K_B\). Suppose, for \(\theta(z)\,dz\)-almost every \(z\),
\begin{equation}
\label{measurable-forward-hall}
 \alpha_z((-\infty,r])\ge\beta_z((-\infty,r])
\end{equation}
at each gap between consecutive classes. These finitely many conditions
are equivalent to forward feasibility. Class masses may vanish, and
the cuts may be tight on arbitrary measurable sets of labels.

Let \(e(a,z;b,w)\) be a calibrated genuine-distance defect, or another
nonnegative cost with the following bounds:
\begin{align}
\label{measurable-forward-cost}
 0\le e&\le B,\qquad
 c_0|w-z|^2\le e(a,z;b,w)\le C_0|w-z|^2 &&(a<b),\\
 e(a,z;b,w)&\ge b_0>0 &&(a>b).
 \nonumber
\end{align}
Suppose also that, uniformly on forward pairs,
\begin{equation}
\label{measurable-forward-expansion}
 \left|e(a,z;b,w)-\tfrac12(w-z)^TA_z(a,b)(w-z)\right|
 \le\omega(|w-z|)|w-z|^2,\qquad \omega(r)\longrightarrow0.
\end{equation}
Here \(A\) is continuous on the compact allowed label--mark product,
with common positive upper and lower ellipticity bounds. Each frozen
coefficient has the common-cumulative form
\[
 A_z(a,b)=[T_z(b)-T_z(a)]^{-1},
 \qquad hI\le T_z'(r)\le HI
\]
on a common compact scalar interval containing the classes, for fixed
\(0<h\le H<\infty\). The comparison hypotheses need only hold on the
fixed label carrier of the sampled inputs. Its boundary may be irregular,
and \(\theta\) may vanish there. The backward gap and bounded diameter
of \(D\) imply a global bound \(e\ge c_*|w-z|^2\).

For \(\Psi\in C(D\times K_A\times K_B)\), put \(G=\|\Psi\|_\infty\).
Sample the two colors independently with respective joint laws
\[
 \theta(z)\,dz\,\alpha_z(da),\qquad
 \theta(w)\,dw\,\beta_w(db).
\]
There are \(n\) iid observations of each color; positions and marks
within one observation need not be independent. Define
\begin{equation}
\label{measurable-forward-finite}
 V_n(t)=\frac1n\min_{\sigma\in S_n}\sum_{i=1}^n
 \left[n^{2/q}e(A_i,Z_i;B_{\sigma(i)},W_{\sigma(i)})
       +t\Psi(Z_i,A_i,B_{\sigma(i)})\right].
\end{equation}
Let \(I_{z,+}(u)\) be the unit-intensity stationary source-Palm
infimum of the frozen cost
\[
                   \tfrac12v^TA_z(a,b)v+u\Psi(z,a,b),
\]
over perfect forward
matchings with finite second displacement moment. Auxiliary stationary
randomness is allowed. Existence, attainment and measurability follow
from Propositions~\ref{prop:forward-poisson-criterion}
and~\ref{prop:forward-weak-continuity}. Set
\begin{equation}
\label{measurable-forward-integral}
 P(t)=\int_D\theta(z)^s I_{z,+}(t\theta(z)^{2/q})\,dz,
\end{equation}
with zero integrand on \(\{\theta=0\}\).

\begin{theorem}[Forward pressure with measurable conditional laws]
\label{measurable-forward-pressure}
Under the preceding hypotheses, if \(q>4\), then
\[
             \E\sup_{|t|\le T}|V_n(t)-P(t)|\longrightarrow0
             \qquad(T<\infty).
\]
No strictly positive uniform Hall margin is required. The comparison
uses all observations with their original iid laws and assigns finite
cost on every realization, including forward-infeasible empirical
class counts.
\end{theorem}

The proof uses the strict-Hall comparison proved in
Proposition~\ref{strict-screen-three-pools}. We recall its precise
interface before the finite-type repair and passage to measurable kernels.

\paragraph{Strict-Hall screening interface.}
\label{measurable-forward-strict-input}
For fixed product mark laws of positive mass on their retained classes,
suppose all proper class-cut surpluses are strictly positive.
For \(q>4\), a stationary forward mean minimizer can be screened in
\(Q_L=[0,L]^q\), using independent reserves of fraction \(\rho>0\),
with quadratic repair cost
\(C\rho^sL^q+o_\rho(L^q)\). The retained main law has intensity
\(\lambda=1-\rho\) and cost density
\(\lambda^s I_+(t\lambda^{2/q})\).
The construction uses universal earliest-source/latest-target reserve
classes. Its exceptional event \(E_{L,\rho}\) satisfies
\begin{equation}
\label{measurable-forward-strict-exception}
 \E\!\left[L^2N_{\rm tot}\mathbf1_{E_{L,\rho}}\right]
                         =o_\rho(L^q).
\end{equation}
On its success event the pool completions are forward, apart from
a possible final pool-count completion discussed below. All put-aside
pool-count vertices have universal classes and their expected number
is \(O_\rho(L^{q/2})\).

Proposition~\ref{strict-screen-three-pools} proves this statement.
Its boundary-sensitive injection is Lemma~\ref{strict-screen-injection},
and \eqref{strict-screen-exceptional-energy} supplies the required
whole-cube exception estimate.
In particular, the input is stronger than merely a bound on the number
of bad edges. Its port injection uses the boundary-depth quantity
\[
 S_L=\sum_{\text{source ports }x}
          \operatorname{dist}(x,\partial Q_L)^2,\qquad
 \E S_L=o(L^q),
\]
and an independent expanded-collar reference process to prove
\eqref{measurable-forward-strict-exception}. No selected port process
is assumed to be Poisson or iid.

\subsubsection{A bounded alternating-path repair}

\begin{lemma}[Short augmentation and original-count deficiency]
\label{measurable-forward-augmentation}
For any finite configuration with the preceding complete class rectangles,
a shortest alternating augmenting path uses at most \(N_A\) new edges and
at most \(N_A-1\) old edges, where \(N_A\) is the number of source classes.
For a spatial cube \(C\) of side \(h\), one augmentation increases the
unnormalized objective in \eqref{measurable-forward-finite} by at most
\begin{equation}
\label{measurable-forward-augmentation-cost}
              C N n^{2/q}h^2+2N|t|G.
\end{equation}

Let \(N_h(C)\) be the original input count of class \(h\), and write
\(\mu_h(C)=\E N_h(C)\). If the mean class vector is forward feasible
and each color has mean total \(\mu(C)\), the number \(R(C)\) of
unmatched vertices in any maximum-cardinality forward matching satisfies
\begin{equation}
\label{measurable-forward-deficiency}
 R(C)\le2\sum_h|N_h(C)-\mu_h(C)|,\qquad
 \E R(C)\le C_N\sqrt{\mu(C)}
\end{equation}
for either independent Poisson inputs or the original iid inputs.
\end{lemma}
\begin{proof}
Write the forward steps of a shortest augmenting path as
\[
 x_0\to y_1,\quad x_1\to y_2,\quad\ldots,\quad x_{r-1}\to y_r,
\]
with \(y_i x_i\) old matching edges. The path is simple, \(x_0\)
is unmatched, and \(y_r\) is unmatched. If \(x_i,x_j\), \(i<j\),
have the same source class, completeness of the class rectangle gives
the shortcut \(x_i\to y_{j+1}\). This is not an old matching edge:
\(x_0\) is unmatched, and for \(i>0\) the unique old partner is
\(y_i\ne y_{j+1}\). The shortcut contradicts minimal path length.
Thus \(r\le N_A\). Each new edge has kinetic cost at most
\(C_0q n^{2/q}h^2\); deleting old nonnegative kinetic costs can only
help. Bounding both removed and inserted probe terms proves
\eqref{measurable-forward-augmentation-cost}.

The max-flow/min-cut formula gives the unmatched source count as
\[
 \max_{U\subset\{\text{source classes}\}}
 \left[\sum_{h\in U}N_h(C)
      -\sum_{k\in\mathcal N(U)}N_k(C)\right]_+ .
\]
Subsets of classes suffice because adding the other vertices of an
occupied source class does not add neighbors. The same expression
with mean counts is nonpositive by population feasibility. Hence the
unmatched source count is at most \(\sum_h|N_h-\mu_h|\), and the
target argument gives the first inequality in
\eqref{measurable-forward-deficiency}. In both sampling models
\(\operatorname{Var}N_h\le\mu_h\), so Cauchy--Schwarz gives the second.
Neither independence among class counts nor a distributional assertion
about unmatched positions is used.
\end{proof}

Integrating \eqref{measurable-forward-hall} against \(\theta\) on
any cell proves mean feasibility in that cell. Thus
\eqref{measurable-forward-deficiency} applies with
\(\mu(C)=n\int_C\theta\), including with measurable mark kernels.

\subsubsection{The dyadic construction on the original observations}

Fix \(L\ge1\) and set
\[
 r_n=\lceil n^{1/q}/L\rceil,\qquad
 \ell_n=r_n^{-1},\qquad L_n=n^{1/q}\ell_n\longrightarrow L.
\]
Partition \(D\) into cubes of side \(\ell_n\). In each leaf choose
a minimum-cost matching among maximum-cardinality forward partial
matchings for the actual objective. This finite optimum exists even
if the leaf does not admit a perfect forward matching.

Pad the grid with empty cells to the next dyadic grid. At each parent
start from the union of the child matchings and apply shortest
augmenting paths until maximum forward cardinality is reached.
An augmentation increases the cardinality by one. If the child
residuals have sizes \(R_1,\ldots,R_{2^q}\), there are at most
\(\frac12\sum_iR_i\) augmentations. Previously matched edges may be
changed; Lemma~\ref{measurable-forward-augmentation} includes their
probe contribution.

At parent side \(h=2^k\ell_n\), summing
\eqref{measurable-forward-deficiency} over the \(O(h^{-q})\)
child cells and using Cauchy--Schwarz gives a bound
\(C\sqrt n\,h^{-q/2}\) for their total expected residual count.
Consequently the normalized expected kinetic increase is at most
\begin{equation}
\label{measurable-forward-dyadic}
 C\sum_{k\ge1}(2^kL_n)^{2-q/2}
                  \le C_qL_n^{2-q/2},
\end{equation}
and the expected absolute normalized probe increase is at most
\begin{equation}
\label{measurable-forward-probe-repair}
                    C_N|t|G L_n^{-q/2}.
\end{equation}
The geometric series in \eqref{measurable-forward-dyadic} converges
because \(q>4\). Its constants do not depend on Hall margins,
conditional mark densities or the positions of selected residuals.

At the final root each color has \(n\) observations. By
\eqref{measurable-forward-deficiency}, the unmatched count per color
has expectation at most \(C_N\sqrt n\). Retain all legal edges and
apply the bounded-cost completion from the proof of
Proposition~\ref{prop:general-upper-energy}. Its normalized contribution is
\begin{equation}
\label{measurable-forward-root}
       CB n^{2/q-1/2}+C|t|G n^{-1/2}=o(1).
\end{equation}
Only this last completion may insert backward edges. The resulting
perfect matching uses all original observations. Starting the same
hierarchy at microscopic unit scale, with any maximum forward matching
in each leaf, also proves a uniform coarse energy upper bound.

\subsubsection{Homogeneous maximum-forward partial values}

Freeze feasible product mark laws, a common-cumulative coefficient
\(A\), and a continuous paired probe \(g\), with \(\|g\|_\infty\le G\).
For independent product-marked Poisson inputs of intensity \(a\ge0\)
per color in \(Q_L\), let \(H_L(a,t)\) be the minimum of
\(\sum(k+tg)\) among maximum-cardinality forward partial matchings.
Set \(b_L(a,t)=L^{-q}\E H_L(a,t)\). The value is finite on every
count realization.

\begin{lemma}[Sharp homogeneous leaf limit]
\label{measurable-forward-homogeneous}
For \(q>4\), using the strict-Hall supporting input above,
\begin{equation}
\label{measurable-forward-leaf-limit}
 b_L(a,t)\longrightarrow a^s I_+(t a^{2/q}),
\end{equation}
with zero value for \(a=0\). Uniformly over feasible mark laws,
including tight cuts and zero class masses,
\begin{equation}
\label{measurable-forward-leaf-bound}
       -|t|Ga\le b_L(a,t)\le C a^s+|t|Ga.
\end{equation}
The constant depends only on dimension, the finite graph, and the
upper ellipticity bound.
\end{lemma}
\begin{proof}
First take \(a=1\), positive class masses and strict proper Hall cuts.
In the strict supporting construction, the first screening pool leaves
its count excess in universal classes, the bridge pool uses only universal
classes, and the ordinary reserve deletes any excess from the earliest
source or latest target class. Every final put-aside vertex is therefore
universal. Opposite-color leftovers can all be paired forward; their
expected number is \(O_\rho(L^{q/2})\), so their expected cost per
volume is \(O_\rho(L^{2-q/2})+O_\rho(|t|L^{-q/2})=o(1)\).
The success-event matching has cardinality \(\min(N_S,N_T)\),
hence maximum forward cardinality, as established in
Proposition~\ref{strict-screen-three-pools}.

On any exception use an arbitrary maximum forward partial matching.
Its absolute cost is at most \(C(L^2+|t|G)N_{\rm tot}\);
\eqref{measurable-forward-strict-exception} pays this fallback.
The retained minimizer and repair estimates give, first as \(L\to\infty\)
at fixed \(\rho\), and then as \(\rho\downarrow0\),
\[
                    \limsup_L b_L(1,t)\le I_+(t).
\]
This uses an energy-weighted exceptional-event estimate; a vanishing
fraction of illegal edges alone would not justify the conclusion.

For a general fixed feasible law, omit zero-mass classes and split at
every proper zero-surplus cut. Each resulting block has equal positive
source and target mass \(p_j\) and strictly positive internal cut
surpluses; there are finitely many blocks and \(\sum_jp_j=1\).
In any unit-intensity stationary forward perfect matching, the intensity
of edges crossing a radial cut \(r\) is
\(F_\alpha(r)-F_\beta(r)\): subtract target mass on the left from source
mass on the left, and use forward orientation. At a tight cut this
nonnegative intensity is zero. Stationarity and a countable spatial
cover then exclude every such edge almost surely. Thus
\begin{equation}
\label{measurable-forward-block-splitting}
 I_+(t)=\sum_j p_j^s I_{j,+}(t p_j^{2/q}).
\end{equation}
For the lower comparison restrict a matching to each block; other blocks
may serve as auxiliary randomness. For the upper comparison superpose
independent block minimizers. Both preserve the prescribed independent
Poisson inputs.

Apply the strict construction separately to the blocks at intensities
\(a p_j\). Their union has expected total residual count
\(O_{\alpha,\beta,a}(L^{q/2})\) by
\eqref{measurable-forward-deficiency}. Augment the union to a global
maximum forward matching. By
\eqref{measurable-forward-augmentation-cost}, its extra expected cost
per volume is
\[
                  O(L^{2-q/2})+O(|t|L^{-q/2})=o(1).
\]
Together with \eqref{measurable-forward-block-splitting} and spatial
scaling, this proves the upper bound in
\eqref{measurable-forward-leaf-limit}. No positive block mass or
internal margin is required to be uniform over varying frozen laws.

For \eqref{measurable-forward-leaf-bound}, start a homogeneous dyadic
hierarchy at side \(a^{-1/q}\). Arbitrary maximum forward leaf
matchings have expected kinetic cost \(C a^sL^q\). At a coarser
scale \(h\), the extra expected kinetic cost per volume is at most
\(C\sqrt a\,h^{2-q/2}\). Summing from \(h=a^{-1/q}\) gives
\(Ca^s\). If \(L<a^{-1/q}\), the whole-cube bound
\(CaL^{q+2}\le Ca^sL^q\) suffices. Padding affects only constants.
The probe has absolute expectation at most \(|t|GaL^q\).
The lower bound follows from nonnegative kinetic cost and at most
\(N_S\) matched edges.

This uniform comparison and lower ellipticity give a uniform second
displacement moment per volume for the partial optima.
Their unmatched density tends to zero by
\eqref{measurable-forward-deficiency}. Uniform spatial rooting and
the two-endpoint compactness argument in
Proposition~\ref{prop:highd-local-limit} retain both endpoints:
long-edge probabilities are \(O(R^{-2})\), and unmatched-root
probabilities vanish. Every limit is a stationary perfect forward
matching of the prescribed independent marked Poisson inputs.
The allowed class relation is closed. Add \(|t|G\) to each matched
edge cost and use compact nonnegative cost truncations. The matched
intensity tends to \(a\), so the added constant cancels. The limiting
stationary cost is at least \(a^s I_+(ta^{2/q})\), proving the
lower bound in \eqref{measurable-forward-leaf-limit}.
\end{proof}

\subsubsection{Fixed microscopic blocks and measurable kernels}

Let \(J_n(C,t)\) denote the leaf's maximum-forward partial objective
before division by \(n\), and let \(C_n(z)\) be its leaf containing \(z\).
At almost every \(z\), Lebesgue differentiation holds simultaneously for
\(\theta\) and
\[
 w\longmapsto\theta(w)\int f\,d\alpha_w,\qquad
 w\longmapsto\theta(w)\int g\,d\beta_w
\]
for countable dense families of continuous mark tests.
After rescaling the leaf by \(n^{1/q}\), its two input clouds converge
jointly to independent marked Poisson clouds in \([0,L]^q\) of
intensity \(\theta(z)\), with product mark laws \(\alpha_z,\beta_z\).
The containing cubes are a bounded-eccentricity differentiation basis.
For a continuous position test times a mark test, the error is bounded
by the normalized integral of the corresponding absolute Lebesgue-point
error on the leaf. The iid rare-event limit gives Poisson counts, and
the originally independent colors remain independent.
At \(\theta(z)=0\), both limiting clouds are empty.
This is a local input limit, not a product-law assertion for the
finite sampled marks and positions.

The classes are clopen relative to their respective compact mark spaces.
Along convergent finite configurations their class assignments therefore
eventually agree. The collection of legal partial matchings and its
maximum cardinality stabilize. Taking a minimum over the finitely many
maximum-cardinality costs is consequently continuous, even for atomic
mark laws or tied optima. Continuity of \(A,\Psi\) and
\eqref{measurable-forward-expansion} give convergence of the rescaled
forward costs on bounded-count configurations. No backward edge enters
a leaf value. For fixed \(L,T\),
\begin{equation}
\label{measurable-forward-leaf-domination}
 |J_n(C,t)|\le C_{L,T}(N_S(C)+N_T(C)),\qquad |t|\le T.
\end{equation}
The count means are bounded by \(\|\theta\|_\infty L_n^q\);
their second moments are uniformly bounded. Expectations thus converge,
and almost everywhere
\[
 \frac{\E J_n(C_n(z),t)}{n|C_n(z)|}
        \longrightarrow b_L(\theta(z),z,t),
\]
where \(z\) records the frozen coefficient, laws and probe.
The domination is uniform in \(z\) at fixed \(L,T\), so
\begin{equation}
\label{measurable-forward-leaf-integral}
 \frac1n\sum_C\E J_n(C,t)
       \longrightarrow\int_D b_L(\theta(z),z,t)\,dz.
\end{equation}
These expectations are measurable by conditioning on finite counts
and integrating finite minima against Borel kernels. This argument
also gives measurability of their limiting integrated stationary value.

\begin{proof}[Proof of Theorem~\ref{measurable-forward-pressure}]
At fixed \(L\), the dyadic construction, genuine final completion,
and \eqref{measurable-forward-leaf-integral} give
\begin{equation}
\label{measurable-forward-upper-fixed}
 \limsup_n\E V_n(t)
 \le\int_D b_L(\theta(z),z,t)\,dz
                +CL^{2-q/2}+C|t|G L^{-q/2}.
\end{equation}
For almost every positive-density root,
Lemma~\ref{measurable-forward-homogeneous} identifies the limit as
\(L\to\infty\). Its bound is dominated by
\(C\theta^s+TG\theta\), integrable on \(D\) and independent of
small Hall margins, tight-cut decompositions, and positive class masses.
Thus, first taking \(n\to\infty\) at fixed \(L\) and then
\(L\to\infty\), we obtain
\begin{equation}
\label{measurable-forward-upper}
                       \limsup_n\E V_n(t)\le P(t).
\end{equation}
All exceptional leaf count configurations are covered by their finite
maximum-forward partial values. No growing-penalty diagonal is used.

The coarse comparator, the global bound \(e\ge c_*|w-z|^2\), and
the fixed backward gap yield, for every exact optimizer \(\sigma\),
\begin{equation}
\label{measurable-forward-optimizer-bounds}
 \sup_n\E\frac1n\sum_i n^{2/q}
       |Z_i-W_{\sigma(i)}|^2<\infty,\qquad
 \E\frac{\#\{\text{backward edges}\}}n=O(n^{-2/q}),
\end{equation}
uniformly for bounded \(t\). Apply the measurable-input local Poisson
argument just proved and the two-endpoint scene compactness in
Propositions~\ref{prop:rough-ray-poisson}
and~\ref{prop:highd-local-limit}. Along a subsequence realizing the
expectation liminf, the limiting conditional scenes at almost every
positive-density root are stationary perfect matchings of the correct
independent Poisson inputs. The second bound in
\eqref{measurable-forward-optimizer-bounds} and closed class legality
make them forward. Compact endpoint cutoffs, the relative expansion,
and the nonnegative shift by \(|t|G\) give the lower mean-cost
comparison. After exhaustion the intensity scaling is exactly
\eqref{measurable-forward-integral}, hence
\begin{equation}
\label{measurable-forward-lower}
                       \liminf_n\E V_n(t)\ge P(t).
\end{equation}
This is the lower compactness and truncation argument in
Proposition~\ref{prop:highd-measurable-pressure}, restricted to the
closed allowed rectangles. The ordered proposition's all-pairs
admissibility was needed for its upper repair, which has been replaced
here. Neither energy uniform integrability nor convergence of
untruncated optimizer energies is required.

For completeness, replacing one original source or target observation
changes the optimum in \eqref{measurable-forward-finite} by at most
\[
                   \frac{Bn^{2/q}+2TG}{n},
\]
by keeping the same permutation for each comparison direction.
There are \(2n\) independent observations. Efron--Stein gives
\begin{equation}
\label{measurable-forward-variance}
             \operatorname{Var}V_n(t)
                 \le C_{B,T,G}n^{4/q-1}\longrightarrow0.
\end{equation}
Together with \eqref{measurable-forward-upper} and
\eqref{measurable-forward-lower}, this proves
\(\E|V_n(t)-P(t)|\to0\). Both \(V_n\) and \(P\) are
\(G\)-Lipschitz in \(t\); for \(P\), the powers of \(\theta\)
sum to one and \(\int\theta=1\).
A finite parameter net proves expected locally uniform convergence.
\end{proof}

\begin{lemma}[Scalar stationary differentiability]
\label{measurable-forward-c1-input}
For the original scalar common-cumulative finite-class model,
including tight cuts, the stationary forward value \(I_+(t)\) is \(C^1\)
in every continuous paired-mark probe direction, with derivative bounded
by the probe norm.
\end{lemma}
\begin{proof}
Theorem~\ref{forward-full-pressure-c1} proves differentiability of the
limiting penalty pressure by a joint large-penalty/vanishing-probe
argument. Theorem~\ref{forward-full-stationary-identity} identifies that
pressure with \(I_+\), including tight Hall cuts. Its probe-norm
Lipschitz bound gives the derivative bound.
\end{proof}

\begin{corollary}[Scalar paired-probe readout]
\label{measurable-forward-readout}
In the scalar model of Theorem~\ref{measurable-forward-pressure},
using the preceding stationary differentiability input, \(P\) is \(C^1\)
and
\[
 P'(t)=\int_D\theta(z)I_{z,+}'(t\theta(z)^{2/q})\,dz.
\]
For every exact unperturbed empirical optimum \(\gamma\),
including fractional optimal couplings,
\[
 \sup_{\gamma\in\operatorname{Opt}_n}
 \left|\int\Psi(z,a,b)\,d\gamma-P'(0)\right|
                           \longrightarrow0
                           \quad\hbox{in probability}.
\]
\end{corollary}
\begin{proof}
The derivative integrand is bounded in absolute value by
\(G\theta(z)\), since \(s+2/q=1\). Dominated differentiation
and then continuity of the derivative give \(P\in C^1\).
The two-sided finite-difference sandwich, using the pressure theorem
at neighboring fixed parameters, controls all exact optimizers
simultaneously. Every optimal fractional coupling has a Birkhoff
decomposition supported on optimal permutations and hence satisfies
the same bounds.
\end{proof}

\begin{proposition}[Compact decorated marks: pressure convergence]
\label{measurable-forward-decorated}
Each class \(K_h\) may instead be a compact metric space with a continuous
scalar radial coordinate \(r:K_h\to\R\), whose images have the same
strict separation and fixed source/target order. The map \(r\) need
not be injective. Replace the kinetic coefficient by
\[
           A_z(a,b)=[T_z(r(b))-T_z(r(a))]^{-1},
\]
and keep the full compact marks in the conditional laws and probe.
If the same continuity, ellipticity, backward gap, forward comparison
and uniform relative expansion hold for this decorated model, the
pressure conclusion of Theorem~\ref{measurable-forward-pressure}
holds with its stationary forward value.
\end{proposition}
\begin{proof}
The graph, deficiency and augmentation arguments use only class
membership. Compact-mark weak convergence and clopen class membership
also give the fixed-leaf continuity argument without a scalar inverse.
Finite-energy stationary competitors split by allowed class pair and
restore the full conditional marks by independent marking.
Attainment and lower compactness use continuous costs on compact marks.

It remains to check the strict homogeneous supporting comparison.
A stationary mean minimizer with full marks satisfies the same boundedly
perturbed two-cycle inequality on each allowed rectangle.
Its kinetic trajectory and phase coordinates depend only on
\(r(a),r(b)\). In the traffic estimate of
Theorem~\ref{thm:highd-kinetic-traffic} and
Lemma~\ref{lem:highd-perturbed-traffic}, the full ambient marked
Poisson occupation integrals can be taken against the radial
pushforward law. These estimates require no inverse radial map,
density of radial marks, or unique mark over a radial value.
The selected-rectangle extension is given explicitly in
Lemma~\ref{forward-full-traffic-centering}.

Full-mark-dependent translations preserve the product marked Poisson
input and its independence from the other color. The same conditional
velocity identities therefore center the kinetic current.
The boundary-current screen, universal class reserves and independent
bridge injection use positions and class intensities, which are
unchanged. The strict supporting comparison, tight-cut block splitting,
and all steps through \eqref{measurable-forward-variance}
therefore apply to the decorated model.
\end{proof}

The compact decorated model retains the common-cumulative representation,
a fixed finite class pattern, uniform ellipticity and \(q>4\).
It permits pressure comparison with the original mark retained beside a
clipped radial coordinate; arbitrary off-support matrix extensions need
not preserve its structure. The scalar readout uses a separate
differentiability input. Degenerate weighted pressures are treated by
Theorem~\ref{thm:general-focal-value-equality}.


% Conversion of general-rough-pressure-gluing.md
% source SHA256 46bd3d24b6abe0d1f9573f0b66ec992f461a6633a74c4d9b792d6203d65a0d33
% independent mathematical review SHA256
% 848b23cff33f414eb263f3aa524f0621e185785bd519cd365a5c312b95947c4f
% The compact relative-expansion proof also converts the accepted source
% rough-relative-expansion-next.md, SHA256
% c7010cee349e9b4099e369c24cd979603084f8d9a904faa9327f97d067eee3d7.
% Insert after the measurable-forward-* fragment. No new macros required.
\section{Global localization and weighted stationary pressure}\label{sec:global-weighted-pressure}

\subsection{General rough-ray localization and pressure gluing}
\label{general-gluing-section}

Retain the assumptions of Proposition~\ref{prop:general-optimal-face}:
\((M,F)\) is connected, complete, smooth, reversible and strongly convex,
and \(\mu,\nu\) have nonnegative Lipschitz densities relative to a smooth
positive volume, with compact disjoint supports \(A,B\). Fix a population
optimal potential \(\phi\), oriented so that transport increases \(\phi\),
and put
\[
 \begin{gathered}
 \delta=\operatorname{dist}_F(A,B)>0,\qquad
 q=d-1>4,\qquad s=1-\frac2q,\\
 e(x,y)=d_F(x,y)-\phi(y)+\phi(x)\ge0.
 \end{gathered}
\]
For a bounded continuous \(f:A\times B\to\R\), fix a bounded continuous
extension to the compact whole-flow endpoint product, preserving its
supremum norm by Tietze's theorem. Only auxiliary off-support evaluations
use that extension. For independent iid empirical probabilities
\(\mu_n,\nu_n\), use the original pressure \(V_n(t)\) from
\eqref{eq:general-finsler-pressure}. Its unperturbed minimizers
are the entire original distance-optimal set.

We prove a lower bound using stationary forward matchings with finite
weighted quadratic energy and an upper bound using finite Euclidean
displacement second moments. Exact radial clipping identifies the
second value. Theorem~\ref{thm:general-focal-value-equality} proves
their equality on the actual conditional populations, including
interleaving supports and focal endpoints. These bounds yield the
general selector in Theorem~\ref{thm:general-finsler-selector}.

\subsubsection{Finite active sheets and compatible compact cores}

Refine the finite whole-ray atlas of
Proposition~\ref{prop:general-upper-atlas} by the intrinsic active-block
index of Theorem~\ref{thm:intrinsic-active-invariance}. There are uniformly
finitely many blocks per ray. Indeed, every positive block contains a
calibrated source-target segment of length at least \(\delta\); the blocks
are disjoint and the population potential range is bounded. Denote the
resulting finite sheets by \(j\), with bounded coordinate cubes \(Q_j\).
Their population restrictions have the form
\begin{equation}
\label{general-gluing-sheet-disintegration}
 \begin{split}
 x&=X_j(a,z),\qquad y=X_j(b,z),\qquad
 \Lambda_j(dz)=\theta_j(z)\,dz,\\
 \mu_j&=\theta_j(z)\,dz\,\alpha_{j,z}(da),\qquad
 \nu_j=\theta_j(z)\,dz\,\beta_{j,z}(db).
 \end{split}
\end{equation}
The densities \(\theta_j\) are bounded and measurable, and the conditional
probabilities need only be measurable. Every population optimum preserves
\((j,z)\), is forward, and stays in its intrinsic active block. The TV
continuity for whole-ray laws in
Proposition~\ref{prop:compact-conditional-whole} need not survive restriction
to a particular active block; it is not required here.

\begin{lemma}[Compact relative estimates in the finite upper-atlas labels]
\label{general-gluing-compatible-cores}
In the same transverse coordinates as
Proposition~\ref{prop:general-upper-atlas}, one can retain finitely many
compact cores with arbitrarily small omitted whole-label quotient mass.
Their characteristic maps have \(C^1\) extensions agreeing with the actual
rays and their full first jets. On every compact subset of the actual ray
interiors, these extensions are nonsingular and have a local \(C^2\)
eikonal potential agreeing with \(\phi\) on the retained rays.

In particular, on a fixed enlarged interior collar and a compact core
\(K\), for forward pairs with \(b-a\ge\delta'>0\),
\begin{equation}
\label{general-gluing-relative-expansion}
 \left|e(X(a,z),X(b,w))-\frac12(w-z)^T S_z(a,b)^{-1}(w-z)\right|
 \le\omega_K(|w-z|)|w-z|^2,
\end{equation}
uniformly for \(z,w\in K\), with \(\omega_K(h)\to0\) as
\(h\to0\), where
\begin{equation}
\label{general-gluing-local-compliance}
 \begin{split}
 J(r,z)&=D_zX(r,z),\\
 H(r,z)&=J(r,z)^TD_{vv}^2F(X(r,z),\partial_rX(r,z))J(r,z),\\
 S_z(a,b)&=\int_a^b H(r,z)^{-1}\,dr.
 \end{split}
\end{equation}
The matrices are continuous and uniformly positive definite on these
compact parameter sets. For such pairs there are also constants
\(0<c_K\le C_K<\infty\) such that
\begin{equation}
\label{general-gluing-core-coercivity}
 c_K|w-z|^2\le e(X(a,z),X(b,w))\le C_K|w-z|^2.
\end{equation}
The same conclusions hold on compact variable interior domains. No
nonsingularity, injectivity or calibration is asserted beyond the actual
ray domains.
\end{lemma}

\begin{proof}
A finite upper-atlas seed is a \(C^{1,1}\) level graph \(\sigma\) in an
ambient coordinate box. Its graph Hessian exists almost everywhere.
Restrict the second jets by Lusin and Egorov so that the Hessian is
continuous and the gradient Taylor remainders are uniform. The values,
gradients and Hessians then satisfy the order-two Whitney conditions:
for retained labels \(z,z+h\), uniformly,
\[
 \begin{split}
 D\sigma(z+h)-D\sigma(z)-D^2\sigma(z)h&=o(|h|),\\
 \sigma(z+h)-\sigma(z)-D\sigma(z)h
                  -\tfrac12D^2\sigma(z)[h,h]&=o(|h|^2).
 \end{split}
\]
For the second identity integrate the gradient remainder on the ambient
seed segment; the original \(C^{1,1}\) graph is defined on its neighborhood.
The finite-order Whitney theorem gives a \(C^2\) graph matching these jets
on an arbitrarily large compact label set. Its oriented unit conormal is
a smooth function of its position and tangent plane. Thus it matches the
original \(P_0,DP_0\), as well as \(\sigma,D\sigma\). Smooth kinetic
Hamiltonian flow gives a \(C^1\) extension of the full phase map with all
these first jets, over the bounded population potential range.

We justify base nonsingularity without requiring ambient twice
differentiability of \(\phi_k\) on the fixed level seed. Use countably
many rational enlarged collars and strain indices \(k\). For each such
choice let \(E\) be the Borel labels whose whole collar lies in the strain
set. Retain differentiability points of the initial phase data and density
points of every such \(E\) to which they belong. This removes only a null
set. On that collar and \(E\),
\[
 P(r,z)=d\phi_k(X(r,z)).
\]
Approaching any parameter direction through \(I\times E\), differentiating
the phase map, and using the Lipschitz differential of \(\phi_k\) give
\begin{equation}
\label{general-gluing-phase-domination}
 |\delta P|\le\Lip(d\phi_k)|\delta X|.
\end{equation}
The phase map \((r,z)\mapsto\Phi_{r-c}(\sigma(z),P_0(z))\) is an
immersion: at the seed the tangent directions and Hamiltonian direction
are independent, since the unit conormal vanishes on seed tangents and
evaluates to one on the unit velocity; Hamiltonian flow preserves this
independence. Hence \(\delta X=0\) would force \(\delta P=0\) and
contradict immersion. Countably many collars prove nonsingularity for
every interior time at each retained good label. The density points used
here belong to the preliminary Borel sets \(E\); we do not require every
point of the final compact core to be its own density-one point.

Select compact subsets of this conull set and include Lusin continuity
of the clipped maximal-ray endpoints, as in
Proposition~\ref{prop:compact-conditional-whole}. On a compact actual
interior domain the extended base map \(Y\) is nonsingular. It is injective
on a neighborhood of that domain: otherwise colliding pairs in shrinking
neighborhoods would converge to a collision of actual core rays; injectivity
there forces equal limiting parameters, contradicting local invertibility.
The propagated covector is \(C^1\), and the canonical one-form pulls back
to \(dr\), because the extended seed is a level surface and the energy is
\(1/2\). Thus \(\psi(Y(r,z))=r\) defines a \(C^2\) local eikonal
potential, with \(d\psi\) the propagated covector. It contains every
selected core ray exactly. It need not be a global Lipschitz calibration.

Here is the uniform relative calculation. The pairs
\((X(a,z),X(b,z))\) with \(b-a\ge\delta'\) in the compact actual
interior domain have unique minimizing segments, by the calibrated initial
velocity and strong convexity. They extend minimizing beyond both endpoints;
the fixed-ray index argument of Proposition~\ref{prop:rough-ray-quadratic}
excludes endpoint conjugacy. The true global distance is smooth near each
pair. Completeness and compactness of bounded-length minimizing segments
exclude a competing branch in a smaller neighborhood, as in that
proposition. Compactness of the pair set gives a common neighborhood.

On it put
\[
 \mathcal E(x,y)=d_F(x,y)-\psi(y)+\psi(x).
\]
This is \(C^2\), and its value and full first derivative vanish at every
diagonal ray pair. Uniform Taylor expansion in ambient coordinates gives
\[
 \mathcal E(X(a,z),Y(b,w))
 =\tfrac12\Delta_y^TD^2_{yy}\mathcal E(X(a,z),X(b,z))\Delta_y
                         +o(|\Delta_y|^2),
\]
where \(\Delta_y\) is the coordinate difference from \(Y(b,z)\).
Uniform \(C^1\) differentiability gives
\(\Delta_y=J(b,z)(w-z)+o(|w-z|)\). For \(w\in K\),
\(\psi(Y(b,w))=b\) exactly, so the left side is the original defect.
The second-variation calculation in
Proposition~\ref{prop:rough-ray-quadratic} identifies its matrix by minimizing
\[
 \int_a^b v'(r)^TH(r,z)v'(r)\,dr,
 \qquad v(a)=0,\quad v(b)=h.
\]
Completing the square gives the minimum \(h^TS_z(a,b)^{-1}h\).
This is an ambient Hessian calculation with the matching first phase jets,
not a second differentiation of a merely \(C^1\) coordinate map. It proves
\eqref{general-gluing-relative-expansion}.

The columns of \(J\) span the kernel of the unit covector, and the length
Hessian has only the radial kernel. Therefore \(H\) is positive definite.
Its continuity and compactness give positive eigenvalue bounds, and the
positive radial gap gives such bounds for \(S\) and \(S^{-1}\).
The relative estimate proves \eqref{general-gluing-core-coercivity} near
equal labels. Away from equal labels, compactness and the global calibration
give the claim: a zero defect between two interior rays would identify
their calibrated geodesics and hence their labels. A finite cover gives
these arguments on compact variable domains too.

Finally, finite compact cores capture any prescribed fraction of the
absolutely continuous quotient measure. Assign overlapping cores by first
membership. Their Borel differences retain every estimate of their
containing cores. No regular boundary or positive distance between the
assigned pieces is required.
\end{proof}

\subsubsection{Actual optimizer localization and retention of both partners}

Proposition~\ref{prop:general-upper-energy}, with its bounded-probe
comparison, implies uniformly for \(t\) in compact intervals and every
optimizer \(\gamma_n\) of \eqref{eq:general-finsler-pressure},
\begin{equation}
\label{general-gluing-energy}
 \E V_n(t)\le C_t+o(1),\qquad
 \E\left[n^{2/q}\int e\,d\gamma_n\right]\le C_t.
\end{equation}
Thus \(\int e\,d\gamma_n\to0\) in probability. Compactness and empirical
marginal convergence imply that every weak limiting plan is a population
optimum.

\begin{lemma}[Invariant finite Borel partitions]
\label{general-gluing-invariant-partitions}
Let \(b_A:A\to\mathcal B\) and \(b_B:B\to\mathcal B\) be fixed finite
Borel labels. If every population optimum preserves these labels, then,
uniformly over optimizers of \eqref{eq:general-finsler-pressure},
\begin{equation}
\label{general-gluing-partition-retention}
 \int\mathbf1_{\{b_A(x)\ne b_B(y)\}}\,d\gamma_n
                         \longrightarrow0
                 \quad\text{in probability}.
\end{equation}
This applies to any fixed finite partition by active sheets, whole-label
cores and their complement.
\end{lemma}

\begin{proof}
For \(\eta>0\), take compact Lusin sets in \(A,B\), each omitting at
most \(\eta\) of its marginal, on which the corresponding labels are
continuous. Extend their disagreement indicator continuously from the
compact product to \(A\times B\), with values in \([0,1]\). Every
population optimum integrates this extension to at most \(2\eta\),
because it vanishes at its retained pairs. Empirical masses outside the
fixed Lusin sets converge by the iid law of large numbers. Along any
subsequence with vanishing slack and weak plan convergence, the limsup
of the disagreement mass is therefore at most \(4\eta\). Let \(\eta\)
decrease to zero and use the subsequence criterion. A measurable choice
of a worst optimizer would contradict the same argument if uniformity
failed. Fractional optima cause no change: the relevant extrema are
attained on a finite optimal permutation face.
\end{proof}

No positive calibration gap across distinct active blocks was used. Such
blocks can have calibrated cross-pairs which cannot occur with positive
mass in any balanced population optimum.

\begin{lemma}[Retention of both microscopic partners]
\label{general-gluing-partner-retention}
For an edge in one active sheet write \(D_n=n^{1/q}(w-z)\), and count a
cross-sheet edge as an escape. Choose either a uniformly sampled source
or a uniformly sampled target and its partner. Then, uniformly over the
empirical optimizers under consideration,
\begin{equation}
\label{general-gluing-no-escape}
 \lim_{R\to\infty}\limsup_{n\to\infty}
  \Pp(\text{escape or }|D_n|>R)=0.
\end{equation}
For fractional plans use a data-measurable Birkhoff mixture of optimal
permutations, drawn before the independent root.
\end{lemma}

\begin{proof}
Fix \(\varepsilon>0\). On almost every ray, compact interior intervals
exhaust both conditional endpoint laws. Choose the first rational interval
capturing at least \(1-\varepsilon\) of each, with an enlarged interval
strictly inside the ray. Finitely many resulting whole-label classes
capture all but \(\varepsilon\) quotient mass. Restrict them to the
cores of Lemma~\ref{general-gluing-compatible-cores}. The omitted mass
of either endpoint is at most \(C\varepsilon\).

Clock-window boundaries are controlled on the original populations.
For their finite set \(T\) of rational thresholds, put
\(C_h=\bigcup_{r\in T}\{x:|\phi(x)-r|<h\}\). The original
conditional clock densities give
\(\mu(C_h)+\nu(C_h)\to0\) as \(h\downarrow0\). Every finite
plan has the prescribed marginals, so its expected mass touching
\(C_h\) is at most this sum; the same bound follows from the exact
mean marginals for a Dirichlet allocation. For a bounded open spatial
window \(W\), Haar averaging bounds either anchored local readout by
\[
 |W|\bigl(\mu(C_h)+\nu(C_h)\bigr).
\]
Open-set lower semicontinuity preserves this bound through both weak
stages. Let \(h\downarrow0\) afterwards. Countability gives one
probability-one event on which all rational endpoint boundaries are null.

On a retained core, backward edges have \(e\ge\delta\) and have vanishing
mass by \eqref{general-gluing-energy}. For forward pairs,
\[
 e(X(a,z),X(b,w))\ge c_\varepsilon|w-z|^2.
\]
Near equal labels the needed radial gap follows from
\(\delta\le |b-a|+L|w-z|\); use
Lemma~\ref{general-gluing-compatible-cores}. Away from equal labels,
compactness and uniqueness of calibrated interior rays give the lower
bound as in its proof. The mass crossing the fixed finite core partition
vanishes by Lemma~\ref{general-gluing-invariant-partitions}. Therefore
\begin{equation}
\label{general-gluing-two-scale-tail}
 \limsup_n\Pp(\text{escape or }|D_n|>R)
       \le C\varepsilon+\frac{C_\varepsilon}{R^2}.
\end{equation}
The same bound holds from the target anchor because every matching edge
has one endpoint of each color. First increase \(R\), then decrease
\(\varepsilon\). This does not assert a global unweighted second-moment
bound. Every positive-weight permutation in a Birkhoff decomposition of
an optimal plan is optimal for the perturbed objective, so all estimates
integrate over the indicated draw.
\end{proof}

\subsubsection{Two stationary classes on the full conditional laws}

At a good active root suppress \(j,z\) and define
\begin{equation}
\label{general-gluing-full-matrix}
 A_z(a,b)=\left[\int_a^b H(r,z)^{-1}\,dr\right]^{-1},
 \qquad a<b,
\end{equation}
using \eqref{general-gluing-local-compliance}. Conditional endpoints are
interior to their maximal ray almost surely. The matrix \(H\) is continuous
and positive definite there. Its upper bound is uniform over the finite
atlas and the bounded population potential range, since the original phase
flows have uniformly bounded first derivatives. Consequently
\(A_z(a,b)\le C\mathrm{Id}\) on supported forward pairs, whose scalar
gap is at least \(\delta\). A lower bound is uniform on each interior
collar, and may degenerate when a source approaches the lower ray endpoint
or a target approaches its upper endpoint.

\begin{definition}[Weighted and second-moment stationary pressures]
\label{general-gluing-two-classes}
Let \(I_{j,\mathrm w}(z,v)\) be the infimum of
\begin{equation}
\label{general-gluing-stationary-cost}
 \E^{\rm S}\left[
    \tfrac12 D^TA_z(a,b)D+v f(X_j(a,z),X_j(b,z))\right]
\end{equation}
over stationary forward perfect matchings of the independent unit-intensity
marked Poisson inputs with laws \(\alpha_{j,z},\beta_{j,z}\), with finite
weighted kinetic mean. Here \(D\) is the paired spatial displacement and
\(\E^{\rm S}\) is source-Palm expectation. Stationary auxiliary randomness
is allowed. Let \(I_{j,2}(z,v)\) be the same infimum with the additional
condition \(\E^{\rm S}|D|^2<\infty\).
\end{definition}

The classes are nonempty and, with constants uniform in the finite atlas,
\begin{equation}
\label{general-gluing-uniform-pressure-bounds}
 -|v|\|f\|_\infty\le I_{j,\mathrm w}(z,v)
   \le I_{j,2}(z,v)\le C_q+|v|\|f\|_\infty.
\end{equation}
Indeed, separated compact scalar supports have finitely many monochromatic
runs, uniformly bounded in number by their scalar span and \(\delta\).
A feasible traffic matrix splits each independent Poisson color into pair
types. Match each pair type by a fixed stationary finite-second-moment
two-color scheme, as in \eqref{eq:forward-weak-reserve}. For type masses
\(m_h\), its energy per volume is at most
\(C\sum_hm_h^s\le C N^{2/q}\), including tight Hall cuts.

The weighted infimum is attained. For fixed input laws, increasing interior
collars omit endpoint probabilities tending to zero. At bounded cost,
inside a collar coercivity gives a partner tail \(C_k/R^2\); outside,
charge the prescribed endpoint marginal probability. Applying this at both
anchors retains a perfect matching in a stationary limit. Compact
nonnegative cost truncations on interior collars prove lower
semicontinuity. The bounded probe and closed forward relation on separated
classes pass to the limit. The limiting law may have infinite unweighted
second moment; attainment does not prove equality of the two values.

The values have universally measurable versions in the root. Stationary
scene laws form a standard Borel space; input, degree and translation
constraints are countably expressible, and the shifted cost is a Borel
nonnegative integral. The constrained infimum is lower semianalytic.
This is enough for Lebesgue integration. Clipped uniformly elliptic values
below also have the usual Borel measurability from
Proposition~\ref{prop:forward-weak-continuity}. Set
\begin{equation}
\label{general-gluing-integrated-pressures}
 \begin{split}
 \mathcal P_{\mathrm w}(t)
   &=\sum_j\int_{Q_j}\theta_j(z)^s
       I_{j,\mathrm w}(z,t\theta_j(z)^{2/q})\,dz,\\
 \mathcal P_2(t)
   &=\sum_j\int_{Q_j}\theta_j(z)^s
       I_{j,2}(z,t\theta_j(z)^{2/q})\,dz.
 \end{split}
\end{equation}
The zero-density integrands are zero. Kinetic cost density scales by
\(\theta^s\), whereas the probe scales by \(\theta\). Bounded cubes,
bounded densities and \eqref{general-gluing-uniform-pressure-bounds} make
both expressions finite.

\subsubsection{A pathwise lower bound for actual empirical optimizers}

\begin{proposition}[The weighted stationary lower value]
\label{general-gluing-weighted-lower}
For every fixed \(t\in\R\) and \(\eta>0\),
\begin{equation}
\label{general-gluing-probability-lower}
 \Pp\bigl(V_n(t)<\mathcal P_{\mathrm w}(t)-\eta\bigr)
                  \longrightarrow0.
\end{equation}
At almost every positive-density root, every subsequential random rooted
limit of actual empirical optimizers has independent marked Poisson inputs
and a perfect forward stationary matching with finite weighted kinetic mean.
\end{proposition}

\begin{proof}
Root uniformly in each transverse cube, independently of the data and of
the Birkhoff draw. Lebesgue differentiation of the bounded measurable
coordinate marginal densities, tested against a countable determining
family of continuous mark functions, gives the independent marked Poisson
input limit. No conditional quantile regularity is used.

This input marginal is also deterministic in each random limiting rooted
law. Replacing one observation changes the root integral of a bounded
input-cylinder test of fixed microscopic radius \(R\) by at most
\(C_R/n\), because only roots in that microscopic window can notice the
change. Efron--Stein therefore gives variance \(O_R(n^{-1})\). Its
expectation has the preceding differentiation and Poisson limit; countably
many tests give the quenched assertion. Bounded input intensities give
local count tightness and moment bounds. The deterministic root-translation
identity, with vanishing outer-cube boundary error, gives stationarity.

Lemma~\ref{general-gluing-invariant-partitions} removes cross-sheet and
cross-core edges in integrated rooted mass, and
Lemma~\ref{general-gluing-partner-retention} retains both partners. The
limit is thus a perfect integer matching. Backward mass vanishes by
\eqref{general-gluing-energy}; the positive gap on supported scalar
classes makes its limiting edges forward. On each compact collar,
\eqref{general-gluing-relative-expansion} identifies its kinetic cost.

To obtain a pathwise value bound, jointly extract the scalar values and
the random empirical laws \(\Gamma_n\) of rooted scenes. The bound
\(V_n(t)\ge-|t|\|f\|_\infty\) and
\eqref{general-gluing-energy} give scalar tightness. For a finite disjoint
assignment of cores, fix a compact nonnegative spatial anchor of integral
one, compact endpoint and mark cutoffs inside the actual ray domains, and
a finite clip \(N\) of the local nonnegative edge sum. Its summand is
frozen kinetic cost plus the frozen probe plus
\(B_t=|t|\|f\|_\infty\). For these fixed cutoffs,
\begin{equation}
\label{general-gluing-finite-domination}
 \Gamma_n(F_{R,N})\le V_n(t)+B_t+o_{\Pp}(1).
\end{equation}
Indeed, integrate first the corresponding nonnegative \emph{actual} edge
cost. Disjoint source-core assignment counts each edge at most once,
and root integration of the anchor contributes one. Its integral is at
most \(V_n(t)+B_t\). Replacing it by the compact frozen summand costs
\(o_{\Pp}(1)\), by the relative expansion, coefficient and probe
continuity, and local count moment bounds. A clip cannot increase this
integral. Choose cutoffs at continuity levels of the subsequential law and
then exhaust a countable family. Measurable core boundaries are handled
by differentiation of their indicators, not a boundary-volume estimate.

Pass jointly in \eqref{general-gluing-finite-domination}, then exhaust
count clips, spatial cutoffs, mark collars and finite core assignments.
Campbell's formula yields, almost surely for the limiting value \(V\)
and random law,
\begin{align*}
 V+B_t
 &\ge\sum_j\int_{Q_j}\theta_j(z)
     \E_z^{\rm S}[\text{kinetic at intensity }\theta_j(z)
                                      +t f+B_t]\,dz\\
 &\ge\mathcal P_{\mathrm w}(t)+B_t.
\end{align*}
The shift cancels since total source mass is one. On realizations with
finite \(V\), weighted energy is finite for almost every positive-density
root. This proves \eqref{general-gluing-probability-lower} by the
subsequence criterion. Only compact-collar energy was used before
exhaustion; no unweighted moment was inferred from weighted energy.
\end{proof}

\subsubsection{Exact radial clipping with all conditional marks retained}

Fix a compact jet core with continuous maximal-ray endpoints clipped to a
bounded interval \(R_0\) strictly containing all supported potential values.
Write these clipped endpoints as \(\ell(z),u(z)\), and take the exterior
buffer larger than \(\tau\), where \(0<\tau<\delta/16\). Define
\begin{equation}
\label{general-gluing-radial-clipping}
 a_\tau(z)=\max\{a,\ell(z)+\tau\},\qquad
 b_\tau(z)=\min\{b,u(z)-\tau\}.
\end{equation}
Almost every source has a forward partner at scalar distance at least
\(\delta\), and almost every target has such a predecessor. Thus source
marks have a \(\delta\) margin from the upper maximal-ray endpoint, and
target marks from the lower one, whenever that endpoint lies in \(R_0\).
The clipping stays on the same actual ray, moves the source forward and
the target backward by at most \(\tau\), and preserves the finite
forward graph. Every original conditional mark and its probability remain.

For actual supported endpoints, their calibrated ray segments and the
triangle inequality give exactly
\begin{align}
 d_F(X(a,z),X(b,w))
 &\le a_\tau(z)-a
       +d_F(X(a_\tau(z),z),X(b_\tau(w),w))
       +b-b_\tau(w),\notag\\
 e(X(a,z),X(b,w))
 &\le d_F(X(a_\tau(z),z),X(b_\tau(w),w))
                         -b_\tau(w)+a_\tau(z).
\label{general-gluing-exact-clipping}
\end{align}
There is no additive \(O(\tau)\) term to amplify by \(n^{2/q}\).
The probe continues to evaluate the original endpoints.

For fixed \(\tau\), these clipped segments lie in a compact actual
interior domain. Lemma~\ref{general-gluing-compatible-cores} gives the
uniform relative expansion there, with frozen matrix
\begin{equation}
\label{general-gluing-clipped-matrix}
 A_{z,\tau}(a,b)
   =\left[\int_{a_\tau(z)}^{b_\tau(z)}H(r,z)^{-1}\,dr\right]^{-1}.
\end{equation}
It is uniformly elliptic on that core for fixed \(\tau\), and its upper
bound is independent of \(\tau\): use \(H\le C\mathrm{Id}\) and the
clipped gap \(b_\tau-a_\tau\ge\delta-2\tau\). Moreover
\begin{equation}
\label{general-gluing-matrix-monotonicity}
 A_{z,\tau}(a,b)\downarrow A_z(a,b)
                  \quad\text{in Loewner order as }\tau\downarrow0.
\end{equation}
For an actual cross-label edge the target clip uses \(w\), not \(z\).
First apply the relative expansion with that actual clipped scalar, then
freeze it using continuity of the endpoint functions and coefficients.
The resulting error remains \(o(|w-z|^2)\).

To use a fixed common-interior model, subdivide the core finitely so that
\(\ell,u\) oscillate by less than \(\tau/4\) on each piece. Shrink its
diameter also so the entire scalar source and target projections are
separated by \(\delta/2\), using
\[
 \delta\le |a-b|+L|z-w|.
\]
Their closures are compact with finitely many monochromatic runs. The
opposite-end margins above and the small endpoint oscillation place all
clipped images in one common interior interval.

Use the full decorated marks \((a,a_\tau(z))\), respectively
\((b,b_\tau(z))\): take the graph closures over the compact core and
the original compact scalar class. The radial coordinate is the second
projection, independent of the label and possibly noninjective. At each
label the conditional law is supported on its graph. Small endpoint
oscillation and cross-label scalar separation ensure uniform separation
of the radial images of the entire fixed decorated classes. The coefficient
is common-cumulative in this second projection, while the probe keeps the
original first coordinate. This is precisely the pressure extension in
Proposition~\ref{measurable-forward-decorated}. No arbitrary off-support
matrix extension is assumed to preserve a common-cumulative representation.

\begin{lemma}[The exact limit of clipping pressure]
\label{general-gluing-clipping-limit}
Let \(I_{j,\tau}(z,v)\) be the stationary pressure with
\(A_{z,\tau}\), the unchanged original input mark laws, and the original
probe. Then
\begin{equation}
\label{general-gluing-clipping-value}
 I_{j,\tau}(z,v)\downarrow I_{j,2}(z,v)
                          \qquad(\tau\downarrow0).
\end{equation}
\end{lemma}

\begin{proof}
Each clipped minimizer has finite second moment by coercivity. Its cost
under \(A_z\) is no larger, so \(I_{j,\tau}\ge I_{j,2}\). Conversely
fix one finite-second-moment competitor for \(I_{j,2}\). The matrices
converge as in \eqref{general-gluing-matrix-monotonicity}, and are bounded
above by \(C\mathrm{Id}\). Its kinetic cost therefore converges by
domination by \(C|D|^2\); its probe is unchanged. First decrease
\(\tau\), then approximate the infimum within this class. Monotonicity
is immediate from the matrix order. No corresponding convergence to
\(I_{j,\mathrm w}\) is inferred for a competitor with infinite
unweighted second moment.
\end{proof}

\subsubsection{Admissible global upper recovery}

The local input is Theorem~\ref{measurable-forward-pressure} and its
compact decorated-mark extension,
Proposition~\ref{measurable-forward-decorated}. It permits bounded measurable
common densities, arbitrary measurable conditional kernels and every
feasible fixed finite forward class graph, including tight Hall cuts.
On the fixed core and clipping pieces just constructed its continuous
uniformly elliptic common-cumulative hypotheses hold. Its finite-type
augmenting repair controls actual counts; it does not assert that adaptive
unmatched points are iid. Cost hypotheses are needed only on the fixed
label support carrying the density, since all samples and repair endpoints
lie there; extending the true metric cost off that support is unnecessary.

\begin{proposition}[The finite-second-moment upper value]
\label{general-gluing-second-moment-upper}
For every fixed \(t\in\R\),
\begin{equation}
\label{general-gluing-expectation-upper}
 \limsup_{n\to\infty}\E V_n(t)\le\mathcal P_2(t).
\end{equation}
\end{proposition}

\begin{proof}
Choose finitely many whole-label cores whose complement has quotient mass
at most \(\varepsilon\), assigning overlaps by first membership. Fix
\(\tau\) and make the finite subdivisions above. All these choices use
only population data and precede sampling.

A piece has equal source and target mass \(p_i\). Conditional on the
original count vectors, each color's observations in original index order
are iid with that piece's normalized restriction. Retain the first
\(m=\min(N_i^A,N_i^B)\) from each color. For \(m>0\), multiply the
size-\(m\) local objective by \((m/n)^s\) and use
\begin{equation}
\label{general-gluing-random-count-scaling}
 t_m=t(m/n)^{2/q},\qquad
 (m/n)^s m^{-s}=n^{-s},\qquad (m/n)^st_m/m=t/n.
\end{equation}
The zero-count contribution is zero. Compact-parameter uniform pressure
convergence, \(m/n\to p_i\), the binomial moment bounds and the local
estimate \(C_i m^s+C_i\sqrt m\) justify these random sizes and parameters.

Normalization affects the point law; counts retain the original cell
mass. At frozen intensity \(\theta>0\), an affine box of side \(\ell\)
has mean count \(n\theta\ell^q\). For its degree-two kinetic kernel \(k\),
\[
 (n\theta\ell^q)^{2/q}\theta^{-2/q}k(v)
                         =n^{2/q}k(\ell v).
\]
Thus the unit-box kernel is \(\theta^{-2/q}k\), with unchanged mark law
and probe. Applying \eqref{eq:highd-variable-intensity-scaling} to the
piece density \(\theta/p_i\), the prefactor \(p_i^s\) and tilt
\(t p_i^{2/q}\) gives the original-density integrand
\[
 \theta^s I_{j,\tau}(z,t\theta^{2/q}).
\]

The local theorem uses
\(\Psi(z,a,b)=f(X_j(a,z),X_j(b,z))\). Its difference from the actual
\(f(X_j(a,z),X_j(b,w))\) contributes \(o(1)\) in expectation. For
fixed core and \(\tau\), the forward recovery edges have bounded
normalized squared-displacement expectation. For each \(\rho>0\), the
expected fraction with \(|w-z|>\rho\) tends to zero; the top-level
completion fraction does too. On the other edges, uniform continuity of
the whole flow and the extended probe bounds the difference by a modulus
vanishing with \(\rho\). First increase \(m\), then decrease \(\rho\).
The same conclusion holds after inserting random counts. The local
recovery and \eqref{general-gluing-exact-clipping} now bound the actual
piece cost. The total expected discrepancy of the two counts over the
fixed finite partition is \(O(\sqrt n)\).

The omitted whole-label set is balanced. Restrict any population optimum
to that set, giving a deterministic subplan of mass at most
\(\varepsilon\). Apply the small-mass construction of
Proposition~\ref{prop:general-upper-energy} to that subplan as part of the
full construction. Equivalently, assign its pair types independently to
original points using the marginal Radon--Nikodym probabilities before
matching. Its attributed normalized expected energy is at most
\begin{equation}
\label{general-gluing-remainder-cost}
 C\varepsilon^s+C'n^{2/q-1/2}\sqrt\varepsilon.
\end{equation}
This is a population thinning construction, not an iid claim about points
left by an optimizer. The remaining finite-piece and finite-type count
leftovers have expected size \(O(\sqrt n)\). Pair them by bounded genuine
distance; the normalized cost is \(O(n^{2/q-1/2})\to0\), since \(q>4\).
Their probe mass vanishes, and the omitted population part contributes at
most \(|t|\|f\|_\infty\varepsilon\) to the probe. This constructs an
actual full empirical coupling.

First let \(n\to\infty\), with all cores, pieces and \(\tau\) fixed.
Next let \(\tau\downarrow0\), using
Lemma~\ref{general-gluing-clipping-limit} and the common comparator bounds
for dominated integration. No core ellipticity constant is required to
stay uniform in this limit. Finally let \(\varepsilon\downarrow0\).
For an omitted label set \(E\), H\"older's inequality gives
\begin{equation}
\label{general-gluing-small-label-integral}
 \int_E\theta^s\,dz
 \le\operatorname{Leb}_q(Q)^{2/q}
                  \left(\int_E\theta\,dz\right)^s.
\end{equation}
Thus the omitted kinetic bound vanishes with its quotient mass. Its
signed pressure is bounded below by
\(-|t|\|f\|_\infty\theta\). The retained integral converges to
\(\mathcal P_2(t)\), and \eqref{general-gluing-remainder-cost} vanishes,
proving \eqref{general-gluing-expectation-upper}.
\end{proof}

The order is \(n\to\infty\), then \(\tau\downarrow0\), then
\(\varepsilon\downarrow0\). There is no term
\(\varepsilon n^{2/q}\), rate assumption on Whitney cores, or uniform
estimate for shrinking ray-end collars.

% Source: degenerate-scalar-pressure-readout.md
% Source SHA256: 45dd0c63c9581a3ea8ec7dc7e1976976bd828323c54ada38fe1335ddadad6bd9
\subsubsection{Differentiability with degenerate ray endpoints}
\label{degenerate-scalar-section}

The second-moment pressure remains differentiable when its quadratic
coefficient degenerates at marginal-null ray endpoints. We prove this
for the original scalar marks and their original paired probe.

Let \(q\ge3\), and let \(\alpha,\beta\) have compact, disjoint
scalar supports \(A,B\) in a bounded interval \(J\), separated by
\(\delta>0\), with finitely many ordered monochromatic classes.
Assume \(F_\alpha\ge F_\beta\); tight Hall cuts are allowed.
Let \(I\) be an open interval such that
\begin{equation}
\label{degenerate-scalar-interior-mass}
                     \alpha(I)=\beta(I)=1.
\end{equation}
The supports may meet the finite endpoints of \(I\). Let
\(D:I\to\operatorname{Sym}_q\) be continuous, with
\(D(r)\ge h\mathrm{Id}\), \(h>0\), on the relevant scalar span;
no upper bound near the endpoints is assumed. On almost-sure legal
pairs \(a<b\), put
\begin{equation}
\label{degenerate-scalar-kernel}
 S(a,b)=\int_a^bD(r)\,dr,\qquad A(a,b)=S(a,b)^{-1},\qquad
 k((x,a),(y,b))=\tfrac12(y-x)^TA(a,b)(y-x).
\end{equation}
These matrices are finite and positive definite, and
\(A(a,b)\le(h\delta)^{-1}\mathrm{Id}\). Values at marginal-null
endpoint marks are immaterial. In the ray geometry, \(D=H^{-1}\);
the upper bound on \(H\) gives the required lower bound on \(D\),
without a lower bound on \(H\).

Let \(\mathcal C_2\) consist of stationary perfect forward integer
matching laws of the independent unit-intensity product-marked Poisson
inputs, with \(\E^{\rm S}|y-x|^2<\infty\); stationary auxiliary
randomness is allowed. The finite-class construction makes this class
nonempty. For \(g\in C(A\times B)\), put \(G=\|g\|_\infty\) and
\begin{equation}
\label{degenerate-scalar-pressure}
 I_2(t)=\inf_{P\in\mathcal C_2}\E_P^{\rm S}[k+tg].
\end{equation}
This is a finite concave \(G\)-Lipschitz function.

\begin{theorem}[Scalar pressure at degenerate endpoints]
\label{degenerate-scalar-c1}
Under the preceding assumptions, \(I_2\in C^1(\R)\) for every
continuous original paired-mark probe \(g\).
\end{theorem}

\begin{proof}
\emph{Compliance regularization.}
Choose continuous cutoffs \(0\le\psi_m\le1\), increasing in \(m\),
vanishing near every finite endpoint of \(I\), and eventually equal
to one on every fixed compact subinterval of \(I\). On the bounded
span \(J\), arrange that their supports stay inside \(I\), and set
\begin{equation}
\label{degenerate-scalar-regularization}
 D_m(r)=
 \begin{cases}
 h\mathrm{Id}+\psi_m(r)(D(r)-h\mathrm{Id}),&r\in I,\\
 h\mathrm{Id},&r\notin I,
 \end{cases}
 \qquad T_m(r)=\int_{r_0}^rD_m(s)\,ds.
\end{equation}
The endpoint definitions are continuous, and
\(h\mathrm{Id}\le T_m'\le H_m\mathrm{Id}\) for some finite \(H_m\).
Define \(S_m(a,b)=T_m(b)-T_m(a)\), \(A_m=S_m^{-1}\), and \(k_m\)
as in \eqref{degenerate-scalar-kernel}. On legal interior pairs,
\begin{equation}
\label{degenerate-scalar-monotonicity}
 S_m\uparrow S,\qquad A_m\downarrow A,\qquad
       0<A\le A_m\le(h\delta)^{-1}\mathrm{Id}.
\end{equation}
For every compact interval \(K\Subset I\), all sufficiently large
\(m\) satisfy \(k_m=k\) on every pair with both marks in \(K\):
the cutoff is one on the entire interval between the marks.
The scalar laws, legality relation and probe are unchanged.

Put \(I_m(t)=\inf_{P\in\mathcal C_2}\E_P^{\rm S}[k_m+tg]\).
Uniform ellipticity at fixed \(m\) makes this the stationary class
in Theorem~\ref{forward-full-stationary-identity}. Moreover,
\begin{equation}
\label{degenerate-scalar-value-limit}
                  I_m(t)\downarrow I_2(t)
                  \quad\text{locally uniformly in }t.
\end{equation}
Indeed \(I_m\ge I_2\); for any fixed \(\varepsilon\)-minimizer
\(P\in\mathcal C_2\), the bound
\(k_m\le|y-x|^2/(2h\delta)\) permits dominated convergence.
Then let \(\varepsilon\downarrow0\). The common Lipschitz constant
gives local uniformity. We do not infer differentiability from this
value limit alone.

\emph{A five-value diagonal and diffuse plans.}
Extend \(k_m\) to forbidden pairs by \(|y-x|^2/2\), denoting this
extension by \(k_{m,\mathrm e}\). The forbidden indicator
\(\chi=\mathbf1_{\{a>b\}}\) is continuous on the separated mark
product. For
\[
                  c_{m,M,t}=k_{m,\mathrm e}+M\chi+tg,
\]
let \(f_{n,m,M}(t)\) be the expected normalized optimum of balanced
iid assignment in \(D_n=(0,n^{1/q})^q\), and let
\(p_{m,M}(t)=\lim_n f_{n,m,M}(t)\). Existence follows from
Proposition~\ref{prop:highd-finite-pressure}. Theorem~
\ref{forward-full-stationary-identity} and Lemma~
\ref{forward-full-penalty-bound} give
\begin{equation}
\label{degenerate-scalar-penalty-limit}
 p_{m,M}(t)\uparrow I_m(t),\qquad
 -|t|G\le p_{m,M}(t)\le C+|t|G.
\end{equation}
Here \(C\) is independent of \(m,M\): the classwise spatial
competitor uses only the common quadratic upper bound in
\eqref{degenerate-scalar-monotonicity} and the fixed number of classes.
The volume limit is taken at each fixed \(m,M\).

Fix \(t_0\) and \(h_j\downarrow0\). Successively choose
\(m_j\to\infty\), \(M_j\to\infty\), and \(n_j\to\infty\), so
that, for \(\ell=-2,-1,0,1,2\),
\begin{equation}
\label{degenerate-scalar-five-values}
 \begin{split}
 |I_{m_j}(t_0+\ell h_j)-I_2(t_0+\ell h_j)|&\le h_j^2,\\
 |p_{m_j,M_j}(t_0+\ell h_j)-I_{m_j}(t_0+\ell h_j)|&\le h_j^2,\\
 |f_{n_j,m_j,M_j}(t_0+\ell h_j)-p_{m_j,M_j}(t_0+\ell h_j)|
    &\le h_j^2.
 \end{split}
\end{equation}
These are successive value limits, with no exchange of limits.
On a common iid target cloud \((Y_l,b_l)_{l=1}^n\), with
independent \(P\sim\operatorname{Dirichlet}(1,\ldots,1)\) and
\(n=n_j\), define
\[
 \lambda_n=\mathbf1_{D_n}(x)\,dx\,\alpha(da),\qquad
 \nu_n=\sum_{l=1}^n nP_l\delta_{(Y_l,b_l)}.
\]
Choose optimal diffuse plans \(\Gamma_j^+,\Gamma_j^-,\Gamma_j^0\)
for these same marginals, at parameters \(t_0+h_j,t_0-h_j,t_0\).
Proposition~\ref{prop:iid-cavity-mean} identifies their expected
normalized optimal values with \(f_{n,m_j,M_j}\). Its atomless
pair-difference hypothesis holds at every index: for a fixed source
mark and distinct spatial targets, the row-cost difference is a
nonconstant polynomial in \(x\). Distinct inverse matrices give a
nonzero quadratic term; identical matrices give a nonzero linear term.
Mark additions affect only the constant. Polynomial level sets are
Lebesgue null, and Fubini permits arbitrary, including atomic, \(\alpha\).

Write \(C_j=k_{m_j,\mathrm e}+M_j\chi+t_0g\) and
\(\overline\Gamma_j=(\Gamma_j^++\Gamma_j^-)/2\). Finite
optimality gives
\begin{equation}
\label{degenerate-scalar-diffuse-excess}
 \begin{split}
 0\le\Delta_j
   &:=\int C_j\,d\overline\Gamma_j-\int C_j\,d\Gamma_j^0
       \le h_jGn,\\
 \int(k_{m_j,\mathrm e}+M_j\chi)\,d\Gamma_j^\pm
   &\le\int C_j\,d\Gamma_j^0+(|t_0|+2h_j)Gn.
 \end{split}
\end{equation}
Consequently the expected kinetic cost per unit volume is uniformly
bounded and the expected forbidden mass per unit volume is
\(O(M_j^{-1})\).

\emph{Both endpoint tails and the rooted limit.}
Root at an independent uniform \(U_j\in D_n\), and translate the
environment and plans by \(-U_j\). For a compact interval
\(K\Subset I\), put \(\varepsilon_K=\alpha(K^c)+\beta(K^c)\).
Boundedness of \(D\) on \(K\) gives \(c_K>0\) such that
\(k_m\ge k\ge c_K|y-x|^2\) on legal pairs with both marks in
\(K\), for every \(m\). Decrease \(c_K\) to at most \(1/2\)
to cover forbidden pairs as well. For any bounded spatial anchor
window \(B_0\), root integration contributes at most \(|B_0|/n\)
for either endpoint. Bad-source mass is exactly \(n\alpha(K^c)\);
expected bad-target mass is \(n\beta(K^c)\), since \(\E[nP_l]=1\).
Thus each sign and either endpoint anchor satisfy
\begin{equation}
\label{degenerate-scalar-both-tails}
 \E[\text{anchored edge mass with }|y-x|>R]
 \le |B_0|\left(\varepsilon_K+\frac{C}{c_KR^2}\right).
\end{equation}
This uses the prescribed marginal masses, without independence of
matched marks. Exhausting \(I\) by \(K\), then choosing \(R\),
gives uniform partner tightness without a uniform second moment.
The source-anchored mass is deterministically at most \(|B_0|\), and
the Dirichlet second moments and uniform target positions give
\begin{equation}
\label{degenerate-scalar-target-ui}
             \E[\nu_n(U_j+B_0)^2]\le2|B_0|+|B_0|^2.
\end{equation}

These local mass bounds and both tails give a joint stationary
subsequential limit \((\Gamma^+,\Gamma^-,\nu)\), retaining both
common marginals
\begin{equation}
\label{degenerate-scalar-rooted-marginals}
 \lambda=dx\,\alpha(da),\qquad
 \nu=\sum_l E_l\delta_{(Y_l,b_l)}.
\end{equation}
Here the targets form a unit-intensity product-marked Poisson process
and \(E_l\) are independent mean-one exponentials, independent of
the targets. Indeed \(nP_l=nE_l/\sum_iE_i\); the binomial local
limit and the law of large numbers give the target marginal limit.
Equations~\eqref{degenerate-scalar-both-tails}--
\eqref{degenerate-scalar-target-ui} prevent loss of either projection.
Compact nonnegative tests against \(\chi\) have expectation
\(O(M_j^{-1})\), so both plans are forward. Their full marginals
and \eqref{degenerate-scalar-interior-mass} place both endpoints in
\(I\) almost surely. No convergence or uniform integrability of
kinetic means is used.

\emph{Legal two-cycles and the diffuse map.}
A strict legal two-cycle improvement of \(k+t_0g\) in the limiting
average has a certificate consisting of bounded open endpoint
rectangles, with mark closures in one \(K\Subset I\), disjoint
same-color rectangles, all four pair products legal, and saving at
least \(\epsilon>0\). A countable family of certificates suffices.
For all sufficiently large \(j\), exact interior agreement gives
the same old and crossed kinetic costs for \(k_{m_j,\mathrm e}\);
all four penalties are zero.

Let \(W\) be the minimum of one and the masses of the two old
rectangles. Restrict and scale the two pieces to common mass \(W\),
then cross their marginals by product couplings. This preserves both
marginals and saves at least \(\epsilon W\). As in the proof of
Theorem~\ref{forward-full-pressure-c1}, repeat the correction on a
translated coarse grid of side \(L\) containing both source
projections in its anchor cell. Different cells remove disjoint
source pieces, and each correction preserves every target mass.
Grid-phase averaging and \eqref{degenerate-scalar-diffuse-excess}
therefore give, with \(\tau_{U_j}\) denoting translation by \(-U_j\),
\begin{equation}
\label{degenerate-scalar-cycle-witness}
 \E W(\tau_{U_j}\overline\Gamma_j)
 \le\frac{L^q\E\Delta_j}{\epsilon n}
 \le\frac{L^qGh_j}{\epsilon}\longrightarrow0.
\end{equation}
The bounded witness is lower semicontinuous for vague convergence.
Portmanteau and the countable certificates show that
\(\overline\Gamma=(\Gamma^++\Gamma^-)/2\) admits no strict legal
two-cycle improvement for the original interior cost.

Apply the legal countable-target rigidity argument from
Theorem~\ref{forward-full-pressure-c1}. Its row-difference hypothesis
holds pointwise in the source mark: the finite positive definite
matrices in \eqref{degenerate-scalar-kernel} give a nonconstant
polynomial in \(x\), with Lebesgue-null level sets. Thus
\(\overline\Gamma\) is a map and
\begin{equation}
\label{degenerate-scalar-diffuse-agreement}
                         \Gamma^+=\Gamma^-
                         \quad\text{almost surely}.
\end{equation}
Only pointwise positive definiteness is needed here, with no uniform
coefficient bound as marks approach the endpoints.

\emph{Bounded probes and derivatives.}
Set \(q_j^\pm=n_j^{-1}\E\int g\,d\Gamma_j^\pm\). The finite
Dirichlet identity makes these supergradients of
\(f=f_{n_j,m_j,M_j}\) at \(t_0\pm h_j\). For example,
\[
 \frac{f(t_0+2h_j)-f(t_0+h_j)}{h_j}
 \le q_j^+\le
 \frac{f(t_0+h_j)-f(t_0)}{h_j}.
\]
Each endpoint value differs from \(I_2\) by at most \(3h_j^2\)
in \eqref{degenerate-scalar-five-values}. Concavity and the
corresponding negative-side inequalities yield
\begin{equation}
\label{degenerate-scalar-one-sided-derivatives}
 q_j^+\longrightarrow I_2'(t_0+),\qquad
 q_j^-\longrightarrow I_2'(t_0-).
\end{equation}

Choose nonnegative \(\eta\in C_c(\R^q)\) with integral one.
The expected rooted integral of \(\eta(x)g(a,b)\) differs from
\(q_j^\pm\) by \(O(Gn_j^{-1/q})\): only a fixed-width source
boundary layer contributes, and the source marginal is deterministic.
A continuous compact target cutoff gives a continuous compact plan
test, bounded in absolute value by a fixed source mass times
\(G\|\eta\|_\infty\). Its expectation therefore passes to the
rooted limit. Removing that cutoff costs at most
\begin{equation}
\label{degenerate-scalar-probe-tail}
                 C_\eta G\left(\varepsilon_K+
                                      \frac{C}{c_KR^2}\right)
\end{equation}
by \eqref{degenerate-scalar-both-tails}. First take the local limit,
then \(R\to\infty\), then exhaust \(I\) by \(K\). This passes
the original bounded probe expectations, including marks outside
each temporary collar. Equation~\eqref{degenerate-scalar-diffuse-agreement}
equates the two limiting expectations in
\eqref{degenerate-scalar-one-sided-derivatives}. Hence \(I_2\) is
differentiable at every \(t_0\). A finite differentiable concave
function on \(\R\) has continuous derivative.
\end{proof}

Any attained minimizer in \(\mathcal C_2\) has paired expectation
\(I_2'(t)\). More generally, if \(P_l\in\mathcal C_2\) have
\(\E_{P_l}^{\rm S}[k+tg]\downarrow I_2(t)\), then
\(\E_{P_l}^{\rm S}g\to I_2'(t)\). Compare each same competitor
at \(t-h,t+h\), let its value excess vanish, and then let
\(h\downarrow0\). No uniform unweighted moment bound on the
near-minimizers is required.

\begin{corollary}[Integrated second-moment pressure]
\label{degenerate-scalar-integrated-c1}
For the atlas and continuous physical probe \(f\) above,
\(\mathcal P_2\) from \eqref{general-gluing-integrated-pressures}
belongs to \(C^1(\R)\), and
\begin{equation}
\label{degenerate-scalar-integrated-derivative}
 \mathcal P_2'(t)=\sum_j\int_{Q_j}\theta_j(z)
   \partial_v I_{j,2}\bigl(z,t\theta_j(z)^{2/q}\bigr)\,dz.
\end{equation}
\end{corollary}
\begin{proof}
At almost every active root use
\(g_z(a,b)=f(X_j(a,z),X_j(b,z))\) and \(D=H^{-1}\).
The conditional endpoints are interior almost surely, and the upper
bound on \(H\) supplies \(D\ge h\mathrm{Id}\); \(h\) may depend
on the root. Theorem~\ref{degenerate-scalar-c1} applies to each
\(I_{j,2}(z,\cdot)\). Since \(s+2/q=1\), its derivative
contribution is bounded by \(\theta_j(z)\|f\|_\infty\), whose
total integral is \(\|f\|_\infty\). The existing atlas bounds
give finiteness and measurability. Dominated difference quotients prove
\eqref{degenerate-scalar-integrated-derivative}, and dominated
continuity of the pointwise derivatives proves continuity.
\end{proof}

Thus pressure differentiability needs neither a uniform endpoint
collar nor an inverse-coefficient moment. Identification with
\(\mathcal P_{\mathrm w}\) and the empirical pressure still uses
the separate equality \(I_{j,\mathrm w}=I_{j,2}\) at almost every
active root for full perturbed pressure convergence.
Once Corollary~\ref{weighted-capped-integrated-c1} is available,
equality only at the central parameter is sufficient for unperturbed
selection: the lower comparisons at \(h\) and \(-h\), together with
convergence of \(V_n(0)\), bound every optimum's paired readout between
the corresponding secants of \(\mathcal P_{\mathrm w}\).
Its derivative at zero gives the common limit.

% Source: degenerate-finite-pressure-existence.md
% Source SHA256: f59e735c91652744406e004742e0765f4b83a652296591f11723ba330007d8de
% Independent source review: degenerate-finite-pressure-existence-review.md, PASS.
\subsubsection{Finite forward pressure without a lower matrix bound}
\label{degenerate-finite-section}

We next identify a homogeneous finite-volume value directly. This does
not identify every weighted stationary competitor with a finite-moment
one. Let \(q>4\), and let two independent unit-intensity marked Poisson
processes have compact scalar laws \(\alpha,\beta\), with finitely many
ordered monochromatic classes, positive separation, and
\(F_\alpha\ge F_\beta\). Write \(J\) for the total number of classes.
An edge is legal when \(a<b\); every allowed class rectangle is complete
bipartite. Fix Borel matrices and a bounded Borel probe satisfying
\[
 0\preceq A(a,b)\preceq C\operatorname{Id},\qquad
 G=\|g\|_\infty,\qquad
 c_t((x,a),(y,b))=\tfrac12(y-x)^TA(a,b)(y-x)+tg(a,b).
\]
In \(Q_L=[0,L)^q\), let \(H_L(t)\) minimize total cost among legal
\emph{maximum-cardinality} partial matchings of the restricted inputs,
and set \(f_L(t)=L^{-q}\E H_L(t)\).

\begin{theorem}[Degenerate finite pressure and monotone coefficients]
\label{degenerate-finite-pressure}
The limit \(p_A(t)=\lim_{L\to\infty}f_L(t)\) exists for all real
\(L\to\infty\), is finite, concave and \(G\)-Lipschitz, and convergence
is locally uniform in \(t\). For \(l\ge1\) and \(|t|\le T\),
\begin{equation}
\label{degenerate-finite-size}
 p_A(t)\le f_l(t)+C_{J,q,C}
       \bigl(l^{2-q/2}+TG l^{-q/2}\bigr).
\end{equation}
The constant is independent of any lower eigenvalue bound. If
\(0\preceq A\preceq A_m\preceq C\operatorname{Id}\) and
\(A_m\downarrow A\) for \(\alpha\otimes\beta\)-almost every legal
pair, then \(p_{A_m}\downarrow p_A\), locally uniformly.
\end{theorem}

\begin{proof}
Apply Lemma~\ref{measurable-forward-augmentation} to the finite
class graph. In a box of diameter \(R\), one augmentation increases cost
by at most
\begin{equation}
\label{degenerate-finite-augmentation}
 C_J(CR^2+TG).
\end{equation}
The lemma counts all removed and inserted probe terms and uses only
nonnegativity of the removed kinetic costs. Adding or deleting \(K\)
vertices therefore changes the optimal maximum-cardinality value in
absolute value by at most \(C_J(CR^2+TG)K\). For deletion, restrict the
larger optimum and augment: at most \(K\) edges were removed and at
most \(K\) augmentations are needed. Compare in the other direction
by extending the smaller matching and augmenting. This also proves
the assertion for signed probes.

For a box \(B\) of volume \(V\), let \(U_B\) count the unmatched
vertices in a maximum-cardinality matching. Population Hall feasibility
and the original-count estimate of the same lemma give
\begin{equation}
\label{degenerate-finite-deficiency}
 U_B\le 2\sum_h|N_h(B)-Vp_h|,\qquad
 \E U_B\le C_J\sqrt V.
\end{equation}
The expectation bound follows from the original Poisson class-count
variances for every \(V>0\); it requires no law for the unmatched points.

Fix \(l\ge1\). For \(L\ge l\), put
\(m=\lfloor L/l\rfloor\) and \(a=L/m\). Tile \(Q_L\) into
\(m^q\) cubes of side \(a\), using a maximum-cardinality optimum in
each. Embed their index grid in the next dyadic grid, with extra cells
empty. At level \(j\), start from the union of the child matchings and
augment each parent to maximum cardinality. Its diameter is at most \(\sqrt q\,2^j a\);
there are at most \(C_qm^q2^{-jq}\) nonempty parents up to the last
level. Each has at most \(2^q\) children of volume at most
\((2^ja)^q\). Equations~\eqref{degenerate-finite-augmentation} and
\eqref{degenerate-finite-deficiency} bound the expected extra cost
per original volume by
\[
 C_{J,q}\bigl(C(2^ja)^{2-q/2}+TG(2^ja)^{-q/2}\bigr).
\]
Every child's unmatched count depends on its original input counts,
so the estimate remains valid after earlier adaptive choices. Since
\(q>4\), summing gives
\begin{equation}
\label{degenerate-finite-hierarchy}
 f_L(t)\le f_a(t)+C_{J,q,C}
             \bigl(a^{2-q/2}+TG a^{-q/2}\bigr).
\end{equation}
The final matching is maximum-cardinality in \(Q_L\).
Also \(f_L\ge-TG\); taking \(l=1\) gives a uniform upper bound.

The function \(a\mapsto f_a(t)\) is continuous on \((0,\infty)\),
locally uniformly for \(|t|\le T\). Couple the inputs on a unit spatial
cube with a Poisson intensity coordinate, retain heights below \(a^q\),
and multiply positions by \(a\). The mean number of added or deleted
vertices between \(a,a'\) is \(2|a^q-a'^q|\), controlled by the vertex
bound above. On common points, scaling changes each quadratic edge cost
by at most \(C_qC|a^2-a'^2|\), with unchanged probe. Expected common
counts are locally bounded. Division by \(a^q\) proves continuity;
no mark continuity is needed.

At fixed \(l\), \(a\to l\) as \(L\to\infty\), so
\eqref{degenerate-finite-hierarchy} bounds \(\limsup_L f_L\) by
\(f_l+C(l^{2-q/2}+TG l^{-q/2})\). Taking \(l\to\infty\) along
a sequence attaining the finite lower limit proves existence and
\eqref{degenerate-finite-size}. Expected normalized matched source count
is at most one. Hence all \(f_L\) are concave and \(G\)-Lipschitz,
which also gives local uniform convergence by a finite parameter net.

For fixed \(L\), independent marks avoid the exceptional null set for
all finitely many input pairs almost surely. The finite set of legal
maximum-cardinality matchings is independent of the matrix coefficients.
Thus \(H_L^{(m)}(t)\downarrow H_L(t)\), dominated in absolute value by
\(C_q(CL^2+TG)(N_S+N_T)\), and \(f_L^{(m)}\downarrow f_L\).
Monotonicity gives \(p_A\le\lim_m p_{A_m}\), whereas the uniform
estimate~\eqref{degenerate-finite-size} gives, for every fixed \(l\),
\[
 \lim_m p_{A_m}(t)\le f_l(t)+C_{J,q,C}
                   \bigl(l^{2-q/2}+TG l^{-q/2}\bigr).
\]
Now let \(l\to\infty\). This proves the reverse inequality and,
again by common Lipschitz bounds, local uniformity. The direction of
coefficient approximation is from above.
\end{proof}

\begin{corollary}[Scalar finite pressure and its full optimal face]
\label{degenerate-finite-readout}
For the endpoint-degenerate scalar cumulative coefficients of
Theorem~\ref{degenerate-scalar-c1}, and a continuous original paired
probe \(g\), one has \(p_A=I_2\in C^1(\R)\). For two independent
lists of \(n\) iid marked points in the cube of side \(n^{1/q}\),
the normalized maximum-cardinality values converge in \(L^2\), locally
uniformly in the probe parameter. At parameter zero the paired average
of every optimum converges in probability to \(I_2'(0)\), uniformly
over the whole optimal face, including fractional matchings.
\end{corollary}

\begin{proof}
Use the compliance regularizations of
Theorem~\ref{degenerate-scalar-c1}. At each fixed index they are
uniformly elliptic, and Lemma~\ref{measurable-forward-homogeneous}
identifies \(p_{A_m}=I_m\). That lemma includes weak Hall inequalities
and tight cuts: split into strict blocks and augment to maximum
cardinality with normalized cost \(O(L^{2-q/2})\). Thus it does not
require strict Hall inequalities for the original full law. Theorem
\ref{degenerate-finite-pressure} gives \(p_{A_m}\downarrow p_A\),
while the fixed-competitor argument in
Theorem~\ref{degenerate-scalar-c1} gives \(I_m\downarrow I_2\).
This proves the asserted identity without identifying \(I_{\rm w}\).

For the fixed-size assertion, couple two infinite iid lists to independent
\(\operatorname{Poisson}(n)\) prefix counts. The expected number of
changed vertices is \(O(\sqrt n)\). The insertion/deletion bound gives
an expected normalized value difference
\(O((Cn^{2/q}+TG)n^{-1/2})=o(1)\). Replacing one labeled sample
changes the optimum by at most \(C_J(Cn^{2/q}+TG)\). Efron--Stein
therefore bounds the variance of the normalized value by
\(C_{J,T}n^{4/q-1}=o(1)\). The deterministic normalized Lipschitz
bound \(G\) gives locally uniform \(L^2\) convergence. Secant bounds
at \(\pm h\), first sending \(n\to\infty\) and then \(h\downarrow0\),
give the readout uniformly over every zero-parameter optimum. Fractional
optima satisfy the same bounds; the maximum-cardinality face of the
bipartite matching polytope is integral. Its unmatched fraction vanishes
by~\eqref{degenerate-finite-deficiency}.
\end{proof}

This is a homogeneous local theorem on the original independent inputs.
Restricting a full geometric optimizer to a microscopic box need not
give a maximum-cardinality matching there. Controlling the cost of
completing those restrictions remains a separate lower-localization
problem; the corollary alone does not close the general-geometric
pressure bracket.

% Source: weighted-pressure-capped-cavity.md
% Source SHA256: ee4f5f16845e3651fd83bdba3b8e6dae63493dddc3e5eadc0870ffef7eab3e6b
% Independent source review: weighted-pressure-capped-cavity-review.md, PASS.
\subsubsection{Weighted pressure from bounded costs and its differentiability}
\label{weighted-capped-section}

We now identify the weighted stationary pressure through bounded finite
costs and prove its differentiability. Retain the scalar setting of
Theorem~\ref{degenerate-scalar-c1}, with \(q>4\): the original
independent unit-intensity marked Poisson inputs have compact separated
scalar laws \(\alpha,\beta\), finitely many ordered monochromatic
classes, and \(F_\alpha\ge F_\beta\). Both laws give full mass to
the open interval \(I\). With continuous \(D(r)\ge h\mathrm{Id}\)
on its relevant interior span, write
\[
 \begin{gathered}
 S(a,b)=\int_a^bD(r)\,dr,\qquad A(a,b)=S(a,b)^{-1},\\
 Q((x,a),(y,b))=\tfrac12(y-x)^TA(a,b)(y-x),\quad a<b.
 \end{gathered}
\]
The matrix \(A\) is positive definite and continuous on interior
legal pairs, with a common upper bound; it may degenerate at
marginal-null endpoints. For
\(g\in C(\operatorname{supp}\alpha\times\operatorname{supp}\beta)\),
put \(G=\|g\|_\infty\). Let \(\mathcal C_{\rm w}\) be the stationary
perfect forward integer matching laws on the prescribed inputs with
\(\E^{\rm S}Q<\infty\), allowing stationary auxiliary randomness, and set
\begin{equation}
\label{weighted-capped-pressure-definition}
 I_{\rm w}(t)=\inf_{P\in\mathcal C_{\rm w}}
                           \E_P^{\rm S}[Q+tg].
\end{equation}
The finite-class construction gives a competitor in \(\mathcal C_2\),
so \(I_{\rm w}\) is finite, concave and \(G\)-Lipschitz.

Extend \(Q\) to forbidden pairs by \(|y-x|^2/2\), denoting the
extension by \(Q_{\rm e}\), and put \(\chi=\mathbf1_{\{a>b\}}\).
For integers \(K,M\ge1\), define
\begin{equation}
\label{weighted-capped-cost}
 \psi_K(u)=\frac{Ku}{K+u},\qquad
 c_{K,M,t}=\psi_K(Q_{\rm e})+M\chi+tg.
\end{equation}
All pairs are permitted in this finite assignment problem. The indicator
\(\chi\) is continuous on the separated mark product.

\begin{theorem}[Bounded-cost representation and weighted differentiability]
\label{weighted-capped-c1}
The fixed-\(K,M\) independent-input thermodynamic pressure
\(p_{K,M}(t)\) exists, and
\begin{equation}
\label{weighted-capped-value-limit}
                   p_{K,M}(t)\uparrow I_{\rm w}(t)
                   \quad(K,M\to\infty),
\end{equation}
locally uniformly in \(t\). The weighted infimum is attained by a
stationary perfect integer matching, and \(I_{\rm w}\in C^1(\R)\).
\end{theorem}

\begin{proof}
\emph{Bounded finite pressures.}
Fix \(K,M,T\), and put \(B_T=K+M+TG\). On \(|t|\le T\),
the costs lie in \([-TG,B_T]\). In \(Q_L=[0,L)^q\), let
\(H_L^{K,M}(t)\) be the optimum over ordinary maximum-cardinality
matchings of the two restricted Poisson inputs, and put
\(f_L^{K,M}(t)=L^{-q}\E H_L^{K,M}(t)\).
Adding or deleting \(k\) vertices changes this optimum in absolute
value by at most \(C B_Tk\). Indeed, restrict an optimum, then
augment to the new maximum cardinality; the complete bipartite graph
requires only a bounded number of changed edges per changed vertex.
Compare in the reverse direction for the other inequality.

In a box of volume \(V\), the unmatched count is
\(|N_{\rm S}-N_{\rm T}|\), of expectation at most \(C\sqrt V\).
For fixed \(l>0\), partition \(Q_L\) into \(m^q\) cubes of side
\(a=L/m\), where \(m=\lfloor L/l\rfloor\). Embed this index
grid in the next dyadic grid, declaring extra cells empty. Starting
with leaf optima, pair opposite unmatched vertices at each merge.
At level \(j\), at most \(C_qm^q2^{-jq}\) parents occur, each
with child unmatched expectation at most \(C_q(2^ja)^{q/2}\).
Each new edge costs at most \(B_T\). Summing the geometric series gives
\begin{equation}
\label{weighted-capped-finite-size}
 f_L^{K,M}(t)\le f_a^{K,M}(t)+C_qB_Ta^{-q/2},\qquad
 \limsup_{L\to\infty}f_L^{K,M}(t)
                  \le f_l^{K,M}(t)+C_qB_Tl^{-q/2}.
\end{equation}
For the last passage, side-length continuity follows by coupling
Poisson intensity prefixes on the unit cube and dilating the common
positions. Changed vertices are controlled by the vertex bound; for
each common finite input the costs vary continuously under spatial
dilation, with bounded per-edge domination. Null endpoint marks cause
no difficulty. The argument is locally uniform in \(t\).

Uniform upper and lower bounds follow from \(|c_{K,M,t}|\le B_T\)
and the expected matched source intensity being at most one. Taking
\(l\to\infty\) along a sequence attaining the lower limit in
\eqref{weighted-capped-finite-size} proves existence of \(p_{K,M}\).
Concavity and the common Lipschitz bound \(G\) give locally uniform
convergence. For two balanced iid lists of size \(n\) in
\(D_n=(0,n^{1/q})^q\), denote the expected normalized full-assignment
value by \(f_{n,K,M}(t)\). Coupling iid prefixes to independent
\(\operatorname{Poisson}(n)\) counts changes \(O(\sqrt n)\)
vertices in expectation, so
\begin{equation}
\label{weighted-capped-iid-limit}
             f_{n,K,M}(t)\longrightarrow p_{K,M}(t).
\end{equation}
Replacing one sample changes the normalized optimum by \(O(B_T/n)\);
Efron--Stein bounds its variance by \(O(B_T^2/n)\). These assertions
are for fixed \(K,M\).

\emph{Comparison with every weighted stationary competitor.}
Fix \(P\in\mathcal C_{\rm w}\), and retain its edges with both
endpoints in \(Q_L\). Every partner has finite displacement almost
surely, so Campbell's formula and bounded convergence show that the
expected fraction of crossing endpoints tends to zero. Complete the
retained matching to ordinary maximum cardinality, at additional cost
at most \(B_T\) times the missing-vertex count. With
\[
 w_L(v)=\prod_{i=1}^q(1-|v_i|/L)_+,
\]
the expected retained cost per volume is exactly
\(\E_P^{\rm S}[(\psi_K(Q)+tg)w_L(y-x)]\). Its integrand is
bounded, and the retained edges are forward. Therefore
\begin{equation}
\label{weighted-capped-stationary-upper}
 p_{K,M}(t)\le\E_P^{\rm S}[\psi_K(Q)+tg]
             \le\E_P^{\rm S}[Q+tg],\qquad
 p_{K,M}(t)\le I_{\rm w}(t).
\end{equation}
Both \(K\) and \(M\) increase the finite costs. Thus their
increasing supremum \(p_\infty(t)\) is finite and at most \(I_{\rm w}(t)\).

\emph{Two-anchor compactness and integer attainment.}
Fix \(t\). Choose \(K_j,M_j\to\infty\), followed by
\(n_j\to\infty\), so that the expected iid optimal values differ
from \(p_{K_j,M_j}(t)\) by at most \(1/j\). Take exact integer
full assignments and root at an independent uniform point of \(D_{n_j}\).
By \eqref{weighted-capped-stationary-upper}, their expected capped
kinetic cost per volume is bounded uniformly, while their forbidden
mass per volume is \(O(M_j^{-1})\).

For a compact scalar interval \(F\Subset I\), put
\(\varepsilon_F=\alpha(F^c)+\beta(F^c)\). On legal pairs with
marks in \(F\), \(Q\ge c_F|y-x|^2\); reduce \(c_F\) to at most
\(1/2\) for the forbidden extension. At fixed \(F,R\), once
\(K_j\ge c_FR^2\),
\[
 \psi_{K_j}(Q_{\rm e})\ge\tfrac12c_FR^2
       \quad\text{if }|y-x|>R\text{ and }a,b\in F.
\]
Root integration contributes at most \(|W|/n_j\) for either endpoint
anchor in a bounded spatial window \(W\). The expected total bad-mark
mass is at most \(n_j\varepsilon_F\), using the original marginals.
Consequently, for either anchor,
\begin{equation}
\label{weighted-capped-both-tails}
 \limsup_{j\to\infty}
 \E[\text{anchored edge mass with }|y-x|>R]
       \le |W|\left(\varepsilon_F+\frac{C}{c_FR^2}\right).
\end{equation}
First choose \(F\) with small excluded mass, then choose \(R\);
the cap threshold concerns only sufficiently large indices. This proves
partner tightness without a uniform unweighted second moment.

The original iid local counts have bounded second moments at both
anchors. Joint rooted compactness therefore retains both full input
projections, whose limit is the prescribed pair of independent marked
unit Poisson processes. Their local simplicity and absence of spatial
collisions preserve the degree-one integer edge relation. The rooted
limit is stationary and perfect; compact tests against \(\chi\)
make it forward. Denote its law by \(P_*\).

On bounded displacement sets with compact interior legal marks,
\(\psi_{K_j}(Q)\to Q\) uniformly. Nonnegative cutoff tests, Fatou
and exhaustion give the kinetic lower bound. For a nonnegative compact
source anchor of integral one, the root-averaging weight is at most
\(1/n_j\), so its rooted capped kinetic expectation is bounded by
the total capped kinetic cost per volume. No kinetic boundary-error
estimate is required. Bounded probe expectations pass separately by
\eqref{weighted-capped-both-tails} and local count integrability;
their source boundary error is \(O(Gn_j^{-1/q})\). Dropping the
nonnegative forbidden penalty yields
\begin{equation}
\label{weighted-capped-attainment}
                    \E_{P_*}^{\rm S}[Q+tg]\le p_\infty(t).
\end{equation}
Together with \eqref{weighted-capped-stationary-upper}, this proves
\(p_\infty=I_{\rm w}\), finite weighted energy of \(P_*\), and
attainment. Any cofinal choice of \(K,M\) has the same limit; common
Lipschitz bounds give local uniformity in \eqref{weighted-capped-value-limit}.

\emph{The exact diffuse value identity.}
For a common iid target cloud \((Y_l,b_l)_{l=1}^n\) and independent
\(P\sim\operatorname{Dirichlet}(1,\ldots,1)\), define
\begin{equation}
\label{weighted-capped-dirichlet-marginals}
 \lambda_n=\mathbf1_{D_n}(x)\,dx\,\alpha(da),\qquad
 \nu_n=\sum_{l=1}^n nP_l\delta_{(Y_l,b_l)}.
\end{equation}
The value identity \eqref{eq:iid-cavity-pressure} applies to
\eqref{weighted-capped-cost}. To verify its atomless-difference
hypothesis, fix an interior source mark and two interior target marks.
The functions \(\psi_K(Q_l(x))\) are real analytic, since their
quadratics are positive definite and \(K+Q_l(x)>0\). If their
difference were constant on a set of positive Lebesgue measure,
analyticity would make it constant everywhere. Both rows tend to
\(K\) as \(|x|\to\infty\), so that constant is zero. Strict
monotonicity of \(\psi_K\) then gives \(Q_l=Q_r\). Equality of
these quadratic polynomials forces equality of their matrices and
centres, contrary to distinct target positions. Penalty and probe
terms change only the constant. Fubini handles arbitrary scalar laws,
including atoms; exceptional endpoint marks are null.

The bounded-row value identity also follows without atomless differences:
add independent continuous source-row noises of amplitude at most
\(\varepsilon\), apply Proposition~\ref{prop:iid-cavity-mean}, then
remove the noises. Each normalized optimum changes by at most
\(\varepsilon\). Projection and independent lifting identify the
unperturbed diffuse feasible sets on the augmented and original source
spaces. Thus measurable bounded rows suffice for the value identity.
For the chosen smooth cap the preceding analytic check also verifies
the original atomless-row theorem directly. Semidiscrete optima exist:
write their allocations as finitely many \(L^\infty\) source densities,
use weak-star compactness, and pair each bounded cost with its allocation.
The expected normalized diffuse values are exactly \(f_{n,K,M}(t)\).

\emph{The five-value diagonal.}
Fix \(t_0\) and \(h_j\downarrow0\). First choose
\(K_j,M_j\to\infty\), then \(n_j\to\infty\), such that, for
\(\ell=-2,-1,0,1,2\),
\begin{equation}
\label{weighted-capped-five-values}
 \begin{split}
 |p_{K_j,M_j}(t_0+\ell h_j)-I_{\rm w}(t_0+\ell h_j)|&\le h_j^2,\\
 |f_{n_j,K_j,M_j}(t_0+\ell h_j)-p_{K_j,M_j}(t_0+\ell h_j)|
                                                    &\le h_j^2.
 \end{split}
\end{equation}
On the common marginals \eqref{weighted-capped-dirichlet-marginals},
choose diffuse optima \(\Gamma_j^+,\Gamma_j^-,\Gamma_j^0\) at
\(t_0+h_j,t_0-h_j,t_0\). Put
\(C_j=\psi_{K_j}(Q_{\rm e})+M_j\chi+t_0g\) and
\(\overline\Gamma_j=(\Gamma_j^++\Gamma_j^-)/2\). Finite
optimality gives
\begin{equation}
\label{weighted-capped-diffuse-excess}
 \begin{split}
 0\le\Delta_j
  &:=\int C_j\,d\overline\Gamma_j-\int C_j\,d\Gamma_j^0
                                                    \le h_jGn_j,\\
 \int(\psi_{K_j}(Q_{\rm e})+M_j\chi)\,d\Gamma_j^\pm
  &\le\int C_j\,d\Gamma_j^0+(|t_0|+2h_j)Gn_j.
 \end{split}
\end{equation}
The exact value identity bounds expected capped kinetic cost per volume
uniformly and forbidden mass by \(O(M_j^{-1})\). The same proof of
\eqref{weighted-capped-both-tails} applies. Source-window mass is
deterministically at most \(|W|\), and the Dirichlet moments give
\[
                  \E[\nu_n(U_j+W)^2]\le2|W|+|W|^2.
\]
Thus a joint stationary rooted limit \((\Gamma^+,\Gamma^-,\nu)\)
retains both common marginals
\begin{equation}
\label{weighted-capped-diffuse-marginals}
 \lambda=dx\,\alpha(da),\qquad
 \nu=\sum_l E_l\delta_{(Y_l,b_l)},
\end{equation}
where the targets are a unit-intensity product-marked Poisson process
and their independent weights \(E_l\) are mean-one exponentials.
This follows from \(nP_l=nE_l/\sum_iE_i\), the local Poisson limit
and the law of large numbers. Both limiting plans are forward.

\emph{Legal cycles and the limiting map.}
A strict legal two-cycle improvement for \(Q+t_0g\) in
\(\overline\Gamma=(\Gamma^++\Gamma^-)/2\) has a certificate of
bounded open endpoint rectangles, with compact interior mark closures,
disjoint same-color rectangles, all four pair products legal, and
saving at least \(\epsilon>0\). On that fixed certificate the
capped costs converge uniformly to \(Q\). For sufficiently large
\(j\), the same exchange for \(C_j\) saves at least
\(\epsilon/2\); all four penalties vanish.

Let \(W\) be the minimum of one and the masses of the two old
rectangles. Scale their restrictions to common mass \(W\) and cross
their marginals. Repeat on a translated coarse grid of side \(L\)
containing both source projections in its anchor cell. Different cells
remove disjoint source pieces; each exchange preserves every target
mass. The grid-phase argument of Theorem~\ref{degenerate-scalar-c1}
and \eqref{weighted-capped-diffuse-excess} imply
\begin{equation}
\label{weighted-capped-cycle-witness}
 \E W(\tau_{U_j}\overline\Gamma_j)
       \le\frac{C L^q\E\Delta_j}{\epsilon n_j}
       \le\frac{C L^qGh_j}{\epsilon}\longrightarrow0,
\end{equation}
where \(\tau_{U_j}\) translates by \(-U_j\). The bounded open-set
witness is lower semicontinuous. Portmanteau and a countable certificate
cover exclude every strict legal improvement of the limiting average.

The same countable-target rigidity argument applies to \(Q\).
At each fixed source mark, differing positive definite matrices give
a nonzero quadratic row difference; identical matrices and distinct
spatial targets give a nonzero linear term. Hence its level sets are
Lebesgue null, and the average is a map. Consequently
\begin{equation}
\label{weighted-capped-diffuse-agreement}
                         \Gamma^+=\Gamma^-
                         \quad\text{almost surely}.
\end{equation}

\emph{Probe passage and derivatives.}
Set \(q_j^\pm=n_j^{-1}\E\int g\,d\Gamma_j^\pm\). The finite
value identity makes them supergradients of \(f=f_{n_j,K_j,M_j}\)
at \(t_0\pm h_j\). In particular,
\[
 \frac{f(t_0+2h_j)-f(t_0+h_j)}{h_j}
 \le q_j^+\le\frac{f(t_0+h_j)-f(t_0)}{h_j}.
\]
The total endpoint-value error in \eqref{weighted-capped-five-values}
is at most \(2h_j^2\). Adjacent secants and their negative-side
counterparts give
\begin{equation}
\label{weighted-capped-one-sided-derivatives}
 q_j^+\longrightarrow I_{\rm w}'(t_0+),\qquad
 q_j^-\longrightarrow I_{\rm w}'(t_0-).
\end{equation}

Take nonnegative \(\eta\in C_c(\R^q)\) with integral one.
The expected rooted \(\eta(x)g(a,b)\) integral differs from
\(q_j^\pm\) by \(O(Gn_j^{-1/q})\), using the deterministic
source marginal and its fixed-width boundary layer. A continuous
compact target cutoff makes this a continuous compact plan test with
a deterministic source-mass bound. Its expectation passes to the
joint limit. Removing the cutoff has limiting expected error at most
\[
                 C_\eta G\left(\varepsilon_F+
                                      \frac{C}{c_FR^2}\right)
\]
by \eqref{weighted-capped-both-tails}. First take the local limit,
then \(R\to\infty\), then exhaust \(I\) by \(F\). This passes
the original probe expectations. Equation~
\eqref{weighted-capped-diffuse-agreement} equates the two limits in
\eqref{weighted-capped-one-sided-derivatives}. Hence \(I_{\rm w}\)
is differentiable everywhere. Its finite concavity makes its derivative
continuous, completing the proof.
\end{proof}

Every weighted minimizer has paired expectation \(I_{\rm w}'(t)\).
The same holds in the limit for any weighted near-minimizing sequence:
evaluate each same competitor at \(t-h,t+h\), let its value excess
vanish, and then let \(h\downarrow0\). No additional energy-tail
bound for such sequences is required.

\begin{corollary}[Integrated weighted pressure]
\label{weighted-capped-integrated-c1}
For the geometric atlas and continuous physical probe \(f\), the
pressure \(\mathcal P_{\rm w}\) in
\eqref{general-gluing-integrated-pressures} belongs to \(C^1(\R)\), with
\begin{equation}
\label{weighted-capped-integrated-derivative}
 \mathcal P_{\rm w}'(t)=\sum_j\int_{Q_j}\theta_j(z)
   \partial_v I_{j,\rm w}\bigl(z,t\theta_j(z)^{2/q}\bigr)\,dz.
\end{equation}
\end{corollary}
\begin{proof}
Apply Theorem~\ref{weighted-capped-c1} at almost every active root,
with \(D=H^{-1}\) and
\(g_z(a,b)=f(X_j(a,z),X_j(b,z))\). The countable representation
by bounded finite pressures gives measurable rootwise values. Existing
\(\mathcal P_2\) bounds give finiteness, and the lower bound is
\(-|t|\|f\|_\infty\). Since \(s+2/q=1\), derivative
contributions are bounded by \(\theta_j(z)\|f\|_\infty\), whose
total integral is \(\|f\|_\infty\). Dominated difference quotients
give \eqref{weighted-capped-integrated-derivative}; dominated continuity
of the pointwise derivatives gives continuity.
\end{proof}

The bounded-cost representation identifies the iterated limit
\(\lim_{K,M\to\infty}\lim_{n\to\infty}f_{n,K,M}=I_{\rm w}\).
Corollary~\ref{degenerate-finite-readout} identifies the original
unbounded maximum-cardinality forward pressure with \(I_2\).
Both pressures are now differentiable. Equality of their values, or an
interchange of these cap and volume limits, remains a separate issue
in the general-geometric pressure bracket.

% Source: degenerate-boundary-functional.md
% Source SHA256: 3609bc3ef6588beeeeb6f3dfcb600602a93a5ab56717b8acc04a8094027e602c
% Scope: homogeneous boundary comparisons; no general-geometric selector claim.
\section{Jacobi geometry at focal endpoints}\label{sec:focal-geometry}

\subsection{Determinant moments through focal ends}
\label{subsec:curved-focal-determinant-moments}

The actual Jacobi frame gives an exact cancellation between endpoint
compliance and conditional volume density. The calculation permits
focusing at supported endpoints and applies to the genuine Finsler
length cost, without a curvature or topology restriction.

\begin{lemma}[A second Jacobi solution through singular endpoints]
\label{lem:curved-jacobi-second-solution}
On a finite closed interval \([\ell,u]\), let
\(\mathsf A,\mathsf B,\mathsf C\) be smooth real \(q\times q\) matrices,
with \(\mathsf A=\mathsf A^\top\) and
\(\mathsf C=\mathsf C^\top>0\). Suppose
\[
 J'=\mathsf BJ+\mathsf CP,\qquad
 P'=-\mathsf AJ-\mathsf B^\top P,\qquad J^\top P=P^\top J,
\]
where \(J\) is invertible on \((\ell,u)\), but may be singular at
either endpoint. Fix \(c\in(\ell,u)\) and, on the interior, set
\[
 H=J^\top\mathsf C^{-1}J,\qquad
 T(t)=\int_c^t H(s)^{-1}\,ds,\qquad
 K=JT,\qquad R=PT+J^{-\top}.
\]
Then \((K,R)\) solves the same system and extends smoothly over
\([\ell,u]\), with \(K(c)=0\) and \(R(c)=J(c)^{-\top}\).
For \(a<c<b\), put \(U(a)=-T(a)\) and \(V(b)=T(b)\). Then
\begin{align}
 \det(I+U(a))|\det J(a)|&=|\det(J(a)-K(a))|,\notag\\
 \det(I+V(b))|\det J(b)|&=|\det(J(b)+K(b))|.
 \label{eq:curved-focal-determinant-cancellation}
\end{align}
The right sides are bounded up to their respective endpoints.
\end{lemma}

\begin{proof}
Since \(T'=J^{-1}\mathsf CJ^{-\top}\), differentiation gives
\(K'=\mathsf BK+\mathsf CR\). The Wronskian identity is equivalent to
\(PJ^{-1}=J^{-\top}P^\top\), and
\[
 (J^{-\top})'=-\mathsf B^\top J^{-\top}
                  -J^{-\top}P^\top\mathsf CJ^{-\top}.
\]
Consequently
\begin{align*}
 R'={}&-\mathsf AJT-\mathsf B^\top PT
       +PJ^{-1}\mathsf CJ^{-\top}
       -\mathsf B^\top J^{-\top}
       -J^{-\top}P^\top\mathsf CJ^{-\top}\\
    ={}&-\mathsf AK-\mathsf B^\top R.
\end{align*}
The two terms containing \(P\) and \(\mathsf C\) cancel in their
written order; no commutativity is used. Uniqueness for the smooth
linear system identifies these formulas with its solution from \(c\),
which extends through both endpoints. Finally,
\(J-K=J(I+U)\) on the left and \(J+K=J(I+V)\) on the right.
Since \(U,V>0\), taking determinants gives
\eqref{eq:curved-focal-determinant-cancellation}. Smooth continuation
of \(J,K\) proves the asserted boundedness.
\end{proof}

\begin{proposition}[Fixed-root determinant moments at focal ends]
\label{prop:curved-focal-determinant-moments}
Under Proposition~\ref{prop:general-optimal-face}, fix a regular
positive-mass ray from Proposition~\ref{prop:rough-ray-quadratic}.
Let \([\ell,u]\) be its finite closed conditional support span,
allowing singularities of the normal Jacobi frame at its endpoints,
and let \(H\) be its actual quadratic matrix from that proposition.
Suppose the source and target conditionals \(\alpha,\beta\) lie on
the respective sides of a cut \(c\in(\ell,u)\). Define
\[
 U(a)=\int_a^c H(r)^{-1}\,dr,\qquad
 V(b)=\int_c^b H(r)^{-1}\,dr.
\]
Then
\[
 \int\det(I+U)\,d\alpha<\infty,\qquad
 \int\det(I+V)\,d\beta<\infty.
\]
In particular, \(\int\|U\|\,d\alpha+\int\|V\|\,d\beta<\infty\).
More generally, without the common-cut ordering, for every interior
\(c\) and \(T(t)=\int_c^t H^{-1}\),
\[
 \int\det(I+|T(t)|)\,d(\alpha+\beta)(t)<\infty.
\]
All conclusions are at the fixed regular root; no uniform bound over
ray labels is asserted.
\end{proposition}

\begin{proof}
We first identify the Hamiltonian system for the actual length cost.
The complete smooth geodesic flow extends the baseline and its first
variation through the finite supported endpoints. The calibrated
segment is embedded because its potential coordinate increases
strictly. Choose smooth tubular coordinates \((t,y)\), with
\(\gamma(t)=(t,0)\). Nearby geodesics are graphs over \(t\) near the
baseline, and their length Lagrangian is
\[
 L(t,y,v)=F((t,y),(1,v)),\qquad
 G=L_{vv}(t,0,0),\quad E=L_{vy}(t,0,0),\quad D=L_{yy}(t,0,0).
\]
These coefficients are smooth on the closed segment. Strong convexity
gives \(G>0\): the full length velocity Hessian has precisely the
radial kernel \(\R(1,0)\), which misses the transverse velocity space.
Write \(J\) for the graph-position variation and
\(P=GJ'+EJ\) for its momentum variation. Linearizing
\(\frac{d}{dt}L_v=L_y\) yields
\[
 J'=-G^{-1}EJ+G^{-1}P,\qquad
 P'=(D-E^\top G^{-1}E)J+E^\top G^{-1}P.
\]
Thus the preceding lemma applies with
\[
 \mathsf C=G^{-1},\qquad \mathsf B=-G^{-1}E,\qquad
 \mathsf A=E^\top G^{-1}E-D.
\]
In particular, \(\mathsf A,\mathsf C\) are symmetric and
\(\mathsf C>0\). A potential-generated variation has
\(P(c)=D_y^2\psi(c,0)J(c)\) on the fixed transverse section.
Hence \(J^\top P=P^\top J\) initially, and the displayed system
preserves this identity. At the typical rough-ray root,
Proposition~\ref{prop:rough-ray-quadratic} supplies exactly such a
smooth auxiliary unit-conormal family with the same position and
cotangent first jets. Symmetry is the symmetry of the extended
potential's second derivative at its regular seed. The smooth flow
then propagates this actual first variation through the closed
segment; no smooth foliation of the original cluster is assumed.

Regraphing at fixed \(t\) subtracts a multiple of the baseline tangent
from each original variation column. The full Hessian \(D^2_{vv}F\)
annihilates that tangent. Therefore the original quadratic matrix is
exactly
\[
 H(t)=J(t)^\top G(t)J(t)=J(t)^\top\mathsf C(t)^{-1}J(t).
\]
Subtracting tangential columns also preserves the full volume
determinant with its longitudinal column. The latter has coordinate
speed one at the baseline. Thus the actual volume Jacobian is
\(w(t)|\det J(t)|\), where \(w\) is the smooth positive ambient
coordinate volume density. The first-jet decomposition gives interior
invertibility of \(J\); its extension may be singular at the ends.

The whole-ray disintegration consequently has the exact form
\[
 \alpha(da)=\frac{\rho_A(\gamma(a))w(a)|\det J(a)|}{\vartheta}\,da,
 \qquad
 \beta(db)=\frac{\rho_B(\gamma(b))w(b)|\det J(b)|}{\vartheta}\,db,
\]
with a common finite positive label density \(\vartheta\).
Insert \eqref{eq:curved-focal-determinant-cancellation}; for example,
\[
 \int\det(I+U)\,d\alpha
 =\vartheta^{-1}\int\rho_A(\gamma(a))w(a)
                    |\det(J(a)-K(a))|\,da<\infty.
\]
Every remaining factor is bounded on a finite interval. Conditional
endpoints have zero mass, so only the interior determinant identity
is used. Vanishing densities and arbitrary conditional support holes
are allowed. For a positive semidefinite matrix,
\(\det(I+U)\ge1+\operatorname{tr}U\ge\|U\|\), giving the first
moments. Finally, \(T\) is negative definite to the left of \(c\)
and positive definite to its right. Applying the corresponding
determinant identity separately on each side proves the last claim.
\end{proof}

These are fixed-root moments, without an endpoint rank assumption or
prescribed vanishing order. The recovery proof below adds legal repair;
no bound uniform over roots is asserted here.

\subsection{Regular focal crossings and mobility moments}
\label{subsec:regular-focal-crossing-moments}

The Hamiltonian identification above also determines the order of every
endpoint singularity. Positivity of its transverse velocity Hessian
gives stronger fixed-root mobility moments without any prescribed
vanishing of the physical density.

\begin{lemma}[Regular crossing at a focal endpoint]
\label{lem:regular-focal-crossing}
Under the Hamiltonian hypotheses of
Lemma~\ref{lem:curved-jacobi-second-solution}, let \(t_0\) be either
endpoint and put \(k=\dim\ker J(t_0)\). If \(k\ge1\), then, as \(t\)
approaches \(t_0\) from the interior,
\begin{equation}
\label{eq:regular-focal-orders}
 |
 \det J(t)|\asymp|t-t_0|^k,\qquad
 \|J(t)^{-1}\|\asymp|t-t_0|^{-1}.
\end{equation}
If \(k=0\), the inverse extends smoothly across that endpoint.
All comparison constants are for the fixed Jacobi system.
\end{lemma}
\begin{proof}
ODE uniqueness and interior invertibility give full endpoint phase rank.
Write
\[
 J_0=J(t_0),\quad P_0=P(t_0),\quad D=J'(t_0),\quad
 N=\ker J_0,\quad E=\operatorname{im}J_0.
\]
The Wronskian gives \(P_0N\perp E\). If \(v\in N\) and
\(Dv\in E\), then
\[
 0=\langle P_0v,Dv\rangle
   =\langle P_0v,\mathsf C(t_0)P_0v\rangle,
\]
so \(v=0\). Rank-nullity makes \(L=\Pi_{E^\perp}D|_N\) invertible.
Use fixed orthonormal bases for \(N^\perp\oplus N\) and
\(E\oplus E^\perp\), and let \(\Pi\) project orthogonally onto \(N\). Then \(A_0=J_0+D\Pi\) is block
triangular with invertible diagonal blocks. For \(s=t-t_0\ne0\), set
\[
 A(s)=J(t_0+s)(I-\Pi)+s^{-1}J(t_0+s)\Pi,\qquad
 J(t_0+s)=A(s)(I-\Pi+s\Pi).
\]
The one-sided derivative gives \(A(s)\to A_0\), hence
\[
 \frac{\det J(t_0+s)}{s^k}\to\det A_0\ne0,\qquad
 sJ(t_0+s)^{-1}\to R=\Pi A_0^{-1},\quad \operatorname{rank}R=k.
\]
For \(k\ge1\), these prove \eqref{eq:regular-focal-orders} at either
endpoint, including \(k=q\). For \(k=0\), smooth inversion applies.

Smoothness also gives \(A(s)=A_0+O(s)\), hence
\(J(t_0+s)^{-1}=R/s+O(1)\). For \(0<k<q\), the fixed blocks have a
smoothly invertible upper-left block, off-diagonal blocks divisible by
\(s\), and Schur complement \(sL(s)\), where \(L(0)=L\).
\end{proof}

\begin{proposition}[Fixed-root moments from regular focal crossings]
\label{prop:regular-focal-crossing-moments}
Retain the genuine Finsler fixed-root geometry and conditional measures
of Proposition~\ref{prop:curved-focal-determinant-moments}. Fix an
interior point \(c\) and set \(T(t)=\int_c^t H(s)^{-1}\,ds\).
At a focal endpoint of nullity \(k\ge1\), every positive moment order
\(r<k+1\) of \(\|T\|\) is integrable against \(\alpha+\beta\) near
that endpoint. A nonfocal endpoint imposes no moment restriction.
In particular,
\begin{equation}
\label{eq:regular-focal-universal-moments}
 \int\|T(t)\|^r\,d(\alpha+\beta)(t)<\infty
 \qquad(0<r<2).
\end{equation}
For \(U=-T\) on source marks before \(c\), and \(V=T\) on target
marks after \(c\), this gives the corresponding two mobility moments.
For every paired law \(\kappa\in\Pi(\alpha,\beta)\) on legal pairs,
\begin{equation}
\label{eq:regular-focal-paired-moment}
 \begin{split}
 \int\|A(a,b)^{-1}\|^r\,\kappa(da,db)
 \le 2^{\max(r-1,0)}\biggl[
 \int\|T(a)\|^r\,\alpha(da)
 +\int\|T(b)\|^r\,\beta(db)\biggr]<\infty
 \end{split}
\end{equation}
for every \(0<r<2\). These bounds are uniform over paired laws at
this fixed root only; no uniform bound over root labels is asserted.
\end{proposition}
\begin{proof}
The actual quadratic matrix and conditional volume density were
identified in Proposition~\ref{prop:curved-focal-determinant-moments}:
\[
 H^{-1}=J^{-1}\mathsf C J^{-\top},\qquad
 d(\alpha+\beta)(t)
 =\frac{(\rho_A+\rho_B)(\gamma(t))w(t)|\det J(t)|}{\vartheta}\,dt,
\]
where \(w\) is smooth positive and \(0<\vartheta<\infty\).
Bounded physical densities and Lemma~\ref{lem:regular-focal-crossing}
give, near a focal endpoint \(t_0\),
\[
 \frac{d(\alpha+\beta)}{dt}\le C|t-t_0|^k,\qquad
 \|H(t)^{-1}\|\le C|t-t_0|^{-2}.
\]
Integration toward the fixed interior point gives
\(\|T(t)\|\le C(1+|t-t_0|^{-1})\). Thus its local \(r\)-moment is
bounded by a constant times \(\int_0^\epsilon s^{k-r}\,ds\), finite
for \(r<k+1\). The same argument applies at the other endpoint;
nonfocal ends and interior compact intervals have bounded \(T\).
Since any focal nullity is at least one, this proves
\eqref{eq:regular-focal-universal-moments}.
Finally,
\(A(a,b)^{-1}=\int_a^b H^{-1}=T(b)-T(a)\).
The norm triangle inequality and
\((x+y)^r\le2^{\max(r-1,0)}(x^r+y^r)\) prove
\eqref{eq:regular-focal-paired-moment} using only the fixed marginals.
This argument allows arbitrary zeros and holes in the physical densities
and does not require common ordering for the bound involving \(T\).
\end{proof}

For a stationary weighted competitor, the Palm paired law has these
same marginals. In particular, if \(\E^S Q<\infty\), H\"older's
inequality gives, for \(1<r<2\) and \(s=2r/(r+1)\),
\[
 \E^S|D|^s
 \le2^{s/2}\bigl(\E^S\|A^{-1}\|^r\bigr)^{(2-s)/2}
                   (\E^S Q)^{s/2}<\infty.
\]
Thus every positive displacement moment below \(4/3\) follows.
The general value recovery below uses full-path traffic and does not
require a displacement second moment of the weighted minimizer.

% Source: minimizer-focal-regularity-attempt.md, Sections 1--3.
% Source SHA256: 2fb73e8abbadc335bba26a776f293b88a7b3ecdaae34a74c60cf499dc6bbe779
% Independent PASS: minimizer-focal-regularity-attempt-review.md
% Review SHA256: c6094002b0e03ba53384531c7da83108384af0ce3e01c64d15eab3c595728a18
\subsection{Simple focal poles}
\label{subsec:focal-simple-pole-dini}

The smoothly continued second Jacobi solution improves the endpoint
mobility expansion: its remainder is bounded, with no logarithmic term.
All statements are at one fixed regular ray.

\begin{lemma}[Smooth remainder after the focal pole]
\label{lem:focal-simple-pole}
Retain the Hamiltonian system, actual quadratic matrix \(H\), and
cumulative compliance \(T\) of
Lemma~\ref{lem:curved-jacobi-second-solution}. At an endpoint \(t_0\),
put \(k=\dim\ker J(t_0)\). There are a fixed symmetric positive
semidefinite matrix \(G_0\) of rank \(k\) and a matrix function
\(B_0\), smooth through zero, such that
\begin{equation}
 T(t_0+s)=-\frac{G_0}{s}+B_0(s).
 \label{eq:focal-simple-pole}
\end{equation}
At a nonfocal endpoint \(G_0=0\). Consequently, on source and target
mark ranges separated from a common interior cut \(c\),
\begin{equation}
 U(a)=\frac{G_L}{a-\ell}+B_L(a),\qquad
 V(b)=\frac{G_R}{u-b}+B_R(b),
 \label{eq:focal-two-residues}
\end{equation}
where the remainders are bounded symmetric matrices. They need not be
positive semidefinite, and \(G_L,G_R\) need not commute.
\end{lemma}

\begin{proof}
Use the fixed domain and codomain decompositions from
Lemma~\ref{lem:regular-focal-crossing}. Its Schur complement has the
exact form \(sL(s)\), with \(L\) smooth and \(L(0)\) invertible.
The upper-left block is smoothly invertible and both off-diagonal
blocks have a factor \(s\). Block inversion therefore strengthens the
bounded inverse remainder to
\[
 J(t_0+s)^{-1}=s^{-1}R+B_J(s),
\]
where \(B_J\) is smooth and \(\operatorname{rank}R=k\).
When \(k=q\), the same statement follows directly by factoring \(s\)
from \(J(t_0+s)\). The second solution \(K=JT\) extends smoothly
by Lemma~\ref{lem:curved-jacobi-second-solution}. Write
\(K(t_0+s)=K_0+sK_1(s)\). Multiplication gives
\[
 T(t_0+s)=\frac{RK_0}{s}+B_0(s)
\]
with smooth remainder. On the interior,
\[
 T'=H^{-1}=J^{-1}\mathsf C J^{-\top},\qquad
 \lim_{s\to0}s^2T'(t_0+s)=R\mathsf C(t_0)R^\top.
\]
Differentiating the pole expression identifies
\(RK_0=-R\mathsf C(t_0)R^\top\). Set
\(G_0=R\mathsf C(t_0)R^\top\); positivity of \(\mathsf C(t_0)\)
gives its positivity and rank. Smooth inversion handles \(k=0\).
On the left \(U=-T\) and \(s=a-\ell>0\); on the right \(V=T\)
and \(s=b-u<0\). This proves~\eqref{eq:focal-two-residues}, with
boundedness elsewhere supplied by the regular compact mark range.
\end{proof}

In particular,
\[
 H(t_0+s)^{-1}=\frac{G_0}{s^2}+B_0'(s).
\]
The possible order-\(s^{-1}\) term cancels in the actual Hamiltonian
geometry. No assumption on physical density vanishing is used here.
\subsection{Conditional centering by endpoint translations}

Retain independent unit-intensity marked Poisson inputs on \(\R^q\),
with source and target laws \(\alpha,\beta\), and the stationary
classes in Definition~\ref{general-gluing-two-classes}. Suppose all sources precede
a common cut and all targets follow it. Write
\[
 \begin{gathered}
 S(a,b)=U(a)+V(b),\qquad A=S^{-1},\qquad U,V\ge hI>0,\\
 Q(a,b,D)=\tfrac12D^TAD,\qquad p=AD.
 \end{gathered}
\]
\begin{proposition}[Conditional neutrality at a weighted minimum]
\label{prop:common-cut-conditional-centering}
If \(P\) attains \(I_w(t)\) for a bounded paired-mark probe, then
\[
 \E_P^S[p\mid a]=0,\qquad \E_P^T[p\mid b]=0.
\]
The same directed momentum \(p=A(y-x)\) is used at both anchors.
If additionally
\begin{equation}
\label{eq:common-cut-first-mobility}
 \int\|U(a)\|\,\alpha(da)+\int\|V(b)\|\,\beta(db)<\infty,
\end{equation}
then \(Up,Vp,D\) are integrable and
\(\E_P^S[Up]=\E_P^S[Vp]=\E_P^S D=0\).
For the invariant sigma field \(\Xi\) of the joint matching law,
the stronger componentwise identities hold:
\[
 \E^S[p\mid a,\Xi]=\E^T[p\mid b,\Xi]=0,\qquad
 \E^S[Up\mid\Xi]=\E^S[Vp\mid\Xi]=\E^S[D\mid\Xi]=0,
\]
where the second group uses \eqref{eq:common-cut-first-mobility}.
\end{proposition}
\begin{proof}
A target translation by \(\epsilon h_T(b)\), with bounded \(h_T\),
preserves the input and paired-mark laws and changes kinetic energy by
\[
 \epsilon\E^S[p\cdot h_T(b)]
 +\tfrac12\epsilon^2\E^S[h_T(b)^TAh_T(b)].
\]
All terms are finite. Minimality for both signs of \(\epsilon\)
gives the target conditional identity. The source variation has the
opposite sign and gives the source identity. Mass transport identifies
the directed Palm edge laws at the two anchors. Furthermore,
\[
 |Up|^2\le\|U\|p^TUp\le\|U\|p^TSp,\qquad
 \E^S|Up|\le(2\E^S\|U\|\,\E^S Q)^{1/2},
\]
and likewise for \(V\). These bounds prove integrability under
\eqref{eq:common-cut-first-mobility}. Multiply the conditional
identities by truncated mobility matrices and pass to the limit using
that integrability. Since \(D=Up+Vp\), all three means vanish.

For the component assertion, the prescribed independent marked Poisson
input law is ergodic under common translations. Its marginal on almost
every ergodic component of a stationary extension is therefore that same
input law. Each component has unit source intensity, so the Palm
decomposition uses the same component weights. Its objective is at
least \(I_w(t)\); their average equals \(I_w(t)\), so almost every
component is itself minimizing. Apply the preceding first variations
on each component. Conditional energies are finite almost surely, and
the conditional mobility means in \eqref{eq:common-cut-first-mobility}
are the same deterministic numbers. The conditional version of the
integrability estimate then gives the asserted componentwise means.
\end{proof}
\subsection{Phase envelopes for the original Poisson inputs}

Consider stationary perfect matchings of the original independent
unit-intensity Poisson inputs with product marks $\alpha,\beta$.
All source marks precede a common interior cut $c$, all targets follow
it, and every crossed source--target pair is legal. For continuous
positive definite $H$ on the open scalar interval, put
\[
 C(r)=\int_c^r H(s)^{-1}\,ds,\quad U(a)=-C(a),\quad V(b)=C(b),
 \quad S=U+V,\quad A=S^{-1}.
\]
Assume $U,V\ge hI$ for $h>0$ and
$M_U=\int\det(I+U(a))\,\alpha(da)<\infty$.
Write $\E^S$ for source Palm expectation. On each edge define
\[
 p=A(y-x),\qquad z=x+Up=y-Vp,\qquad X(r)=z+C(r)p,
 \qquad Q=\tfrac12p^TSp.
\]
Suppose $\E^S Q<\infty$, so $\E^S|p|^2\le h^{-1}\E^S Q$.
We allow a bounded two-cycle defect
\begin{equation}\label{eq:degenerate-central-cycle-defect}
 Q_{ii}+Q_{jj}\le Q_{ij}+Q_{ji}+D_0,\qquad D_0\ge0.
\end{equation}
Here $Q_{ij}$ is the cost from source $i$ to target $j$.
Local optimality for $Q+tg$, with $|g|\le G$, gives $D_0=4|t|G$.
Campbell's formula, integrating the source anchor first, gives
\begin{equation}\label{eq:degenerate-cut-campbell}
 \E\sum_e\mathbf1_{\{z_e\in B\}}f(a_e,b_e,p_e)
 =|B|\E^S f(a,b,p)
\end{equation}
for nonnegative $f$ and finite-volume $B$. The shift $z=x+U(a)p$
requires no integrability of $Up$ in this nonnegative identity.

\begin{lemma}[Approximate momentum field and Poisson envelopes]
\label{degenerate-central-phase-envelope}
Set $L=(1+\sqrt2)/h$ and $b=2\sqrt{D_0/h}$.
The edge momenta satisfy $|p_i-p_j|\le L|z_i-z_j|+b$.
There is a measurable translation-equivariant field $\Pi$ with
$\Lip(\Pi)\le\sqrt qL$ and $|\Pi(z_i)-p_i|\le\sqrt q b$.
A union of $m$ unit phase cells in $(z,p)$ has a containing full-input
marked Poisson envelope of mean at most $C M_Um$ for its source vertices.
For a phase ball of radius $W\ge1$, its envelope count $N_W$ satisfies
$\E N_W^4\le C(1+W^{4q})$.
\end{lemma}

\begin{proof}
The common-cut midpoint calculation uses only $U^{-1},V^{-1}\le h^{-1}I$.
For $e=z_i-z_j$, the own-pair excess at the midpoint is at most
$|e|^2/(2h)$. Minimizing the first crossed quadratic saves at least
$h|r|^2/4$, where
$r=p_i-p_j+(V_j^{-1}-U_i^{-1})e/2$.
Equation~\eqref{eq:degenerate-central-cycle-defect} therefore gives
$h|r|^2/4\le |e|^2/(2h)+D_0$. Since
$|r-(p_i-p_j)|\le|e|/h$, the asserted bound follows.
Define the coordinatewise McShane extension by
\[
 \Pi_k(w)=\inf_i\{p_{i,k}+L|w-z_i|\},\qquad k=1,\ldots,q.
\]
A fixed edge bounds each infimum above; the momentum bound and the
triangle inequality bound it below. At $z_i$ its value lies between
$p_{i,k}-b$ and $p_{i,k}$. This proves the extension bounds; countability
of the edges gives measurability, and the formula commutes with translations.

For a fixed-size phase cell centered at $(k,l)$, all relevant sources lie in
\[
 \{(x,a):x\in k-U(a)l+(I+U(a))B_R\},
\]
with fixed $R$. Indeed $(I+U)^{-1}$ and $(I+U)^{-1}U$ have norm at most
one. Its marked intensity is at most $C_R\int\det(I+U)\,d\alpha$.
The count in this deterministic marked set is Poisson. Taking unions
proves the $m$-cell bound, without independence between envelopes.
A phase ball of radius $W$ is contained in the marked envelope
$\{(x,a):x\in(I+U(a))B_{CW}\}$, of mean at most $CM_UW^q$.
The Poisson fourth-moment formula proves the last assertion.
\end{proof}
\section{Complete-path traffic and interleaved populations}\label{sec:complete-path-traffic}

\subsection{Full kinetic-path traffic through focal endpoints}
\label{subsec:focal-traffic}

Fix a regular positive-mass root with all source marks before one
interior cut $c$ and all target marks after it. Retain the actual
conditional laws $\alpha,\beta$ of the bounded physical densities and
the genuine transverse matrix $H$. Write
\[
 C(r)=\int_c^r H(s)^{-1}\,ds,\quad U(a)=-C(a),\quad V(b)=C(b),
 \quad S=U+V,\quad A=S^{-1}.
\]
For a matched edge let
\[
 D=y-x,\qquad p=AD,\qquad z=x+Up=y-Vp,\qquad
 X(r)=z+C(r)p,\qquad Q=\tfrac12p^\top Sp.
\]
All expectations of edge variables below are source-Palm expectations.
The common cut gives $U,V\ge hI$ for a fixed $h>0$.

\begin{theorem}[Full-path traffic at an actual focal root]
\label{thm:focal-full-path-traffic}
Let $q>4$. Suppose a stationary perfect matching on the original
independent marked Poisson inputs has $\E Q<\infty$ and, for every two
of its edges,
\[
 Q_{ii}+Q_{jj}\le Q_{ij}+Q_{ji}+D_0
\]
with a fixed $D_0\ge0$. For every fixed nonnegative smooth compactly
supported kernel $\eta$, its full absolute kinetic traffic
\[
 \mathcal T_\eta(w)=\sum_e\int_{a_e}^{b_e}
       |H(r)^{-1}p_e|\eta(w-X_e(r))\,dr
\]
satisfies
\begin{equation}\label{eq:focal-traffic-full-bound}
 \E[\mathcal T_\eta(0)^2]\le C(1+\E Q)<\infty.
\end{equation}
The constant can depend on the fixed root, its coefficient and density
bounds, $D_0$, and $\eta$. If stationary edge indicators
$\mathbf1_{B_m}$ tend to zero on every edge almost surely, then the
corresponding absolute traffic tends to zero in $L^2$. The theorem
applies in particular to a weighted mean minimizer of $Q+t g$ for
bounded $g$, with $D_0=4|t|\|g\|_\infty$.
\end{theorem}

\begin{proof}
We first prove the result for a kernel supported in the physical unit
ball, of integral one. All estimates use the original physical source
anchor in Campbell's formula. Constants never depend on the dyadic
indices introduced below.

\emph{Focal coefficients and a physical one-edge bound.}
Lemma~\ref{lem:focal-simple-pole} and the regular crossing imply, at a
focal endpoint of nullity $k\ge1$, in its inward distance coordinate $s$,
\begin{equation}\label{eq:focal-traffic-pole}
 H^{-1}=G/s^2+R(s),\qquad
 c_0(s^{-2}\Pi+I-\Pi)\le H^{-1}
                  \le C_0(s^{-2}\Pi+I-\Pi).
\end{equation}
Here $G$ is fixed positive semidefinite of rank $k$, $\Pi$ projects
onto $\operatorname{range}G$, and $R$ is smooth through zero.
The conditional density is at most $Cs^k$.
The anisotropic comparison follows directly from the regular-crossing
block inverse: $(s\Pi+I-\Pi)J(s)^{-1}$ has an invertible limit, so its
product with the positive Hamiltonian coefficient and its transpose
has a positive definite limit. No simultaneous diagonalization is used.
At a nonfocal end ordinary smooth positive bounds suffice.
The bounded mobility remainders also give
\[
 e_L=\frac{p^\top G_Lp}{a-\ell},\qquad
 e_R=\frac{p^\top G_Rp}{u-b},\qquad
 \mathcal E:=1+\E Q+\E|p|^2+\E(e_L+e_R)\le C(1+\E Q).
\]
Zero residues give zero focal energies; $\E|p|^2\le\E Q/h$.

Translate the left endpoint to zero and abbreviate $W(r)=H(r)^{-1}$.
The block decomposition $\operatorname{range}G\oplus\ker G$ gives
\[
 \|W(r_0)^{-1}\|\le C,\qquad
 \|G W(r_0)^{-1}\|\le Cr_0^2.
\]
Indeed the complementary block of $R$ is uniformly positive, and its
Schur complement is $G|_{\operatorname{range}G}/r_0^2+O(1)$;
the first block row of the inverse is $O(r_0^2)$. The full-rank case
is immediate. Consequently, if $|r-r_0|\le\epsilon r_0$,
\[
 \|W(r)W(r_0)^{-1}-I\|\le C\epsilon.
\]
Choose $\epsilon$ so that this is at most $1/4$. For $p\ne0$ the
velocity remains within $|W(r_0)p|/4$ of $W(r_0)p$. Projection onto
that fixed direction is increasing at speed at least
$3|W(r_0)p|/4$, and total speed is at most $5|W(r_0)p|/4$.
The time spent in the unit ball is at most $8/(3|W(r_0)p|)$.
Every one-edge contribution on such a window therefore obeys
\begin{equation}\label{eq:focal-traffic-one-edge}
 u_i:=\int_I|W(r)p_i|\eta(-X_i(r))\,dr\le C,
 \qquad u_i\le C|I|\sup_I|W(r)p_i|.
\end{equation}
This is a physical bound, without an anisotropic condition-number loss.

\emph{A short-window estimate at a coarse phase scale.}
Consider a window about $r_0$ of length $O(\epsilon s)$ and a
stationary family active there, whose sources belong to a deterministic
mark set $B$. Choose a positive matrix $\Lambda$ and $L\ge1$, and put
\[
 p'=\Lambda^{1/2}p,\quad
 U'=\Lambda^{-1/2}[C(r_0)-C(a)]\Lambda^{-1/2},\quad
 V'=\Lambda^{-1/2}[C(b)-C(r_0)]\Lambda^{-1/2}.
\]
Suppose $U',V'\ge h_0I$,
$s\Lambda^{-1/2}W(r)\Lambda^{-1/2}$ is uniformly bounded on the window,
and
\begin{equation}\label{eq:focal-traffic-scale-conditions}
 L^{-1}\|\Lambda^{-1/2}\|\le1,\qquad
 (\det\Lambda)^{1/2}L^q
       \int_B\det(I+U'(a))\,\alpha(da)\le C_0.
\end{equation}
Assume also \eqref{eq:focal-traffic-one-edge} with uniform constants.
Additional adaptive restrictions on the family are allowed: envelopes
always count the full original source input on $B$.

Use coarse phase coordinates
$\bar z=L^{-1}\Lambda^{-1/2}X(r_0)$, $\bar p=p'/L$.
The common-cut midpoint inequality gives a momentum extension
$\mathcal P$ with uniformly bounded Lipschitz constant and bounded
error at the actual phase points; its cycle defect is $D_0/L^2$.
Fix once per band and window the coordinatewise McShane extension,
which is measurable and translation equivariant, before making any
phase or exceptional-event split. For small fixed $\epsilon$, the maps
\[
 F_r(\bar z)=\bar z+R_r\mathcal P(\bar z),\qquad
 R_r=\Lambda^{-1/2}[C(r)-C(r_0)]\Lambda^{-1/2}
\]
are uniformly bi-Lipschitz by contraction. Their inverse image of zero
traces a curve in phase space of length $O(W_0)$ and radius $O(1+W_0)$,
where $W_0=|\mathcal P(0)|$: subtracting the inverse equations gives
$|F_r^{-1}(0)-F_v^{-1}(0)|\le C W_0|r-v|/s$.
Every contributor to the physical kernel has coarse position norm at
most one by \eqref{eq:focal-traffic-scale-conditions}, so its phase point
lies in a fixed-radius tube about this curve. For large $W_0$, its
$|p'|/L$ is comparable to $W_0$; bounded $W_0$ gives a fixed phase ball.

A unit coarse phase cell pulls back to a marked source envelope with
spatial volume at mark $a$ at most
$C(\det\Lambda)^{1/2}L^q\det(I+U'(a))$. Its count is Poisson with
uniformly bounded mean. Unions of $m$ cells have mean at most $Cm$.
A phase ball of radius $W\ge1$ has an ellipsoid envelope of mean
$CW^q$ and fourth count moment at most $C(1+W^{4q})$.
For $W_0\in[W,2W]$, a candidate tube is covered by an animal of
$m=\lceil C'W\rceil$ cells with root in a ball of radius $C'W$.
There are at most $CW^{2q}a_q^m$ such animals. A Poisson tail and
union bound give a data-only exceptional event of probability
$Ce^{-cKW}$, $K\ge K_0$, outside which every such animal contains
at most $CKW$ source points. For bounded $W_0$ use one fixed phase
envelope and probability $Ce^{-cK}$. These events control all adaptive
curves; no selected cloud is declared Poisson.

Split the traffic into large-phase and bounded-phase good-event parts,
$T_{\rm lg},T_{\rm sm}$, and an exceptional part. Pointwise,
\[
 T_{\rm lg}^2\le CK\sum_i (|p_i'|/L)u_i,
 \qquad T_{\rm sm}^2\le CK\sum_i u_i.
\]
Integrating the original physical source anchor first in Campbell gives
\begin{align}
 \E T_{\rm lg}^2&\le CK L^{-1}
        \int_I\E[|p'|\,|W(r)p|\mathbf1_B]\,dr,
       \label{eq:focal-traffic-large}\\
 \E T_{\rm sm}^2&\le CK\int_I\E[|W(r)p|\mathbf1_B]\,dr.
       \label{eq:focal-traffic-small-first}
\end{align}
If also $U'\le CI$, bounded-phase contributors have original sources
in $\Lambda^{1/2}L B_C$, independently of their marks. If
$u_i\le CM|p_i'|$, Cauchy--Schwarz over at most $CK$ contributors
therefore gives the alternative bound
\begin{equation}\label{eq:focal-traffic-small-second}
 \E T_{\rm sm}^2\le CK M^2(\det\Lambda)^{1/2}L^q
                                      \E[|p'|^2\mathbf1_B].
\end{equation}
The Palm momentum is not conditioned to be independent of this source
set: Campbell supplies its deterministic spatial volume factor.
On an exceptional event, \eqref{eq:focal-traffic-one-edge} bounds the
traffic by $CN_W$. Cauchy--Schwarz and the fourth count moment give
$CW^{2q}e^{-cKW}$. Summing dyadic $W$, including the bounded case,
shows that the exceptional part has $L^2$ norm at most $Ce^{-cK}$,
after changing $c$. All these constants are uniform under the stated
scale conditions.

\emph{Old sources near the focus.}
For a small reference time $t$, take a window of width $\epsilon t$,
sources $a<\gamma t$ with fixed $\gamma=3/4$, and
\[
 \Lambda_t=t^{-1}\Pi+t(I-\Pi),\qquad L_t=t^{-1/2-1/q}.
\]
Both normalized legs have fixed positive lower bounds. Integration of
\eqref{eq:focal-traffic-pole} gives
$\det(I+U'_t(a))\le C(t/a)^k$. Since $\alpha(da)\le Ca^kda$,
\[
 (\det\Lambda_t)^{1/2}\int_0^{\gamma t}
       \det(I+U'_t(a))\,\alpha(da)\le Ct^{q/2+1}.
\]
Multiplication by $L_t^q$ gives a uniform envelope; the physical coarse
kernel radius is at most $t^{1/q}$. All short-window hypotheses hold.
Set $e(t)=\E[e_L\mathbf1_{a<\gamma t}]$,
$v(t)=\E[|p|^2\mathbf1_{a<\gamma t}]$.
The bound $\E[|\Pi p|^2\mathbf1_{a<\gamma t}]\le Ct e(t)$ and
$\int_I|W(r)p|dr\le Ct^{-1/2}|p'|$ give
\[
 \E T_{\rm lg}^2\le CK\{t^{1/q}e(t)+t^{1+1/q}v(t)\}.
\]
For the small-phase term use
$\int_I|W(r)p|dr\le C(|Gp|/t+t|p|)$ and
\[
 \alpha\{a<\gamma t\}\le Ct^{k+1},\qquad
 \int_{a<\gamma t}a\,\alpha(da)\le Ct^{k+2}.
\]
Cauchy--Schwarz in \eqref{eq:focal-traffic-small-first} yields
\[
 \E T_{\rm sm}^2\le CK\{t^{k/2}\sqrt{e(t)}
                              +t^{(k+3)/2}\sqrt{v(t)}\}.
\]
Thus, since $q>4$ and $k\ge1$,
\begin{equation}\label{eq:focal-traffic-old-bound}
 \|T_{{\rm old},t}\|_2
       \le C\sqrt K\sqrt{\mathcal E}\,t^{1/(2q)}+Ce^{-cK}.
\end{equation}
A fixed number of relative windows covers each dyadic time slab.
Every portion with $a\le r/2$ is included by choosing a nearby window
center, since $a<3t/4$ there. For $t=t_*2^{-n}$ choose
$K=K_0+C_1(n+1)$; both terms in \eqref{eq:focal-traffic-old-bound}
are summable.

\emph{Births comparable to the focal time.}
For the remaining portions, $r/2<a<r$. Partition source marks into
$t\le a<2t$, and elapsed times into $s<r-a\le2s$, with
$s=t2^{-j}$, $j\ge0$. These half-open intervals partition $(0,2t]$,
which contains every $r-a<a$ in this source shell. Partition the shell
into $N\le Ct/s$ mark bands $B_i$ of width at most $s/8$.
Their relevant time ranges are covered by a fixed number of windows
of width $\epsilon s$, with centers satisfying
$c_1s\le r_0-a\le C_1s$ for every mark in the band and $r_0\asymp t$.
Use
\[
 \Lambda_{t,s}=(s/t^2)\Pi+s(I-\Pi),\qquad
 L_{t,s}=s^{-1/2}t^{-1/q}.
\]
Here $cI\le U'\le CI$, $V'\ge cI$, and the normalized coefficient
condition holds at time scale $s$. The physical coarse kernel radius
is at most $t^{1/q}$. Moreover
\[
 \alpha(B_i)\le Ct^k s,\qquad
 (\det\Lambda_{t,s})^{1/2}L_{t,s}^q=t^{-k-1}.
\]
Thus each phase-cell envelope has mean at most $Cs/t\le C$.
The one-edge bound holds because the window width is at most a fixed
small multiple of $r_0$. All event constants remain uniform.

Put $e_i=\E[e_L\mathbf1_{B_i}]$, $v_i=\E[|p|^2\mathbf1_{B_i}]$.
Their sums are at most $\mathcal E$, and
\[
 |p'|^2=(s/t^2)|\Pi p|^2+s|(I-\Pi)p|^2,
 \qquad \E[|\Pi p|^2\mathbf1_{B_i}]\le Ct e_i.
\]
Since $\int_I|W(r)p|dr\le C\sqrt{s}|p'|/t$,
\eqref{eq:focal-traffic-large} gives
\[
 \E T_{{\rm lg},i}^2
      \le CK\{s^2t^{-2+1/q}e_i+s^2t^{-1+1/q}v_i\}.
\]
Minkowski and Cauchy--Schwarz over the $N\le Ct/s$ bands imply
\begin{equation}\label{eq:focal-traffic-birth-large}
 \Big\|\sum_iT_{{\rm lg},i}\Big\|_2
      \le C\sqrt K\sqrt{\mathcal E}\,t^{1/(2q)}\sqrt{s/t}.
\end{equation}

For the same bounded-phase good-event part we need two estimates.
First, \eqref{eq:focal-traffic-small-first} gives
\[
 \E T_{{\rm sm},i}^2\le CK\{
 s^{3/2}t^{(k-3)/2}\sqrt{e_i}
             +s^{3/2}t^{k/2}\sqrt{v_i}\}.
\]
Indeed $\E[|Gp|\mathbf1_{B_i}]\le C\sqrt{e_i t\alpha(B_i)}$,
and the focal speed integrated over the window is at most $Cs|Gp|/t^2$.
Using $\sum_i e_i^{1/4}\le N^{3/4}(\sum_i e_i)^{1/4}$ and the
corresponding bound for $v_i$, we obtain
\begin{equation}\label{eq:focal-traffic-birth-first}
 \Big\|\sum_iT_{{\rm sm},i}\Big\|_2
                 \le C\sqrt K\sqrt{\mathcal E}\,t^{k/4}.
\end{equation}
This alone has no summable gain in the birth-distance scale.
Instead use $U'\le CI$, $u_i\le C\sqrt{s}|p_i'|/t$, and
\eqref{eq:focal-traffic-small-second} to get
\[
 \E T_{{\rm sm},i}^2
           \le CK\{s^2t^{-k-4}e_i+s^2t^{-k-3}v_i\}.
\]
Summing square roots yields the second estimate
\begin{equation}\label{eq:focal-traffic-birth-second}
 \Big\|\sum_iT_{{\rm sm},i}\Big\|_2
      \le C\sqrt K\sqrt{\mathcal E}\,t^{-(k+2)/2}\sqrt{s/t}.
\end{equation}
Both estimates concern the same random contribution, so take their
minimum. With $\theta=k/[2(3k+4)]$ and $b_k=\theta/2>0$, the inequality
$\min(A,B)\le A^{1-\theta}B^\theta$ gives
\[
 \Big\|\sum_iT_{{\rm sm},i}\Big\|_2
       \le C\sqrt K\sqrt{\mathcal E}\,t^{k/8}(s/t)^{b_k}.
\]
Including the fixed number of windows per band and all exceptional
parts proves
\begin{equation}\label{eq:focal-traffic-birth-sum}
 \begin{split}
 \|T_{{\rm birth},t,s}\|_2\le{}&
 C\sqrt K\sqrt{\mathcal E}
       \{t^{1/(2q)}(s/t)^{1/2}+t^{k/8}(s/t)^{b_k}\}\\
 &+C(t/s)e^{-cK}.
 \end{split}
\end{equation}
For $t=t_*2^{-n}$, $s=t2^{-j}$, choose
$K=K_0+C_2(n+j+1)$. With $C_2$ large, the last term is at most
$C2^{-2n-2j}$. The two main terms are summable over both indices,
since their geometric exponents are strictly positive; the factor
$\sqrt{n+j+1}$ is harmless. This sums every mark band and every
exceptional event without asserting their independence.

\emph{The remaining path pieces.}
Reverse time and exchange colors for the right focal region. The
reversed compliance derivative is positive, the same pole expansion
holds, and the original target PPP provides its envelopes. Mass
transport identifies the required target-Palm edge moments with their
source-Palm counterparts.
In the compact regular remainder choose a fixed small birth/death
distance below the common source-target separation. At a source birth,
partition elapsed time into scales $s$, marks into bands of width $s/8$,
and time into windows of length $\epsilon s$. Take
$\Lambda=sI$, $L=s^{-1/2}$. The source normalized leg is bounded
above and below, the target leg is bounded below, and the source
envelope has bounded mean because $(\det\Lambda)^{1/2}L^q=1$.
Both \eqref{eq:focal-traffic-large} and
\eqref{eq:focal-traffic-small-second} give
$CKs^2\E[|p|^2\mathbf1_B]$. Summing norms over $O(s^{-1})$ bands
gives $C\sqrt{Ks}\sqrt{\E|p|^2}+Cs^{-1}e^{-cK}$, summable for
dyadic $s$ and $K$ proportional to $1+\log(1/s)$. Target deaths are
handled in the same way.
Away from births and deaths both legs have fixed positive lower
bounds. Finitely many short windows with $\Lambda=I,L=1$ suffice.
Their source determinant envelopes are finite: changing the interior
cut adds a fixed bounded matrix to the common-cut mobility, so
$\det(I+U_{\rm new})\le C\det(I+U)$ on sources preceding the new cut.
Proposition~\ref{prop:curved-focal-determinant-moments} gives its
integrability. Additional restrictions on the matched target only
decrease counts and Palm moments; the envelope is the full original
source restriction. These finitely many norms are at most
$C\sqrt{\mathcal E}$.

Minkowski and the preceding convergent series prove
\eqref{eq:focal-traffic-full-bound}. At a fixed point the full traffic
is therefore finite almost surely. Edgewise convergence of indicators
gives $\mathcal T_{\eta,B_m}\to0$ there, by dominated convergence in
the countable sum of one-edge integrals. Since it is bounded by
$\mathcal T_\eta$, a second dominated convergence proves the $L^2$ tail
claim. Finally, any fixed compact kernel, or the absolute values of
its first derivatives, is bounded by a constant times finitely many
translates of a smooth nonnegative bump positive on a smaller unit
ball. Minkowski and stationarity extend the result to these kernels.
\end{proof}

\begin{corollary}[Original-input divergence and mean current]
\label{cor:focal-traffic-current}
For a kernel of integral one, define the full signed current
\[
 J_\eta(w)=\sum_e\int_{a_e}^{b_e}
                 H(r)^{-1}p_e\,\eta(w-X_e(r))\,dr.
\]
It belongs locally to $L^2$ in expectation and, in distributions,
\begin{equation}\label{eq:focal-traffic-divergence}
 \operatorname{div}J_\eta=\eta*\mu_S-\eta*\mu_T,
\end{equation}
where $\mu_S,\mu_T$ are the actual original endpoint PPPs.
For a stationary weighted mean minimizer the full current has zero
conditional mean on every ergodic component. Selected edge tails
have vanishing $L^2$ current when their indicators vanish edgewise,
but are not asserted individually centered.
\end{corollary}
\begin{proof}
The absolute traffic dominates the signed current and justifies
integration of the derivative along every complete path, proving
\eqref{eq:focal-traffic-divergence}. The first marginal mobility moments
and finite weighted energy give $\E|Up|+\E|Vp|<\infty$.
For a mean minimizer, bounded markwise endpoint translations preserve
the exact inputs and paired probe; their first variations give
$\E[p\mid a]=\E[p\mid b]=0$ on each ergodic component. Each component
has the same original ergodic product PPP input marginal. Truncation
and the preceding integrability imply $\E(Up+Vp)=0$ componentwise.
Campbell, justified by the full absolute-traffic bound, identifies
the mean signed current with this mean displacement. Tail domination
gives the last assertion without imposing centering on a selected
family.
\end{proof}

\subsection{Whole-path screening and equality of the two stationary values}
\label{subsec:general-path-weighted-screening}

The screening argument uses square-integrable traffic of complete paths.
It does not require a second moment of their endpoint displacements.
We first isolate this distinction, then apply the focal traffic theorem.

\begin{proposition}[Screening arbitrary finite paths]
\label{prop:general-finite-path-screening}
Let $q\ge3$. On the original independent marked Poisson inputs of common
intensity $\lambda>0$, let a stationary perfect matching carry covariant
finite absolutely continuous paths from each source to its target.
Suppose their smoothed full absolute traffic $T$ belongs to $L^2$, and
the full signed current $J$ has zero invariant conditional mean.
The same traffic bound may be assumed for a smooth compact kernel
dominating the first derivatives of the chosen kernel; an additional
convolution also supplies the required smooth current.

Choose $\ell_L\to\infty$ with $\ell_L/L\to0$. Write
$\mathcal L_e=\operatorname{length}(X_e)$ and retain exactly the edges
with source in $Q_{L-4\ell_L}$ and $\mathcal L_e\le\ell_L$.
This is a measurable rule, and every retained path lies in
$Q_{L-2\ell_L}$. The remaining endpoints of the original inputs in
$Q_L$ have expected number $o(L^q)$, and their
mass-preserving reflected smoothings $a_L,b_L$ satisfy
\begin{equation}
 \E\|a_L-b_L\|_{H_N^{-1}(Q_L)}^2=o(L^q).
 \label{eq:general-path-port-bound}
\end{equation}
When $\lambda=1-\rho$, adding independent reserves of intensity $\rho$
and the same mark laws gives a maximum-cardinality matching of the full
unit-intensity input in $Q_L$, preserving every retained main edge, with
\begin{equation}
 \E[\text{additional squared-distance cost}]
 \le C_q\rho^{1-2/q}L^q+o_{\rho,P}(L^q).
 \label{eq:general-path-reserve-bound}
\end{equation}
Neither conclusion assumes a finite displacement second moment.
\end{proposition}

\begin{proof}
Let $T_{>R}$ be the traffic of paths with $\mathcal L_e>R$.
Every $\mathcal L_e$ is finite. Dominated convergence, first over the
edge integrals and then under $T(0)^2$, gives
$\E T_{>R}(0)^2\to0$. Let $r_\eta$ bound the kernel support and write
$K_L$ for the current of retained paths. Outside a boundary collar of
width $3\ell_L+r_\eta$, every contributing path with
$\mathcal L_e\le\ell_L$ is retained. Stationarity therefore gives
\begin{equation}
 \frac1{L^q}\E\int_{Q_L}|J-K_L|^2
 \le C\min\left\{1,\frac{\ell_L+r_\eta}{L}\right\}\E T(0)^2
                   +\E T_{>\ell_L}(0)^2\longrightarrow0.
 \label{eq:general-path-current-truncation}
\end{equation}
A source port lies in the collar outside $Q_{L-4\ell_L}$ or belongs
to a path longer than $\ell_L$. A target port lies outside
$Q_{L-6\ell_L}$ or belongs to such a long path. Campbell and
mark-dependent spatial translation give, for either color,
\[
 \frac{\E[\text{ports}]}{L^q}
 \le C\lambda\frac{\ell_L}{L}
       +\lambda\P^S(\mathcal L_e>\ell_L)\longrightarrow0.
\]
Only almost sure finiteness of the path length is used here.

Integrating along the complete paths gives the exact identity
$\operatorname{div}J=\eta*(\mu_S-\mu_T)$.
Absolute traffic justifies the local distributional exchanges.
Lemma~\ref{lem:highd-current-boundary} therefore applies to $J$: its
hypotheses concern this divergence, square integrability, and invariant
mean, with no restriction on path shape. The reflected full-input error
$Z_L$ from \eqref{eq:highd-reflected-input-error} still satisfies
$\E\|Z_L\|_{H_N^{-1}}^2\le C\lambda L^{q-1}$, since the complete input
is the original independent Poisson pair. For the reflected port
smoothings, the exact identity is
\[
 a_L-b_L=\operatorname{div}(J-K_L)+Z_L.
\]
The kept paths and their kernels are interior for large $L$.
Integration by parts against mean-zero $H^1$ tests gives the bulk
term in \eqref{eq:general-path-current-truncation}, the full-current
boundary term, and $Z_L$. Their expected squared norms are all
$o(L^q)$, proving \eqref{eq:general-path-port-bound}.

The proof of Corollary~\ref{cor:highd-kinetic-screening} now uses only
this port bound and the vanishing port count. Explicitly, if $H_S,H_T$
are reserve counts and $N_S,N_T$ full counts, use the event
\[
 \min(H_S,H_T)\ge\rho L^q/2,\qquad
 |N_S-N_T|\le\rho L^q/4.
\]
Its complement has probability $C\exp(-c_\rho L^q)$. Delete the full
color excess from reserves using auxiliary indices, then move each
remaining iid uniform reserve list to its constant density. The expected
cost is $C_q\rho^{1-2/q}L^q$, and both constant densities are at least
$\rho/4$. After smoothing the ports at cost at most their number,
the Neumann-field construction transports the two residual densities
at cost at most $(4/\rho)\|a_L-b_L\|_{H_N^{-1}}^2$.
Equal integer atomic masses admit a permutation, so composing these
transports gives an actual residual matching. On the count exception,
an arbitrary residual matching costs at most $qL^2(N_S+N_T)$, whose
expected exceptional contribution is $o(L^q)$. This proves
\eqref{eq:general-path-reserve-bound}. No cost of the retained edges
is included in this additional squared-distance bound, and no selected
port cloud is assigned an independent Poisson law.
\end{proof}
\subsection{Full focal currents for interleaved conditional supports}
\label{subsec:interleaved-focal-current}

Fix one regular positive-mass ray on a finite conditional support span
\([\ell,u]\), with the genuine Hamiltonian Jacobi frame and conditional
laws \(\alpha,\beta\) of bounded physical densities. The matrix \(H\)
is positive definite in \((\ell,u)\) and extends as a bounded positive
semidefinite matrix to the closed span. The two conditional supports may
interleave. Choose any \(c\in(\ell,u)\), and write
\[
 \begin{gathered}
 T(r)=\int_c^rH(s)^{-1}\,ds,\qquad
 S(a,b)=T(b)-T(a),\qquad A(a,b)=S(a,b)^{-1},\\
 D=y-x,\qquad p=AD,\qquad Q=\tfrac12p^TSp,\qquad
 X_e(r)=x+[T(r)-T(a)]p\quad(a\le r\le b).
 \end{gathered}
\]
The spatial inputs are the original independent unit-intensity marked
Poisson processes, with laws \(\alpha,\beta\). Edge expectations below
are source-Palm expectations; expectations of fields refer to their
stationary spatial law. No sign is prescribed for \(T(a)\) or \(T(b)\).

\begin{theorem}[Full traffic for legal interleaved matchings]
\label{thm:interleaved-focal-traffic}
Let \(q>4\). Suppose a stationary perfect matching of these inputs is
forward, every actual edge satisfies \(b-a\ge\delta>0\), and
\(\E^S Q<\infty\). Assume that for a fixed \(D_0\ge0\),
\[
 Q_{ii}+Q_{jj}\le Q_{ij}+Q_{ji}+D_0
\]
whenever both crossed pairs of two matched edges are forward legal.
For every nonnegative \(\eta\in C_c^\infty(\R^q)\), the full absolute
traffic satisfies
\begin{equation}
\label{eq:interleaved-focal-traffic}
 \mathcal T_\eta(w)=\sum_e\int_{a_e}^{b_e}
 |H(r)^{-1}p_e|\eta(w-X_e(r))\,dr,
 \qquad
 \E[\mathcal T_\eta(0)^2]\le C(1+\E^S Q)<\infty.
\end{equation}
Constants may depend on this fixed root, \(\delta,D_0\), and \(\eta\).
If stationary edge indicators tend to zero on every edge almost surely,
their corresponding absolute traffic tends to zero in \(L^2\).
The full signed current
\[
 J_\eta(w)=\sum_e\int_{a_e}^{b_e}
 H(r)^{-1}p_e\eta(w-X_e(r))\,dr
\]
exists locally and obeys
\begin{equation}
\label{eq:interleaved-current-divergence}
 \operatorname{div}J_\eta=\eta*\mu_S-\eta*\mu_T.
\end{equation}
The theorem applies to actual weighted mean minimizers of \(Q+v g\),
for a bounded paired-mark probe, with \(D_0=4|v|\|g\|_\infty\).
\end{theorem}

\begin{proof}
We verify each changed use of common ordering in the proof of
Theorem~\ref{thm:focal-full-path-traffic}. Its short-window estimates
apply to incomplete stationary edge families, with counts bounded by
the original deterministic mark restrictions. We first recover their
global moment and envelope inputs without a common cut.

\emph{Momentum and focal energy.}
Put \(M=\sup_{[\ell,u]}\|H\|\). On every actual edge,
\begin{equation}
\label{eq:interleaved-momentum-bound}
 S(a,b)\ge\frac\delta M I,
 \qquad \E^S|p|^2\le\frac{2M}{\delta}\E^S Q.
\end{equation}
At a left focal endpoint of nullity \(k\ge1\), use its inward
coordinate \(s\). Lemma~\ref{lem:focal-simple-pole} and the regular
crossing give
\[
 H(\ell+s)^{-1}=G_L/s^2+R_L(s),\qquad
 \alpha(d(\ell+s))\le Cs^k\,ds,
\]
with \(G_L\ge0\) of rank \(k\) and \(R_L\) smooth through zero.
If \(a-\ell<\delta/4\), its actual target lies beyond
\(\ell+\delta\). Positivity of the integrated matrix therefore gives
\[
 S(a,b)\ge\int_a^{\ell+\delta/2}H(r)^{-1}\,dr
       =\frac{G_L}{a-\ell}-\frac{2G_L}{\delta}+B_L(a),
\]
where \(B_L\) is bounded on this collar. Thus
\(p^TG_Lp/(a-\ell)\le2Q+C|p|^2\) there; away from the collar
it is at most \(C|p|^2\). Time reversal gives the corresponding
right-target bound. With zero residues at nonfocal ends, set
\begin{equation}
\label{eq:interleaved-focal-energy}
 e_L=\frac{p^TG_Lp}{a-\ell},\qquad
 e_R=\frac{p^TG_Rp}{u-b},\qquad
 \mathcal E_*:=1+\E^S Q+\E^S|p|^2+\E^S(e_L+e_R)
                  \le C(1+\E^S Q).
\end{equation}
Perfection and the actual-edge gap also imply
\begin{equation}
\label{eq:interleaved-opposite-end-margins}
 \alpha\{a>u-\delta\}=0,\qquad
 \beta\{b<\ell+\delta\}=0.
\end{equation}
Otherwise such an input could have no forward partner in the span.

\emph{Signed determinant envelopes.}
Let \(\mathsf K=JT\) be the smooth second Jacobi solution of
Lemma~\ref{lem:curved-jacobi-second-solution}. Since \(T\) is negative
definite below \(c\) and positive definite above it,
\[
 \det(I+|T(r)|)|\det J(r)|=
 \begin{cases}
 |\det(J(r)-\mathsf K(r))|,&r<c,\\
 |\det(J(r)+\mathsf K(r))|,&r>c.
 \end{cases}
\]
The actual conditional volume factor, bounded physical densities and
smooth positive volume therefore give
\begin{equation}
\label{eq:interleaved-signed-determinant}
 \int\det(I+|T(r)|)\,d(\alpha+\beta)(r)<\infty.
\end{equation}
The determinant identity is used only at interior marks; endpoints
have zero conditional mass. In particular the first absolute
cumulative-mobility moments are finite.

For an edge active at a local cut \(r_0\), its positive legs are
\[
 U_{r_0}(a)=T(r_0)-T(a),\qquad
 V_{r_0}(b)=T(b)-T(r_0).
\]
For \(a<r_0\), in Loewner order,
\[
 I+U_{r_0}(a)
 \le(1+\|T(r_0)\|)(I+|T(a)|).
\]
Determinant monotonicity and
\eqref{eq:interleaved-signed-determinant} bound the original source
envelope \(\int_{a<r_0}\det(I+U_{r_0}(a))\,\alpha(da)\), uniformly
when \(r_0\) stays in a compact interior interval. The target
statement follows in the same way.

\emph{Legal short windows and their inherited estimates.}
If two edges are active throughout a window \(I\), then
\[
 a_i<\inf I\le\sup I<b_j,
 \qquad a_j<\inf I\le\sup I<b_i.
\]
Both crossed pairs are legal. Hence the midpoint argument from
Theorem~\ref{thm:focal-full-path-traffic} applies to this family.
For clarity, its precise local inputs at time scale \(s\), with
\(W(r)=H(r)^{-1}\), are
\[
 \begin{gathered}
 p'=\Lambda^{1/2}p,\qquad
 U'=\Lambda^{-1/2}U_{r_0}\Lambda^{-1/2}\ge h_1I,
 \quad V'=\Lambda^{-1/2}V_{r_0}\Lambda^{-1/2}\ge h_1I,\\
 s\|\Lambda^{-1/2}W(r)\Lambda^{-1/2}\|\le C\quad(r\in I),
 \qquad L\ge1,\quad L^{-1}\|\Lambda^{-1/2}\|\le1,\\
 (\det\Lambda)^{1/2}L^q
       \int_B\det(I+U'(a))\,\alpha(da)\le C,
 \end{gathered}
\]
where \(B\) is a deterministic original source-mark restriction.
Also require the physical one-edge bound
\(u_i=\int_I|W(r)p_i|\eta(-X_i(r))\,dr\le C\).
Additional adaptive edge restrictions are permitted.

The coarse phase coordinates are
\(\bar z=L^{-1}\Lambda^{-1/2}X(r_0)\) and \(\bar p=p'/L\).
The legal midpoint comparison gives the same measurable coarse
Lipschitz momentum extension, with defect \(D_0/L^2\). Each unit
phase cell pulls back to an ORIGINAL marked source set of mean at
most \(C\), by the last displayed envelope bound. The predecessor's
Poisson-animal events control all candidate random phase curves at
once. They do not assign a Poisson law to selected points.
Its physical-source Campbell calculation consequently gives
\begin{align}
 \E T_{\rm lg}^2
 &\le CK L^{-1}\int_I\E^S[|p'|\,|W(r)p|\mathbf1_B]\,dr,
 \label{eq:interleaved-window-large}\\
 \E T_{\rm sm}^2
 &\le CK\int_I\E^S[|W(r)p|\mathbf1_B]\,dr.
 \label{eq:interleaved-window-small}
\end{align}
If also \(U'\le CI\) and \(u_i\le C M_1|p_i'|\), the same
bounded-phase contribution obeys
\begin{equation}
\label{eq:interleaved-window-confined}
 \E T_{\rm sm}^2\le CK M_1^2(\det\Lambda)^{1/2}L^q
                                      \E^S[|p'|^2\mathbf1_B].
\end{equation}
The exceptional part has \(L^2\) norm at most \(Ce^{-cK}\).
These are precisely
\eqref{eq:focal-traffic-large}--\eqref{eq:focal-traffic-small-second},
with the legality and all full-input hypotheses now checked locally.
We next verify their scale conditions over every part of every path.

\emph{Old sources near the left focus.}
Translate \(\ell\) to zero and choose \(t_*<\delta/64\) so small
that all enlarged windows and their relative target comparison
intervals remain in the focal asymptotic neighborhood. The physical
one-edge estimate remains valid: on a short relative window,
\[
 \|W(r)W(r_0)^{-1}-I\|\le C|r-r_0|/r_0.
\]
For a sufficiently small relative width, the velocity stays in a
fixed cone around \(W(r_0)p\), so projection on that vector is
strictly increasing. The integral of speed while in a fixed physical
ball is uniformly bounded, as in
\eqref{eq:focal-traffic-one-edge}. This uses no mark ordering.

For a cut at time \(t\le t_*\), restrict to \(a<3t/4\). Every
actual target obeys \(b\ge a+\delta\ge\delta\), so the target
leg contains \([t,2t]\). Set
\[
 \Lambda_t=t^{-1}\Pi+t(I-\Pi),\qquad
 L_t=t^{-1/2-1/q},
\]
where \(\Pi\) projects onto the left residue range, of dimension
\(k\). Both normalized legs have uniform positive lower bounds.
The pole estimates give
\(\det(I+U'_t(a))\le C(t/a)^k\). Integrating the ORIGINAL
bound \(\alpha(da)\le Ca^k\,da\) gives phase-cell mean
\(C L_t^q t^{q/2+1}\le C\); the coarse kernel radius is at most
\(t^{1/q}\). The normalized coefficient condition also holds.

Put \(e(t)=\E^S[e_L\mathbf1_{a<3t/4}]\) and
\(v(t)=\E^S[|p|^2\mathbf1_{a<3t/4}]\).
The unchanged physical estimates are
\[
 \begin{split}
 \E T_{\rm lg}^2&\le CK
       \{t^{1/q}e(t)+t^{1+1/q}v(t)\},\\
 \E T_{\rm sm}^2&\le CK
       \{t^{k/2}\sqrt{e(t)}+t^{(k+3)/2}\sqrt{v(t)}\}.
 \end{split}
\]
Indeed \(\E^S[|\Pi p|^2\mathbf1_{a<3t/4}]\le Ct e(t)\),
the integrated speed is at most
\(C(|G_Lp|/t+t|p|)\), and the original marginal bounds are
\(\alpha\{a<3t/4\}\le Ct^{k+1}\) and
\(\int_{a<3t/4}a\,\alpha(da)\le Ct^{k+2}\).
Equations~\eqref{eq:interleaved-window-large} and
\eqref{eq:interleaved-window-small}, followed by Cauchy--Schwarz,
give the displayed estimates. Thus
\begin{equation}
\label{eq:interleaved-old-sum}
 \|T_{{\rm old},t}\|_2
 \le C\sqrt{K\mathcal E_*}\,t^{1/(2q)}+Ce^{-cK}.
\end{equation}
A fixed number of relative windows covers each dyadic time slab and
all portions \(a\le r/2\). Taking \(t=t_*2^{-n}\) and
\(K=K_0+C_1(n+1)\) makes both terms summable.

\emph{Births near the left focus.}
For the remaining portions \(r/2<a<r\), use source shells
\(t\le a<2t\) and elapsed-time bands \(s<r-a\le2s\), where
\(s=t2^{-j}\), \(j\ge0\). These bands cover every positive
elapsed time in question. Divide the source shell into at most
\(Ct/s\) bands \(B_i\) of width \(s/8\). A fixed number of
windows per band, of width \(\epsilon s\), has cuts satisfying
\(r_0-a\asymp s\) for every mark in that band and \(r_0\asymp t\).
The windows can be enlarged while remaining after all these births.
All actual targets still lie beyond \(\delta\). With
\[
 \Lambda_{t,s}=(s/t^2)\Pi+s(I-\Pi),\qquad
 L_{t,s}=s^{-1/2}t^{-1/q},
\]
the source normalized leg is bounded above and below. The target
leg contains a relative focal interval beyond the cut, whose
compliance dominates \(\Lambda_{t,s}\), since \(s\le t\).
The coefficient condition holds at time scale \(s\), and the
physical coarse kernel radius is at most \(t^{1/q}\). Moreover,
\[
 (\det\Lambda_{t,s})^{1/2}L_{t,s}^q=t^{-k-1},
 \qquad \alpha(B_i)\le Ct^k s.
\]
Thus the original phase-envelope mean is at most \(Cs/t\le C\).
Every two contributors active throughout these windows cross legally.

Write \(e_i=\E^S[e_L\mathbf1_{B_i}]\) and
\(v_i=\E^S[|p|^2\mathbf1_{B_i}]\); their sums are at most
\(\mathcal E_*\). The momentum and speed bounds are
\[
 |p'|^2=(s/t^2)|\Pi p|^2+s|(I-\Pi)p|^2,
 \quad \E^S[|\Pi p|^2\mathbf1_{B_i}]\le Ct e_i,
 \quad \int_I|W(r)p|\,dr\le C\sqrt{s}|p'|/t.
\]
Equations~\eqref{eq:interleaved-window-large}--
\eqref{eq:interleaved-window-confined} give, respectively,
\[
 \begin{split}
 \E T_{{\rm lg},i}^2
 &\le CK\{s^2t^{-2+1/q}e_i+s^2t^{-1+1/q}v_i\},\\
 \E T_{{\rm sm},i}^2
 &\le CK\{s^{3/2}t^{(k-3)/2}\sqrt{e_i}
                         +s^{3/2}t^{k/2}\sqrt{v_i}\},\\
 \E T_{{\rm sm},i}^2
 &\le CK\{s^2t^{-k-4}e_i+s^2t^{-k-3}v_i\}.
 \end{split}
\]
Summing norms over the bands gives the large-phase bound
\(C\sqrt{K\mathcal E_*}t^{1/(2q)}\sqrt{s/t}\).
For the same bounded-phase good-event part the two bounds are
\[
 C\sqrt{K\mathcal E_*}t^{k/4},\qquad
 C\sqrt{K\mathcal E_*}t^{-(k+2)/2}\sqrt{s/t}.
\]
The first uses \(\sum_i e_i^{1/4}\le
(Ct/s)^{3/4}(\sum_i e_i)^{1/4}\), and likewise for \(v_i\).
Interpolate their minimum with exponent \(k/[2(3k+4)]\).
With \(b_k=k/[4(3k+4)]>0\), the complete level bound is
\begin{equation}
\label{eq:interleaved-birth-sum}
 \begin{split}
 \|T_{{\rm birth},t,s}\|_2
 \le{}&C\sqrt{K\mathcal E_*}
       \{t^{1/(2q)}(s/t)^{1/2}+t^{k/8}(s/t)^{b_k}\}\\
      &+C(t/s)e^{-cK}.
 \end{split}
\end{equation}
For \(t=t_*2^{-n}\), \(s=t2^{-j}\), choose
\(K=K_0+C_2(n+j+1)\). The positive powers in both indices are
summable, as is the entire exceptional term for sufficiently large
\(C_2\). No independence of the selected bands or events is used.

\emph{Right endpoints and the compact remainder.}
Reverse time and exchange colors for right-focal deaths. The reversed
compliance derivative is positive, the original target PPP supplies
the envelopes, and target-Palm edge moments equal source-Palm moments.
The opposite-end margin in
\eqref{eq:interleaved-opposite-end-margins} supplies the required
partner separation. The two residue matrices need not commute.

In the compact regular region choose a small birth-distance cutoff
below \(\delta/16\), also small enough to keep the corresponding
windows away from the focal ends. At a source birth of scale \(s\),
actual partners satisfy \(b-r\ge\delta-Cs\ge\delta/2\).
Use bands of width \(s/8\), \(\Lambda=sI\), and \(L=s^{-1/2}\).
The source normalized leg is bounded above and below, the target
leg is bounded below, and \((\det\Lambda)^{1/2}L^q=1\).
Both the large-phase and confined small-phase squared estimates are
\(CKs^2\E^S[|p|^2\mathbf1_B]\). Summing over \(O(1/s)\) bands
gives \(C\sqrt{Ks}\sqrt{\E^S|p|^2}\), plus
\(Cs^{-1}e^{-cK}\). Take \(K\) proportional to
\(1+\log(1/s)\) and sum dyadic \(s\). Regular deaths are identical
after time reversal.

The remaining compact middle pieces stay a fixed positive time away
from both actual endpoints. Both local legs then have positive lower
bounds from \(W\ge M^{-1}I\); upper leg bounds are unnecessary.
The signed determinant estimate supplies a uniform source envelope
over their compact cuts. Additional target restrictions only reduce
the selected contributions. Finitely many windows cover these pieces.

Minkowski over this complete coverage and
\eqref{eq:interleaved-old-sum}--\eqref{eq:interleaved-birth-sum}
proves \eqref{eq:interleaved-focal-traffic} first for a nonnegative
unit-mass kernel supported in the unit ball. A finite collection of
translated positive bumps dominates any prescribed compact kernel and
its first derivatives. Stationarity and Minkowski give the asserted
kernel scope. Local absolute integrability justifies summing the
endpoint identity along complete paths, proving
\eqref{eq:interleaved-current-divergence}. Finally, dominated convergence
first in the countable edge sum and then under
\(\mathcal T_\eta(0)^2\) proves the adaptive-tail conclusion.
\end{proof}

\begin{corollary}[Centered prefixes and cut occupation]
\label{cor:interleaved-prefix-current}
Under Theorem~\ref{thm:interleaved-focal-traffic}, suppose the matching
is an actual weighted mean minimizer of \(Q+v g\) for a bounded
paired-mark probe. Its full signed current and each deterministic
time-restricted current
\[
 J_{\eta,I}(w)=\sum_e\int_{I\cap(a_e,b_e)}
 H(r)^{-1}p_e\eta(w-X_e(r))\,dr
\]
have zero mean conditional on the translation-invariant sigma field
\(\Xi\). For each fixed \(r\in(\ell,u)\), the cut measure
\(\nu_r=\sum_{a_e<r<b_e}\delta_{X_e(r)}\), counting multiplicities,
is locally finite with intensity
\(\alpha\{a<r\}-\beta\{b\le r\}\), and
\begin{equation}
\label{eq:interleaved-prefix-divergence}
 \operatorname{div}J_{\eta,(\ell,r)}
  =\eta*\mu_S^{a<r}-\eta*\mu_T^{b\le r}-\eta*\nu_r.
\end{equation}
Its local counts have finite second moments. No centering or independent
Poisson law is asserted for adaptively selected edge families or cut points.
\end{corollary}

\begin{proof}
Bounded deterministic markwise translations preserve the original
Poisson inputs, paired marks, forward legality, stationarity and the
probe. Since \(A\le(M/\delta)I\), the exact quadratic first
variation is integrable. Minimality for both signs gives
\[
 \E^S[p\mid a]=0,\qquad \E^T[p\mid b]=0.
\]
The same holds componentwise: the prescribed product Poisson law is
ergodic, so every stationary ergodic component retains this input
law and unit intensity. Almost every component is itself minimizing,
since its value is at least the global infimum and their average
attains it. Mass transport identifies the two directed Palm edge laws.

By \eqref{eq:interleaved-opposite-end-margins} and the simple pole,
\(T(a)=-G_L/(a-\ell)+O(1)\) on the source law and
\(T(b)=G_R/(u-b)+O(1)\) on the target law. The residues are zero
at nonfocal ends. For example,
\[
 \E^S\frac{|G_Lp|}{a-\ell}
 \le\left(\|G_L\|\E^S e_L
                   \int\frac{\alpha(da)}{a-\ell}\right)^{1/2}
 <\infty
\]
at a focal left end, by its density bound. The bounded remainders
are controlled by \eqref{eq:interleaved-momentum-bound}. Hence
\(\E^S|T(a)p|+\E^S|T(b)p|<\infty\), componentwise as well.
Bounded truncations in the conditional momentum identities give
\(\E^S[D\mid\Xi]=0\). Campbell therefore centers the full current.

For each fixed interior time, forward legality gives the exact identity
\[
 \mathbf1_{\{a<r<b\}}=
       \mathbf1_{\{a<r\}}-\mathbf1_{\{b\le r\}},
\]
outside a null endpoint event. The conditional momentum identities
therefore imply \(\E^S[p\mathbf1_{\{a<r<b\}}\mid\Xi]=0\).
The full absolute traffic bound justifies Campbell and integration
over any deterministic interval, yielding
\(\E[J_{\eta,I}(0)\mid\Xi]=0\).

Re-anchoring each active edge at \(X_e(r)\) in Campbell's formula
gives the stated intensity of \(\nu_r\); finite intensity implies
local finiteness almost surely. Integrating each path only up to time
\(r\) gives its source minus its target if completed, or its source
minus its cut point if still active. Summing gives
\eqref{eq:interleaved-prefix-divergence}. Its two endpoint terms are
smoothings of deterministic mark restrictions of the original PPPs,
and have finite second moments. The derivative-kernel traffic bound
controls the squared divergence term. Thus \(\eta*\nu_r(0)\) has
finite second moment. Taking a nonnegative bump bounded below on a
unit ball and using stationarity proves the local count assertion.
\end{proof}

These current estimates retain the full genuine interleaved conditional
laws and permit focal supported endpoints. Legal reserve completion and
stationary value identification are separate steps; no displacement
second moment is inferred here.

% Source: interleaved-screening-without-d2.md
% SHA256: a4d1d2f54d2102a6a4c09f8c5ea9996f665546de7938a914f1eaf93f964b5aef
% Full source proof: independently reviewed; conversion checked separately.
\section{Legal recovery and equality of stationary values}\label{sec:legal-recovery}

\subsection{Legal stagewise screening for interleaved populations}
\label{subsec:interleaved-path-screening}

The omitted paths supply nonnegative intermediate populations at the
successive mark cuts. Independent reserves keep the densities positive,
so a Neumann transport at each class can be composed into a legal
endpoint matching. The argument below assumes square-integrable full
absolute traffic and does not assume a displacement second moment.
The interleaved traffic theorem supplies that premise separately.

\begin{lemma}[A common orientation for boundary energy and flux]
\label{lem:generic-current-boundary}
Let $q>1$, and let $J^1,\ldots,J^s$ be finitely many stationary smooth
square-integrable vector fields on $\R^q$, each with zero invariant
conditional mean. There is a deterministic orthogonal matrix $O$ such
that, on $Q_L=O[-L/2,L/2]^q$, simultaneously for every $j$,
\begin{equation}
 \E\mathcal B_L(J^j)=o(L^q),\qquad
 \E\left|\int_{\partial Q_L}J^j\cdot n\right|^2
       =o(L^{2q-2}),
 \label{eq:interleaved-path-boundary}
\end{equation}
where
\[
 \mathcal B_L(J)=
 \sup_{\substack{\phi\in H^1(Q_L),\ \int_{Q_L}\phi=0\\
                         \int_{Q_L}|\nabla\phi|^2\le1}}
 \left|\int_{\partial Q_L}(J\cdot n)\phi\right|^2.
\]
No condition on the divergence or on mixing is required.
\end{lemma}
\begin{proof}
Let $\Sigma_j$ be the finite matrix spectral measure of $J^j$, and put
$\mu=\sum_j\operatorname{tr}\Sigma_j$. The zero-frequency atom vanishes:
its trace is the squared $L^2$ norm of the invariant conditional mean.
For almost every unit vector $n$,
\[
                    \mu(\R n\setminus\{0\})=0.
\]
Indeed, Tonelli applied to integration over sphere directions reduces
this assertion to the fact that a fixed nonzero frequency lies on the
normal line of only two sphere directions. Almost every orthogonal
matrix therefore has this property for all its finitely many axes.
Fix such an $O$. The scalar spectral measure of $J^j\cdot n$ is
dominated by $\operatorname{tr}\Sigma_j$.

On a face normal to an axis $n$, trace and Poincar\'e inequalities bound
the squared dual norm of the datum $e^{ik\cdot x}$ by $CL^q$.
After rescaling to the unit cube, that squared norm divided by $L^q$
equals the squared norm of the adjoint trace applied to
$e^{iLk_{\rm tan}\cdot x}$. If $k_{\rm tan}\ne0$, these functions
converge weakly to zero in the face's $L^2$. The trace from mean-zero
$H^1(Q_1)$, with the gradient norm, to the face's $L^2$ is compact,
so the adjoint-trace norms tend to zero.
Expand the random face functional in an orthonormal basis of this
Hilbert test space. Scalar spectral isometry followed by Tonelli
expresses its expected squared dual norm as the integral of the
plane-wave squared norms. The uniform bound and the absence of spectral
mass on the normal line give convergence by dominated convergence.
Sum the finitely many faces with the squared triangle inequality.
This proves the first assertion in
\eqref{eq:interleaved-path-boundary}.

For the second assertion, divide the integral on each face by its area
$L^{q-1}$. Its plane-wave multiplier, in the chosen coordinates, is
\[
 e^{\pm iLk_n/2}
       \prod_{a\ne n}\operatorname{sinc}(Lk_a/2),
 \qquad \operatorname{sinc}(u)=\frac{\sin u}{u},\quad
 \operatorname{sinc}(0)=1.
\]
The multiplier has modulus at most one and tends to zero off the
normal line. Spectral isometry, dominated convergence, and the
finite-face squared triangle inequality give the claimed flux bound.
\end{proof}

\begin{proposition}[Legal screening from full-path traffic]
\label{prop:interleaved-path-screening}
Let $q>4$. Let $C_1,\ldots,C_m$ be finitely many ordered monochromatic
mark classes, each of positive mass for its color, with gaps between
consecutive classes. Let $c_i$ lie in those gaps. For the prescribed
probability laws $\alpha,\beta$, suppose
\begin{equation}
 H_i=\alpha(( -\infty,c_i))-
              \beta(( -\infty,c_i))\ge\eta>0,
 \qquad 1\le i<m.
 \label{eq:interleaved-path-hall}
\end{equation}
The first class is a source class and the last is a target class.
Every source--target class rectangle is wholly legal or wholly illegal,
and legality means $a<b$.

Write
\[
 \begin{gathered}
 F(r)=H(r)^{-1},\qquad S(a,b)=\int_a^bF(r)\,dr,\\
 A(a,b)=S(a,b)^{-1},\qquad
 Q_A(D,a,b)=\tfrac12D^\top A(a,b)D.
 \end{gathered}
\]
Assume $Q_A(D,a,b)\le\Lambda|D|^2$ on legal rectangles.
On the original independent product-marked Poisson inputs of common
intensity $\lambda$, let $P$ be a stationary forward mean minimizer
of $Q_A+t g$, with finite weighted energy and bounded paired-mark
probe $g$. Its kinetic paths are
\[
 p_e=A(a_e,b_e)(y_e-x_e),\qquad
 X_e(r)=x_e+\int_{a_e}^rF(u)p_e\,du.
\]
Assume the first marginal cumulative-compliance moments are finite,
and the smoothed full absolute traffic belongs to $L^2$, also for
a compact kernel dominating the necessary kernel derivatives.
Every individual path is finite.

Fix $0<\rho<1$ and set $\lambda=1-\rho$.
There is a deterministic cube orientation, depending on $P$, such that
independent reserves of intensity $\rho$ and the respective original
mark laws give the following construction on $Q_L$.
Choose $\ell_L\to\infty$ with $\ell_L/L\to0$ and retain the main
edges whose sources lie in $Q_{L-4\ell_L}$ and whose complete path
lengths are at most $\ell_L$, as in
Proposition~\ref{prop:general-finite-path-screening}.
There is an event $G_L$ on which
these edges extend to a legal matching of the smaller full color count,
using the exact original unit-intensity inputs. The only put-aside
points are the total color excess, selected from the earliest source
or latest target reserve class. Moreover,
\begin{align}
 \E[\mathbf1_{G_L}\,\text{additional squared cost}]
     &\le C\rho^{1-2/q}L^q+o_{\rho,P}(L^q),
       \label{eq:interleaved-path-screening-cost}\\
 \E[L^2N_{\rm tot}\mathbf1_{G_L^c}]&=o_{\rho,P}(L^q),
       \label{eq:interleaved-path-screening-exception}
\end{align}
where $N_{\rm tot}$ is the full original input count. The leading
constant depends only on $q,m$, independently of $\rho$ and the chosen
minimizing law. No selected port population is
assigned an iid law.
\end{proposition}
\begin{proof}
\emph{Centering every prefix and slab.}
Bounded deterministic spatial shifts depending on one endpoint mark
preserve each original marked Poisson law, their independence, all
paired marks, and the condition $a<b$. The same endpoint first-variation
argument as in Proposition~\ref{prop:common-cut-conditional-centering}
therefore gives
\begin{equation}
 \E[p\mid a,\mathcal I]=0,\qquad
 \E[p\mid b,\mathcal I]=0,
 \label{eq:interleaved-path-momentum-centering}
\end{equation}
where $\mathcal I$ is the invariant sigma-field. Bounded shifts have
finite quadratic cost because $A$ is bounded. Almost every stationary
component has the same prescribed ergodic product Poisson input law;
its objective is at least the global infimum, so components of an
attained minimizer are themselves minimizing. This justifies the
componentwise statement. Palm Cauchy--Schwarz and the first compliance
moments give integrable path length, using
\[
 \operatorname{length}(X_e)
       \le\sqrt{\operatorname{tr}S(a_e,b_e)\,
                         p_e^\top S(a_e,b_e)p_e}.
\]

For a legal edge,
$\mathbf1_{\{a<r<b\}}=\mathbf1_{\{a<r\}}-
\mathbf1_{\{b\le r\}}$ away from endpoint times. Thus
\eqref{eq:interleaved-path-momentum-centering} implies
$\E[p\mathbf1_{\{a<r<b\}}\mid\mathcal I]=0$ for almost every $r$.
For a fixed smooth radial unit-mass kernel $\eta_0$ supported in the
unit ball, define
\[
 J_{\le i}(w)=\sum_e\int_{a_e}^{\min(b_e,c_i)}
                 F(r)p_e\eta_0(w-X_e(r))\,dr,
\]
with zero contribution if $a_e\ge c_i$.
The empty and full prefixes give the end conventions. Integration is
justified by integrable absolute traffic. Campbell identifies the
invariant mean of every $J_{\le i}$ as zero; the same holds for
$J_i=J_{\le i}-J_{\le i-1}$. Their absolute values are bounded by
full traffic, so they are square-integrable. This is the deterministic
prefix/slab centering in Corollary~\ref{cor:interleaved-prefix-current}.
Apply Lemma~\ref{lem:generic-current-boundary} simultaneously to all
prefixes and slabs and fix the resulting orientation.

\emph{Retained paths and the classwise discrepancy.}
Put $V=L^q$, and let $K_i$ be the retained contribution to $J_i$.
The length cutoff in
Proposition~\ref{prop:general-finite-path-screening} gives
\begin{equation}
 \E[\text{number of main ports}]=o(V),\qquad
 \E\int_{Q_L}|J_i-K_i|^2=o(V).
 \label{eq:interleaved-path-retention}
\end{equation}
Indeed, each slab's omitted current is bounded in the deep cube by
the same full long-path traffic $T_{>\ell_L}$. Applying
\eqref{eq:general-path-current-truncation} proves the current estimate;
the source and target port estimates follow from the same retained
set. Both conclusions use a single limit $L\to\infty$ and require
no path-length moment.

At each proper cut let the nonnegative omitted occupation be
\begin{equation}
 \nu_i(w)=\sum_{\substack{e\text{ unretained}\\a_e<c_i<b_e}}
                         \eta_0(w-X_e(c_i)),
 \qquad \nu_0=\nu_m=0.
 \label{eq:interleaved-path-occupation}
\end{equation}
Include cut points whose endpoints lie outside $Q_L$.
The full cut process has intensity $\lambda H_i$ by nonnegative
Campbell and the forward indicator identity. It is therefore locally
finite almost surely. Its compact-kernel smoothing, and hence
$\nu_i$, is a bounded smooth density on each finite cube. No Poisson
or square-integrability assertion about this cut process is needed.

Use the reflected kernel of
Proposition~\ref{prop:highd-signed-ports} to smooth the source and
target ports in class $i$ into $a_i,b_i$. It preserves mass and constant
density, agrees with the original kernel in the interior, and costs
at most one per endpoint. At most one of $a_i,b_i$ is nonzero.
Let $Z_i$ be the reflected-kernel discrepancy of the \emph{full}
original main input in this class, with its color sign. Its mean is
zero even for a single color, since reflection preserves constant
density. The full-input Poisson variance calculation in
\eqref{eq:highd-reflected-input-error} gives
\begin{equation}
             \E\|Z_i\|_{H_N^{-1}(Q_L)}^2\le CL^{q-1}.
 \label{eq:interleaved-path-reflection}
\end{equation}
Integrating the complete path segments in the class slab gives the
exact identity on $Q_L$
\begin{equation}
 \nu_{i-1}+a_i-\nu_i-b_i
           =\operatorname{div}(J_i-K_i)+Z_i.
 \label{eq:interleaved-path-stage-divergence}
\end{equation}
The retained currents have zero boundary trace for large $L$.
Integration by parts against mean-zero tests, the boundary lemma,
\eqref{eq:interleaved-path-retention}, and
\eqref{eq:interleaved-path-reflection} give
\begin{equation}
 \E\|\nu_{i-1}+a_i-\nu_i-b_i\|_{H_N^{-1}(Q_L)}^2=o(V).
 \label{eq:interleaved-path-stage-neumann}
\end{equation}

\emph{Cut-mass errors and the reserve event.}
Let $M_i^{\rm main}$ be the full main source prefix count minus target
prefix count in $Q_L$. Let $\bar\nu_i$ be the full cut occupation
before retained paths are subtracted, and put
\[
             \delta_i=M_i^{\rm main}-\int_{Q_L}\bar\nu_i.
\]
The full prefix version of
\eqref{eq:interleaved-path-stage-divergence} implies
\begin{equation}
 \delta_i=\int_{\partial Q_L}J_{\le i}\cdot n
                              +\int_{Q_L}Z_{\le i}.
 \label{eq:interleaved-path-cut-error}
\end{equation}
The last integral is a full-input Poisson integral with bounded weight
\[
 w_L(x)=\mathbf1_{Q_L}(x)-\int_{Q_L}\eta_0(w-x)\,dw.
\]
This weight has integral zero and is supported in a fixed-width boundary
layer. Its variance is $O(L^{q-1})$. The flux assertion in
\eqref{eq:interleaved-path-boundary} therefore gives
\begin{equation}
                  \E\delta_i^2=o(L^{2q-2}).
 \label{eq:interleaved-path-cut-variance}
\end{equation}

Add one independent reserve process per color with intensity $\rho$
and the original respective mark law. Their union with the main input
has exactly the prescribed unit-intensity product Poisson law.
Let $\Delta$ be the full source-minus-target count in $Q_L$. Put aside
$\Delta$ reserve sources from $C_1$ if $\Delta>0$, or $-\Delta$
reserve targets from $C_m$ if $\Delta<0$, using marks and auxiliary
indices only. Define
\[
 r=\frac{\rho}{16}
        \min\{\eta,\alpha(C_1),\beta(C_m),1\}>0.
\]
Except on an exponentially unlikely count event, the extreme reserves
are sufficient, every remaining proper reserve prefix is at least
$4rV$, and both remaining extreme class counts are at least $4rV$.
Indeed their respective means are at least $16rV$, and finitely many
Poisson count deviations, together with $|\Delta|\le rV$, suffice.
Write the remaining class counts as $h_i^S,h_i^T$ and their signed
prefixes as $H_i^{\rm res}$.
Intersect this count event with $|\delta_i|\le rV$ at every proper
cut, and call the result $G_L$. By
\eqref{eq:interleaved-path-cut-variance} and Chebyshev,
\[
                         \Pp(G_L^c)=o(L^{-2}).
\]
For a fixed $C$ larger than the total input intensity, split at
$N_{\rm tot}\le CV$ to obtain
\[
 \E[L^2N_{\rm tot}\mathbf1_{G_L^c}]
 \le CL^2V\Pp(G_L^c)
       +L^2\E[N_{\rm tot}\mathbf1_{\{N_{\rm tot}>CV\}}]
 =o(V).
\]
The second term has an exponential Poisson tail. This proves
\eqref{eq:interleaved-path-screening-exception} without requiring
independence of the count and the exceptional event.

\emph{Positive densities and stagewise gluing.}
Let $P_i$ be the source prefix port count minus target prefix port count.
On $G_L$ define
\begin{equation}
 \gamma_i=V^{-1}\left(P_i+H_i^{\rm res}-\int_{Q_L}\nu_i\right),
 \qquad\gamma_0=\gamma_m=0.
 \label{eq:interleaved-path-gamma}
\end{equation}
The final equality uses the deletion of precisely the full color excess.
Each retained path crossing $c_i$ contributes one both to its retained
endpoint prefix difference and to its smoothed cut mass: its entire
path and kernel are interior. These contributions cancel exactly,
giving
\begin{equation}
 \gamma_i=V^{-1}(H_i^{\rm res}+\delta_i)\ge3r,
                         \qquad 1\le i<m.
 \label{eq:interleaved-path-gamma-cancellation}
\end{equation}
At class $i$ use the two nonnegative densities
\begin{equation}
 f_i=\nu_{i-1}+\gamma_{i-1}+a_i+h_i^S/V,
 \qquad
 g_i=\nu_i+\gamma_i+b_i+h_i^T/V.
 \label{eq:interleaved-path-stage-densities}
\end{equation}
Their total masses agree by \eqref{eq:interleaved-path-gamma}.
At interior stages both densities are at least $3r$ by
\eqref{eq:interleaved-path-gamma-cancellation}; at the first starting
side use $h_1^S/V\ge4r$, and at the last ending side use
$h_m^T/V\ge4r$. Constant terms vanish on mean-zero tests. Linear
interpolation with a Neumann flux therefore yields
\begin{equation}
 W_2^2(f_i\,dw,g_i\,dw)
 \le(3r)^{-1}
             \|\nu_{i-1}+a_i-\nu_i-b_i\|_{H_N^{-1}(Q_L)}^2.
 \label{eq:interleaved-path-stage-cost}
\end{equation}
This is the same positive-density interpolation as in
Corollary~\ref{cor:highd-kinetic-screening}. Its expected sum over
the finitely many stages, restricted to $G_L$, is $o_{\rho,P}(V)$.

First move each reserve class cloud to its constant density.
Conditional on marks, counts, and main data, its selected positions
are iid uniform, because the excess selection used no positions.
The sharp iid quadratic estimate and concavity give
\begin{equation}
 C L^2\sum_i\E[(h_i^S)^{1-2/q}+(h_i^T)^{1-2/q}]
                      \le C\rho^{1-2/q}V.
 \label{eq:interleaved-path-reserve-spreading}
\end{equation}
For this unconditional bound extend the count-based deletion
arbitrarily on its count-failure event; the actual construction is
used only on $G_L$. Restriction to $G_L$ only decreases a nonnegative
cost. Smoothing the ports costs at most their number, hence $o(V)$
in expectation.

Compose the stage couplings through the common intermediate measures
$\nu_i+\gamma_i$. At a source class, new source mass joins the incoming
intermediate mass and moves to the next intermediate state. At a target
class, part of the incoming mass terminates in the target population
and the rest continues. Disintegration gives finite Markov gluing.
The first incoming and last outgoing intermediate masses are zero, so
every resulting trajectory begins at an original source and ends at an
original target in a strictly later class. No intermediate mass is
created without an original source ancestor. Thus every endpoint pair
is legal.

Include the atomic smoothing and reserve-spreading couplings at the
two ends. Each trajectory has at most $m+2$ spatial stages, so its
squared endpoint distance is at most $m+2$ times the sum of their squared
lengths. The resulting finite atomic coupling has integer endpoint
masses and uses only legal bipartite edges. The corresponding
transportation polytope is integral, so a minimum-cost integral
matching costs no more than this fractional coupling. Adding back the
retained main edges preserves all capacities and all retained pairs.
All put-aside points have the excess color, and every other point is
matched. This proves \eqref{eq:interleaved-path-screening-cost}.

\emph{Exceptional configurations.}
On $G_L^c$, an arbitrary maximum-forward matching of the full finite
input costs at most $CL^2N_{\rm tot}$. Its entire expected cost is
$o(V)$ by \eqref{eq:interleaved-path-screening-exception}; it need not
retain the main edges. Finite Hall failures can prevent a legal matching
from reaching the smaller full color count on this event. Any additional
unmatched vertices can be charged at bounded multiples of $L^2$ each
at the same $o(V)$ cost. Equivalently, a full maximum-color-cardinality
comparison can pair its unmatched endpoints arbitrarily on this event
and pay a backward penalty $M_L\le CL^2$. This does not assert a
forward perfect matching on every finite configuration.
\end{proof}

\begin{remark}[The fixed-root value interface]
\label{rem:interleaved-path-value-interface}
For a homogeneous comparison, the orientation in
Lemma~\ref{lem:generic-current-boundary} causes no change of stationary
value: rotating space preserves the independent Poisson input laws and
conjugates all coefficient matrices, bijecting both stationary classes.
The finite-box theorem can be applied to the conjugated coefficients
on the standard cube. No measurable orientation across root labels is
required for this fixed-root scalar comparison.

Let $I_{w,+}$ denote the finite-weighted-energy stationary forward
infimum. Start from an attained unit-intensity minimizer at parameter
$t\lambda^{2/q}$, and dilate space by $\lambda^{-1/q}$.
Its objective density becomes
$\lambda^{1-2/q}I_{w,+}(t\lambda^{2/q})$.
Fixed dilation preserves the traffic premise after rescaling the compact
kernel. Every added legal kinetic cost is at most $\Lambda$ times its
squared cost. Removing main edges costs only their possible negative
probe part, at most $|t|\|g\|_\infty$ per omitted source. New probe
terms cost at most $|t|\|g\|_\infty$ times the reserve and port counts.
Thus, for the finite comparator having the charged-exception convention
above,
\begin{equation}
 \limsup_{L\to\infty}\frac{\E[\text{finite value}]}{L^q}
 \le \lambda^{1-2/q}I_{w,+}(t\lambda^{2/q})
          +C\rho^{1-2/q}+C|t|\|g\|_\infty\rho.
 \label{eq:interleaved-path-value-comparison}
\end{equation}
Here the leading $C$ depends only on $q,m,\Lambda$, independently of
$\rho$ and the chosen minimizing law. Density-floor reciprocals, current
norms, and convergence rates occur only in the fixed-$\rho$ remainder.
The volume limit is taken first for each fixed $\rho$, stationary law,
and orientation; only then is $\rho\downarrow0$ taken.
Bounded probes give the required Lipschitz continuity in $t$.
There is no uniform traffic-tail or ergodic-rate requirement.

Equation~\eqref{eq:interleaved-path-value-comparison} supplies the
legal screening input in the form of
\eqref{strict-screen-poisson-upper}. Identification of that finite
comparison with the finite-second-moment forward stationary value,
including its exact maximum-cardinality and penalty conventions, is
a separate finite-pressure interface. Likewise, the interleaved
whole-path traffic assertion remains the stated premise of
Proposition~\ref{prop:interleaved-path-screening}; it is supplied by
Theorem~\ref{thm:interleaved-focal-traffic}, not by the screening
argument. Tight Hall cuts are handled by the existing balanced
active-block decomposition. The proof here permits root-dependent
positive slack and does not add strict Hall inequalities to the final
geometric target.
\end{remark}

% Source: general-finsler-selector-assembly.md
% Source SHA256: a3fc8f8b254a01e38713f7daf567cab5f8afba928701a762e85a66f41dcb76b0
% Complete unified source review: general-finsler-selector-assembly-review.md, PASS.
% Review SHA256: a468859ca28e5bc6e4316bb4272da9a3e61999bd863d99cdf4929af69514a1c2
\subsection{General Finsler distance selection}
\label{subsec:general-finsler-selector}

The legal interleaved traffic and screening results identify the two
stationary pressures on every actual active population block. This
removes the remaining relaxation premise from the general geometric
pressure bracket. We first give the exact finite-value comparison,
including tight cuts, and then assemble the full empirical selector.

\begin{theorem}[Weighted value recovery for interleaved geometric roots]
\label{thm:general-focal-value-equality}
Let \(q>4\), and fix a regular positive-mass conditional ray, or a
normalized positive active block, supplied by
Proposition~\ref{prop:general-optimal-face} and the genuine first-jet
geometry of Proposition~\ref{prop:rough-ray-quadratic}.
The conditional laws come from bounded physical densities relative
to smooth positive volume, and the opposite-color supports have a
positive scalar separation. Supported endpoints may focus, and the
conditional populations may interleave. For every continuous bounded
original paired-mark probe \(g\) and \(t\in\R\),
\begin{equation}
 I_{\mathrm w}(t)=I_2(t)=p_A(t),\qquad
 A(a,b)=\left[\int_a^b H(r)^{-1}\,dr\right]^{-1},
 \label{eq:interleaved-weighted-value-equality}
\end{equation}
on the original independent marked Poisson inputs.
Here \(I_{\mathrm w},I_2\) are the classes in
Definition~\ref{general-gluing-two-classes}, and \(p_A\) is the legal
maximum-cardinality finite pressure of
Theorem~\ref{degenerate-finite-pressure}.
The conclusion is an equality of values and does not assert that a
weighted minimizer has a finite unweighted displacement second moment.
\end{theorem}

\begin{proof}
\emph{Finite classes and strict blocks.}
By Propositions~\ref{prop:ray-finite-bands}
and~\ref{prop:ray-forward-criterion}, the compact conditional supports
have finitely many monochromatic order-convex classes. Each legal class
rectangle is complete bipartite. If
\[
 h(r)=\alpha((-\infty,r])-\beta((-\infty,r]),
\]
then \(h\ge0\), both marginals vanish on \(\{h=0\}\), and
the finitely many positive components have equal positive marginal
masses. Every forward coupling stays within these components.

First normalize one such component. Delete its zero-mass classes.
Choose a cut \(c_i\) in each proper inter-class gap. The normalized
surplus at each cut is strictly positive, so their finite minimum
is a number \(\eta>0\). The first class is a source class and the
last a target class. These properties are consequences of the active
decomposition, not assumptions on the original populations.

There is no omitted within-class Hall constraint. The surplus is
nondecreasing within a source class, nonincreasing within a target
class, and constant in each empty gap. Thus its signs at the adjacent
inter-class cuts control its sign everywhere. Equivalently, every
mark in one class has exactly the same allowed opposite-color
classes. Density holes inside a class do not change that graph.

\emph{The actual matrix and endpoint laws.}
Suppress the fixed root in the notation. On its finite support span
\([\ell,u]\), the matrix \(H\) is continuous and positive definite
in the interior and extends as a bounded positive semidefinite
matrix at either endpoint. If \(H\le M\mathrm{Id}\) and
the supported opposite-color gap is \(\delta>0\), then
\begin{equation}
 S(a,b)\succeq(\delta/M)\mathrm{Id},\qquad
 0\prec A(a,b)\preceq(M/\delta)\mathrm{Id}.
 \label{eq:interleaved-root-matrix-bound}
\end{equation}
The matrices are continuous and uniformly positive definite on
each compact interior collar. Supported endpoints themselves have
zero conditional mass.

Proposition~\ref{prop:curved-focal-determinant-moments} identifies
the actual first variation with the smooth Hamiltonian system of
Lemma~\ref{lem:curved-jacobi-second-solution},
\[
 H=J^\top\mathsf C^{-1}J,\qquad
 J'=\mathsf BJ+\mathsf CP,\qquad
 P'=-\mathsf AJ-\mathsf B^\top P,
 \qquad J^\top P=P^\top J,
\]
where \(\mathsf C=\mathsf C^\top>0\) and
\(\mathsf A=\mathsf A^\top\). Its phase stack has full rank and
its base \(J\) is invertible at interior times. The actual conditional
densities, after restriction and normalization, are
\begin{equation}
 \begin{split}
 \alpha(dr)&=\frac{\mathbf1_{\rm block}(r)
      \rho_A(X(r))w(r)|\det J(r)|}{\vartheta}\,dr,\\
 \beta(dr)&=\frac{\mathbf1_{\rm block}(r)
      \rho_B(X(r))w(r)|\det J(r)|}{\vartheta}\,dr,
              \qquad 0<\vartheta<\infty,
 \end{split}
 \label{eq:interleaved-root-density}
\end{equation}
with \(w\) smooth positive. They are therefore bounded above by a
root-dependent multiple of \(|\det J|\). Restriction changes only
the finite positive normalizer and introduces no density lower bound.

At a singular supported end,
Lemma~\ref{lem:regular-focal-crossing},
Proposition~\ref{prop:regular-focal-crossing-moments}, and
Lemma~\ref{lem:focal-simple-pole} give a fixed positive semidefinite
residue of rank \(k\ge1\), with
\[
 H(\ell+r)^{-1}=G_L/r^2+R_L(r),\qquad
 \frac{d\alpha}{dr},\frac{d\beta}{dr}\le C_{\rm root}r^k
\]
near that end, where \(R_L\) is smooth and bounded. The right end
has the corresponding expansion; a nonfocal end has bounded inverse
and zero residue. If an active block ends inside a larger regular
ray, that cut is simply a nonfocal end. No commuting-matrix or
endpoint-rank hypothesis is imposed.

For any interior reference \(c\), write
\(T(r)=\int_c^r H^{-1}\). The non-common-cut conclusion of
Proposition~\ref{prop:curved-focal-determinant-moments} gives
\begin{equation}
 \int\det(I+|T(r)|)\,d(\alpha+\beta)(r)<\infty.
 \label{eq:interleaved-root-moment}
\end{equation}
It follows by using the determinant of \(J-JT\) to the left of
\(c\) and that of \(J+JT\) to the right. In particular the first
absolute cumulative-compliance moments are finite. These statements
use the full original conditional laws, including their supported
endpoint collars.

\emph{Minimizing currents and legal screening.}
Theorem~\ref{weighted-capped-c1} supplies an attained weighted
integer minimum of \(Q+tg\). Its scalar mobility is \(H^{-1}\);
\eqref{eq:interleaved-root-matrix-bound}, interior continuity, the
finite separated class graph, and Hall feasibility verify that
theorem's hypotheses. The finite-second-moment class is nonempty
by the finite-class construction in
\eqref{general-gluing-uniform-pressure-bounds}.

At a weighted minimum, two-edge comparisons with legal crossed
pairs have pure-quadratic defect at most \(4|t|\|g\|_\infty\).
Every pair of edges active throughout a common short window has
both crossed pairs forward. Thus
Theorem~\ref{thm:interleaved-focal-traffic} applies with its original
deterministic Poisson envelopes; the selected window families are
not assigned their own independent input law. With
\[
 p=A(a,b)(y-x),\qquad
 X_e(r)=x+[T(r)-T(a)]p,
\]
it gives full absolute traffic in \(L^2\), derivative-kernel
majorants, and vanishing adaptive traffic tails. Every individual
path has finite length and diameter because both of its marks are
interior.

The same minimizer has the conditional momentum identities
\[
 \E[p\mid a,\mathcal I]=\E[p\mid b,\mathcal I]=0
\]
on each invariant component \(\mathcal I\). Bounded markwise
translations preserve the exact independent input law, marks, and
legality; bounded \(A\) and finite \(\E|p|^2\) justify their
first variations. Each component retains the ergodic prescribed
Poisson marginal and itself minimizes.
Moreover, by \eqref{eq:interleaved-root-moment} and mass transport,
\[
 \E\operatorname{tr}S(a,b)<\infty,\qquad
 \operatorname{length}(X_e)
 \le\sqrt{\operatorname{tr}S(a,b)\,p^\top S(a,b)p},
 \qquad \E\operatorname{length}(X_e)<\infty.
\]
The identity
\(\mathbf1_{\{a<r<b\}}=\mathbf1_{\{a<r\}}-
\mathbf1_{\{b\le r\}}\) and Campbell's formula consequently
center every deterministic prefix and slab current, with the exact
original-input divergences. These are precisely the conclusions of
Corollary~\ref{cor:interleaved-prefix-current}. No adaptive
omitted-edge current is asserted centered.

\emph{The exact finite comparator.}
Let \(H_L(t)\) minimize \(Q+tg\) among legal maximum-cardinality
partial matchings of the original restricted unit-intensity inputs
in a cube of side \(L\). Theorem~\ref{degenerate-finite-pressure}
gives
\[
 p_A(t)=\lim_{L\to\infty}L^{-q}\E H_L(t).
\]
The coefficient is the genuine scalar cumulative coefficient with
null endpoint degeneration of Theorem~\ref{degenerate-scalar-c1}.
Hence Corollary~\ref{degenerate-finite-readout} identifies this
same finite value with \(I_2(t)\). In particular the comparator
allows unequal finite color counts and Hall deficits; it is not a
finite perfect-matching problem or a fixed-unmatched-penalty value.

Fix \(0<\rho<1\), put \(\lambda=1-\rho\), and take a unit-input
weighted minimum at parameter \(t\lambda^{2/q}\). Dilate its full
joint scene by \(\lambda^{-1/q}\). Its input intensity is
\(\lambda\) per color and its mean objective density at parameter
\(t\) is
\begin{equation}
                  \lambda^{1-2/q}I_{\mathrm w}(t\lambda^{2/q}).
 \label{eq:interleaved-root-dilation}
\end{equation}
The traffic conclusions hold for this fixed dilation because they
hold for each fixed rescaled smoothing kernel. Add independent
reserves of intensity \(\rho\) with the same original mark laws.
Their union with the main process is exactly the prescribed
independent unit-intensity marked pair.

Proposition~\ref{prop:interleaved-path-screening} applies on this
strict block. Its deterministic generic cube orientation may depend
on the chosen stationary law. Rotating the spatial coordinates
conjugates \(H,A\), preserves their cumulative representation, and
bijectionally preserves both stationary competitor classes and the
independent Poisson law. The preceding finite-pressure identity
therefore applies in that orientation with unchanged scalar values.
No measurable orientation over root labels is needed.

For clarity, the screening comparison has the exact finite
cardinality required here. Write \(V=L^q\). Retain precisely the
main edges with source in \(Q_{L-4\ell_L}\) and path length at most
\(\ell_L\); their endpoint ports have expected number \(o(V)\).
The omitted paths' nonnegative cut occupations
and the centered prefix currents give
\[
 \E|\delta_i|^2=o(L^{2q-2}),
\]
where \(\delta_i\) is the full original prefix-count discrepancy
against the full smoothed cut occupation.
Only the total full color excess is put aside from the first source
reserve class or the last target reserve class. The auxiliary
selection uses counts and marks, preserving iid uniform reserve
positions. On the screening good event the intermediate constants
are exactly
\[
 \gamma_i=\frac{\text{remaining reserve prefix}_i+\delta_i}{V}
                 \ge3r,
 \qquad r=r(\rho)>0.
\]
These identities include omitted paths whose endpoints are outside
the cube. The positive densities at every class have equal masses;
stagewise Neumann transport and integrality of the legal bipartite
transportation polytope match all residual mass to strictly later
target classes. With the retained edges this is a legal matching
of cardinality \(\min(N_{\rm S},N_{\rm T})\), and hence is
maximum-cardinality. Density holes within a class create no further
legality constraint.

The exception satisfies
\begin{equation}
 \Pp(G_L^c)=o(L^{-2}),\qquad
 \E[L^2N_{\rm tot}\mathbf1_{G_L^c}]=o(V).
 \label{eq:interleaved-root-exception}
\end{equation}
The second estimate follows by truncating to \(N_{\rm tot}\le CV\)
and using the exponential original Poisson tail on the complement.
On the exception choose any legal maximum-cardinality matching.
Its kinetic cost is at most \(CL^2N_{\rm tot}\); its absolute probe
cost is at most \(|t|\|g\|_\infty N_{\rm tot}\). Thus it is an
admissible comparator for \(H_L(t)\) on every configuration and
costs \(o(V)\) in expectation, even when Hall deficits prevent
attaining the smaller color count.

No unmatched penalty is needed. In a formulation charging at most
\(CL^2\) per unmatched point, the good-event total excess has mean
\(O(\sqrt V)\) and normalized charge \(O(L^{2-q/2})=o(1)\);
the exception is already covered by
\eqref{eq:interleaved-root-exception}. Neither an infinite-penalty
limit nor a simultaneous limit is being used here.

The expected additional good-event squared-distance cost is
\(C\rho^{1-2/q}V+o_\rho(V)\). The coefficient upper bound pays
its kinetic cost, and new probe terms are charged to
\(O(\rho V)+o(V)\) reserve and port counts.
For signed probes, the retained main cost is at most the full
source-anchored main cost in the cube plus
\(|t|\|g\|_\infty\) times its source-port count: omitted kinetic
cost is nonnegative. Discarding this comparison on the bad event
costs at most its negative part, bounded by
\(|t|\|g\|_\infty N_{\rm main,S}\mathbf1_{G_L^c}\), also \(o(V)\).
No uniform integrability of a selected large kinetic cost is needed.
It follows that
\begin{equation}
 I_2(t)\le
 \lambda^{1-2/q}I_{\mathrm w}(t\lambda^{2/q})
       +C\rho^{1-2/q}+C|t|\|g\|_\infty\rho.
 \label{eq:interleaved-root-recovery}
\end{equation}
The displayed \(C\) is independent of \(\rho\) and of the chosen
minimum: it depends only on \(q\), the fixed class count and the
coefficient upper bound, through reserve spreading and the bounded
number of stages. The law, slack, orientation, current norms and
\(\rho\) may affect the discarded fixed-\(\rho\) \(o(V)\)
remainders and their thresholds. Every inverse density-floor factor
multiplies such an \(o(V)\) term and has disappeared before
\eqref{eq:interleaved-root-recovery}.

First let volume tend to infinity, then let \(\rho\downarrow0\).
The common \(\|g\|_\infty\)-Lipschitz bound in the probe parameter
gives \(I_2(t)\le I_{\mathrm w}(t)\). Class inclusion gives the
opposite inequality, proving
\eqref{eq:interleaved-weighted-value-equality} on a strict block.

\emph{Tight cuts and the full original law.}
At a tight scalar cut, a stationary perfect forward matching has
source-Palm crossing probability
\[
 \Pp^{\rm S}(a\le c<b)
       =\alpha((-\infty,c])-\beta((-\infty,c])=0.
\]
Perfection and mass transport supply the target marginal in this
identity. Thus all edges stay in their finitely many balanced
blocks. Original marked Poisson restrictions to the deterministic
block mark sets are independent, with equal color intensities
\(w_b\). A restricted matching may depend on the other blocks,
which are allowed stationary auxiliary randomness. Conversely,
independent block schemes superpose to the exact original joint
law. Restriction and near-minimizing superposition give, for
\(\star\in\{\mathrm w,2\}\),
\begin{equation}
 I_\star^{\rm whole}(t)
       =\sum_b w_b^{1-2/q}I_\star^b(tw_b^{2/q}).
 \label{eq:interleaved-root-block-scaling}
\end{equation}
Zero-mass blocks are omitted, and finitely many finite-moment schemes
remain finite-moment. Blockwise equality proves the whole-ray equality.

Finite samples need not balance these blocks. The already used
Corollary~\ref{degenerate-finite-readout} handles their count defects
with \eqref{degenerate-finite-deficiency} and
\eqref{degenerate-finite-augmentation}. Their expected unmatched
counts are \(O(\sqrt V)\); each bounded-class augmentation costs
at most \(C(L^2+|t|\|g\|_\infty)\). The normalized error
\(O(L^{2-q/2})+O(|t|L^{-q/2})\) vanishes because \(q>4\).
No predetermined empirical mass balance has been introduced.
\end{proof}

\section{Pressure readout and the general selector}\label{sec:finsler-extension}

\subsection{Pressure convergence and derivative readout}

The proved general-geometric bracket is
\begin{equation}
\label{general-gluing-pressure-bracket}
 \mathcal P_{\mathrm w}(t)
 \le \text{every subsequential pathwise value lower limit},
 \qquad
 \limsup_n\E V_n(t)\le\mathcal P_2(t).
\end{equation}
The local equality needed for readout is
\begin{equation}
\label{general-gluing-relaxation-condition}
 I_{j,\mathrm w}(z,v)=I_{j,2}(z,v)
       \quad\text{for almost every active root}.
\end{equation}
By Lemma~\ref{general-gluing-clipping-limit}, this is exactly convergence
of clipping pressure to the weighted stationary infimum. It asserts
approximation of that value by finite-unweighted-second-moment matchings;
pointwise positive definiteness and a uniform upper matrix bound alone
have not established it.

\begin{proposition}[Pressure convergence and whole-optimal-set readout]
\label{general-gluing-conditional-readout}
If \eqref{general-gluing-relaxation-condition} holds at the needed probe
parameters, then
\[
 V_n(t)\longrightarrow\mathcal P_{\mathrm w}(t)=\mathcal P_2(t)
               \quad\text{in probability and in }L^1.
\]
If the common pressure \(\mathcal P_f\) is differentiable at zero and
this convergence holds at neighboring finite-difference parameters, then
\begin{equation}
\label{general-gluing-all-optimal-readout}
 \sup_{\gamma\in\operatorname{Opt}_n}
     \left|\int f\,d\gamma-\mathcal P_f'(0)\right|
               \longrightarrow0\quad\text{in probability},
\end{equation}
where \(\operatorname{Opt}_n\) is the entire original length-optimal set.
\end{proposition}

\begin{proof}
The probability lower bound and deterministic bound
\(V_n\ge-|t|\|f\|_\infty\) make the negative parts relative to the
common limit converge to zero in mean. The expectation upper bound then
forces the positive parts to converge to zero in mean. For every
unperturbed optimum and \(h>0\), optimality gives
\[
 \frac{V_n(h)-V_n(0)}h
 \le\int f\,d\gamma
 \le\frac{V_n(0)-V_n(-h)}h.
\]
Pass to the pressure limit for fixed \(h\), then let \(h\downarrow0\).
This bounds all optima at once, including fractional ones.
\end{proof}

Corollaries~\ref{degenerate-scalar-integrated-c1}
and~\ref{weighted-capped-integrated-c1} supply differentiability of
both atlas pressures without an endpoint nondegeneracy condition.
Corollary~\ref{degenerate-finite-readout} also identifies the homogeneous
finite-box value with the second-moment value. Theorem~\ref{weighted-capped-c1}
identifies bounded-cost limits with the attained weighted value.
Theorem~\ref{thm:general-focal-value-equality} establishes equality of
these values for the actual geometric populations.
\begin{proof}[Proof of Theorem~\ref{thm:general-finsler-selector}]
Put \(q=d-1>4\), \(s=1-2/q\),
\(\delta=\operatorname{dist}_F(A,B)>0\), and
\(D=\operatorname{diam}_F(A\cup B)<\infty\).
Fix a population optimal \(1\)-Lipschitz potential \(\phi\) with
transport increasing \(\phi\), and write
\[
                     e(x,y)=d_F(x,y)-\phi(y)+\phi(x)\ge0.
\]

\emph{The whole continuous face and actual local geometry.}
Propositions~\ref{prop:general-regular-pairs}
and~\ref{prop:general-optimal-face} give a common full-marginal-measure
Borel interior-ray carrier, a standard Borel quotient with common
measure \(\Lambda\), absolutely continuous scalar conditional laws
\(\alpha_z,\beta_z\), and injective coordinates \(\gamma_z(r)\).
The entire continuous optimal face is exactly
\[
 \pi=\int(\gamma_z\times\gamma_z)_\#\kappa_z\,\Lambda(dz),
 \qquad
 \kappa_z\in\Pi(\alpha_z,\beta_z),\qquad
 \kappa_z\{a<b\}=1.
\]
The equality of quotient measures follows from optimality, and
\(\delta\le b-a\le D\) on supported forward pairs. It imposes no
equality on a previously chosen spatial projection.

Split the finite order classes into their positive-flux blocks using
Propositions~\ref{prop:ray-finite-bands}
and~\ref{prop:ray-forward-criterion}. Theorem~
\ref{thm:intrinsic-active-invariance} makes the measured active arcs
intrinsic up to null sets and Borel relabeling. This yields strict
finite Hall blocks without a final Hall-slack hypothesis.

Refine Proposition~\ref{prop:general-upper-atlas} by these finitely
many block indices. In the notation of
\eqref{general-gluing-sheet-disintegration}, the finite sheets have
bounded coordinate cubes \(Q_j\), bounded measurable densities
\(\theta_j\), and full conditional laws
\[
 \mu_j=\theta_j(z)\,dz\,\alpha_{j,z}(da),\qquad
 \nu_j=\theta_j(z)\,dz\,\beta_{j,z}(db),\qquad
 x=X_j(a,z),\quad y=X_j(b,z).
\]
The kernels need only be measurable in \(z\).

At almost every positive root,
Proposition~\ref{prop:rough-ray-quadratic} and
Lemma~\ref{general-gluing-compatible-cores} identify the actual
quadratic coefficient by \eqref{general-gluing-local-compliance}
and~\eqref{general-gluing-full-matrix}. Its \(H\) is positive
definite at interior times and has a bounded upper matrix on the
finite atlas. The actual global distance is smooth near the
retained calibrated interior pairs: each such minimizing segment
extends past its endpoints and lies before its cut and conjugate
times, as proved in Proposition~\ref{prop:general-regular-pairs}.
Completeness and compactness exclude competing minimizing branches
near a fixed compact retained pair set.
Thus \eqref{general-gluing-relative-expansion} concerns \(d_F\)
itself, without any global cut-locus exclusion.

The smooth first variation extends through the finite closed
conditional support span. Proposition~
\ref{prop:curved-focal-determinant-moments} identifies its Hamiltonian
Jacobi system and the volume density in
\eqref{eq:interleaved-root-density}; its non-common-cut determinant
identity gives \eqref{eq:interleaved-root-moment}. Supported singular
ends are regular focal crossings, and an interior active-block cut
is nonfocal. These facts verify all geometric and conditional-law
hypotheses of Theorem~\ref{thm:general-focal-value-equality}
at almost every positive root. Arbitrary density zeros, disconnected
supports and both possible focal ends are permitted. Constants may
depend on the root; no such constants will be integrated.

\emph{The common measurable pressure.}
Fix \(f\in C(A\times B)\) and use its bounded continuous extension
to the compact whole-flow endpoint product when needed in clipping.
It preserves its supremum norm and changes no sampled paired value.
Use the original iid pressure \(V_n(t)\) defined in
\eqref{eq:general-finsler-pressure}.

Use the original continuous root probe
\[
 g_{j,z}(a,b)=f(X_j(a,z),X_j(b,z)),\qquad
                         \|g_{j,z}\|_\infty\le\|f\|_\infty.
\]
Theorem~\ref{thm:general-focal-value-equality} gives
\(I_{j,\mathrm w}(z,v)=I_{j,2}(z,v)\) for all real \(v\) at
almost every positive root. If necessary, intersect the conull
sets for rational \(v\); the common Lipschitz bound extends equality
to every real parameter on one conull set. In particular it holds
at \(v=t\theta_j(z)^{2/q}\).

The universally measurable versions following
Definition~\ref{general-gluing-two-classes} suffice for the completed
Lebesgue integrals. Clipped elliptic values and their monotone limit
give the corresponding measurable second-moment version. Thus
\begin{equation}
 \begin{split}
 \mathcal P_f(t)
 &=\sum_j\int_{Q_j}\theta_j(z)^s
       I_{j,\mathrm w}(z,t\theta_j(z)^{2/q})\,dz\\
 &=\mathcal P_{\mathrm w}(t)=\mathcal P_2(t).
 \end{split}
 \label{eq:general-finsler-common-pressure}
\end{equation}
By \eqref{general-gluing-uniform-pressure-bounds}, the absolute
integrands are bounded by
\[
                   C\theta_j(z)^s+|t|\|f\|_\infty\theta_j(z).
\]
This is integrable on the bounded finite atlas. No measurable
choice of root minimizers, cube orientations, or screening scales
is necessary: the pointwise comparison has already established a
scalar value identity before integration. In particular, Hall-slack
reciprocals, focal constants, current norms and boundary rates do
not enter the integrable majorant.

\emph{Global iid convergence and differentiation.}
Proposition~\ref{general-gluing-weighted-lower} gives, for each
fixed \(t\) and \(\varepsilon>0\),
\[
 \Pp\{V_n(t)<\mathcal P_{\mathrm w}(t)-\varepsilon\}\longrightarrow0.
\]
Its rooted scenes use the original independent observations and retain
both microscopic partners. Fractional optimizers are represented by
an optimal Birkhoff draw before the independent root; arbitrary
selected point clouds are not assigned independent input laws.
Proposition~\ref{general-gluing-second-moment-upper} gives
\[
                   \limsup_n\E V_n(t)\le\mathcal P_2(t).
\]
Its radial clipping retains all original marks and uses the exact
distance inequality \eqref{general-gluing-exact-clipping}, without
an amplified additive clipping error. Its compact decorated model
satisfies Theorem~\ref{measurable-forward-pressure} and
Proposition~\ref{measurable-forward-decorated}. The proved order is
sample size first, clipping removal second, and omitted core mass
last. Count excesses are repaired at the existing \(q>4\) scale.
None of these estimates imposes an endpoint interior-support condition.

Since \(V_n(t)\ge-|t|\|f\|_\infty\), the probability lower bound
makes its negative part relative to the common limit vanish in
mean. The expectation upper bound then makes its positive part
vanish in mean. Thus
\begin{equation}
 V_n(t)\longrightarrow\mathcal P_f(t)
                 \quad\text{in probability and in }L^1.
 \label{eq:general-finsler-pressure-limit}
\end{equation}
This is precisely Proposition~
\ref{general-gluing-conditional-readout} with its relaxation premise
now supplied by Theorem~\ref{thm:general-focal-value-equality}.
Corollary~\ref{weighted-capped-integrated-c1}, or the equal
second-moment version in
Corollary~\ref{degenerate-scalar-integrated-c1}, gives
\(\mathcal P_f\in C^1(\R)\). Its derivative integrand is bounded
by \(\theta_j(z)\|f\|_\infty\), whose total integral is
\(\|f\|_\infty\). No additional uniform endpoint estimate is needed.

\emph{Readout of the whole fractional face.}
Proposition~\ref{general-gluing-conditional-readout}, applied to
\eqref{eq:general-finsler-pressure-limit} and the differentiable
common pressure, gives
\begin{equation}
 \sup_{\gamma\in\operatorname{Opt}_n}
       \left|\int f\,d\gamma-\mathcal P_f'(0)\right|
                   \longrightarrow0\quad\text{in probability}.
 \label{eq:general-finsler-readout}
\end{equation}
The estimate is simultaneous over all empirical optima, including
fractional ones.

\emph{Deterministic intrinsic reconstruction.}
Choose a countable uniformly dense family \((f_k)\) in
\(C(A\times B)\), including \(1\), and put
\(c_k=\mathcal P_{f_k}'(0)\). Select a measurable empirical
optimum, for example the first minimizing permutation in a fixed
finite ordering. A deterministic subsequence makes all its countably
many readouts converge almost surely simultaneously, and both
empirical marginals converge almost surely.

The calibrated slack of these optima also tends to zero almost
surely. Continuity of \(d_F\) on the fixed compact product gives
convergence of empirical optimal costs to the population optimum,
while the two potential integrals converge by the law of large
numbers. Their difference is the nonnegative slack \(e\).
Compactness now gives a subsequential limiting probability measure
with marginals \(\mu,\nu\), integrals \(c_k\), and zero slack.
It is a population optimum. The dense family determines it uniquely;
since every \(c_k\) is deterministic, every realization in the
full-probability subsequence event gives the same measure \(\pi_*\).
No measurable field of local minimizing laws has been postulated.

One compatible weak metric is
\[
 d_0(\pi,\pi')=\sum_{k\ge1}2^{-k}
 \frac{\left|\int f_k\,d\pi-\int f_k\,d\pi'\right|}
      {1+2\|f_k\|_\infty}.
\]
Truncate the sum, use the finitely many estimates
\eqref{eq:general-finsler-readout} uniformly over the entire optimal
set, and bound the remaining geometric tail. This proves
\eqref{eq:general-finsler-entire-selector} for \(d_0\).
Compactness makes all compatible weak metrics uniformly equivalent.
The empirical optimal polytope and the compact maximum theorem
also give measurability of the displayed suprema.

Different auxiliary potentials or coordinates yield deterministic
limits of the same original empirical optimization problem.
The triangle inequality forces those limits to agree. Hence the
physical plan is intrinsic, and the proof is complete.
\end{proof}

\begin{remark}[Geometric and dimensional scope]
\label{rem:general-finsler-scope}
Theorem~\ref{thm:general-finsler-selector} covers all admissible
populations in Euclidean space and every complete smooth Riemannian
manifold, as well as smooth complete reversible strongly convex
Finsler tori. No global nonvanishing field or tangent-bundle
trivialization is chosen: local seeds and normal frames are handled
by the finite measurable atlas. The theorem requires no curvature
sign, global convex normal neighborhood, global cut-locus avoidance,
or additional homotopy condition. Its local regularity conclusions
concern the actual calibrated pairs and are derived from the
distance problem itself.

The conclusion is convergence in probability of the full optimal
set. It does not identify the limit with an entropy selector or a
Monge map. Dimensions \(d=4,5\), incomplete or nonreversible
metrics, touching or noncompact supports, and singular marginal
populations are outside the proved statement.
\end{remark}

\begin{corollary}[Isometric transport of empirical selection]
\label{cor:isometry}
Let \(J:(M,d)\to(M',d')\) be an isometry onto its image and let the
two marginals have compact support. If their empirical optimal sets
converge to a deterministic plan \(\pi_*\), then the pushed-forward
problem has deterministic selector \((J\times J)_\#\pi_*\).
\end{corollary}

\begin{proof}
Under the coupling of samples \(X'_i=J(X_i)\), \(Y'_j=J(Y_j)\),
the finite cost matrices agree entry by entry. Pushforward maps the
entire optimal set bijectively to the new one. On the compact product
supports, pushforward by \(J\times J\) is continuous for weak
convergence, hence uniformly continuous on their compact probability
spaces. This transfers convergence uniformly over the optimal sets.
\end{proof}


\subsection{Scope and formalization}

The theorem uses the original length distance and independently sampled
marginals throughout. The temporary Euclidean coordinates describe local
geometry and sampling; no global trivialization, prescribed marginal
balance or sign of curvature is required. The conclusion is a selected
population plan. An entropy characterization, Monge representation and
the general cases \(d=4,5\) remain separate questions.

The Lean development follows the general Finsler proof. Its physical
mean-pressure selector endpoint in
\texttt{FinslerPhysicalMeanSelector.lean} constructs one common potential,
pressure derivatives and the entire-face selector from the original
geometric and population data, conditional on local pointwise convergence
of the actual mean pressures. The conditional geometric, finite-optimizer,
energy and readout witnesses are constructed internally. Mean-pressure
convergence remains an explicit analytic premise of that endpoint, so the
complete general theorem is not yet formally verified.

\begin{thebibliography}{99}

\bibitem{Lyu2026EmpiricalW1}
W.~Lyu.
Deterministic Selection in Empirical Wasserstein-One Matching.
\emph{Preprints.org}, version~1, 12 October 2026.
\href{https://doi.org/10.20944/preprints202610.1018.v1}{doi:10.20944/preprints202610.1018.v1}.

\bibitem{GoldmanTrevisan2023}
M.~Goldman and D.~Trevisan.
Optimal transport methods for combinatorial optimization over two random point sets.
\emph{Probability Theory and Related Fields} \textbf{188} (2024), 1315--1384.
\href{https://doi.org/10.1007/s00440-023-01245-1}{doi:10.1007/s00440-023-01245-1}.
\href{https://arxiv.org/abs/2209.14615v2}{arXiv:2209.14615v2}.

\bibitem{GoldmanTrevisan2021}
M.~Goldman and D.~Trevisan.
Convergence of asymptotic costs for random Euclidean matching problems.
\emph{Probability and Mathematical Physics} \textbf{2} (2021), 341--362.
\href{https://arxiv.org/abs/2009.04128v2}{arXiv:2009.04128v2}.

\bibitem{AmirAngelHolroyd}
G.~Amir, O.~Angel, and A.~E.~Holroyd.
Multicolour Poisson matching.
\emph{Annales de l'Institut Henri Poincar\'e, Probabilit\'es et Statistiques}
\textbf{57} (2021), 1811--1833.
\href{https://doi.org/10.1214/21-AIHP1147}{doi:10.1214/21-AIHP1147}.
\href{https://arxiv.org/abs/1605.06485v2}{arXiv:1605.06485v2}.

\bibitem{Assouline}
R.~Assouline.
Curvature-Dimension for Autonomous Lagrangians.
Preprint, 2024, revised 2026.
\href{https://arxiv.org/abs/2409.08001v6}{arXiv:2409.08001v6}.


\bibitem{BhowmickPrasad}
A.~Bhowmick and S.~Prasad.
On the cut locus of submanifolds of a Finsler manifold.
\emph{The Journal of Geometric Analysis} \textbf{34} (2024), article~308.
\href{https://doi.org/10.1007/s12220-024-01751-1}{doi:10.1007/s12220-024-01751-1}.

\bibitem{OhtaNeedles}
S.-i.~Ohta.
Needle decompositions and isoperimetric inequalities in Finsler geometry.
\emph{J. Math. Soc. Japan} \textbf{70} (2018), no.~2, 651--693.
\href{https://doi.org/10.2969/jmsj/07027604}{doi:10.2969/jmsj/07027604}.

\bibitem{FeldmanMcCannDensity}
M.~Feldman and R.~J.~McCann.
Uniqueness and transport density in Monge's mass transportation problem.
\emph{Calculus of Variations and Partial Differential Equations}
\textbf{15} (2002), 81--113.
\href{https://doi.org/10.1007/s005260100119}{doi:10.1007/s005260100119}.

\bibitem{NutzZhong}
M.~Nutz and C.~Zhong.
Entropic regularization of Monge's problem.
Preprint, 2026.
\href{https://arxiv.org/abs/2604.21578}{arXiv:2604.21578}, version~2.

\bibitem{delBarrioEtAl2025}
E.~del Barrio, A.~Gonz\'alez-Sanz, J.-M.~Loubes, and
D.~Rodr\'iguez-V\'itores.
Distributional limit theory for optimal transport.
Preprint, 2025, revised 2026.
\href{https://arxiv.org/abs/2505.19104}{arXiv:2505.19104}.

\bibitem{RichterGebertSturmfelsTheobald}
J.~Richter-Gebert, B.~Sturmfels, and T.~Theobald.
First steps in tropical geometry.
\href{https://arxiv.org/abs/math/0306366v2}{arXiv:math/\allowbreak0306366v2}, 2003.

\bibitem{CaraccioloSicuroSportiello}
S.~Caracciolo, G.~Sicuro, and A.~Sportiello.
Spanning trees in the Assignment Problem: two theorems and two conjectures.
\href{https://arxiv.org/abs/2608.22308v1}{arXiv:2608.22308v1}, 2026.

\end{thebibliography}
\end{document}
